\documentclass[14pt,a4paper,oneside]{extbook}
\usepackage[utf8]{vietnam}
\usepackage{amsmath,amssymb,mathrsfs,fancyhdr,enumerate,multirow,makecell,venndiagram}
\usepackage{tkz-euclide,tikz-3dplot,tikz,tkz-tab}
\usetikzlibrary{arrows,calc,intersections,angles,quotes,backgrounds,shapes.geometric,patterns,through}
\usepackage{yhmath}
\usepackage{pgfplots}
\pgfplotsset{compat=1.9}
\usepgfplotslibrary{fillbetween}
\usepackage[tikz]{bclogo}
\usepackage[top=1.8cm, bottom=1.8cm, left=1.8cm, right=1.8cm, headheight=16pt]{geometry}
\usepackage[hidelinks,unicode]{hyperref}
\usepackage{qrcode}
\usepackage[many]{tcolorbox}

\usepackage[loigiai]{ex_test}
\providecommand{\shortans}[1]{\par\vspace{1.5mm}\noindent\textbf{Đáp số: } \framebox{\bfseries\color{red!85!black}#1}\par\vspace{2mm}}

\def\colorEX{\color{blue}}
\everymath{\color{blue}}

\usepackage{fontawesome}
\usepackage{framed}
\usepackage{polynom}

\pagestyle{fancy}
\fancyhf{}
\fancyhead[L]{\small\itshape\color{gray} Tuyển tập Bài tập Toán 3 theo từng bài}
\fancyhead[R]{\small\bfseries\color{blue!80!black} Lớp Toán Thầy Lĩnh -- loptoanthaylinh.vn}
\fancyfoot[C]{\bfseries\thepage}
\renewcommand{\headrulewidth}{0.4pt}

\setcounter{secnumdepth}{4}
\setcounter{tocdepth}{2}
\usepackage{titlesec}
\usepackage{titletoc}

\makeatletter
\renewcommand\thesection{\@arabic\c@section}
\makeatother

\titleformat{\section}{\centering\Large\bfseries\color{blue!80!black}}{Chủ đề~\thesection.}{1ex}{
    \setcounter{ex}{0}
    \MakeUppercase
}
\titleformat{\subsection}{\large\bfseries\color{blue!70!black}}{}{0ex}{}

% KHUNG TIÊU ĐỀ PHIẾU BÀI TẬP KÈM QR CODE
\newcommand{\TaoPhieuBaiTap}[4]{%
  \refstepcounter{subsection}%
  \addcontentsline{toc}{subsection}{Bài #2. #3}%
  \setcounter{ex}{0}%
  \noindent
  \begin{tcolorbox}[
    colback=white,
    colframe=blue!85!black,
    arc=2.5mm,
    boxrule=1pt,
    left=2.5mm, right=2.5mm, top=1.8mm, bottom=1.8mm
  ]
    \begin{minipage}[c]{0.27\textwidth}
      \centering
      \begin{tcolorbox}[colback=blue!8, colframe=blue!80!black, boxrule=0.6pt, arc=1.5mm, halign=center, top=1.2mm, bottom=1.2mm, left=1mm, right=1mm]
        {\bfseries\color{blue!95!black}\small LỚP TOÁN THẦY LĨNH}\\[1.5pt]
        {\scriptsize\color{blue!70!black}\faGlobe~loptoanthaylinh.vn}\\[1.5pt]
        {\scriptsize\bfseries\color{blue!85!black}\faPhone~Hotline: 0866976878}
      \end{tcolorbox}
    \end{minipage}%
    \hfill
    \begin{minipage}[c]{0.56\textwidth}
      \centering
      {\footnotesize\bfseries\color{blue!80!black}\MakeUppercase{#1}}\\[2pt]
      {\large\bfseries\color{red!85!black}\MakeUppercase{BÀI #2: #3}}\\[2pt]
      {\scriptsize\itshape Phiếu Bài Tập Luyện Tập -- Hệ Thống Bài Tập Theo Buổi}
    \end{minipage}%
    \hfill
    \begin{minipage}[c]{0.14\textwidth}
      \centering
      \qrcode[height=1.35cm]{#4}\\[2pt]
      {\tiny\bfseries\color{blue!90!black}\faQrcode~Quét QR học}
    \end{minipage}
  \end{tcolorbox}
  \vspace{2mm}
}

\begin{document}

\tableofcontents
\newpage


% =========================================================================
% Ôn tập và bổ sung
% =========================================================================
\section{Ôn tập và bổ sung}

\TaoPhieuBaiTap{Chủ đề 1 -- Ôn tập và bổ sung}{01}{Ôn tập các số trong phạm vi 1000}{https://loptoanthaylinh.vn/bai-tap/toan-3/bai-01}
\begin{ex}
Số liền sau của số 999 là số nào?
\shortans{1000}
\loigiai{
Số liền sau của 999 là $999 + 1 = 1000$.
}
\end{ex}

\begin{ex}
Số liền trước của số 500 là số nào?
\shortans{499}
\loigiai{
Số liền trước của 500 là $500 - 1 = 499$.
}
\end{ex}

\begin{ex}
Số gồm 6 trăm, 3 chục và 8 đơn vị viết là gì?
\shortans{638}
\loigiai{
Số gồm 6 trăm, 3 chục, 8 đơn vị là 638.
}
\end{ex}

\begin{ex}
Số 705 có chữ số hàng chục là bao nhiêu?
\shortans{0}
\loigiai{
Trong số 705, chữ số hàng chục là 0.
}
\end{ex}

\begin{ex}
Số lớn nhất có ba chữ số khác nhau là số nào?
\shortans{987}
\loigiai{
Số lớn nhất có 3 chữ số khác nhau là 987.
}
\end{ex}

\begin{ex}
Số bé nhất có ba chữ số khác nhau là số nào?
\shortans{102}
\loigiai{
Số bé nhất có 3 chữ số khác nhau là 102.
}
\end{ex}

\begin{ex}
Số tròn trăm liền sau của số 400 là số nào?
\shortans{500}
\loigiai{
Số tròn trăm liền sau của 400 là 500.
}
\end{ex}

\begin{ex}
Số tròn chục liền trước của số 320 là số nào?
\shortans{310}
\loigiai{
Số tròn chục liền trước của 320 là 310.
}
\end{ex}

\begin{ex}
Trong các số $485, 548, 458, 584$, số lớn nhất là số nào?
\shortans{584}
\loigiai{
So sánh hàng trăm rồi đến hàng chục: 584 là số lớn nhất.
}
\end{ex}

\begin{ex}
Trong các số $612, 621, 609, 619$, số bé nhất là số nào?
\shortans{609}
\loigiai{
Số bé nhất trong các số đã cho là 609.
}
\end{ex}

\begin{ex}
Cho các chữ số 2, 5, 9. Số lớn nhất có ba chữ số ghép từ cả ba chữ số này là số nào?
\shortans{952}
\loigiai{
Sắp xếp theo thứ tự giảm dần: 952.
}
\end{ex}

\begin{ex}
Cho các chữ số 0, 4, 7. Số bé nhất có ba chữ số khác nhau ghép từ ba chữ số này là số nào?
\shortans{407}
\loigiai{
Chữ số hàng trăm khác 0 nhỏ nhất là 4, tiếp theo là 0 và 7: 407.
}
\end{ex}

\begin{ex}
Tìm $x$ biết: $x$ là số liền trước của 780.
\shortans{779}
\loigiai{
Số liền trước của 780 là $780 - 1 = 779$.
}
\end{ex}

\begin{ex}
Tìm $y$ biết $y$ là số tròn chục nằm giữa 855 và 865.
\shortans{860}
\loigiai{
Số tròn chục duy nhất nằm giữa 855 và 865 là 860.
}
\end{ex}

\begin{ex}
Có bao nhiêu số tròn chục có hai chữ số?
\shortans{9}
\loigiai{
Các số đó là 10, 20, 30, 40, 50, 60, 70, 80, 90: có 9 số.
}
\end{ex}

\begin{ex}
Có bao nhiêu số tròn trăm từ 100 đến 1000?
\shortans{10}
\loigiai{
Các số: 100, 200, 300, 400, 500, 600, 700, 800, 900, 1000 (10 số).
}
\end{ex}

\begin{ex}
Viết số thích hợp vào chỗ chấm: $600 + 40 + 5 = \dots$
\shortans{645}
\loigiai{
$600 + 40 + 5 = 645$.
}
\end{ex}

\begin{ex}
Viết số thích hợp vào chỗ chấm: $900 + 8 = \dots$
\shortans{908}
\loigiai{
$900 + 8 = 908$.
}
\end{ex}

\begin{ex}
Số 354 gồm mấy chục nếu chỉ xét giá trị hàng chục?
\shortans{5}
\loigiai{
Chữ số hàng chục của 354 là 5.
}
\end{ex}

\begin{ex}
Điền số thích hợp: $299 < \dots < 301$.
\shortans{300}
\loigiai{
Số nằm giữa 299 và 301 là 300.
}
\end{ex}

\begin{ex}
Tổng của số lớn nhất có ba chữ số và số bé nhất có một chữ số là bao nhiêu?
\shortans{999}
\loigiai{
Số lớn nhất có ba chữ số là 999, số bé nhất có một chữ số là 0: $999 + 0 = 999$.
}
\end{ex}

\begin{ex}
Hiệu giữa số bé nhất có ba chữ số và số lớn nhất có hai chữ số là bao nhiêu?
\shortans{1}
\loigiai{
Số bé nhất có ba chữ số là 100, số lớn nhất có hai chữ số là 99: $100 - 99 = 1$.
}
\end{ex}

\begin{ex}
Từ ba chữ số 1, 2, 3 có thể lập được bao nhiêu số có ba chữ số khác nhau?
\shortans{6}
\loigiai{
Các số: 123, 132, 213, 231, 312, 321 (tổng cộng 6 số).
}
\end{ex}

\begin{ex}
Một đoàn tàu có 8 toa chở khách, mỗi toa có 100 chỗ ngồi. Hỏi đoàn tàu chở được tối đa bao nhiêu hành khách?
\shortans{800}
\loigiai{
$8 \times 100 = 800$ chỗ ngồi.
}
\end{ex}

\begin{ex}
Số nào cộng với 350 thì được kết quả là 500?
\shortans{150}
\loigiai{
$500 - 350 = 150$.
}
\end{ex}

\newpage

\TaoPhieuBaiTap{Chủ đề 1 -- Ôn tập và bổ sung}{02}{Ôn tập phép cộng, phép trừ trong phạm vi 1000}{https://loptoanthaylinh.vn/bai-tap/toan-3/bai-02}
\begin{ex}
Tính: $350 + 240 = \dots$
\shortans{590}
\loigiai{
$350 + 240 = 590$.
}
\end{ex}

\begin{ex}
Tính: $780 - 450 = \dots$
\shortans{330}
\loigiai{
$780 - 450 = 330$.
}
\end{ex}

\begin{ex}
Tính: $425 + 138 = \dots$
\shortans{563}
\loigiai{
$425 + 138 = 563$.
}
\end{ex}

\begin{ex}
Tính: $652 - 237 = \dots$
\shortans{415}
\loigiai{
$652 - 237 = 415$.
}
\end{ex}

\begin{ex}
Tính giá trị của biểu thức: $500 + 200 - 300$.
\shortans{400}
\loigiai{
$500 + 200 - 300 = 700 - 300 = 400$.
}
\end{ex}

\begin{ex}
Tìm kết quả của phép tính: $846 - 325$.
\shortans{521}
\loigiai{
$846 - 325 = 521$.
}
\end{ex}

\begin{ex}
Tính tổng của 415 và 279.
\shortans{694}
\loigiai{
$415 + 279 = 694$.
}
\end{ex}

\begin{ex}
Tính hiệu của 900 và 456.
\shortans{444}
\loigiai{
$900 - 456 = 444$.
}
\end{ex}

\begin{ex}
Đội Một trồng được 245 cây, đội Hai trồng được 318 cây. Hỏi cả hai đội trồng được bao nhiêu cây?
\shortans{563}
\loigiai{
$245 + 318 = 563$ cây.
}
\end{ex}

\begin{ex}
Một cửa hàng buổi sáng bán được 450 kg gạo, buổi chiều bán được ít hơn buổi sáng 120 kg gạo. Hỏi buổi chiều bán được bao nhiêu ki-lô-gam gạo?
\shortans{330}
\loigiai{
$450 - 120 = 330$ kg.
}
\end{ex}

\begin{ex}
Một thư viện có 850 quyển sách, đã cho mượn 375 quyển. Hỏi thư viện còn lại bao nhiêu quyển sách?
\shortans{475}
\loigiai{
$850 - 375 = 475$ quyển sách.
}
\end{ex}

\begin{ex}
Điền số thích hợp: $240 + \dots = 670$.
\shortans{430}
\loigiai{
$670 - 240 = 430$.
}
\end{ex}

\begin{ex}
Điền số thích hợp: $\dots - 150 = 420$.
\shortans{570}
\loigiai{
$420 + 150 = 570$.
}
\end{ex}

\begin{ex}
Điền số thích hợp: $800 - \dots = 340$.
\shortans{460}
\loigiai{
$800 - 340 = 460$.
}
\end{ex}

\begin{ex}
Tính nhanh: $250 + 180 + 50$.
\shortans{480}
\loigiai{
$250 + 50 + 180 = 300 + 180 = 480$.
}
\end{ex}

\begin{ex}
Tính nhanh: $670 - 240 - 70$.
\shortans{360}
\loigiai{
$670 - 70 - 240 = 600 - 240 = 360$.
}
\end{ex}

\begin{ex}
Quãng đường từ Hà Nội đến Phủ Lý dài 60 km, từ Phủ Lý đến Ninh Bình dài 35 km. Hỏi quãng đường từ Hà Nội đến Ninh Bình dài bao nhiêu ki-lô-mét?
\shortans{95}
\loigiai{
$60 + 35 = 95$ km.
}
\end{ex}

\begin{ex}
Một tấm vải dài 300 m, người ta cắt bán đi 125 m. Hỏi tấm vải còn lại dài bao nhiêu mét?
\shortans{175}
\loigiai{
$300 - 125 = 175$ m.
}
\end{ex}

\begin{ex}
Bác An thu hoạch được 560 kg thóc nếp và 380 kg thóc tẻ. Hỏi bác An thu hoạch tất cả bao nhiêu ki-lô-gam thóc?
\shortans{940}
\loigiai{
$560 + 380 = 940$ kg.
}
\end{ex}

\begin{ex}
Tính tổng của số tròn trăm lớn nhất có ba chữ số và số tròn trăm bé nhất có ba chữ số.
\shortans{1000}
\loigiai{
$900 + 100 = 1000$.
}
\end{ex}

\begin{ex}
Một xe ô tô chở 650 thùng hàng, đã dỡ xuống 280 thùng. Hỏi trên xe còn lại bao nhiêu thùng hàng?
\shortans{370}
\loigiai{
$650 - 280 = 370$ thùng.
}
\end{ex}

\begin{ex}
Giá trị của chữ số 7 trong số 745 lớn hơn giá trị của chữ số 4 trong số đó bao nhiêu đơn vị?
\shortans{660}
\loigiai{
Giá trị chữ số 7 là 700, giá trị chữ số 4 là 40: $700 - 40 = 660$.
}
\end{ex}

\begin{ex}
Tổng của hai số là 450, nếu số hạng thứ nhất tăng thêm 50 đơn vị và giữ nguyên số hạng thứ hai thì tổng mới là bao nhiêu?
\shortans{500}
\loigiai{
$450 + 50 = 500$.
}
\end{ex}

\begin{ex}
Hiệu của hai số là 200, nếu số bị trừ bớt đi 30 đơn vị và giữ nguyên số trừ thì hiệu mới là bao nhiêu?
\shortans{170}
\loigiai{
$200 - 30 = 170$.
}
\end{ex}

\begin{ex}
Tìm một số biết rằng lấy số đó trừ đi 235 thì được 465.
\shortans{700}
\loigiai{
$465 + 235 = 700$.
}
\end{ex}

\newpage

\TaoPhieuBaiTap{Chủ đề 1 -- Ôn tập và bổ sung}{03}{Tìm thành phần trong phép cộng, phép trừ}{https://loptoanthaylinh.vn/bai-tap/toan-3/bai-03}
\begin{ex}
Tìm $x$, biết: $x + 125 = 458$.
\shortans{333}
\loigiai{
$x = 458 - 125 = 333$.
}
\end{ex}

\begin{ex}
Tìm $x$, biết: $246 + x = 689$.
\shortans{443}
\loigiai{
$x = 689 - 246 = 443$.
}
\end{ex}

\begin{ex}
Tìm $x$, biết: $x - 314 = 275$.
\shortans{589}
\loigiai{
$x = 275 + 314 = 589$.
}
\end{ex}

\begin{ex}
Tìm $x$, biết: $758 - x = 423$.
\shortans{335}
\loigiai{
$x = 758 - 423 = 335$.
}
\end{ex}

\begin{ex}
Tìm số hạng chưa biết trong phép cộng có tổng là 560 và một số hạng là 230.
\shortans{330}
\loigiai{
$560 - 230 = 330$.
}
\end{ex}

\begin{ex}
Tìm số bị trừ trong phép trừ có hiệu là 145 và số trừ là 320.
\shortans{465}
\loigiai{
$145 + 320 = 465$.
}
\end{ex}

\begin{ex}
Tìm số trừ trong phép trừ có số bị trừ là 800 và hiệu là 350.
\shortans{450}
\loigiai{
$800 - 350 = 450$.
}
\end{ex}

\begin{ex}
Tìm $y$, biết: $y + 400 = 920$.
\shortans{520}
\loigiai{
$y = 920 - 400 = 520$.
}
\end{ex}

\begin{ex}
Tìm $y$, biết: $650 - y = 180$.
\shortans{470}
\loigiai{
$y = 650 - 180 = 470$.
}
\end{ex}

\begin{ex}
Tìm $y$, biết: $y - 250 = 350$.
\shortans{600}
\loigiai{
$y = 350 + 250 = 600$.
}
\end{ex}

\begin{ex}
Một người có một số quả bưởi. Sau khi bán đi 120 quả thì còn lại 80 quả. Hỏi ban đầu người đó có bao nhiêu quả bưởi?
\shortans{200}
\loigiai{
$80 + 120 = 200$ quả.
}
\end{ex}

\begin{ex}
Mẹ mua một bó hoa gồm cả hoa hồng và hoa cúc tổng cộng 45 bông, trong đó có 28 bông hoa hồng. Hỏi có bao nhiêu bông hoa cúc?
\shortans{17}
\loigiai{
$45 - 28 = 17$ bông.
}
\end{ex}

\begin{ex}
Hùng có 65 viên bi. Hùng cho bạn một số viên bi thì còn lại 38 viên bi. Hỏi Hùng đã cho bạn bao nhiêu viên bi?
\shortans{27}
\loigiai{
$65 - 38 = 27$ viên bi.
}
\end{ex}

\begin{ex}
Một cửa hàng nhập về một số bao gạo, sau khi bán đi 75 bao thì còn lại 125 bao. Hỏi ban đầu cửa hàng nhập về bao nhiêu bao gạo?
\shortans{200}
\loigiai{
$125 + 75 = 200$ bao gạo.
}
\end{ex}

\begin{ex}
Tìm số bị trừ biết số trừ là 215 và hiệu là 485.
\shortans{700}
\loigiai{
$485 + 215 = 700$.
}
\end{ex}

\begin{ex}
Tìm số trừ biết số bị trừ là 900 và hiệu là 540.
\shortans{360}
\loigiai{
$900 - 540 = 360$.
}
\end{ex}

\begin{ex}
Tìm $x$, biết: $x + 45 = 100 - 20$.
\shortans{35}
\loigiai{
$x + 45 = 80 \Rightarrow x = 80 - 45 = 35$.
}
\end{ex}

\begin{ex}
Tìm $x$, biết: $x - 30 = 40 + 50$.
\shortans{120}
\loigiai{
$x - 30 = 90 \Rightarrow x = 90 + 30 = 120$.
}
\end{ex}

\begin{ex}
Tìm $x$, biết: $150 - x = 20 + 60$.
\shortans{70}
\loigiai{
$150 - x = 80 \Rightarrow x = 150 - 80 = 70$.
}
\end{ex}

\begin{ex}
Tìm một số biết rằng nếu lấy 500 trừ đi số đó rồi cộng với 100 thì được 300.
\shortans{300}
\loigiai{
$500 - x + 100 = 300 \Rightarrow 500 - x = 200 \Rightarrow x = 300$.
}
\end{ex}

\begin{ex}
Tổng của hai số là 800. Nếu thêm vào số hạng thứ nhất 50 đơn vị thì số hạng thứ hai phải bớt đi bao nhiêu đơn vị để tổng không đổi?
\shortans{50}
\loigiai{
Để tổng không đổi, số hạng thứ hai cũng phải bớt đi 50 đơn vị.
}
\end{ex}

\begin{ex}
Tìm $x$, biết: $x + 312 = 500 + 120$.
\shortans{308}
\loigiai{
$x + 312 = 620 \Rightarrow x = 620 - 312 = 308$.
}
\end{ex}

\begin{ex}
Tìm $x$, biết: $x - 145 = 600 - 250$.
\shortans{495}
\loigiai{
$x - 145 = 350 \Rightarrow x = 350 + 145 = 495$.
}
\end{ex}

\begin{ex}
Biết số bị trừ hơn hiệu 128 đơn vị. Số trừ là bao nhiêu?
\shortans{128}
\loigiai{
Số trừ = Số bị trừ - Hiệu = 128.
}
\end{ex}

\begin{ex}
Biết tổng hơn số hạng thứ nhất 240 đơn vị. Số hạng thứ hai là bao nhiêu?
\shortans{240}
\loigiai{
Số hạng thứ hai = Tổng - Số hạng thứ nhất = 240.
}
\end{ex}

\newpage

\TaoPhieuBaiTap{Chủ đề 1 -- Ôn tập và bổ sung}{04}{Ôn tập bảng nhân 2, bảng nhân 5}{https://loptoanthaylinh.vn/bai-tap/toan-3/bai-04}
\begin{ex}
Tính: $2 \times 8 = \dots$
\shortans{16}
\loigiai{
$2 \times 8 = 16$.
}
\end{ex}

\begin{ex}
Tính: $5 \times 7 = \dots$
\shortans{35}
\loigiai{
$5 \times 7 = 35$.
}
\end{ex}

\begin{ex}
Tính: $18 : 2 = \dots$
\shortans{9}
\loigiai{
$18 : 2 = 9$.
}
\end{ex}

\begin{ex}
Tính: $45 : 5 = \dots$
\shortans{9}
\loigiai{
$45 : 5 = 9$.
}
\end{ex}

\begin{ex}
Tính: $2 \times 9 + 15 = \dots$
\shortans{33}
\loigiai{
$2 \times 9 + 15 = 18 + 15 = 33$.
}
\end{ex}

\begin{ex}
Tính: $5 \times 8 - 18 = \dots$
\shortans{22}
\loigiai{
$5 \times 8 - 18 = 40 - 18 = 22$.
}
\end{ex}

\begin{ex}
Tính: $20 : 2 + 35 = \dots$
\shortans{45}
\loigiai{
$10 + 35 = 45$.
}
\end{ex}

\begin{ex}
Tính: $50 : 5 - 6 = \dots$
\shortans{4}
\loigiai{
$10 - 6 = 4$.
}
\end{ex}

\begin{ex}
Mỗi đôi đũa có 2 chiếc đũa. Hỏi 9 đôi đũa có bao nhiêu chiếc đũa?
\shortans{18}
\loigiai{
$2 \times 9 = 18$ chiếc.
}
\end{ex}

\begin{ex}
Có 16 học sinh chia đều thành 2 nhóm. Hỏi mỗi nhóm có bao nhiêu học sinh?
\shortans{8}
\loigiai{
$16 : 2 = 8$ học sinh.
}
\end{ex}

\begin{ex}
Mỗi tuần lễ mẹ đi làm 5 ngày. Hỏi trong 6 tuần lễ mẹ đi làm tất cả bao nhiêu ngày?
\shortans{30}
\loigiai{
$5 \times 6 = 30$ ngày.
}
\end{ex}

\begin{ex}
Có 35 lít dầu chia đều vào các can, mỗi can chứa 5 lít. Hỏi rót được bao nhiêu can dầu?
\shortans{7}
\loigiai{
$35 : 5 = 7$ can.
}
\end{ex}

\begin{ex}
Có 20 cái kẹo chia đều cho 5 bạn. Hỏi mỗi bạn được mấy cái kẹo?
\shortans{4}
\loigiai{
$20 : 5 = 4$ cái kẹo.
}
\end{ex}

\begin{ex}
Điền số thích hợp: $2 \times \dots = 14$.
\shortans{7}
\loigiai{
$14 : 2 = 7$.
}
\end{ex}

\begin{ex}
Điền số thích hợp: $5 \times \dots = 40$.
\shortans{8}
\loigiai{
$40 : 5 = 8$.
}
\end{ex}

\begin{ex}
Điền số thích hợp: $\dots : 2 = 6$.
\shortans{12}
\loigiai{
$6 \times 2 = 12$.
}
\end{ex}

\begin{ex}
Điền số thích hợp: $\dots : 5 = 5$.
\shortans{25}
\loigiai{
$5 \times 5 = 25$.
}
\end{ex}

\begin{ex}
Lớp 3A có 5 tổ, mỗi tổ có 6 bạn. Hỏi lớp 3A có bao nhiêu bạn?
\shortans{30}
\loigiai{
$5 \times 6 = 30$ bạn.
}
\end{ex}

\begin{ex}
Mỗi bàn học có 2 học sinh ngồi. Có tất cả 10 bàn học. Hỏi có bao nhiêu học sinh ngồi?
\shortans{20}
\loigiai{
$2 \times 10 = 20$ học sinh.
}
\end{ex}

\begin{ex}
Tính: $2 \times 6 : 2 = \dots$
\shortans{6}
\loigiai{
$12 : 2 = 6$.
}
\end{ex}

\begin{ex}
Tính: $5 \times 4 : 2 = \dots$
\shortans{10}
\loigiai{
$20 : 2 = 10$.
}
\end{ex}

\begin{ex}
Một con gà có 2 chân. Hỏi 8 con gà có tất cả bao nhiêu chân?
\shortans{16}
\loigiai{
$2 \times 8 = 16$ chân.
}
\end{ex}

\begin{ex}
Mỗi bàn tay có 5 ngón tay. Hỏi 2 bàn tay có bao nhiêu ngón tay?
\shortans{10}
\loigiai{
$5 \times 2 = 10$ ngón tay.
}
\end{ex}

\begin{ex}
Có 14 quả cam xếp đều vào 2 đĩa. Hỏi mỗi đĩa có bao nhiêu quả cam?
\shortans{7}
\loigiai{
$14 : 2 = 7$ quả cam.
}
\end{ex}

\begin{ex}
Có 30 bông hoa cắm đều vào 5 lọ. Hỏi mỗi lọ có bao nhiêu bông hoa?
\shortans{6}
\loigiai{
$30 : 5 = 6$ bông hoa.
}
\end{ex}

\newpage

\TaoPhieuBaiTap{Chủ đề 1 -- Ôn tập và bổ sung}{05}{Ôn tập bảng chia 2, bảng chia 5}{https://loptoanthaylinh.vn/bai-tap/toan-3/bai-05}
\begin{ex}
Tính: $3 \times 4 = \dots$
\shortans{12}
\loigiai{
$3 \times 4 = 12$.
}
\end{ex}

\begin{ex}
Tính: $3 \times 7 = \dots$
\shortans{21}
\loigiai{
$3 \times 7 = 21$.
}
\end{ex}

\begin{ex}
Tính: $3 \times 9 = \dots$
\shortans{27}
\loigiai{
$3 \times 9 = 27$.
}
\end{ex}

\begin{ex}
Tính: $15 : 3 = \dots$
\shortans{5}
\loigiai{
$15 : 3 = 5$.
}
\end{ex}

\begin{ex}
Tính: $24 : 3 = \dots$
\shortans{8}
\loigiai{
$24 : 3 = 8$.
}
\end{ex}

\begin{ex}
Tính: $27 : 3 = \dots$
\shortans{9}
\loigiai{
$27 : 3 = 9$.
}
\end{ex}

\begin{ex}
Tính: $3 \times 6 + 18 = \dots$
\shortans{36}
\loigiai{
$3 \times 6 + 18 = 18 + 18 = 36$.
}
\end{ex}

\begin{ex}
Tính: $3 \times 8 - 14 = \dots$
\shortans{10}
\loigiai{
$3 \times 8 - 14 = 24 - 14 = 10$.
}
\end{ex}

\begin{ex}
Tính: $21 : 3 + 25 = \dots$
\shortans{32}
\loigiai{
$7 + 25 = 32$.
}
\end{ex}

\begin{ex}
Tính: $30 : 3 - 7 = \dots$
\shortans{3}
\loigiai{
$10 - 7 = 3$.
}
\end{ex}

\begin{ex}
Mỗi tam giác có 3 cạnh. Hỏi 8 hình tam giác như thế có tất cả bao nhiêu cạnh?
\shortans{24}
\loigiai{
$3 \times 8 = 24$ cạnh.
}
\end{ex}

\begin{ex}
Có 18 quyển vở chia đều cho 3 bạn. Hỏi mỗi bạn nhận được bao nhiêu quyển vở?
\shortans{6}
\loigiai{
$18 : 3 = 6$ quyển vở.
}
\end{ex}

\begin{ex}
Mỗi bạn gấp được 3 con hạc giấy. Hỏi 9 bạn gấp được tất cả bao nhiêu con hạc giấy?
\shortans{27}
\loigiai{
$3 \times 9 = 27$ con hạc.
}
\end{ex}

\begin{ex}
Cô giáo có 21 bút chì thưởng đều cho 3 tổ. Hỏi mỗi tổ nhận được bao nhiêu bút chì?
\shortans{7}
\loigiai{
$21 : 3 = 7$ bút chì.
}
\end{ex}

\begin{ex}
Điền số thích hợp: $3 \times \dots = 18$.
\shortans{6}
\loigiai{
$18 : 3 = 6$.
}
\end{ex}

\begin{ex}
Điền số thích hợp: $3 \times \dots = 30$.
\shortans{10}
\loigiai{
$30 : 3 = 10$.
}
\end{ex}

\begin{ex}
Điền số thích hợp: $\dots : 3 = 4$.
\shortans{12}
\loigiai{
$4 \times 3 = 12$.
}
\end{ex}

\begin{ex}
Điền số thích hợp: $\dots : 3 = 7$.
\shortans{21}
\loigiai{
$7 \times 3 = 21$.
}
\end{ex}

\begin{ex}
Tính: $3 \times 5 + 3 \times 2 = \dots$
\shortans{21}
\loigiai{
$15 + 6 = 21$ (hoặc $3 \times 7 = 21$).
}
\end{ex}

\begin{ex}
Một bình hoa cắm 3 cành hoa. Hỏi để cắm 5 bình hoa như thế cần bao nhiêu cành hoa?
\shortans{15}
\loigiai{
$3 \times 5 = 15$ cành hoa.
}
\end{ex}

\begin{ex}
Có 12 quả táo chia đều cho 3 anh em. Hỏi mỗi người nhận được mấy quả táo?
\shortans{4}
\loigiai{
$12 : 3 = 4$ quả táo.
}
\end{ex}

\begin{ex}
Tìm $x$, biết: $x \times 3 = 12 + 6$.
\shortans{6}
\loigiai{
$x \times 3 = 18 \Rightarrow x = 18 : 3 = 6$.
}
\end{ex}

\begin{ex}
Tìm $y$, biết: $y : 3 = 10 - 2$.
\shortans{24}
\loigiai{
$y : 3 = 8 \Rightarrow y = 8 \times 3 = 24$.
}
\end{ex}

\begin{ex}
Một sợi dây đồng dài 27 cm được cắt thành 3 đoạn bằng nhau. Hỏi mỗi đoạn dài bao nhiêu xăng-ti-mét?
\shortans{9}
\loigiai{
$27 : 3 = 9$ cm.
}
\end{ex}

\begin{ex}
Tính: $3 \times 10 - 15 = \dots$
\shortans{15}
\loigiai{
$30 - 15 = 15$.
}
\end{ex}

\newpage

\TaoPhieuBaiTap{Chủ đề 1 -- Ôn tập và bổ sung}{06}{Gấp một số lên một số lần}{https://loptoanthaylinh.vn/bai-tap/toan-3/bai-06}
\begin{ex}
Tính: $4 \times 3 = \dots$
\shortans{12}
\loigiai{
$4 \times 3 = 12$.
}
\end{ex}

\begin{ex}
Tính: $4 \times 6 = \dots$
\shortans{24}
\loigiai{
$4 \times 6 = 24$.
}
\end{ex}

\begin{ex}
Tính: $4 \times 8 = \dots$
\shortans{32}
\loigiai{
$4 \times 8 = 32$.
}
\end{ex}

\begin{ex}
Tính: $4 \times 9 = \dots$
\shortans{36}
\loigiai{
$4 \times 9 = 36$.
}
\end{ex}

\begin{ex}
Tính: $16 : 4 = \dots$
\shortans{4}
\loigiai{
$16 : 4 = 4$.
}
\end{ex}

\begin{ex}
Tính: $28 : 4 = \dots$
\shortans{7}
\loigiai{
$28 : 4 = 7$.
}
\end{ex}

\begin{ex}
Tính: $36 : 4 = \dots$
\shortans{9}
\loigiai{
$36 : 4 = 9$.
}
\end{ex}

\begin{ex}
Tính: $40 : 4 = \dots$
\shortans{10}
\loigiai{
$40 : 4 = 10$.
}
\end{ex}

\begin{ex}
Tính giá trị của biểu thức: $4 \times 7 + 25$.
\shortans{53}
\loigiai{
$4 \times 7 + 25 = 28 + 25 = 53$.
}
\end{ex}

\begin{ex}
Tính giá trị của biểu thức: $4 \times 5 - 12$.
\shortans{8}
\loigiai{
$4 \times 5 - 12 = 20 - 12 = 8$.
}
\end{ex}

\begin{ex}
Tính giá trị của biểu thức: $32 : 4 + 19$.
\shortans{27}
\loigiai{
$8 + 19 = 27$.
}
\end{ex}

\begin{ex}
Một con bàn học có 4 chân. Hỏi 7 cái bàn như thế có bao nhiêu cái chân?
\shortans{28}
\loigiai{
$4 \times 7 = 28$ chân.
}
\end{ex}

\begin{ex}
Một hình vuông có 4 cạnh bằng nhau. Nếu mỗi cạnh dài 6 cm thì tổng độ dài 4 cạnh của hình vuông là bao nhiêu xăng-ti-mét?
\shortans{24}
\loigiai{
$4 \times 6 = 24$ cm.
}
\end{ex}

\begin{ex}
Có 32 học sinh xếp đều thành 4 hàng. Hỏi mỗi hàng có bao nhiêu học sinh?
\shortans{8}
\loigiai{
$32 : 4 = 8$ học sinh.
}
\end{ex}

\begin{ex}
Có 20 cái bánh chia đều vào 4 hộp. Hỏi mỗi hộp có bao nhiêu cái bánh?
\shortans{5}
\loigiai{
$20 : 4 = 5$ cái bánh.
}
\end{ex}

\begin{ex}
Điền số thích hợp: $4 \times \dots = 20$.
\shortans{5}
\loigiai{
$20 : 4 = 5$.
}
\end{ex}

\begin{ex}
Điền số thích hợp: $4 \times \dots = 32$.
\shortans{8}
\loigiai{
$32 : 4 = 8$.
}
\end{ex}

\begin{ex}
Điền số thích hợp: $\dots : 4 = 6$.
\shortans{24}
\loigiai{
$6 \times 4 = 24$.
}
\end{ex}

\begin{ex}
Điền số thích hợp: $\dots : 4 = 9$.
\shortans{36}
\loigiai{
$9 \times 4 = 36$.
}
\end{ex}

\begin{ex}
Một ô tô con có 4 bánh xe. Hỏi 9 ô tô con như thế có tất cả bao nhiêu bánh xe?
\shortans{36}
\loigiai{
$4 \times 9 = 36$ bánh xe.
}
\end{ex}

\begin{ex}
Có 24 quả cam chia đều cho 4 người. Hỏi mỗi người được mấy quả cam?
\shortans{6}
\loigiai{
$24 : 4 = 6$ quả cam.
}
\end{ex}

\begin{ex}
Tìm $x$, biết: $x \times 4 = 28$.
\shortans{7}
\loigiai{
$x = 28 : 4 = 7$.
}
\end{ex}

\begin{ex}
Tìm $x$, biết: $x : 4 = 5$.
\shortans{20}
\loigiai{
$x = 5 \times 4 = 20$.
}
\end{ex}

\begin{ex}
Tính: $4 \times 4 : 2 = \dots$
\shortans{8}
\loigiai{
$16 : 2 = 8$.
}
\end{ex}

\begin{ex}
Tính: $40 : 4 \times 3 = \dots$
\shortans{30}
\loigiai{
$10 \times 3 = 30$.
}
\end{ex}

\newpage

\TaoPhieuBaiTap{Chủ đề 1 -- Ôn tập và bổ sung}{07}{Giảm đi một số lần}{https://loptoanthaylinh.vn/bai-tap/toan-3/bai-07}
\begin{ex}
Một hình tam giác có độ dài ba cạnh lần lượt là 12 cm, 15 cm và 18 cm. Tính chu vi hình tam giác đó theo đơn vị xăng-ti-mét.
\shortans{45}
\loigiai{
$12 + 15 + 18 = 45$ cm.
}
\end{ex}

\begin{ex}
Một hình tứ giác có độ dài các cạnh lần lượt là 10 cm, 20 cm, 15 cm và 25 cm. Chu vi hình tứ giác đó là bao nhiêu xăng-ti-mét?
\shortans{70}
\loigiai{
$10 + 20 + 15 + 25 = 70$ cm.
}
\end{ex}

\begin{ex}
Một đoạn thẳng AB dài 15 cm, đoạn thẳng CD dài hơn đoạn AB là 7 cm. Hỏi đoạn thẳng CD dài bao nhiêu xăng-ti-mét?
\shortans{22}
\loigiai{
$15 + 7 = 22$ cm.
}
\end{ex}

\begin{ex}
Một sợi dây dài 80 cm, người ta cắt đi 35 cm. Hỏi đoạn dây còn lại dài bao nhiêu xăng-ti-mét?
\shortans{45}
\loigiai{
$80 - 35 = 45$ cm.
}
\end{ex}

\begin{ex}
Đổi: $1\text{ m} = \dots\text{ cm}$. Điền số thích hợp.
\shortans{100}
\loigiai{
$1\text{ m} = 100\text{ cm}$.
}
\end{ex}

\begin{ex}
Đổi: $1\text{ m} = \dots\text{ dm}$. Điền số thích hợp.
\shortans{10}
\loigiai{
$1\text{ m} = 10\text{ dm}$.
}
\end{ex}

\begin{ex}
Đổi: $3\text{ m } 20\text{ cm} = \dots\text{ cm}$. Điền số thích hợp.
\shortans{320}
\loigiai{
$3\text{ m } 20\text{ cm} = 300 + 20 = 320\text{ cm}$.
}
\end{ex}

\begin{ex}
Đổi: $5\text{ dm } 4\text{ cm} = \dots\text{ cm}$. Điền số thích hợp.
\shortans{54}
\loigiai{
$5\text{ dm } 4\text{ cm} = 50 + 4 = 54\text{ cm}$.
}
\end{ex}

\begin{ex}
Một can đựng 10 lít nước, một xô đựng 5 lít nước. Hỏi cả can và xô đựng bao nhiêu lít nước?
\shortans{15}
\loigiai{
$10 + 5 = 15$ lít.
}
\end{ex}

\begin{ex}
Mẹ mua 25 kg gạo nếp và 30 kg gạo tẻ. Hỏi mẹ mua tất cả bao nhiêu ki-lô-gam gạo?
\shortans{55}
\loigiai{
$25 + 30 = 55$ kg.
}
\end{ex}

\begin{ex}
Một túi đường nặng 2 kg. Hỏi 5 túi đường như thế nặng tất cả bao nhiêu ki-lô-gam?
\shortans{10}
\loigiai{
$2 \times 5 = 10$ kg.
}
\end{ex}

\begin{ex}
Một thùng chứa 40 lít dầu, người ta đã lấy ra 16 lít. Hỏi trong thùng còn lại bao nhiêu lít dầu?
\shortans{24}
\loigiai{
$40 - 16 = 24$ lít.
}
\end{ex}

\begin{ex}
Hình vuông có cạnh dài 5 cm. Tính tổng độ dài 4 cạnh của hình vuông đó theo đơn vị xăng-ti-mét.
\shortans{20}
\loigiai{
$5 \times 4 = 20$ cm.
}
\end{ex}

\begin{ex}
Hình chữ nhật có chiều dài 8 cm, chiều rộng 5 cm. Nửa chu vi hình chữ nhật đó là bao nhiêu xăng-ti-mét?
\shortans{13}
\loigiai{
Nửa chu vi = chiều dài + chiều rộng = $8 + 5 = 13$ cm.
}
\end{ex}

\begin{ex}
Có bao nhiêu cạnh trong một hình tứ giác?
\shortans{4}
\loigiai{
Hình tứ giác có 4 cạnh.
}
\end{ex}

\begin{ex}
Có bao nhiêu đỉnh trong một hình tam giác?
\shortans{3}
\loigiai{
Hình tam giác có 3 đỉnh.
}
\end{ex}

\begin{ex}
Một khối lập phương có bao nhiêu mặt?
\shortans{6}
\loigiai{
Khối lập phương có 6 mặt.
}
\end{ex}

\begin{ex}
Một khối lập phương có bao nhiêu đỉnh?
\shortans{8}
\loigiai{
Khối lập phương có 8 đỉnh.
}
\end{ex}

\begin{ex}
Một khối hộp chữ nhật có bao nhiêu cạnh?
\shortans{12}
\loigiai{
Khối hộp chữ nhật có 12 cạnh.
}
\end{ex}

\begin{ex}
Đường gấp khúc gồm 3 đoạn thẳng có độ dài lần lượt là 14 cm, 16 cm và 20 cm. Tính độ dài đường gấp khúc đó theo đơn vị xăng-ti-mét.
\shortans{50}
\loigiai{
$14 + 16 + 20 = 50$ cm.
}
\end{ex}

\begin{ex}
Đoạn thẳng MN dài 40 cm. Điểm P là trung điểm của MN. Hỏi đoạn thẳng MP dài bao nhiêu xăng-ti-mét?
\shortans{20}
\loigiai{
$40 : 2 = 20$ cm.
}
\end{ex}

\begin{ex}
Rót hết 18 lít mật ong vào các can loại 2 lít. Hỏi rót được bao nhiêu can mật ong?
\shortans{9}
\loigiai{
$18 : 2 = 9$ can.
}
\end{ex}

\begin{ex}
Có 3 gói kẹo, mỗi gói cân nặng 200 g. Hỏi cả 3 gói kẹo cân nặng bao nhiêu gam?
\shortans{600}
\loigiai{
$200 \times 3 = 600$ g.
}
\end{ex}

\begin{ex}
Một chiếc thước kẻ dài 30 cm, bút chì dài 18 cm. Chiếc thước kẻ dài hơn bút chì bao nhiêu xăng-ti-mét?
\shortans{12}
\loigiai{
$30 - 18 = 12$ cm.
}
\end{ex}

\begin{ex}
Chu vi hình tam giác đều có cạnh 7 cm là bao nhiêu xăng-ti-mét?
\shortans{21}
\loigiai{
$7 \times 3 = 21$ cm.
}
\end{ex}

\newpage

\TaoPhieuBaiTap{Chủ đề 1 -- Ôn tập và bổ sung}{08}{Luyện tập chung}{https://loptoanthaylinh.vn/bai-tap/toan-3/bai-08}
\begin{ex}
Tính: $354 + 215 = \dots$
\shortans{569}
\loigiai{
$354 + 215 = 569$.
}
\end{ex}

\begin{ex}
Tính: $687 - 342 = \dots$
\shortans{345}
\loigiai{
$687 - 342 = 345$.
}
\end{ex}

\begin{ex}
Tính: $3 \times 7 + 120 = \dots$
\shortans{141}
\loigiai{
$21 + 120 = 141$.
}
\end{ex}

\begin{ex}
Tính: $4 \times 8 - 15 = \dots$
\shortans{17}
\loigiai{
$32 - 15 = 17$.
}
\end{ex}

\begin{ex}
Tính: $36 : 4 + 45 = \dots$
\shortans{54}
\loigiai{
$9 + 45 = 54$.
}
\end{ex}

\begin{ex}
Tính: $27 : 3 + 128 = \dots$
\shortans{137}
\loigiai{
$9 + 128 = 137$.
}
\end{ex}

\begin{ex}
Tìm số lớn nhất trong các số: $879, 897, 789, 798$.
\shortans{897}
\loigiai{
897 là số lớn nhất.
}
\end{ex}

\begin{ex}
Tìm số bé nhất trong các số: $512, 521, 502, 520$.
\shortans{502}
\loigiai{
502 là số bé nhất.
}
\end{ex}

\begin{ex}
Số liền trước của 1000 là số nào?
\shortans{999}
\loigiai{
$1000 - 1 = 999$.
}
\end{ex}

\begin{ex}
Số liền sau của 689 là số nào?
\shortans{690}
\loigiai{
$689 + 1 = 690$.
}
\end{ex}

\begin{ex}
Tìm $x$, biết: $x + 234 = 568$.
\shortans{334}
\loigiai{
$x = 568 - 234 = 334$.
}
\end{ex}

\begin{ex}
Tìm $x$, biết: $x - 145 = 320$.
\shortans{465}
\loigiai{
$x = 320 + 145 = 465$.
}
\end{ex}

\begin{ex}
Tìm $x$, biết: $4 \times x = 36$.
\shortans{9}
\loigiai{
$x = 36 : 4 = 9$.
}
\end{ex}

\begin{ex}
Tìm $x$, biết: $x : 3 = 8$.
\shortans{24}
\loigiai{
$x = 8 \times 3 = 24$.
}
\end{ex}

\begin{ex}
Một trường tiểu học có 420 học sinh nữ và 390 học sinh nam. Hỏi trường đó có tất cả bao nhiêu học sinh?
\shortans{810}
\loigiai{
$420 + 390 = 810$ học sinh.
}
\end{ex}

\begin{ex}
Một cửa hàng có 500 kg đường, đã bán 235 kg. Hỏi cửa hàng còn lại bao nhiêu ki-lô-gam đường?
\shortans{265}
\loigiai{
$500 - 235 = 265$ kg.
}
\end{ex}

\begin{ex}
Mỗi bàn học có 2 bạn ngồi. Hỏi 8 bàn học như thế có tất cả bao nhiêu bạn?
\shortans{16}
\loigiai{
$2 \times 8 = 16$ bạn.
}
\end{ex}

\begin{ex}
Có 28 quả cam chia đều vào 4 giỏ. Hỏi mỗi giỏ có bao nhiêu quả cam?
\shortans{7}
\loigiai{
$28 : 4 = 7$ quả cam.
}
\end{ex}

\begin{ex}
Một tam giác có độ dài 3 cạnh đều bằng 8 cm. Tính chu vi tam giác đó theo đơn vị xăng-ti-mét.
\shortans{24}
\loigiai{
$8 \times 3 = 24$ cm.
}
\end{ex}

\begin{ex}
Một hình tứ giác có độ dài 4 cạnh đều bằng 5 cm. Tính chu vi tứ giác đó theo đơn vị xăng-ti-mét.
\shortans{20}
\loigiai{
$5 \times 4 = 20$ cm.
}
\end{ex}

\begin{ex}
Tính: $2 \times 9 + 4 \times 5 = \dots$
\shortans{38}
\loigiai{
$18 + 20 = 38$.
}
\end{ex}

\begin{ex}
Tính: $30 : 5 + 24 : 3 = \dots$
\shortans{14}
\loigiai{
$6 + 8 = 14$.
}
\end{ex}

\begin{ex}
Số tròn chục lớn nhất có ba chữ số là số nào?
\shortans{990}
\loigiai{
Số đó là 990.
}
\end{ex}

\begin{ex}
Số gồm 5 trăm và 7 đơn vị viết là số nào?
\shortans{507}
\loigiai{
Số đó là 507.
}
\end{ex}

\begin{ex}
Một sợi dây dài 1 m, người ta cắt bớt đi 30 cm. Đoạn dây còn lại dài bao nhiêu xăng-ti-mét?
\shortans{70}
\loigiai{
$1\text{ m} = 100\text{ cm}$; $100 - 30 = 70$ cm.
}
\end{ex}

\newpage


% =========================================================================
% Bảng nhân, bảng chia
% =========================================================================
\section{Bảng nhân, bảng chia}

\TaoPhieuBaiTap{Chủ đề 2 -- Bảng nhân, bảng chia}{09}{Bảng nhân 6, bảng chia 6}{https://loptoanthaylinh.vn/bai-tap/toan-3/bai-09}
\begin{ex}
Tính: $6 \times 3 = \dots$
\shortans{18}
\loigiai{
$6 \times 3 = 18$.
}
\end{ex}

\begin{ex}
Tính: $6 \times 6 = \dots$
\shortans{36}
\loigiai{
$6 \times 6 = 36$.
}
\end{ex}

\begin{ex}
Tính: $6 \times 7 = \dots$
\shortans{42}
\loigiai{
$6 \times 7 = 42$.
}
\end{ex}

\begin{ex}
Tính: $6 \times 8 = \dots$
\shortans{48}
\loigiai{
$6 \times 8 = 48$.
}
\end{ex}

\begin{ex}
Tính: $6 \times 9 = \dots$
\shortans{54}
\loigiai{
$6 \times 9 = 54$.
}
\end{ex}

\begin{ex}
Tính: $18 : 6 = \dots$
\shortans{3}
\loigiai{
$18 : 6 = 3$.
}
\end{ex}

\begin{ex}
Tính: $30 : 6 = \dots$
\shortans{5}
\loigiai{
$30 : 6 = 5$.
}
\end{ex}

\begin{ex}
Tính: $42 : 6 = \dots$
\shortans{7}
\loigiai{
$42 : 6 = 7$.
}
\end{ex}

\begin{ex}
Tính: $54 : 6 = \dots$
\shortans{9}
\loigiai{
$54 : 6 = 9$.
}
\end{ex}

\begin{ex}
Tính: $60 : 6 = \dots$
\shortans{10}
\loigiai{
$60 : 6 = 10$.
}
\end{ex}

\begin{ex}
Tính giá trị biểu thức: $6 \times 4 + 26$.
\shortans{50}
\loigiai{
$24 + 26 = 50$.
}
\end{ex}

\begin{ex}
Tính giá trị biểu thức: $6 \times 5 - 15$.
\shortans{15}
\loigiai{
$30 - 15 = 15$.
}
\end{ex}

\begin{ex}
Tính giá trị biểu thức: $48 : 6 + 12$.
\shortans{20}
\loigiai{
$8 + 12 = 20$.
}
\end{ex}

\begin{ex}
Tính giá trị biểu thức: $36 : 6 \times 4$.
\shortans{24}
\loigiai{
$6 \times 4 = 24$.
}
\end{ex}

\begin{ex}
Mỗi hộp bánh có 6 chiếc bánh. Hỏi 5 hộp bánh như thế có bao nhiêu chiếc bánh?
\shortans{30}
\loigiai{
$6 \times 5 = 30$ chiếc bánh.
}
\end{ex}

\begin{ex}
Một tuần bạn Mai học bơi 6 buổi. Hỏi trong 4 tuần bạn Mai học bơi bao nhiêu buổi?
\shortans{24}
\loigiai{
$6 \times 4 = 24$ buổi.
}
\end{ex}

\begin{ex}
Có 48 chiếc ghế xếp đều thành 6 hàng. Hỏi mỗi hàng có bao nhiêu chiếc ghế?
\shortans{8}
\loigiai{
$48 : 6 = 8$ chiếc ghế.
}
\end{ex}

\begin{ex}
Một cuộn vải dài 54 m, người ta cắt thành 6 đoạn bằng nhau. Hỏi mỗi đoạn dài bao nhiêu mét?
\shortans{9}
\loigiai{
$54 : 6 = 9$ m.
}
\end{ex}

\begin{ex}
Điền số thích hợp: $6 \times \dots = 42$.
\shortans{7}
\loigiai{
$42 : 6 = 7$.
}
\end{ex}

\begin{ex}
Điền số thích hợp: $\dots : 6 = 6$.
\shortans{36}
\loigiai{
$6 \times 6 = 36$.
}
\end{ex}

\begin{ex}
Tìm $x$, biết: $x \times 6 = 24$.
\shortans{4}
\loigiai{
$x = 24 : 6 = 4$.
}
\end{ex}

\begin{ex}
Tìm $x$, biết: $x : 6 = 8$.
\shortans{48}
\loigiai{
$x = 8 \times 6 = 48$.
}
\end{ex}

\begin{ex}
Một lớp có 36 học sinh chia đều thành 6 tổ. Hỏi mỗi tổ có bao nhiêu học sinh?
\shortans{6}
\loigiai{
$36 : 6 = 6$ học sinh.
}
\end{ex}

\begin{ex}
Mỗi con côn trùng có 6 chân. Hỏi 7 con côn trùng như vậy có tất cả bao nhiêu chân?
\shortans{42}
\loigiai{
$6 \times 7 = 42$ chân.
}
\end{ex}

\begin{ex}
Tính: $6 \times 2 + 6 \times 8 = \dots$
\shortans{60}
\loigiai{
$12 + 48 = 60$ (hoặc $6 \times 10 = 60$).
}
\end{ex}

\newpage

\TaoPhieuBaiTap{Chủ đề 2 -- Bảng nhân, bảng chia}{10}{Bảng nhân 7, bảng chia 7}{https://loptoanthaylinh.vn/bai-tap/toan-3/bai-10}
\begin{ex}
Tính: $7 \times 4 = \dots$
\shortans{28}
\loigiai{
$7 \times 4 = 28$.
}
\end{ex}

\begin{ex}
Tính: $7 \times 6 = \dots$
\shortans{42}
\loigiai{
$7 \times 6 = 42$.
}
\end{ex}

\begin{ex}
Tính: $7 \times 7 = \dots$
\shortans{49}
\loigiai{
$7 \times 7 = 49$.
}
\end{ex}

\begin{ex}
Tính: $7 \times 8 = \dots$
\shortans{56}
\loigiai{
$7 \times 8 = 56$.
}
\end{ex}

\begin{ex}
Tính: $7 \times 9 = \dots$
\shortans{63}
\loigiai{
$7 \times 9 = 63$.
}
\end{ex}

\begin{ex}
Tính: $21 : 7 = \dots$
\shortans{3}
\loigiai{
$21 : 7 = 3$.
}
\end{ex}

\begin{ex}
Tính: $35 : 7 = \dots$
\shortans{5}
\loigiai{
$35 : 7 = 5$.
}
\end{ex}

\begin{ex}
Tính: $42 : 7 = \dots$
\shortans{6}
\loigiai{
$42 : 7 = 6$.
}
\end{ex}

\begin{ex}
Tính: $56 : 7 = \dots$
\shortans{8}
\loigiai{
$56 : 7 = 8$.
}
\end{ex}

\begin{ex}
Tính: $63 : 7 = \dots$
\shortans{9}
\loigiai{
$63 : 7 = 9$.
}
\end{ex}

\begin{ex}
Tính giá trị biểu thức: $7 \times 5 + 15$.
\shortans{50}
\loigiai{
$35 + 15 = 50$.
}
\end{ex}

\begin{ex}
Tính giá trị biểu thức: $7 \times 8 - 26$.
\shortans{30}
\loigiai{
$56 - 26 = 30$.
}
\end{ex}

\begin{ex}
Tính giá trị biểu thức: $49 : 7 + 33$.
\shortans{40}
\loigiai{
$7 + 33 = 40$.
}
\end{ex}

\begin{ex}
Tính giá trị biểu thức: $70 : 7 \times 5$.
\shortans{50}
\loigiai{
$10 \times 5 = 50$.
}
\end{ex}

\begin{ex}
Một tuần lễ có 7 ngày. Hỏi 6 tuần lễ có bao nhiêu ngày?
\shortans{42}
\loigiai{
$7 \times 6 = 42$ ngày.
}
\end{ex}

\begin{ex}
Một tuần lễ có 7 ngày. Hỏi 9 tuần lễ có bao nhiêu ngày?
\shortans{63}
\loigiai{
$7 \times 9 = 63$ ngày.
}
\end{ex}

\begin{ex}
Có 28 bông hoa cắm đều vào 7 bình hoa. Hỏi mỗi bình có bao nhiêu bông hoa?
\shortans{4}
\loigiai{
$28 : 7 = 4$ bông hoa.
}
\end{ex}

\begin{ex}
Có 35 kg gạo chia đều vào 7 túi. Hỏi mỗi túi có bao nhiêu ki-lô-gam gạo?
\shortans{5}
\loigiai{
$35 : 7 = 5$ kg gạo.
}
\end{ex}

\begin{ex}
Điền số thích hợp: $7 \times \dots = 49$.
\shortans{7}
\loigiai{
$49 : 7 = 7$.
}
\end{ex}

\begin{ex}
Điền số thích hợp: $\dots : 7 = 4$.
\shortans{28}
\loigiai{
$4 \times 7 = 28$.
}
\end{ex}

\begin{ex}
Tìm $x$, biết: $x \times 7 = 14 + 21$.
\shortans{5}
\loigiai{
$x \times 7 = 35 \Rightarrow x = 35 : 7 = 5$.
}
\end{ex}

\begin{ex}
Tìm $x$, biết: $x : 7 = 10 - 3$.
\shortans{49}
\loigiai{
$x : 7 = 7 \Rightarrow x = 7 \times 7 = 49$.
}
\end{ex}

\begin{ex}
Bác An mua 7 gói bánh, mỗi gói giá 8 nghìn đồng. Bác An phải trả tất cả bao nhiêu nghìn đồng?
\shortans{56}
\loigiai{
$8 \times 7 = 56$ nghìn đồng.
}
\end{ex}

\begin{ex}
Có 56 quyển vở chia đều cho 7 học sinh giỏi. Hỏi mỗi bạn được bao nhiêu quyển vở?
\shortans{8}
\loigiai{
$56 : 7 = 8$ quyển vở.
}
\end{ex}

\begin{ex}
Tính: $7 \times 3 + 7 \times 7 = \dots$
\shortans{70}
\loigiai{
$21 + 49 = 70$ (hoặc $7 \times 10 = 70$).
}
\end{ex}

\newpage

\TaoPhieuBaiTap{Chủ đề 2 -- Bảng nhân, bảng chia}{11}{Bảng nhân 8, bảng chia 8}{https://loptoanthaylinh.vn/bai-tap/toan-3/bai-11}
\begin{ex}
Tính: $8 \times 3 = \dots$
\shortans{24}
\loigiai{
$8 \times 3 = 24$.
}
\end{ex}

\begin{ex}
Tính: $8 \times 5 = \dots$
\shortans{40}
\loigiai{
$8 \times 5 = 40$.
}
\end{ex}

\begin{ex}
Tính: $8 \times 6 = \dots$
\shortans{48}
\loigiai{
$8 \times 6 = 48$.
}
\end{ex}

\begin{ex}
Tính: $8 \times 7 = \dots$
\shortans{56}
\loigiai{
$8 \times 7 = 56$.
}
\end{ex}

\begin{ex}
Tính: $8 \times 8 = \dots$
\shortans{64}
\loigiai{
$8 \times 8 = 64$.
}
\end{ex}

\begin{ex}
Tính: $8 \times 9 = \dots$
\shortans{72}
\loigiai{
$8 \times 9 = 72$.
}
\end{ex}

\begin{ex}
Tính: $24 : 8 = \dots$
\shortans{3}
\loigiai{
$24 : 8 = 3$.
}
\end{ex}

\begin{ex}
Tính: $40 : 8 = \dots$
\shortans{5}
\loigiai{
$40 : 8 = 5$.
}
\end{ex}

\begin{ex}
Tính: $56 : 8 = \dots$
\shortans{7}
\loigiai{
$56 : 8 = 7$.
}
\end{ex}

\begin{ex}
Tính: $64 : 8 = \dots$
\shortans{8}
\loigiai{
$64 : 8 = 8$.
}
\end{ex}

\begin{ex}
Tính: $72 : 8 = \dots$
\shortans{9}
\loigiai{
$72 : 8 = 9$.
}
\end{ex}

\begin{ex}
Tính giá trị biểu thức: $8 \times 4 + 18$.
\shortans{50}
\loigiai{
$32 + 18 = 50$.
}
\end{ex}

\begin{ex}
Tính giá trị biểu thức: $8 \times 6 - 28$.
\shortans{20}
\loigiai{
$48 - 28 = 20$.
}
\end{ex}

\begin{ex}
Tính giá trị biểu thức: $48 : 8 + 34$.
\shortans{40}
\loigiai{
$6 + 34 = 40$.
}
\end{ex}

\begin{ex}
Một con bạch tuộc có 8 xúc tu. Hỏi 5 con bạch tuộc có bao nhiêu xúc tu?
\shortans{40}
\loigiai{
$8 \times 5 = 40$ xúc tu.
}
\end{ex}

\begin{ex}
Mỗi cuộn len đan được một chiếc khăn dài 8 dm. Hỏi 6 cuộn len như thế đan được chiếc khăn dài bao nhiêu đề-xi-mét?
\shortans{48}
\loigiai{
$8 \times 6 = 48$ dm.
}
\end{ex}

\begin{ex}
Có 32 học sinh xếp đều vào 8 hàng. Hỏi mỗi hàng có bao nhiêu học sinh?
\shortans{4}
\loigiai{
$32 : 8 = 4$ học sinh.
}
\end{ex}

\begin{ex}
Một can dầu có 64 lít chia đều vào 8 thùng nhỏ. Hỏi mỗi thùng nhỏ có bao nhiêu lít dầu?
\shortans{8}
\loigiai{
$64 : 8 = 8$ lít.
}
\end{ex}

\begin{ex}
Điền số thích hợp: $8 \times \dots = 48$.
\shortans{6}
\loigiai{
$48 : 8 = 6$.
}
\end{ex}

\begin{ex}
Điền số thích hợp: $\dots : 8 = 9$.
\shortans{72}
\loigiai{
$9 \times 8 = 72$.
}
\end{ex}

\begin{ex}
Tìm $x$, biết: $x \times 8 = 56$.
\shortans{7}
\loigiai{
$x = 56 : 8 = 7$.
}
\end{ex}

\begin{ex}
Tìm $x$, biết: $x : 8 = 3$.
\shortans{24}
\loigiai{
$x = 3 \times 8 = 24$.
}
\end{ex}

\begin{ex}
Có 40 viên bi chia đều vào 8 túi. Hỏi mỗi túi có bao nhiêu viên bi?
\shortans{5}
\loigiai{
$40 : 8 = 5$ viên bi.
}
\end{ex}

\begin{ex}
Một tổ công nhân làm việc 8 giờ mỗi ngày. Hỏi sau 9 ngày tổ đó làm được tất cả bao nhiêu giờ?
\shortans{72}
\loigiai{
$8 \times 9 = 72$ giờ.
}
\end{ex}

\begin{ex}
Tính: $8 \times 8 : 8 = \dots$
\shortans{8}
\loigiai{
$64 : 8 = 8$.
}
\end{ex}

\newpage

\TaoPhieuBaiTap{Chủ đề 2 -- Bảng nhân, bảng chia}{12}{Bảng nhân 9, bảng chia 9}{https://loptoanthaylinh.vn/bai-tap/toan-3/bai-12}
\begin{ex}
Tính: $9 \times 2 = \dots$
\shortans{18}
\loigiai{
$9 \times 2 = 18$.
}
\end{ex}

\begin{ex}
Tính: $9 \times 4 = \dots$
\shortans{36}
\loigiai{
$9 \times 4 = 36$.
}
\end{ex}

\begin{ex}
Tính: $9 \times 5 = \dots$
\shortans{45}
\loigiai{
$9 \times 5 = 45$.
}
\end{ex}

\begin{ex}
Tính: $9 \times 7 = \dots$
\shortans{63}
\loigiai{
$9 \times 7 = 63$.
}
\end{ex}

\begin{ex}
Tính: $9 \times 8 = \dots$
\shortans{72}
\loigiai{
$9 \times 8 = 72$.
}
\end{ex}

\begin{ex}
Tính: $9 \times 9 = \dots$
\shortans{81}
\loigiai{
$9 \times 9 = 81$.
}
\end{ex}

\begin{ex}
Tính: $18 : 9 = \dots$
\shortans{2}
\loigiai{
$18 : 9 = 2$.
}
\end{ex}

\begin{ex}
Tính: $36 : 9 = \dots$
\shortans{4}
\loigiai{
$36 : 9 = 4$.
}
\end{ex}

\begin{ex}
Tính: $54 : 9 = \dots$
\shortans{6}
\loigiai{
$54 : 9 = 6$.
}
\end{ex}

\begin{ex}
Tính: $63 : 9 = \dots$
\shortans{7}
\loigiai{
$63 : 9 = 7$.
}
\end{ex}

\begin{ex}
Tính: $81 : 9 = \dots$
\shortans{9}
\loigiai{
$81 : 9 = 9$.
}
\end{ex}

\begin{ex}
Tính: $90 : 9 = \dots$
\shortans{10}
\loigiai{
$90 : 9 = 10$.
}
\end{ex}

\begin{ex}
Tính giá trị biểu thức: $9 \times 6 + 16$.
\shortans{70}
\loigiai{
$54 + 16 = 70$.
}
\end{ex}

\begin{ex}
Tính giá trị biểu thức: $9 \times 3 - 17$.
\shortans{10}
\loigiai{
$27 - 17 = 10$.
}
\end{ex}

\begin{ex}
Tính giá trị biểu thức: $72 : 9 + 42$.
\shortans{50}
\loigiai{
$8 + 42 = 50$.
}
\end{ex}

\begin{ex}
Một phòng học có 9 dãy bàn, mỗi dãy có 4 bạn ngồi. Hỏi phòng học đó có tất cả bao nhiêu bạn?
\shortans{36}
\loigiai{
$9 \times 4 = 36$ bạn.
}
\end{ex}

\begin{ex}
Cô giáo có 45 quyển vở thưởng đều cho 9 bạn học sinh. Hỏi mỗi bạn được bao nhiêu quyển vở?
\shortans{5}
\loigiai{
$45 : 9 = 5$ quyển vở.
}
\end{ex}

\begin{ex}
Có 54 lít nước mắm rót đều vào 9 can. Hỏi mỗi can có bao nhiêu lít nước mắm?
\shortans{6}
\loigiai{
$54 : 9 = 6$ lít.
}
\end{ex}

\begin{ex}
Điền số thích hợp: $9 \times \dots = 63$.
\shortans{7}
\loigiai{
$63 : 9 = 7$.
}
\end{ex}

\begin{ex}
Điền số thích hợp: $\dots : 9 = 8$.
\shortans{72}
\loigiai{
$8 \times 9 = 72$.
}
\end{ex}

\begin{ex}
Tìm $x$, biết: $x \times 9 = 81$.
\shortans{9}
\loigiai{
$x = 81 : 9 = 9$.
}
\end{ex}

\begin{ex}
Tìm $x$, biết: $x : 9 = 3$.
\shortans{27}
\loigiai{
$x = 3 \times 9 = 27$.
}
\end{ex}

\begin{ex}
Mỗi túi có 9 kg gạo. Hỏi 8 túi như vậy có tất cả bao nhiêu ki-lô-gam gạo?
\shortans{72}
\loigiai{
$9 \times 8 = 72$ kg.
}
\end{ex}

\begin{ex}
Một sợi dây dài 27 dm cắt thành 9 đoạn bằng nhau. Hỏi mỗi đoạn dài bao nhiêu đề-xi-mét?
\shortans{3}
\loigiai{
$27 : 9 = 3$ dm.
}
\end{ex}

\begin{ex}
Tính: $9 \times 5 + 9 \times 5 = \dots$
\shortans{90}
\loigiai{
$45 + 45 = 90$.
}
\end{ex}

\newpage

\TaoPhieuBaiTap{Chủ đề 2 -- Bảng nhân, bảng chia}{13}{Tìm thành phần trong phép nhân, phép chia}{https://loptoanthaylinh.vn/bai-tap/toan-3/bai-13}
\begin{ex}
Tìm thừa số chưa biết $x$: $x \times 6 = 48$.
\shortans{8}
\loigiai{
$x = 48 : 6 = 8$.
}
\end{ex}

\begin{ex}
Tìm thừa số chưa biết $x$: $7 \times x = 56$.
\shortans{8}
\loigiai{
$x = 56 : 7 = 8$.
}
\end{ex}

\begin{ex}
Tìm số bị chia $x$: $x : 8 = 5$.
\shortans{40}
\loigiai{
$x = 5 \times 8 = 40$.
}
\end{ex}

\begin{ex}
Tìm số chia $x$: $45 : x = 9$.
\shortans{5}
\loigiai{
$x = 45 : 9 = 5$.
}
\end{ex}

\begin{ex}
Tìm thừa số chưa biết trong tích có một thừa số bằng 9 và tích bằng 72.
\shortans{8}
\loigiai{
$72 : 9 = 8$.
}
\end{ex}

\begin{ex}
Tìm số bị chia biết thương là 7 và số chia là 6.
\shortans{42}
\loigiai{
$7 \times 6 = 42$.
}
\end{ex}

\begin{ex}
Tìm số chia biết số bị chia là 63 và thương là 9.
\shortans{7}
\loigiai{
$63 : 9 = 7$.
}
\end{ex}

\begin{ex}
Tìm $y$, biết: $y \times 8 = 64$.
\shortans{8}
\loigiai{
$y = 64 : 8 = 8$.
}
\end{ex}

\begin{ex}
Tìm $y$, biết: $y : 7 = 6$.
\shortans{42}
\loigiai{
$y = 6 \times 7 = 42$.
}
\end{ex}

\begin{ex}
Tìm $y$, biết: $54 : y = 6$.
\shortans{9}
\loigiai{
$y = 54 : 6 = 9$.
}
\end{ex}

\begin{ex}
Một người xếp 36 quả trứng vào các vỉ, mỗi vỉ có 6 quả. Hỏi xếp được bao nhiêu vỉ trứng?
\shortans{6}
\loigiai{
$36 : 6 = 6$ vỉ.
}
\end{ex}

\begin{ex}
Có một số lít dầu rót đều vào 8 can, mỗi can có 4 lít. Hỏi có tất cả bao nhiêu lít dầu?
\shortans{32}
\loigiai{
$4 \times 8 = 32$ lít.
}
\end{ex}

\begin{ex}
Một thanh gỗ dài được cưa thành 7 đoạn, mỗi đoạn dài 5 dm. Hỏi thanh gỗ ban đầu dài bao nhiêu đề-xi-mét?
\shortans{35}
\loigiai{
$5 \times 7 = 35$ dm.
}
\end{ex}

\begin{ex}
Có 48 học sinh chia đều thành các đội, mỗi đội có 8 học sinh. Hỏi chia được bao nhiêu đội?
\shortans{6}
\loigiai{
$48 : 8 = 6$ đội.
}
\end{ex}

\begin{ex}
Tìm $x$, biết: $x \times 4 = 15 + 17$.
\shortans{8}
\loigiai{
$x \times 4 = 32 \Rightarrow x = 32 : 4 = 8$.
}
\end{ex}

\begin{ex}
Tìm $x$, biết: $x : 5 = 20 - 13$.
\shortans{35}
\loigiai{
$x : 5 = 7 \Rightarrow x = 7 \times 5 = 35$.
}
\end{ex}

\begin{ex}
Tìm $x$, biết: $42 : x = 14 - 7$.
\shortans{6}
\loigiai{
$42 : x = 7 \Rightarrow x = 42 : 7 = 6$.
}
\end{ex}

\begin{ex}
Tích của hai số là 40, nếu một thừa số là 5 thì thừa số còn lại là bao nhiêu?
\shortans{8}
\loigiai{
$40 : 5 = 8$.
}
\end{ex}

\begin{ex}
Trong phép chia có thương là 9 và số chia là 4. Số bị chia là bao nhiêu?
\shortans{36}
\loigiai{
$9 \times 4 = 36$.
}
\end{ex}

\begin{ex}
Số bị chia là 56, thương là 7. Số chia là bao nhiêu?
\shortans{8}
\loigiai{
$56 : 7 = 8$.
}
\end{ex}

\begin{ex}
Tìm $y$, biết: $9 \times y = 30 + 33$.
\shortans{7}
\loigiai{
$9 \times y = 63 \Rightarrow y = 63 : 9 = 7$.
}
\end{ex}

\begin{ex}
Tìm $y$, biết: $y : 6 = 2 \times 4$.
\shortans{48}
\loigiai{
$y : 6 = 8 \Rightarrow y = 8 \times 6 = 48$.
}
\end{ex}

\begin{ex}
Nếu chia 24 chiếc kẹo cho một số bạn, mỗi bạn được 3 chiếc kẹo. Hỏi có bao nhiêu bạn được chia kẹo?
\shortans{8}
\loigiai{
$24 : 3 = 8$ bạn.
}
\end{ex}

\begin{ex}
Một cuộn dây thép chia đều cho 8 tổ, mỗi tổ nhận 7 mét thép. Hỏi cuộn dây dài bao nhiêu mét?
\shortans{56}
\loigiai{
$7 \times 8 = 56$ m.
}
\end{ex}

\begin{ex}
Thương của hai số là 8, số chia là 3. Hỏi số bị chia là bao nhiêu?
\shortans{24}
\loigiai{
$8 \times 3 = 24$.
}
\end{ex}

\newpage

\TaoPhieuBaiTap{Chủ đề 2 -- Bảng nhân, bảng chia}{14}{Một phần mấy (1/2, 1/3, 1/4, 1/5...)}{https://loptoanthaylinh.vn/bai-tap/toan-3/bai-14}
\begin{ex}
Một lớp có 32 học sinh, trong đó $\frac{1}{4}$ số học sinh là học sinh giỏi. Hỏi lớp đó có bao nhiêu học sinh giỏi?
\shortans{8}
\loigiai{
$32 : 4 = 8$ học sinh.
}
\end{ex}

\begin{ex}
Một khúc gỗ dài 24 m. Người ta cưa lấy $\frac{1}{3}$ khúc gỗ. Hỏi đoạn gỗ được cưa lấy dài bao nhiêu mét?
\shortans{8}
\loigiai{
$24 : 3 = 8$ m.
}
\end{ex}

\begin{ex}
Một giỏ có 35 quả cam. Mẹ lấy ra $\frac{1}{5}$ số cam để vắt nước. Hỏi mẹ đã lấy ra bao nhiêu quả cam?
\shortans{7}
\loigiai{
$35 : 5 = 7$ quả cam.
}
\end{ex}

\begin{ex}
Một ngày có 24 giờ. Một phần tư ($\frac{1}{4}$) của một ngày là bao nhiêu giờ?
\shortans{6}
\loigiai{
$24 : 4 = 6$ giờ.
}
\end{ex}

\begin{ex}
Một giờ có 60 phút. Một phần sáu ($\frac{1}{6}$) của một giờ là bao nhiêu phút?
\shortans{10}
\loigiai{
$60 : 6 = 10$ phút.
}
\end{ex}

\begin{ex}
Một sợi dây dài 42 cm. Hỏi $\frac{1}{7}$ sợi dây đó dài bao nhiêu xăng-ti-mét?
\shortans{6}
\loigiai{
$42 : 7 = 6$ cm.
}
\end{ex}

\begin{ex}
Một hộp có 48 viên bi. Nam lấy $\frac{1}{8}$ số bi trong hộp. Hỏi Nam lấy được bao nhiêu viên bi?
\shortans{6}
\loigiai{
$48 : 8 = 6$ viên bi.
}
\end{ex}

\begin{ex}
Đàn gà có 54 con, trong đó $\frac{1}{9}$ là gà trống. Hỏi có bao nhiêu con gà trống?
\shortans{6}
\loigiai{
$54 : 9 = 6$ con gà trống.
}
\end{ex}

\begin{ex}
Một cuộn vải dài 36 m. Người ta đã bán đi $\frac{1}{6}$ cuộn vải. Hỏi người ta đã bán bao nhiêu mét vải?
\shortans{6}
\loigiai{
$36 : 6 = 6$ m.
}
\end{ex}

\begin{ex}
Tìm $\frac{1}{2}$ của 18 kg.
\shortans{9}
\loigiai{
$18 : 2 = 9$ kg.
}
\end{ex}

\begin{ex}
Tìm $\frac{1}{3}$ của 27 lít.
\shortans{9}
\loigiai{
$27 : 3 = 9$ lít.
}
\end{ex}

\begin{ex}
Tìm $\frac{1}{7}$ của 49 chiếc bút.
\shortans{7}
\loigiai{
$49 : 7 = 7$ chiếc.
}
\end{ex}

\begin{ex}
Tìm $\frac{1}{8}$ của 64 bạn học sinh.
\shortans{8}
\loigiai{
$64 : 8 = 8$ bạn.
}
\end{ex}

\begin{ex}
Tìm $\frac{1}{9}$ của 72 quả bưởi.
\shortans{8}
\loigiai{
$72 : 9 = 8$ quả.
}
\end{ex}

\begin{ex}
Một cửa hàng có 40 bao gạo, đã bán $\frac{1}{5}$ số bao đó. Hỏi cửa hàng còn lại bao nhiêu bao gạo?
\shortans{32}
\loigiai{
Đã bán: $40 : 5 = 8$ bao. Còn lại: $40 - 8 = 32$ bao gạo.
}
\end{ex}

\begin{ex}
Một rổ táo có 24 quả. Hà ăn $\frac{1}{4}$ số táo đó. Hỏi rổ táo còn lại bao nhiêu quả?
\shortans{18}
\loigiai{
Hà ăn: $24 : 4 = 6$ quả. Còn lại: $24 - 6 = 18$ quả.
}
\end{ex}

\begin{ex}
Bác nông dân thu được 63 kg bí ngô và đã bán đi $\frac{1}{7}$ số bí ngô đó. Hỏi bác còn lại bao nhiêu ki-lô-gam bí ngô?
\shortans{54}
\loigiai{
Đã bán: $63 : 7 = 9$ kg. Còn lại: $63 - 9 = 54$ kg.
}
\end{ex}

\begin{ex}
Có 81 bông hoa xếp vào các bình, mỗi bình có $\frac{1}{9}$ số hoa. Hỏi mỗi bình có bao nhiêu bông hoa?
\shortans{9}
\loigiai{
$81 : 9 = 9$ bông hoa.
}
\end{ex}

\begin{ex}
Một mảnh vườn hình vuông có chu vi 28 m. Độ dài mỗi cạnh của mảnh vườn bằng $\frac{1}{4}$ chu vi. Hỏi cạnh mảnh vườn dài bao nhiêu mét?
\shortans{7}
\loigiai{
$28 : 4 = 7$ m.
}
\end{ex}

\begin{ex}
Tìm một số biết rằng $\frac{1}{6}$ của số đó bằng 7.
\shortans{42}
\loigiai{
Số đó là $7 \times 6 = 42$.
}
\end{ex}

\begin{ex}
Tìm một số biết rằng $\frac{1}{8}$ của số đó bằng 5.
\shortans{40}
\loigiai{
Số đó là $5 \times 8 = 40$.
}
\end{ex}

\begin{ex}
Tìm một số biết rằng $\frac{1}{9}$ của số đó bằng 9.
\shortans{81}
\loigiai{
Số đó là $9 \times 9 = 81$.
}
\end{ex}

\begin{ex}
Một bao gạo nặng 50 kg. Sau khi dùng $\frac{1}{5}$ bao gạo thì còn lại bao nhiêu ki-lô-gam?
\shortans{40}
\loigiai{
Đã dùng: $50 : 5 = 10$ kg. Còn lại: $50 - 10 = 40$ kg.
}
\end{ex}

\begin{ex}
Đội văn nghệ có 20 bạn, trong đó $\frac{1}{4}$ là bạn nam. Hỏi đội văn nghệ có bao nhiêu bạn nữ?
\shortans{15}
\loigiai{
Số bạn nam: $20 : 4 = 5$ bạn. Số bạn nữ: $20 - 5 = 15$ bạn.
}
\end{ex}

\begin{ex}
Biết $\frac{1}{3}$ của số $A$ là 8. Hỏi số $A$ là bao nhiêu?
\shortans{24}
\loigiai{
Số $A$ là $8 \times 3 = 24$.
}
\end{ex}

\newpage

\TaoPhieuBaiTap{Chủ đề 2 -- Bảng nhân, bảng chia}{15}{Luyện tập chung}{https://loptoanthaylinh.vn/bai-tap/toan-3/bai-15}
\begin{ex}
Tính: $6 \times 7 + 18 = \dots$
\shortans{60}
\loigiai{
$42 + 18 = 60$.
}
\end{ex}

\begin{ex}
Tính: $7 \times 8 - 26 = \dots$
\shortans{30}
\loigiai{
$56 - 26 = 30$.
}
\end{ex}

\begin{ex}
Tính: $8 \times 9 - 32 = \dots$
\shortans{40}
\loigiai{
$72 - 32 = 40$.
}
\end{ex}

\begin{ex}
Tính: $9 \times 5 + 45 = \dots$
\shortans{90}
\loigiai{
$45 + 45 = 90$.
}
\end{ex}

\begin{ex}
Tính: $54 : 6 + 41 = \dots$
\shortans{50}
\loigiai{
$9 + 41 = 50$.
}
\end{ex}

\begin{ex}
Tính: $56 : 7 + 62 = \dots$
\shortans{70}
\loigiai{
$8 + 62 = 70$.
}
\end{ex}

\begin{ex}
Tính: $64 : 8 + 32 = \dots$
\shortans{40}
\loigiai{
$8 + 32 = 40$.
}
\end{ex}

\begin{ex}
Tính: $81 : 9 + 51 = \dots$
\shortans{60}
\loigiai{
$9 + 51 = 60$.
}
\end{ex}

\begin{ex}
Tìm $x$, biết: $7 \times x = 63$.
\shortans{9}
\loigiai{
$x = 63 : 7 = 9$.
}
\end{ex}

\begin{ex}
Tìm $x$, biết: $x : 8 = 6$.
\shortans{48}
\loigiai{
$x = 6 \times 8 = 48$.
}
\end{ex}

\begin{ex}
Tìm $x$, biết: $72 : x = 8$.
\shortans{9}
\loigiai{
$x = 72 : 8 = 9$.
}
\end{ex}

\begin{ex}
Tìm $x$, biết: $6 \times x = 36$.
\shortans{6}
\loigiai{
$x = 36 : 6 = 6$.
}
\end{ex}

\begin{ex}
Tìm $\frac{1}{6}$ của 48 kg.
\shortans{8}
\loigiai{
$48 : 6 = 8$ kg.
}
\end{ex}

\begin{ex}
Tìm $\frac{1}{7}$ của 35 lít.
\shortans{5}
\loigiai{
$35 : 7 = 5$ lít.
}
\end{ex}

\begin{ex}
Tìm $\frac{1}{8}$ của 56 mét.
\shortans{7}
\loigiai{
$56 : 8 = 7$ m.
}
\end{ex}

\begin{ex}
Tìm $\frac{1}{9}$ của 63 con vịt.
\shortans{7}
\loigiai{
$63 : 9 = 7$ con vịt.
}
\end{ex}

\begin{ex}
Một cửa hàng có 72 gói kẹo, đã bán đi $\frac{1}{8}$ số gói kẹo. Hỏi cửa hàng đã bán bao nhiêu gói kẹo?
\shortans{9}
\loigiai{
$72 : 8 = 9$ gói kẹo.
}
\end{ex}

\begin{ex}
Mỗi xe ô tô chở được 8 hành khách. Hỏi 9 xe ô tô như vậy chở được tất cả bao nhiêu hành khách?
\shortans{72}
\loigiai{
$8 \times 9 = 72$ hành khách.
}
\end{ex}

\begin{ex}
Có 42 quyển vở chia đều cho 7 bạn. Hỏi mỗi bạn được bao nhiêu quyển vở?
\shortans{6}
\loigiai{
$42 : 7 = 6$ quyển vở.
}
\end{ex}

\begin{ex}
Một tuần lễ có 7 ngày. Hỏi 8 tuần lễ có bao nhiêu ngày?
\shortans{56}
\loigiai{
$7 \times 8 = 56$ ngày.
}
\end{ex}

\begin{ex}
Tính giá trị biểu thức: $6 \times 8 : 8$.
\shortans{6}
\loigiai{
$48 : 8 = 6$.
}
\end{ex}

\begin{ex}
Tính giá trị biểu thức: $7 \times 9 - 7 \times 4$.
\shortans{35}
\loigiai{
$63 - 28 = 35$ (hoặc $7 \times 5 = 35$).
}
\end{ex}

\begin{ex}
Mẹ mua 54 quả cam và chia đều vào 6 rổ. Hỏi mỗi rổ có bao nhiêu quả cam?
\shortans{9}
\loigiai{
$54 : 6 = 9$ quả cam.
}
\end{ex}

\begin{ex}
Một sợi dây dài 63 dm được cắt thành các đoạn bằng nhau dài 7 dm. Hỏi cắt được bao nhiêu đoạn?
\shortans{9}
\loigiai{
$63 : 7 = 9$ đoạn.
}
\end{ex}

\begin{ex}
Tính: $9 \times 4 : 6 = \dots$
\shortans{6}
\loigiai{
$36 : 6 = 6$.
}
\end{ex}

\newpage


% =========================================================================
% Làm quen với hình phẳng, hình khối
% =========================================================================
\section{Làm quen với hình phẳng, hình khối}

\TaoPhieuBaiTap{Chủ đề 3 -- Làm quen với hình phẳng, hình khối}{16}{Điểm ở giữa, trung điểm của đoạn thẳng}{https://loptoanthaylinh.vn/bai-tap/toan-3/bai-16}
\begin{ex}
Đoạn thẳng AB dài 12 cm. Điểm M là trung điểm của đoạn thẳng AB. Độ dài đoạn thẳng AM là bao nhiêu xăng-ti-mét?
\shortans{6}
\loigiai{
$AM = 12 : 2 = 6$ cm.
}
\end{ex}

\begin{ex}
Đoạn thẳng CD dài 20 cm. Điểm I là trung điểm của đoạn thẳng CD. Độ dài đoạn thẳng ID là bao nhiêu xăng-ti-mét?
\shortans{10}
\loigiai{
$ID = 20 : 2 = 10$ cm.
}
\end{ex}

\begin{ex}
Điểm O là trung điểm của đoạn thẳng MN. Biết MO = 8 cm. Độ dài đoạn thẳng MN là bao nhiêu xăng-ti-mét?
\shortans{16}
\loigiai{
$MN = 8 \times 2 = 16$ cm.
}
\end{ex}

\begin{ex}
Điểm E là trung điểm của đoạn thẳng HK. Biết EK = 15 cm. Độ dài đoạn thẳng HK là bao nhiêu xăng-ti-mét?
\shortans{30}
\loigiai{
$HK = 15 \times 2 = 30$ cm.
}
\end{ex}

\begin{ex}
Cho ba điểm thẳng hàng A, B, C theo thứ tự đó. Biết AB = 7 cm, BC = 7 cm. Điểm B cách đều A và C. Hỏi độ dài đoạn AC là bao nhiêu xăng-ti-mét?
\shortans{14}
\loigiai{
$AC = 7 + 7 = 14$ cm.
}
\end{ex}

\begin{ex}
Cho ba điểm thẳng hàng A, B, C theo thứ tự. B có phải là điểm ở giữa A và C không? (Nếu đúng ghi 1, sai ghi 0).
\shortans{1}
\loigiai{
Ba điểm thẳng hàng A, B, C thì B là điểm ở giữa A và C.
}
\end{ex}

\begin{ex}
Đoạn thẳng PQ dài 18 cm. Điểm K là trung điểm của PQ. Hỏi PK dài bao nhiêu xăng-ti-mét?
\shortans{9}
\loigiai{
$18 : 2 = 9$ cm.
}
\end{ex}

\begin{ex}
Một sợi dây dài 24 cm, gập đôi sợi dây lại thì mỗi nửa sợi dây dài bao nhiêu xăng-ti-mét?
\shortans{12}
\loigiai{
$24 : 2 = 12$ cm.
}
\end{ex}

\begin{ex}
Đoạn thẳng AB dài 10 cm, M là trung điểm AB. Đoạn thẳng BC dài 6 cm, N là trung điểm BC (biết A, B, C thẳng hàng theo thứ tự). Tính MN theo xăng-ti-mét.
\shortans{8}
\loigiai{
$MB = 5$ cm, $BN = 3$ cm, $MN = 5 + 3 = 8$ cm.
}
\end{ex}

\begin{ex}
Đoạn thẳng AC dài 16 cm. B là điểm ở giữa A và C sao cho AB = 9 cm. Tính độ dài đoạn thẳng BC theo xăng-ti-mét.
\shortans{7}
\loigiai{
$BC = 16 - 9 = 7$ cm.
}
\end{ex}

\begin{ex}
Nếu I là trung điểm của đoạn thẳng AB dài 28 cm thì IA bằng bao nhiêu xăng-ti-mét?
\shortans{14}
\loigiai{
$28 : 2 = 14$ cm.
}
\end{ex}

\begin{ex}
Biết M là trung điểm của AB, và AM = 25 cm. Độ dài AB bằng bao nhiêu xăng-ti-mét?
\shortans{50}
\loigiai{
$25 \times 2 = 50$ cm.
}
\end{ex}

\begin{ex}
Một thanh gỗ dài 40 dm, bác thợ mộc cưa đôi thanh gỗ tại đúng trung điểm. Hỏi mỗi đoạn gỗ dài bao nhiêu đề-xi-mét?
\shortans{20}
\loigiai{
$40 : 2 = 20$ dm.
}
\end{ex}

\begin{ex}
Đoạn thẳng AB dài 36 cm. Điểm C là trung điểm AB, điểm D là trung điểm AC. Hỏi AD dài bao nhiêu xăng-ti-mét?
\shortans{9}
\loigiai{
$AC = 36 : 2 = 18$ cm, $AD = 18 : 2 = 9$ cm.
}
\end{ex}

\begin{ex}
Cho đoạn thẳng MN dài 48 cm, I là trung điểm MN, K là trung điểm IN. Tính độ dài IK theo xăng-ti-mét.
\shortans{12}
\loigiai{
$IN = 48 : 2 = 24$ cm, $IK = 24 : 2 = 12$ cm.
}
\end{ex}

\begin{ex}
Một đoạn dây thép dài 50 cm. Người ta uốn cong tại chính giữa. Mỗi nửa đoạn dây dài bao nhiêu xăng-ti-mét?
\shortans{25}
\loigiai{
$50 : 2 = 25$ cm.
}
\end{ex}

\begin{ex}
Trên một đường thẳng có ba điểm A, B, C. Biết AB = 5 cm, AC = 10 cm, B nằm giữa A và C. B có phải là trung điểm của AC không? (Nếu đúng ghi 1, sai ghi 0).
\shortans{1}
\loigiai{
$BC = 10 - 5 = 5$ cm = $AB$, nên B là trung điểm của AC.
}
\end{ex}

\begin{ex}
Đoạn thẳng AB dài 14 cm. Điểm M nằm giữa A và B sao cho AM = 6 cm. Tính độ dài đoạn MB theo xăng-ti-mét.
\shortans{8}
\loigiai{
$MB = 14 - 6 = 8$ cm.
}
\end{ex}

\begin{ex}
Nếu O là trung điểm của đoạn PQ dài 60 cm thì OQ dài bao nhiêu xăng-ti-mét?
\shortans{30}
\loigiai{
$60 : 2 = 30$ cm.
}
\end{ex}

\begin{ex}
Cho đoạn thẳng CD dài 32 cm. Điểm M là trung điểm của CD. Tính tổng độ dài CM + MD theo xăng-ti-mét.
\shortans{32}
\loigiai{
$CM + MD = CD = 32$ cm.
}
\end{ex}

\begin{ex}
Một đoạn thẳng có bao nhiêu trung điểm?
\shortans{1}
\loigiai{
Mỗi đoạn thẳng chỉ có duy nhất 1 trung điểm.
}
\end{ex}

\begin{ex}
Một thước dây dài 100 cm. Vạch chỉ trung điểm của thước dây là vạch số bao nhiêu xăng-ti-mét?
\shortans{50}
\loigiai{
$100 : 2 = 50$ cm.
}
\end{ex}

\begin{ex}
Đoạn thẳng MN dài 1 m. Trung điểm của đoạn MN cách điểm M bao nhiêu đề-xi-mét?
\shortans{5}
\loigiai{
$1\text{ m} = 10\text{ dm}$, $10 : 2 = 5$ dm.
}
\end{ex}

\begin{ex}
Điểm K là trung điểm của đoạn thẳng AB dài 16 dm. Điểm K cách đầu mút B bao nhiêu đề-xi-mét?
\shortans{8}
\loigiai{
$16 : 2 = 8$ dm.
}
\end{ex}

\begin{ex}
Đoạn thẳng AB dài 22 cm, M là trung điểm AB. Hỏi độ dài AM là bao nhiêu xăng-ti-mét?
\shortans{11}
\loigiai{
$22 : 2 = 11$ cm.
}
\end{ex}

\newpage

\TaoPhieuBaiTap{Chủ đề 3 -- Làm quen với hình phẳng, hình khối}{17}{Hình tròn, tâm, bán kính, đường kính}{https://loptoanthaylinh.vn/bai-tap/toan-3/bai-17}
\begin{ex}
Một hình tròn có bán kính là 5 cm. Đường kính của hình tròn đó là bao nhiêu xăng-ti-mét?
\shortans{10}
\loigiai{
Đường kính gấp đôi bán kính: $5 \times 2 = 10$ cm.
}
\end{ex}

\begin{ex}
Một hình tròn có bán kính là 8 cm. Đường kính của hình tròn đó là bao nhiêu xăng-ti-mét?
\shortans{16}
\loigiai{
$8 \times 2 = 16$ cm.
}
\end{ex}

\begin{ex}
Một hình tròn có đường kính là 14 cm. Bán kính của hình tròn đó là bao nhiêu xăng-ti-mét?
\shortans{7}
\loigiai{
$14 : 2 = 7$ cm.
}
\end{ex}

\begin{ex}
Một hình tròn có đường kính là 18 cm. Bán kính của hình tròn đó là bao nhiêu xăng-ti-mét?
\shortans{9}
\loigiai{
$18 : 2 = 9$ cm.
}
\end{ex}

\begin{ex}
Một hình tròn có bán kính là 12 cm. Đường kính của hình tròn đó là bao nhiêu xăng-ti-mét?
\shortans{24}
\loigiai{
$12 \times 2 = 24$ cm.
}
\end{ex}

\begin{ex}
Một hình tròn có đường kính là 26 cm. Bán kính của hình tròn đó là bao nhiêu xăng-ti-mét?
\shortans{13}
\loigiai{
$26 : 2 = 13$ cm.
}
\end{ex}

\begin{ex}
Trong một hình tròn, đường kính gấp mấy lần bán kính?
\shortans{2}
\loigiai{
Đường kính dài gấp 2 lần bán kính.
}
\end{ex}

\begin{ex}
Trong một hình tròn, tâm có phải là trung điểm của mọi đường kính không? (Nếu đúng ghi 1, sai ghi 0).
\shortans{1}
\loigiai{
Tâm hình tròn là trung điểm của mọi đường kính.
}
\end{ex}

\begin{ex}
Một compa có độ mở hai đầu bằng 6 cm để vẽ một đường tròn. Đường kính của hình tròn vẽ được là bao nhiêu xăng-ti-mét?
\shortans{12}
\loigiai{
Độ mở compa chính là bán kính: $R = 6$ cm $\Rightarrow$ Đường kính = $6 \times 2 = 12$ cm.
}
\end{ex}

\begin{ex}
Một mặt đồng hồ hình tròn có bán kính 10 cm. Đường kính của mặt đồng hồ đó là bao nhiêu xăng-ti-mét?
\shortans{20}
\loigiai{
$10 \times 2 = 20$ cm.
}
\end{ex}

\begin{ex}
Bánh xe đạp có đường kính là 60 cm. Bán kính của bánh xe đạp đó là bao nhiêu xăng-ti-mét?
\shortans{30}
\loigiai{
$60 : 2 = 30$ cm.
}
\end{ex}

\begin{ex}
Một chiếc đĩa tròn có đường kính 22 cm. Bán kính của chiếc đĩa đó là bao nhiêu xăng-ti-mét?
\shortans{11}
\loigiai{
$22 : 2 = 11$ cm.
}
\end{ex}

\begin{ex}
Hình tròn tâm O có bán kính OA = 7 cm. Đoạn thẳng OB cũng là một bán kính. Đoạn OB dài bao nhiêu xăng-ti-mét?
\shortans{7}
\loigiai{
Tất cả các bán kính trong hình tròn đều bằng nhau: $OB = OA = 7$ cm.
}
\end{ex}

\begin{ex}
Hình tròn tâm O có đường kính AB = 16 cm. Bán kính OM dài bao nhiêu xăng-ti-mét?
\shortans{8}
\loigiai{
$16 : 2 = 8$ cm.
}
\end{ex}

\begin{ex}
Nếu bán kính hình tròn tăng thêm 3 cm thì đường kính của hình tròn đó tăng thêm bao nhiêu xăng-ti-mét?
\shortans{6}
\loigiai{
$3 \times 2 = 6$ cm.
}
\end{ex}

\begin{ex}
Một hình tròn có đường kính là 40 cm. Bán kính của nó là bao nhiêu đề-xi-mét?
\shortans{2}
\loigiai{
$40 : 2 = 20\text{ cm} = 2\text{ dm}$.
}
\end{ex}

\begin{ex}
Một hình tròn có bán kính là 15 dm. Đường kính của hình tròn đó là bao nhiêu đề-xi-mét?
\shortans{30}
\loigiai{
$15 \times 2 = 30$ dm.
}
\end{ex}

\begin{ex}
Đường kính của hình tròn bằng 50 cm. Bán kính của hình tròn đó là bao nhiêu xăng-ti-mét?
\shortans{25}
\loigiai{
$50 : 2 = 25$ cm.
}
\end{ex}

\begin{ex}
Hình tròn có bao nhiêu tâm?
\shortans{1}
\loigiai{
Hình tròn chỉ có duy nhất 1 tâm.
}
\end{ex}

\begin{ex}
Hình tròn tâm O có bán kính 4 cm. Độ dài đường kính MN của hình tròn đó là bao nhiêu xăng-ti-mét?
\shortans{8}
\loigiai{
$4 \times 2 = 8$ cm.
}
\end{ex}

\begin{ex}
Biết bán kính một hình tròn bằng chu vi một hình vuông có cạnh 2 cm. Đường kính hình tròn đó là bao nhiêu xăng-ti-mét?
\shortans{16}
\loigiai{
Chu vi hình vuông = $2 \times 4 = 8$ cm. Bán kính = 8 cm $\Rightarrow$ đường kính = $8 \times 2 = 16$ cm.
}
\end{ex}

\begin{ex}
Độ mở hai đầu compa là 9 cm. Khi quay compa vẽ được hình tròn có đường kính bao nhiêu xăng-ti-mét?
\shortans{18}
\loigiai{
$9 \times 2 = 18$ cm.
}
\end{ex}

\begin{ex}
Đoạn thẳng nối tâm với một điểm trên hình tròn gọi là gì? (Bán kính hay đường kính? Nếu là bán kính ghi 1, đường kính ghi 2).
\shortans{1}
\loigiai{
Đoạn thẳng nối tâm với một điểm trên hình tròn là bán kính.
}
\end{ex}

\begin{ex}
Đoạn thẳng đi qua tâm và cắt hai điểm trên đường tròn gọi là gì? (Bán kính ghi 1, đường kính ghi 2).
\shortans{2}
\loigiai{
Đó là đường kính.
}
\end{ex}

\begin{ex}
Bán kính hình tròn là 35 cm. Đường kính hình tròn là bao nhiêu xăng-ti-mét?
\shortans{70}
\loigiai{
$35 \times 2 = 70$ cm.
}
\end{ex}

\newpage

\TaoPhieuBaiTap{Chủ đề 3 -- Làm quen với hình phẳng, hình khối}{18}{Góc, góc vuông, góc không vuông}{https://loptoanthaylinh.vn/bai-tap/toan-3/bai-18}
\begin{ex}
Hình chữ nhật có bao nhiêu góc vuông?
\shortans{4}
\loigiai{
Hình chữ nhật có 4 góc vuông.
}
\end{ex}

\begin{ex}
Hình vuông có bao nhiêu góc vuông?
\shortans{4}
\loigiai{
Hình vuông có 4 góc vuông.
}
\end{ex}

\begin{ex}
Dụng cụ nào trong các dụng cụ học tập thường dùng để kiểm tra góc vuông? (Nếu là ê ke ghi 1, compa ghi 2).
\shortans{1}
\loigiai{
Ê-ke dùng để kiểm tra góc vuông.
}
\end{ex}

\begin{ex}
Kim giờ và kim phút của đồng hồ tạo thành một góc vuông lúc mấy giờ đúng vào buổi sáng?
\shortans{9}
\loigiai{
Lúc 3 giờ hoặc 9 giờ tạo thành góc vuông. Lúc 9 giờ hoặc 3 giờ.
}
\end{ex}

\begin{ex}
Lúc 3 giờ đúng, góc tạo bởi kim giờ và kim phút là góc gì? (Góc vuông ghi 1, góc không vuông ghi 0).
\shortans{1}
\loigiai{
Lúc 3 giờ đúng góc tạo bởi 2 kim là góc vuông.
}
\end{ex}

\begin{ex}
Lúc 6 giờ đúng, góc tạo bởi kim giờ và kim phút có phải là góc vuông không? (Phải ghi 1, không ghi 0).
\shortans{0}
\loigiai{
Lúc 6 giờ đúng góc bẹt (không vuông).
}
\end{ex}

\begin{ex}
Một hình tam giác vuông có bao nhiêu góc vuông?
\shortans{1}
\loigiai{
Hình tam giác vuông có đúng 1 góc vuông.
}
\end{ex}

\begin{ex}
Một hình tứ giác có nhiều nhất bao nhiêu góc vuông?
\shortans{4}
\loigiai{
Một hình tứ giác có thể có nhiều nhất 4 góc vuông (như hình chữ nhật).
}
\end{ex}

\begin{ex}
Một hình tam giác có nhiều nhất bao nhiêu góc vuông?
\shortans{1}
\loigiai{
Hình tam giác chỉ có nhiều nhất 1 góc vuông.
}
\end{ex}

\begin{ex}
Tờ giấy hình chữ nhật có 4 góc. Cắt đi 1 góc thì tờ giấy mới có bao nhiêu góc?
\shortans{5}
\loigiai{
Khi cắt 1 góc của tờ giấy hình chữ nhật ta được hình ngũ giác có 5 góc.
}
\end{ex}

\begin{ex}
Một góc gồm có một đỉnh và mấy cạnh?
\shortans{2}
\loigiai{
Một góc gồm 1 đỉnh và 2 cạnh xuất phát từ đỉnh đó.
}
\end{ex}

\begin{ex}
Mỗi mặt bàn học hình chữ nhật có bao nhiêu góc vuông?
\shortans{4}
\loigiai{
Có 4 góc vuông.
}
\end{ex}

\begin{ex}
Một hình chữ nhật có 4 góc vuông, 2 hình chữ nhật có tất cả bao nhiêu góc vuông?
\shortans{8}
\loigiai{
$4 \times 2 = 8$ góc vuông.
}
\end{ex}

\begin{ex}
Trong hình vuông ABCD, có bao nhiêu góc vuông đỉnh A?
\shortans{1}
\loigiai{
Tại đỉnh A có đúng 1 góc vuông của hình vuông.
}
\end{ex}

\begin{ex}
Lúc 9 giờ đúng, góc tạo bởi kim giờ và kim phút là góc gì? (Góc vuông ghi 1, không vuông ghi 0).
\shortans{1}
\loigiai{
Lúc 9 giờ đúng là góc vuông.
}
\end{ex}

\begin{ex}
Hình tam giác có 3 góc, trong đó có 1 góc vuông thì có mấy góc không vuông?
\shortans{2}
\loigiai{
$3 - 1 = 2$ góc không vuông.
}
\end{ex}

\begin{ex}
Một hình chữ nhật bị cắt mất một góc vuông ở một đỉnh thì còn lại bao nhiêu góc vuông ban đầu?
\shortans{3}
\loigiai{
$4 - 1 = 3$ góc vuông ban đầu còn lại.
}
\end{ex}

\begin{ex}
Khung cửa sổ hình chữ nhật có mấy góc vuông?
\shortans{4}
\loigiai{
Khung cửa sổ hình chữ nhật có 4 góc vuông.
}
\end{ex}

\begin{ex}
Mỗi chiếc ê-ke có bao nhiêu góc vuông?
\shortans{1}
\loigiai{
Một chiếc ê-ke có đúng 1 góc vuông.
}
\end{ex}

\begin{ex}
3 chiếc ê-ke có tất cả bao nhiêu góc vuông?
\shortans{3}
\loigiai{
$1 \times 3 = 3$ góc vuông.
}
\end{ex}

\begin{ex}
Góc tạo bởi hai cạnh liên tiếp của quyển vở ô li là góc gì? (Góc vuông ghi 1, không vuông ghi 0).
\shortans{1}
\loigiai{
Góc của quyển vở là góc vuông.
}
\end{ex}

\begin{ex}
Một hình chữ nhật có 4 đỉnh. Mỗi đỉnh tạo thành bao nhiêu góc vuông?
\shortans{1}
\loigiai{
Mỗi đỉnh tạo thành 1 góc vuông.
}
\end{ex}

\begin{ex}
Đồng hồ chỉ 12 giờ đúng. Góc tạo bởi hai kim có phải là góc vuông không? (Phải ghi 1, không ghi 0).
\shortans{0}
\loigiai{
Lúc 12 giờ hai kim trùng nhau, góc không vuông (góc 0 độ).
}
\end{ex}

\begin{ex}
Hình ngũ giác có 5 góc. Nếu có 2 góc vuông thì có bao nhiêu góc không vuông?
\shortans{3}
\loigiai{
$5 - 2 = 3$ góc không vuông.
}
\end{ex}

\begin{ex}
Một chiếc khăn quàng đỏ hình tam giác có 1 góc tù và 2 góc nhọn. Có góc vuông nào không? (Có ghi số lượng góc vuông).
\shortans{0}
\loigiai{
Chiếc khăn quàng đỏ không có góc vuông nào.
}
\end{ex}

\newpage

\TaoPhieuBaiTap{Chủ đề 3 -- Làm quen với hình phẳng, hình khối}{19}{Hình tam giác, hình tứ giác, hình chữ nhật, hình vuông}{https://loptoanthaylinh.vn/bai-tap/toan-3/bai-19}
\begin{ex}
Hình tam giác có bao nhiêu cạnh?
\shortans{3}
\loigiai{
Hình tam giác có 3 cạnh.
}
\end{ex}

\begin{ex}
Hình tứ giác có bao nhiêu đỉnh?
\shortans{4}
\loigiai{
Hình tứ giác có 4 đỉnh.
}
\end{ex}

\begin{ex}
Hình chữ nhật có mấy cặp cạnh đối diện dài bằng nhau?
\shortans{2}
\loigiai{
Hình chữ nhật có 2 cặp cạnh đối diện bằng nhau.
}
\end{ex}

\begin{ex}
Hình vuông có bao nhiêu cạnh bằng nhau?
\shortans{4}
\loigiai{
Hình vuông có 4 cạnh dài bằng nhau.
}
\end{ex}

\begin{ex}
Một hình vuông có cạnh dài 7 cm. Tính chu vi hình vuông đó theo đơn vị xăng-ti-mét.
\shortans{28}
\loigiai{
$7 \times 4 = 28$ cm.
}
\end{ex}

\begin{ex}
Một hình chữ nhật có chiều dài 9 cm và chiều rộng 5 cm. Chu vi hình chữ nhật đó là bao nhiêu xăng-ti-mét?
\shortans{28}
\loigiai{
$(9 + 5) \times 2 = 28$ cm.
}
\end{ex}

\begin{ex}
Một hình chữ nhật có chu vi 24 cm. Nửa chu vi hình chữ nhật đó là bao nhiêu xăng-ti-mét?
\shortans{12}
\loigiai{
$24 : 2 = 12$ cm.
}
\end{ex}

\begin{ex}
Một hình chữ nhật có nửa chu vi là 15 cm, chiều rộng là 6 cm. Chiều dài hình chữ nhật đó là bao nhiêu xăng-ti-mét?
\shortans{9}
\loigiai{
$15 - 6 = 9$ cm.
}
\end{ex}

\begin{ex}
Một hình tam giác có độ dài 3 cạnh là 6 cm, 8 cm, 10 cm. Chu vi hình tam giác đó là bao nhiêu xăng-ti-mét?
\shortans{24}
\loigiai{
$6 + 8 + 10 = 24$ cm.
}
\end{ex}

\begin{ex}
Một hình tứ giác có độ dài các cạnh là 5 cm, 7 cm, 8 cm, 10 cm. Chu vi hình tứ giác đó là bao nhiêu xăng-ti-mét?
\shortans{30}
\loigiai{
$5 + 7 + 8 + 10 = 30$ cm.
}
\end{ex}

\begin{ex}
Một hình vuông có chu vi 36 cm. Độ dài một cạnh của hình vuông đó là bao nhiêu xăng-ti-mét?
\shortans{9}
\loigiai{
$36 : 4 = 9$ cm.
}
\end{ex}

\begin{ex}
Một hình vuông có chu vi 20 cm. Độ dài cạnh hình vuông đó là bao nhiêu xăng-ti-mét?
\shortans{5}
\loigiai{
$20 : 4 = 5$ cm.
}
\end{ex}

\begin{ex}
Ghép hai hình vuông cạnh 3 cm lại với nhau tạo thành một hình chữ nhật. Chiều dài hình chữ nhật đó là bao nhiêu xăng-ti-mét?
\shortans{6}
\loigiai{
$3 + 3 = 6$ cm.
}
\end{ex}

\begin{ex}
Ghép hai hình vuông cạnh 4 cm lại thành hình chữ nhật. Chu vi hình chữ nhật tạo thành là bao nhiêu xăng-ti-mét?
\shortans{24}
\loigiai{
Chiều dài = 8 cm, chiều rộng = 4 cm. Chu vi = $(8 + 4) \times 2 = 24$ cm.
}
\end{ex}

\begin{ex}
Một hình chữ nhật có chiều dài 12 cm, chiều rộng bằng $\frac{1}{3}$ chiều dài. Chiều rộng hình chữ nhật đó là bao nhiêu xăng-ti-mét?
\shortans{4}
\loigiai{
$12 : 3 = 4$ cm.
}
\end{ex}

\begin{ex}
Một hình tam giác đều có 3 cạnh bằng nhau, mỗi cạnh dài 9 cm. Chu vi hình tam giác đó là bao nhiêu xăng-ti-mét?
\shortans{27}
\loigiai{
$9 \times 3 = 27$ cm.
}
\end{ex}

\begin{ex}
Một hình tam giác có chu vi 30 cm, biết 3 cạnh bằng nhau. Độ dài mỗi cạnh là bao nhiêu xăng-ti-mét?
\shortans{10}
\loigiai{
$30 : 3 = 10$ cm.
}
\end{ex}

\begin{ex}
Một hình chữ nhật có chiều rộng 4 cm, chiều dài gấp 3 lần chiều rộng. Chiều dài hình chữ nhật đó là bao nhiêu xăng-ti-mét?
\shortans{12}
\loigiai{
$4 \times 3 = 12$ cm.
}
\end{ex}

\begin{ex}
Một hình chữ nhật có chiều rộng 4 cm, chiều dài 12 cm. Chu vi hình chữ nhật đó là bao nhiêu xăng-ti-mét?
\shortans{32}
\loigiai{
$(12 + 4) \times 2 = 32$ cm.
}
\end{ex}

\begin{ex}
Hình nào có 4 cạnh bằng nhau và 4 góc vuông? (Hình vuông ghi 1, hình chữ nhật ghi 2).
\shortans{1}
\loigiai{
Đó là hình vuông.
}
\end{ex}

\begin{ex}
Tổng số cạnh của 3 hình tam giác là bao nhiêu?
\shortans{9}
\loigiai{
$3 \times 3 = 9$ cạnh.
}
\end{ex}

\begin{ex}
Tổng số đỉnh của 4 hình tứ giác là bao nhiêu?
\shortans{16}
\loigiai{
$4 \times 4 = 16$ đỉnh.
}
\end{ex}

\begin{ex}
Một hình chữ nhật có chiều dài 15 cm, chiều rộng kém chiều dài 5 cm. Chiều rộng hình chữ nhật đó là bao nhiêu xăng-ti-mét?
\shortans{10}
\loigiai{
$15 - 5 = 10$ cm.
}
\end{ex}

\begin{ex}
Một miếng bìa hình vuông có chu vi 32 cm. Cạnh miếng bìa dài bao nhiêu xăng-ti-mét?
\shortans{8}
\loigiai{
$32 : 4 = 8$ cm.
}
\end{ex}

\begin{ex}
Một khu vườn hình vuông có chu vi là 40 m. Cạnh khu vườn đó dài bao nhiêu mét?
\shortans{10}
\loigiai{
$40 : 4 = 10$ m.
}
\end{ex}

\newpage

\TaoPhieuBaiTap{Chủ đề 3 -- Làm quen với hình phẳng, hình khối}{20}{Thực hành vẽ hình và trang trí}{https://loptoanthaylinh.vn/bai-tap/toan-3/bai-20}
\begin{ex}
Để vẽ một hình tròn bán kính 4 cm bằng compa, ta mở hai đầu compa cách nhau bao nhiêu xăng-ti-mét?
\shortans{4}
\loigiai{
Mở compa đúng bằng độ dài bán kính: 4 cm.
}
\end{ex}

\begin{ex}
Để vẽ hình tròn có đường kính 10 cm, ta mở hai đầu compa cách nhau bao nhiêu xăng-ti-mét?
\shortans{5}
\loigiai{
Bán kính = $10 : 2 = 5$ cm.
}
\end{ex}

\begin{ex}
Bạn Mai vẽ 4 hình tròn bằng nhau, mỗi hình có bán kính 3 cm. Tổng đường kính của 4 hình tròn đó là bao nhiêu xăng-ti-mét?
\shortans{24}
\loigiai{
Mỗi hình tròn có đường kính $3 \times 2 = 6$ cm. Tổng đường kính = $6 \times 4 = 24$ cm.
}
\end{ex}

\begin{ex}
Vẽ một hình vuông có cạnh 5 cm trên giấy kẻ ô vuông. Chu vi hình vuông đó là bao nhiêu xăng-ti-mét?
\shortans{20}
\loigiai{
$5 \times 4 = 20$ cm.
}
\end{ex}

\begin{ex}
Vẽ một hình chữ nhật có chiều dài 6 cm, chiều rộng 4 cm. Chu vi hình chữ nhật vẽ được là bao nhiêu xăng-ti-mét?
\shortans{20}
\loigiai{
$(6 + 4) \times 2 = 20$ cm.
}
\end{ex}

\begin{ex}
Dùng 3 que tính bằng nhau có thể xếp thành hình gì? (Hình tam giác ghi 1, tứ giác ghi 2).
\shortans{1}
\loigiai{
3 que tính xếp thành hình tam giác.
}
\end{ex}

\begin{ex}
Dùng 4 que tính bằng nhau có thể xếp thành hình gì có 4 góc vuông? (Hình vuông ghi 1, hình tròn ghi 2).
\shortans{1}
\loigiai{
Xếp thành hình vuông.
}
\end{ex}

\begin{ex}
Cần ít nhất bao nhiêu que tính bằng nhau để xếp thành một hình chữ nhật (không phải hình vuông)?
\shortans{6}
\loigiai{
Mỗi chiều dài 2 que, mỗi chiều rộng 1 que: $2 \times 2 + 1 \times 2 = 6$ que tính.
}
\end{ex}

\begin{ex}
Cần bao nhiêu que tính bằng nhau để ghép thành 2 hình vuông riêng biệt?
\shortans{8}
\loigiai{
$4 \times 2 = 8$ que.
}
\end{ex}

\begin{ex}
Nếu ghép chung 1 cạnh, dùng bao nhiêu que tính bằng nhau để tạo thành 2 hình vuông liền nhau?
\shortans{7}
\loigiai{
$4 + 3 = 7$ que.
}
\end{ex}

\begin{ex}
Bạn Lan dùng dây kim tuyến viền quanh một khung ảnh hình vuông cạnh 15 cm. Bạn Lan cần bao nhiêu xăng-ti-mét dây?
\shortans{60}
\loigiai{
$15 \times 4 = 60$ cm.
}
\end{ex}

\begin{ex}
Một dải ruy băng trang trí hình chữ nhật dài 20 cm, rộng 10 cm. Chu vi dải ruy băng đó là bao nhiêu xăng-ti-mét?
\shortans{60}
\loigiai{
$(20 + 10) \times 2 = 60$ cm.
}
\end{ex}

\begin{ex}
Bạn Nam dùng compa vẽ 3 hình tròn đồng tâm có bán kính lần lượt là 2 cm, 4 cm, 6 cm. Đường kính của hình tròn lớn nhất là bao nhiêu xăng-ti-mét?
\shortans{12}
\loigiai{
$6 \times 2 = 12$ cm.
}
\end{ex}

\begin{ex}
Một hình vuông có cạnh 6 cm được chia thành các ô vuông nhỏ cạnh 1 cm. Có tất cả bao nhiêu ô vuông nhỏ?
\shortans{36}
\loigiai{
$6 \times 6 = 36$ ô vuông.
}
\end{ex}

\begin{ex}
Một hình chữ nhật có chiều dài 5 cm, chiều rộng 3 cm được chia thành các ô vuông cạnh 1 cm. Có tất cả bao nhiêu ô vuông nhỏ?
\shortans{15}
\loigiai{
$5 \times 3 = 15$ ô vuông.
}
\end{ex}

\begin{ex}
Để gấp một chiếc thuyền giấy, bạn An dùng một tờ giấy hình vuông cạnh 20 cm. Chu vi tờ giấy đó là bao nhiêu xăng-ti-mét?
\shortans{80}
\loigiai{
$20 \times 4 = 80$ cm.
}
\end{ex}

\begin{ex}
Một họa tiết hình tam giác đều có cạnh dài 8 cm. Chu vi họa tiết đó là bao nhiêu xăng-ti-mét?
\shortans{24}
\loigiai{
$8 \times 3 = 24$ cm.
}
\end{ex}

\begin{ex}
Dùng 12 que tính bằng nhau xếp thành một hình vuông. Mỗi cạnh của hình vuông có bao nhiêu que tính?
\shortans{3}
\loigiai{
$12 : 4 = 3$ que tính.
}
\end{ex}

\begin{ex}
Dùng 9 que tính bằng nhau xếp thành 3 hình tam giác rời nhau. Mỗi tam giác có mấy que tính?
\shortans{3}
\loigiai{
$9 : 3 = 3$ que tính.
}
\end{ex}

\begin{ex}
Bạn Hoa cắt một hình chữ nhật có chu vi 30 cm thành hai hình vuông nhỏ bằng nhau cạnh 5 cm. Đúng hay sai? (Đúng ghi 1, sai ghi 0).
\shortans{1}
\loigiai{
Chu vi hình chữ nhật ghép từ 2 hình vuông cạnh 5 cm: chiều dài 10 cm, rộng 5 cm $\Rightarrow (10+5) \times 2 = 30$ cm.
}
\end{ex}

\begin{ex}
Bác thợ mộc đóng khung ảnh gỗ hình vuông cạnh 25 cm. Tổng chiều dài 4 thanh gỗ là bao nhiêu xăng-ti-mét?
\shortans{100}
\loigiai{
$25 \times 4 = 100$ cm.
}
\end{ex}

\begin{ex}
Để vẽ đường tròn có bán kính 8 cm, đường kính của nó là bao nhiêu xăng-ti-mét?
\shortans{16}
\loigiai{
$8 \times 2 = 16$ cm.
}
\end{ex}

\begin{ex}
Một bức tranh hình chữ nhật có chiều dài 40 cm, chiều rộng 30 cm. Nửa chu vi khung tranh là bao nhiêu xăng-ti-mét?
\shortans{70}
\loigiai{
$40 + 30 = 70$ cm.
}
\end{ex}

\begin{ex}
Một góc vuông được vẽ bằng ê-ke có số đo bao nhiêu độ? (Điền số).
\shortans{90}
\loigiai{
Góc vuông có số đo $90^\circ$.
}
\end{ex}

\begin{ex}
Có bao nhiêu góc vuông trong một hình vuông?
\shortans{4}
\loigiai{
Hình vuông có 4 góc vuông.
}
\end{ex}

\newpage

\TaoPhieuBaiTap{Chủ đề 3 -- Làm quen với hình phẳng, hình khối}{21}{Khối lập phương, khối hộp chữ nhật}{https://loptoanthaylinh.vn/bai-tap/toan-3/bai-21}
\begin{ex}
Khối lập phương có bao nhiêu mặt?
\shortans{6}
\loigiai{
Khối lập phương có 6 mặt.
}
\end{ex}

\begin{ex}
Khối lập phương có bao nhiêu đỉnh?
\shortans{8}
\loigiai{
Khối lập phương có 8 đỉnh.
}
\end{ex}

\begin{ex}
Khối lập phương có bao nhiêu cạnh?
\shortans{12}
\loigiai{
Khối lập phương có 12 cạnh.
}
\end{ex}

\begin{ex}
Các mặt của khối lập phương là hình gì? (Hình vuông ghi 1, hình chữ nhật ghi 2).
\shortans{1}
\loigiai{
Các mặt của khối lập phương đều là hình vuông.
}
\end{ex}

\begin{ex}
Khối hộp chữ nhật có bao nhiêu mặt?
\shortans{6}
\loigiai{
Khối hộp chữ nhật có 6 mặt.
}
\end{ex}

\begin{ex}
Khối hộp chữ nhật có bao nhiêu đỉnh?
\shortans{8}
\loigiai{
Khối hộp chữ nhật có 8 đỉnh.
}
\end{ex}

\begin{ex}
Khối hộp chữ nhật có bao nhiêu cạnh?
\shortans{12}
\loigiai{
Khối hộp chữ nhật có 12 cạnh.
}
\end{ex}

\begin{ex}
Tổng số mặt của 2 khối lập phương là bao nhiêu?
\shortans{12}
\loigiai{
$6 \times 2 = 12$ mặt.
}
\end{ex}

\begin{ex}
Tổng số cạnh của 2 khối lập phương là bao nhiêu?
\shortans{24}
\loigiai{
$12 \times 2 = 24$ cạnh.
}
\end{ex}

\begin{ex}
Tổng số đỉnh của 3 khối hộp chữ nhật là bao nhiêu?
\shortans{24}
\loigiai{
$8 \times 3 = 24$ đỉnh.
}
\end{ex}

\begin{ex}
Một con xúc xắc hình khối lập phương có các mặt được chấm từ 1 đến 6. Tổng số chấm trên 6 mặt là bao nhiêu?
\shortans{21}
\loigiai{
$1 + 2 + 3 + 4 + 5 + 6 = 21$.
}
\end{ex}

\begin{ex}
Một khối lập phương có cạnh dài 5 cm. Tổng độ dài tất cả 12 cạnh của khối lập phương đó là bao nhiêu xăng-ti-mét?
\shortans{60}
\loigiai{
$5 \times 12 = 60$ cm.
}
\end{ex}

\begin{ex}
Một khối lập phương có cạnh dài 3 cm. Diện tích một mặt của khối lập phương đó là bao nhiêu xăng-ti-mét vuông?
\shortans{9}
\loigiai{
$3 \times 3 = 9\text{ cm}^2$.
}
\end{ex}

\begin{ex}
Dùng các khối lập phương nhỏ cạnh 1 cm để xếp thành một khối lập phương lớn có cạnh dài 2 cm. Cần bao nhiêu khối lập phương nhỏ?
\shortans{8}
\loigiai{
$2 \times 2 \times 2 = 8$ khối.
}
\end{ex}

\begin{ex}
Dùng các khối lập phương nhỏ cạnh 1 cm xếp thành khối hộp chữ nhật dài 3 cm, rộng 2 cm, cao 1 cm. Cần bao nhiêu khối lập phương nhỏ?
\shortans{6}
\loigiai{
$3 \times 2 \times 1 = 6$ khối.
}
\end{ex}

\begin{ex}
Một hộp phấn hình khối lập phương có bao nhiêu cặp mặt đối diện song song và bằng nhau?
\shortans{3}
\loigiai{
Có 3 cặp mặt đối diện.
}
\end{ex}

\begin{ex}
Một thùng các-tông hình khối hộp chữ nhật có bao nhiêu đỉnh ở mặt đáy trên?
\shortans{4}
\loigiai{
Mỗi mặt đáy có 4 đỉnh.
}
\end{ex}

\begin{ex}
Một khối lập phương có tổng độ dài 12 cạnh là 48 cm. Độ dài mỗi cạnh của khối lập phương đó là bao nhiêu xăng-ti-mét?
\shortans{4}
\loigiai{
$48 : 12 = 4$ cm.
}
\end{ex}

\begin{ex}
Một khối lập phương có tổng độ dài 12 cạnh là 84 cm. Độ dài mỗi cạnh của khối lập phương đó là bao nhiêu xăng-ti-mét?
\shortans{7}
\loigiai{
$84 : 12 = 7$ cm.
}
\end{ex}

\begin{ex}
Xếp 4 khối lập phương nhỏ bằng nhau thành một hàng ngang. Khối hình tạo thành là khối gì? (Khối hộp chữ nhật ghi 1, khối lập phương ghi 2).
\shortans{1}
\loigiai{
Đó là khối hộp chữ nhật.
}
\end{ex}

\begin{ex}
Một bể cá hình hộp chữ nhật có 4 chân chống tại 4 đỉnh đáy. Bể cá đó có tất cả bao nhiêu đỉnh?
\shortans{8}
\loigiai{
Khối hộp chữ nhật có 8 đỉnh.
}
\end{ex}

\begin{ex}
Có bao nhiêu cạnh ở mặt đáy của khối hộp chữ nhật?
\shortans{4}
\loigiai{
Mặt đáy có 4 cạnh.
}
\end{ex}

\begin{ex}
Có bao nhiêu cạnh đứng nối hai mặt đáy của khối hộp chữ nhật?
\shortans{4}
\loigiai{
Có 4 cạnh bên đứng.
}
\end{ex}

\begin{ex}
Xếp các khối lập phương nhỏ cạnh 1 cm thành một khối lập phương có cạnh 3 cm. Cần bao nhiêu khối lập phương nhỏ?
\shortans{27}
\loigiai{
$3 \times 3 \times 3 = 27$ khối.
}
\end{ex}

\begin{ex}
Khối Rubik $3 \times 3 \times 3$ gồm bao nhiêu khối lập phương nhỏ?
\shortans{27}
\loigiai{
$3 \times 3 \times 3 = 27$ khối nhỏ.
}
\end{ex}

\newpage

\TaoPhieuBaiTap{Chủ đề 3 -- Làm quen với hình phẳng, hình khối}{22}{Luyện tập chung hình học}{https://loptoanthaylinh.vn/bai-tap/toan-3/bai-22}
\begin{ex}
Đoạn thẳng MN dài 18 cm, I là trung điểm của MN. Hỏi đoạn thẳng MI dài bao nhiêu xăng-ti-mét?
\shortans{9}
\loigiai{
$18 : 2 = 9$ cm.
}
\end{ex}

\begin{ex}
Hình tròn có bán kính 6 cm thì đường kính là bao nhiêu xăng-ti-mét?
\shortans{12}
\loigiai{
$6 \times 2 = 12$ cm.
}
\end{ex}

\begin{ex}
Hình tròn có đường kính 24 cm thì bán kính là bao nhiêu xăng-ti-mét?
\shortans{12}
\loigiai{
$24 : 2 = 12$ cm.
}
\end{ex}

\begin{ex}
Một hình vuông có cạnh dài 8 cm. Chu vi của hình vuông đó là bao nhiêu xăng-ti-mét?
\shortans{32}
\loigiai{
$8 \times 4 = 32$ cm.
}
\end{ex}

\begin{ex}
Một hình chữ nhật có chiều dài 14 cm, chiều rộng 6 cm. Chu vi hình chữ nhật đó là bao nhiêu xăng-ti-mét?
\shortans{40}
\loigiai{
$(14 + 6) \times 2 = 40$ cm.
}
\end{ex}

\begin{ex}
Một hình tam giác có độ dài 3 cạnh lần lượt là 7 cm, 9 cm, 11 cm. Chu vi hình tam giác đó là bao nhiêu xăng-ti-mét?
\shortans{27}
\loigiai{
$7 + 9 + 11 = 27$ cm.
}
\end{ex}

\begin{ex}
Khối lập phương có bao nhiêu đỉnh?
\shortans{8}
\loigiai{
Có 8 đỉnh.
}
\end{ex}

\begin{ex}
Khối hộp chữ nhật có bao nhiêu cạnh?
\shortans{12}
\loigiai{
Có 12 cạnh.
}
\end{ex}

\begin{ex}
Một hình vuông có chu vi 28 cm. Cạnh hình vuông đó dài bao nhiêu xăng-ti-mét?
\shortans{7}
\loigiai{
$28 : 4 = 7$ cm.
}
\end{ex}

\begin{ex}
Một hình chữ nhật có chu vi 36 cm, chiều rộng là 8 cm. Chiều dài hình chữ nhật đó là bao nhiêu xăng-ti-mét?
\shortans{10}
\loigiai{
Nửa chu vi = $36 : 2 = 18$ cm. Chiều dài = $18 - 8 = 10$ cm.
}
\end{ex}

\begin{ex}
Một hình chữ nhật có nửa chu vi là 25 cm, chiều dài là 15 cm. Chiều rộng hình chữ nhật đó là bao nhiêu xăng-ti-mét?
\shortans{10}
\loigiai{
$25 - 15 = 10$ cm.
}
\end{ex}

\begin{ex}
Một khối lập phương có cạnh 6 cm. Tổng độ dài các cạnh của khối lập phương đó là bao nhiêu xăng-ti-mét?
\shortans{72}
\loigiai{
$6 \times 12 = 72$ cm.
}
\end{ex}

\begin{ex}
Hình chữ nhật có mấy góc vuông?
\shortans{4}
\loigiai{
Hình chữ nhật có 4 góc vuông.
}
\end{ex}

\begin{ex}
Có bao nhiêu góc vuông trong một hình vuông?
\shortans{4}
\loigiai{
Hình vuông có 4 góc vuông.
}
\end{ex}

\begin{ex}
Đoạn thẳng AB dài 30 cm, điểm M là trung điểm của AB. Đoạn MB dài bao nhiêu xăng-ti-mét?
\shortans{15}
\loigiai{
$30 : 2 = 15$ cm.
}
\end{ex}

\begin{ex}
Compa mở hai đầu 7 cm vẽ một đường tròn. Đường kính đường tròn đó bằng bao nhiêu xăng-ti-mét?
\shortans{14}
\loigiai{
$7 \times 2 = 14$ cm.
}
\end{ex}

\begin{ex}
Một hình tam giác đều có chu vi 24 cm. Mỗi cạnh của hình tam giác đó dài bao nhiêu xăng-ti-mét?
\shortans{8}
\loigiai{
$24 : 3 = 8$ cm.
}
\end{ex}

\begin{ex}
Một hình tứ giác có chu vi 35 cm, độ dài 3 cạnh là 8 cm, 9 cm, 10 cm. Cạnh còn lại dài bao nhiêu xăng-ti-mét?
\shortans{8}
\loigiai{
$35 - (8 + 9 + 10) = 35 - 27 = 8$ cm.
}
\end{ex}

\begin{ex}
Đồng hồ chỉ mấy giờ thì hai kim vuông góc với nhau vào buổi sáng trước 6 giờ?
\shortans{3}
\loigiai{
Lúc 3 giờ đúng.
}
\end{ex}

\begin{ex}
Một khối hộp chữ nhật có bao nhiêu mặt?
\shortans{6}
\loigiai{
Có 6 mặt.
}
\end{ex}

\begin{ex}
Tính: $6 \times 8 + 15 = \dots$
\shortans{63}
\loigiai{
$48 + 15 = 63$.
}
\end{ex}

\begin{ex}
Tính: $7 \times 9 - 25 = \dots$
\shortans{38}
\loigiai{
$63 - 25 = 38$.
}
\end{ex}

\begin{ex}
Tính: $64 : 8 + 42 = \dots$
\shortans{50}
\loigiai{
$8 + 42 = 50$.
}
\end{ex}

\begin{ex}
Tìm $x$: $x \times 7 = 49$.
\shortans{7}
\loigiai{
$x = 49 : 7 = 7$.
}
\end{ex}

\begin{ex}
Tìm $x$: $x : 6 = 9$.
\shortans{54}
\loigiai{
$x = 9 \times 6 = 54$.
}
\end{ex}

\newpage


% =========================================================================
% Phép nhân, phép chia trong phạm vi 100
% =========================================================================
\section{Phép nhân, phép chia trong phạm vi 100}

\TaoPhieuBaiTap{Chủ đề 4 -- Phép nhân, phép chia trong phạm vi 100}{23}{Nhân số có hai chữ số với số có một chữ số}{https://loptoanthaylinh.vn/bai-tap/toan-3/bai-23}
\begin{ex}
Tính: $12 \times 3 = \dots$
\shortans{36}
\loigiai{
$12 \times 3 = 36$.
}
\end{ex}

\begin{ex}
Tính: $14 \times 2 = \dots$
\shortans{28}
\loigiai{
$14 \times 2 = 28$.
}
\end{ex}

\begin{ex}
Tính: $23 \times 3 = \dots$
\shortans{69}
\loigiai{
$23 \times 3 = 69$.
}
\end{ex}

\begin{ex}
Tính: $31 \times 3 = \dots$
\shortans{93}
\loigiai{
$31 \times 3 = 93$.
}
\end{ex}

\begin{ex}
Tính: $13 \times 4 = \dots$
\shortans{52}
\loigiai{
$13 \times 4 = 52$.
}
\end{ex}

\begin{ex}
Tính: $16 \times 3 = \dots$
\shortans{48}
\loigiai{
$16 \times 3 = 48$.
}
\end{ex}

\begin{ex}
Tính: $24 \times 3 = \dots$
\shortans{72}
\loigiai{
$24 \times 3 = 72$.
}
\end{ex}

\begin{ex}
Tính: $15 \times 5 = \dots$
\shortans{75}
\loigiai{
$15 \times 5 = 75$.
}
\end{ex}

\begin{ex}
Tính: $18 \times 4 = \dots$
\shortans{72}
\loigiai{
$18 \times 4 = 72$.
}
\end{ex}

\begin{ex}
Tính: $25 \times 3 = \dots$
\shortans{75}
\loigiai{
$25 \times 3 = 75$.
}
\end{ex}

\begin{ex}
Tính: $17 \times 5 = \dots$
\shortans{85}
\loigiai{
$17 \times 5 = 85$.
}
\end{ex}

\begin{ex}
Tính: $19 \times 4 = \dots$
\shortans{76}
\loigiai{
$19 \times 4 = 76$.
}
\end{ex}

\begin{ex}
Tính: $26 \times 3 = \dots$
\shortans{78}
\loigiai{
$26 \times 3 = 78$.
}
\end{ex}

\begin{ex}
Tính: $48 \times 2 = \dots$
\shortans{96}
\loigiai{
$48 \times 2 = 96$.
}
\end{ex}

\begin{ex}
Mỗi hộp bút có 12 chiếc bút chì. Hỏi 4 hộp bút như thế có tất cả bao nhiêu chiếc bút chì?
\shortans{48}
\loigiai{
$12 \times 4 = 48$ chiếc.
}
\end{ex}

\begin{ex}
Mỗi cuộn vải dài 15 m. Hỏi 5 cuộn vải như thế dài tất cả bao nhiêu mét?
\shortans{75}
\loigiai{
$15 \times 5 = 75$ m.
}
\end{ex}

\begin{ex}
Một lớp học có 3 dãy bàn, mỗi dãy bàn có 11 học sinh. Hỏi lớp học có tất cả bao nhiêu học sinh?
\shortans{33}
\loigiai{
$11 \times 3 = 33$ học sinh.
}
\end{ex}

\begin{ex}
Mỗi gói kẹo có 24 cái kẹo. Hỏi 3 gói kẹo như vậy có tất cả bao nhiêu cái kẹo?
\shortans{72}
\loigiai{
$24 \times 3 = 72$ cái kẹo.
}
\end{ex}

\begin{ex}
Một ngày mẹ đi làm 8 giờ. Hỏi trong 12 ngày mẹ đi làm tất cả bao nhiêu giờ?
\shortans{96}
\loigiai{
$12 \times 8 = 96$ giờ.
}
\end{ex}

\begin{ex}
Tính giá trị biểu thức: $14 \times 3 + 28$.
\shortans{70}
\loigiai{
$42 + 28 = 70$.
}
\end{ex}

\begin{ex}
Tính giá trị biểu thức: $16 \times 4 - 34$.
\shortans{30}
\loigiai{
$64 - 34 = 30$.
}
\end{ex}

\begin{ex}
Tìm $x$, biết: $x : 3 = 24$.
\shortans{72}
\loigiai{
$x = 24 \times 3 = 72$.
}
\end{ex}

\begin{ex}
Tìm $x$, biết: $x : 4 = 18$.
\shortans{72}
\loigiai{
$x = 18 \times 4 = 72$.
}
\end{ex}

\begin{ex}
Một trang sách có khoảng 32 dòng, mỗi dòng có 3 từ tiếng Anh cần chú ý. Trang sách đó có bao nhiêu từ chú ý?
\shortans{96}
\loigiai{
$32 \times 3 = 96$ từ.
}
\end{ex}

\begin{ex}
Tính: $12 \times 5 + 15 = \dots$
\shortans{75}
\loigiai{
$60 + 15 = 75$.
}
\end{ex}

\newpage

\TaoPhieuBaiTap{Chủ đề 4 -- Phép nhân, phép chia trong phạm vi 100}{24}{Gấp một số lên nhiều lần}{https://loptoanthaylinh.vn/bai-tap/toan-3/bai-24}
\begin{ex}
Gấp 6 lên 5 lần ta được số nào?
\shortans{30}
\loigiai{
$6 \times 5 = 30$.
}
\end{ex}

\begin{ex}
Gấp 8 lên 4 lần ta được số nào?
\shortans{32}
\loigiai{
$8 \times 4 = 32$.
}
\end{ex}

\begin{ex}
Gấp 9 lên 7 lần ta được số nào?
\shortans{63}
\loigiai{
$9 \times 7 = 63$.
}
\end{ex}

\begin{ex}
Gấp 12 lên 3 lần ta được số nào?
\shortans{36}
\loigiai{
$12 \times 3 = 36$.
}
\end{ex}

\begin{ex}
Gấp 15 lên 4 lần ta được số nào?
\shortans{60}
\loigiai{
$15 \times 4 = 60$.
}
\end{ex}

\begin{ex}
Gấp 16 lên 5 lần ta được số nào?
\shortans{80}
\loigiai{
$16 \times 5 = 80$.
}
\end{ex}

\begin{ex}
Gấp 24 lên 3 lần ta được số nào?
\shortans{72}
\loigiai{
$24 \times 3 = 72$.
}
\end{ex}

\begin{ex}
Gấp 18 lên 4 lần ta được số nào?
\shortans{72}
\loigiai{
$18 \times 4 = 72$.
}
\end{ex}

\begin{ex}
Đoạn thẳng AB dài 7 cm. Đoạn thẳng CD dài gấp 4 lần đoạn thẳng AB. Đoạn thẳng CD dài bao nhiêu xăng-ti-mét?
\shortans{28}
\loigiai{
$7 \times 4 = 28$ cm.
}
\end{ex}

\begin{ex}
Bạn Lan có 6 cái nhãn vở. Bạn Mai có số nhãn vở gấp 3 lần số nhãn vở của bạn Lan. Hỏi bạn Mai có bao nhiêu cái nhãn vở?
\shortans{18}
\loigiai{
$6 \times 3 = 18$ cái.
}
\end{ex}

\begin{ex}
Năm nay em 6 tuổi. Tuổi mẹ gấp 5 lần tuổi em. Hỏi năm nay mẹ bao nhiêu tuổi?
\shortans{30}
\loigiai{
$6 \times 5 = 30$ tuổi.
}
\end{ex}

\begin{ex}
Năm nay con 7 tuổi, tuổi bố gấp 5 lần tuổi con. Hỏi bố bao nhiêu tuổi?
\shortans{35}
\loigiai{
$7 \times 5 = 35$ tuổi.
}
\end{ex}

\begin{ex}
Bao gạo thứ nhất nặng 12 kg. Bao gạo thứ hai nặng gấp 3 lần bao thứ nhất. Hỏi bao thứ hai nặng bao nhiêu ki-lô-gam?
\shortans{36}
\loigiai{
$12 \times 3 = 36$ kg.
}
\end{ex}

\begin{ex}
Một con vịt nặng 2 kg, một con lợn nặng gấp 25 lần con vịt. Hỏi con lợn nặng bao nhiêu ki-lô-gam?
\shortans{50}
\loigiai{
$2 \times 25 = 50$ kg.
}
\end{ex}

\begin{ex}
Hùng có 8 viên bi, số bi của Dũng gấp 6 lần số bi của Hùng. Hỏi Dũng có bao nhiêu viên bi?
\shortans{48}
\loigiai{
$8 \times 6 = 48$ viên bi.
}
\end{ex}

\begin{ex}
Một số gấp lên 6 lần thì được 42. Số đó là số nào?
\shortans{7}
\loigiai{
$42 : 6 = 7$.
}
\end{ex}

\begin{ex}
Một số gấp lên 8 lần thì được 72. Số đó là số nào?
\shortans{9}
\loigiai{
$72 : 8 = 9$.
}
\end{ex}

\begin{ex}
Một số gấp lên 4 lần rồi cộng với 12 thì được 48. Số đó là số nào?
\shortans{9}
\loigiai{
Số đó gấp 4 lần: $48 - 12 = 36$. Số đó là $36 : 4 = 9$.
}
\end{ex}

\begin{ex}
Thùng thứ nhất có 15 lít dầu. Thùng thứ hai có số lít dầu gấp 4 lần thùng thứ nhất. Hỏi thùng thứ hai có bao nhiêu lít dầu?
\shortans{60}
\loigiai{
$15 \times 4 = 60$ lít.
}
\end{ex}

\begin{ex}
Mẹ hái được 8 quả bưởi, bà hái được gấp 3 lần số bưởi của mẹ. Hỏi cả hai người hái được bao nhiêu quả bưởi?
\shortans{32}
\loigiai{
Bà hái: $8 \times 3 = 24$ quả. Cả hai người: $8 + 24 = 32$ quả.
}
\end{ex}

\begin{ex}
Gấp 14 lên 5 lần rồi bớt đi 20 thì được bao nhiêu?
\shortans{50}
\loigiai{
$14 \times 5 - 20 = 70 - 20 = 50$.
}
\end{ex}

\begin{ex}
Gấp đôi của 35 là bao nhiêu?
\shortans{70}
\loigiai{
$35 \times 2 = 70$.
}
\end{ex}

\begin{ex}
Gấp ba của 24 là bao nhiêu?
\shortans{72}
\loigiai{
$24 \times 3 = 72$.
}
\end{ex}

\begin{ex}
Gấp bốn của 16 là bao nhiêu?
\shortans{64}
\loigiai{
$16 \times 4 = 64$.
}
\end{ex}

\begin{ex}
Một hình vuông có cạnh 9 cm. Nếu gấp cạnh của nó lên 2 lần thì cạnh mới dài bao nhiêu xăng-ti-mét?
\shortans{18}
\loigiai{
$9 \times 2 = 18$ cm.
}
\end{ex}

\newpage

\TaoPhieuBaiTap{Chủ đề 4 -- Phép nhân, phép chia trong phạm vi 100}{25}{Phép chia hết và phép chia có dư}{https://loptoanthaylinh.vn/bai-tap/toan-3/bai-25}
\begin{ex}
Trong phép chia $17 : 3$, số dư là bao nhiêu?
\shortans{2}
\loigiai{
$17 = 3 \times 5 + 2$, dư 2.
}
\end{ex}

\begin{ex}
Trong phép chia $23 : 4$, số dư là bao nhiêu?
\shortans{3}
\loigiai{
$23 = 4 \times 5 + 3$, dư 3.
}
\end{ex}

\begin{ex}
Trong phép chia $29 : 5$, thương là bao nhiêu?
\shortans{5}
\loigiai{
$29 = 5 \times 5 + 4$, thương là 5.
}
\end{ex}

\begin{ex}
Trong phép chia $38 : 6$, thương là bao nhiêu?
\shortans{6}
\loigiai{
$38 = 6 \times 6 + 2$, thương là 6.
}
\end{ex}

\begin{ex}
Trong phép chia $45 : 7$, số dư là bao nhiêu?
\shortans{3}
\loigiai{
$45 = 7 \times 6 + 3$, dư 3.
}
\end{ex}

\begin{ex}
Trong phép chia $53 : 8$, số dư là bao nhiêu?
\shortans{5}
\loigiai{
$53 = 8 \times 6 + 5$, dư 5.
}
\end{ex}

\begin{ex}
Trong phép chia $67 : 9$, số dư là bao nhiêu?
\shortans{4}
\loigiai{
$67 = 9 \times 7 + 4$, dư 4.
}
\end{ex}

\begin{ex}
Trong một phép chia có số chia là 6, số dư lớn nhất có thể có là bao nhiêu?
\shortans{5}
\loigiai{
Số dư luôn nhỏ hơn số chia, lớn nhất là $6 - 1 = 5$.
}
\end{ex}

\begin{ex}
Trong một phép chia có số chia là 8, số dư lớn nhất có thể có là bao nhiêu?
\shortans{7}
\loigiai{
$8 - 1 = 7$.
}
\end{ex}

\begin{ex}
Trong một phép chia cho 5, số dư có thể nhận bao nhiêu giá trị khác nhau (kể cả số dư bằng 0)?
\shortans{5}
\loigiai{
Các số dư có thể là 0, 1, 2, 3, 4 (5 giá trị).
}
\end{ex}

\begin{ex}
Có 19 bông hoa cắm đều vào các lọ, mỗi lọ 4 bông hoa. Hỏi cắm được nhiều nhất bao nhiêu lọ hoa?
\shortans{4}
\loigiai{
$19 : 4 = 4$ dư 3 lọ.
}
\end{ex}

\begin{ex}
Có 25 cái kẹo chia đều cho 4 bạn. Hỏi mỗi bạn được nhiều nhất bao nhiêu cái kẹo?
\shortans{6}
\loigiai{
$25 : 4 = 6$ dư 1.
}
\end{ex}

\begin{ex}
Có 31 học sinh đi đò, mỗi chuyến chở được 6 học sinh. Hỏi cần ít nhất bao nhiêu chuyến đò để chở hết số học sinh?
\shortans{6}
\loigiai{
$31 : 6 = 5$ dư 1, cần ít nhất $5 + 1 = 6$ chuyến.
}
\end{ex}

\begin{ex}
Một lớp có 34 học sinh. Mỗi bàn ngồi 2 học sinh. Hỏi cần ít nhất bao nhiêu bàn cho cả lớp ngồi?
\shortans{17}
\loigiai{
$34 : 2 = 17$ bàn (chia hết).
}
\end{ex}

\begin{ex}
Một phép chia có số chia là 7, thương là 8 và số dư là 5. Số bị chia là bao nhiêu?
\shortans{61}
\loigiai{
$8 \times 7 + 5 = 56 + 5 = 61$.
}
\end{ex}

\begin{ex}
Một phép chia có số chia là 9, thương là 6 và số dư lớn nhất có thể. Số bị chia là bao nhiêu?
\shortans{62}
\loigiai{
Số dư lớn nhất là 8: Số bị chia = $6 \times 9 + 8 = 54 + 8 = 62$.
}
\end{ex}

\begin{ex}
Một phép chia có số chia là 5, thương là 12 và số dư là 3. Số bị chia là bao nhiêu?
\shortans{63}
\loigiai{
$12 \times 5 + 3 = 60 + 3 = 63$.
}
\end{ex}

\begin{ex}
Một cuộn dây dài 43 m được cắt thành các đoạn dài 5 m. Hỏi cắt được bao nhiêu đoạn và thừa ra mấy mét? (Điền số mét thừa).
\shortans{3}
\loigiai{
$43 : 5 = 8$ dư 3 m.
}
\end{ex}

\begin{ex}
Có 50 chiếc bút xếp vào các hộp, mỗi hộp đựng 8 chiếc. Hỏi còn thừa mấy chiếc bút?
\shortans{2}
\loigiai{
$50 : 8 = 6$ dư 2 chiếc.
}
\end{ex}

\begin{ex}
Tính thương trong phép chia $49 : 6$.
\shortans{8}
\loigiai{
$49 = 6 \times 8 + 1$, thương là 8.
}
\end{ex}

\begin{ex}
Số nào chia cho 4 được thương là 9 và dư 2?
\shortans{38}
\loigiai{
$9 \times 4 + 2 = 38$.
}
\end{ex}

\begin{ex}
Số nào chia cho 7 được thương là 5 và dư 6?
\shortans{41}
\loigiai{
$5 \times 7 + 6 = 41$.
}
\end{ex}

\begin{ex}
Trong phép chia cho 9, có bao nhiêu số dư khác 0 có thể có?
\shortans{8}
\loigiai{
Các số dư từ 1 đến 8 (8 giá trị).
}
\end{ex}

\begin{ex}
May mỗi bộ quần áo cần 3 m vải. Có 26 m vải thì may được nhiều nhất bao nhiêu bộ quần áo?
\shortans{8}
\loigiai{
$26 : 3 = 8$ dư 2 m, may được 8 bộ.
}
\end{ex}

\begin{ex}
Số dư của phép chia $75 : 8$ là bao nhiêu?
\shortans{3}
\loigiai{
$75 = 8 \times 9 + 3$, số dư là 3.
}
\end{ex}

\newpage

\TaoPhieuBaiTap{Chủ đề 4 -- Phép nhân, phép chia trong phạm vi 100}{26}{Chia số có hai chữ số cho số có một chữ số}{https://loptoanthaylinh.vn/bai-tap/toan-3/bai-26}
\begin{ex}
Tính: $48 : 2 = \dots$
\shortans{24}
\loigiai{
$48 : 2 = 24$.
}
\end{ex}

\begin{ex}
Tính: $69 : 3 = \dots$
\shortans{23}
\loigiai{
$69 : 3 = 23$.
}
\end{ex}

\begin{ex}
Tính: $84 : 4 = \dots$
\shortans{21}
\loigiai{
$84 : 4 = 21$.
}
\end{ex}

\begin{ex}
Tính: $55 : 5 = \dots$
\shortans{11}
\loigiai{
$55 : 5 = 11$.
}
\end{ex}

\begin{ex}
Tính: $72 : 3 = \dots$
\shortans{24}
\loigiai{
$72 : 3 = 24$.
}
\end{ex}

\begin{ex}
Tính: $56 : 4 = \dots$
\shortans{14}
\loigiai{
$56 : 4 = 14$.
}
\end{ex}

\begin{ex}
Tính: $65 : 5 = \dots$
\shortans{13}
\loigiai{
$65 : 5 = 13$.
}
\end{ex}

\begin{ex}
Tính: $78 : 6 = \dots$
\shortans{13}
\loigiai{
$78 : 6 = 13$.
}
\end{ex}

\begin{ex}
Tính: $84 : 6 = \dots$
\shortans{14}
\loigiai{
$84 : 6 = 14$.
}
\end{ex}

\begin{ex}
Tính: $91 : 7 = \dots$
\shortans{13}
\loigiai{
$91 : 7 = 13$.
}
\end{ex}

\begin{ex}
Tính: $96 : 8 = \dots$
\shortans{12}
\loigiai{
$96 : 8 = 12$.
}
\end{ex}

\begin{ex}
Tính: $98 : 7 = \dots$
\shortans{14}
\loigiai{
$98 : 7 = 14$.
}
\end{ex}

\begin{ex}
Tính: $96 : 6 = \dots$
\shortans{16}
\loigiai{
$96 : 6 = 16$.
}
\end{ex}

\begin{ex}
Tính: $92 : 4 = \dots$
\shortans{23}
\loigiai{
$92 : 4 = 23$.
}
\end{ex}

\begin{ex}
Có 72 kg gạo chia đều vào 6 bao. Hỏi mỗi bao có bao nhiêu ki-lô-gam gạo?
\shortans{12}
\loigiai{
$72 : 6 = 12$ kg.
}
\end{ex}

\begin{ex}
Một cuộn dây dài 84 m được cắt thành 4 đoạn bằng nhau. Hỏi mỗi đoạn dài bao nhiêu mét?
\shortans{21}
\loigiai{
$84 : 4 = 21$ m.
}
\end{ex}

\begin{ex}
Mẹ mua 68 quả táo chia đều cho 4 người. Hỏi mỗi người nhận được bao nhiêu quả táo?
\shortans{17}
\loigiai{
$68 : 4 = 17$ quả táo.
}
\end{ex}

\begin{ex}
Lớp 3B có 36 học sinh chia đều thành 3 tổ. Hỏi mỗi tổ có bao nhiêu học sinh?
\shortans{12}
\loigiai{
$36 : 3 = 12$ học sinh.
}
\end{ex}

\begin{ex}
Tìm $x$, biết: $x \times 3 = 87$.
\shortans{29}
\loigiai{
$x = 87 : 3 = 29$.
}
\end{ex}

\begin{ex}
Tìm $x$, biết: $5 \times x = 95$.
\shortans{19}
\loigiai{
$x = 95 : 5 = 19$.
}
\end{ex}

\begin{ex}
Tính giá trị biểu thức: $96 : 3 - 15$.
\shortans{17}
\loigiai{
$32 - 15 = 17$.
}
\end{ex}

\begin{ex}
Tính giá trị biểu thức: $85 : 5 + 27$.
\shortans{44}
\loigiai{
$17 + 27 = 44$.
}
\end{ex}

\begin{ex}
Một người nuôi 96 con thỏ nhốt đều vào 6 chuồng. Hỏi mỗi chuồng có bao nhiêu con thỏ?
\shortans{16}
\loigiai{
$96 : 6 = 16$ con thỏ.
}
\end{ex}

\begin{ex}
Một thanh sắt dài 90 dm cắt thành các đoạn dài 5 dm. Hỏi được bao nhiêu đoạn?
\shortans{18}
\loigiai{
$90 : 5 = 18$ đoạn.
}
\end{ex}

\begin{ex}
Tính: $75 : 5 + 4 = \dots$
\shortans{19}
\loigiai{
$15 + 4 = 19$.
}
\end{ex}

\newpage

\TaoPhieuBaiTap{Chủ đề 4 -- Phép nhân, phép chia trong phạm vi 100}{27}{Giảm một số đi nhiều lần}{https://loptoanthaylinh.vn/bai-tap/toan-3/bai-27}
\begin{ex}
Giảm 24 đi 3 lần ta được số nào?
\shortans{8}
\loigiai{
$24 : 3 = 8$.
}
\end{ex}

\begin{ex}
Giảm 36 đi 4 lần ta được số nào?
\shortans{9}
\loigiai{
$36 : 4 = 9$.
}
\end{ex}

\begin{ex}
Giảm 45 đi 5 lần ta được số nào?
\shortans{9}
\loigiai{
$45 : 5 = 9$.
}
\end{ex}

\begin{ex}
Giảm 48 đi 6 lần ta được số nào?
\shortans{8}
\loigiai{
$48 : 6 = 8$.
}
\end{ex}

\begin{ex}
Giảm 56 đi 7 lần ta được số nào?
\shortans{8}
\loigiai{
$56 : 7 = 8$.
}
\end{ex}

\begin{ex}
Giảm 64 đi 8 lần ta được số nào?
\shortans{8}
\loigiai{
$64 : 8 = 8$.
}
\end{ex}

\begin{ex}
Giảm 72 đi 9 lần ta được số nào?
\shortans{8}
\loigiai{
$72 : 9 = 8$.
}
\end{ex}

\begin{ex}
Giảm 80 đi 4 lần ta được số nào?
\shortans{20}
\loigiai{
$80 : 4 = 20$.
}
\end{ex}

\begin{ex}
Giảm 90 đi 3 lần ta được số nào?
\shortans{30}
\loigiai{
$90 : 3 = 30$.
}
\end{ex}

\begin{ex}
Một sợi dây dài 28 cm. Sau khi cắt đi thì độ dài sợi dây giảm đi 4 lần. Hỏi sợi dây còn lại dài bao nhiêu xăng-ti-mét?
\shortans{7}
\loigiai{
$28 : 4 = 7$ cm.
}
\end{ex}

\begin{ex}
Bác An thu hoạch được 60 kg cam. Bác Ba thu hoạch được số cam giảm đi 3 lần so với bác An. Hỏi bác Ba thu hoạch được bao nhiêu ki-lô-gam cam?
\shortans{20}
\loigiai{
$60 : 3 = 20$ kg.
}
\end{ex}

\begin{ex}
Năm nay mẹ 36 tuổi, tuổi con bằng tuổi mẹ giảm đi 4 lần. Hỏi năm nay con bao nhiêu tuổi?
\shortans{9}
\loigiai{
$36 : 4 = 9$ tuổi.
}
\end{ex}

\begin{ex}
Bố 42 tuổi, con có tuổi bằng tuổi bố giảm đi 6 lần. Hỏi con bao nhiêu tuổi?
\shortans{7}
\loigiai{
$42 : 6 = 7$ tuổi.
}
\end{ex}

\begin{ex}
Một cuộn vải dài 70 m. Sau khi bán, số vải giảm đi 5 lần. Hỏi cuộn vải còn lại bao nhiêu mét?
\shortans{14}
\loigiai{
$70 : 5 = 14$ m.
}
\end{ex}

\begin{ex}
Giảm 60 đi 5 lần rồi cộng thêm 18 ta được số nào?
\shortans{30}
\loigiai{
$60 : 5 + 18 = 12 + 18 = 30$.
}
\end{ex}

\begin{ex}
Giảm 84 đi 4 lần rồi bớt đi 11 ta được số nào?
\shortans{10}
\loigiai{
$84 : 4 - 11 = 21 - 11 = 10$.
}
\end{ex}

\begin{ex}
Một số giảm đi 6 lần thì được 8. Số đó là số nào?
\shortans{48}
\loigiai{
$8 \times 6 = 48$.
}
\end{ex}

\begin{ex}
Một số giảm đi 7 lần thì được 9. Số đó là số nào?
\shortans{63}
\loigiai{
$9 \times 7 = 63$.
}
\end{ex}

\begin{ex}
Thùng thứ nhất có 48 lít dầu. Thùng thứ hai có số dầu giảm đi 4 lần so với thùng thứ nhất. Hỏi thùng thứ hai có bao nhiêu lít dầu?
\shortans{12}
\loigiai{
$48 : 4 = 12$ lít.
}
\end{ex}

\begin{ex}
Một rổ có 54 quả lê. Số quả táo ít hơn số quả lê giảm đi 6 lần. Hỏi có bao nhiêu quả táo?
\shortans{9}
\loigiai{
$54 : 6 = 9$ quả táo.
}
\end{ex}

\begin{ex}
Một số giảm đi 3 lần thì bằng 15. Tìm số đó.
\shortans{45}
\loigiai{
$15 \times 3 = 45$.
}
\end{ex}

\begin{ex}
Một người có 96 con gà. Sau khi bán thì số gà giảm đi 3 lần. Hỏi người đó còn lại bao nhiêu con gà?
\shortans{32}
\loigiai{
$96 : 3 = 32$ con gà.
}
\end{ex}

\begin{ex}
Giảm 50 đi 5 lần được bao nhiêu?
\shortans{10}
\loigiai{
$50 : 5 = 10$.
}
\end{ex}

\begin{ex}
Giảm 100 đi 4 lần được bao nhiêu?
\shortans{25}
\loigiai{
$100 : 4 = 25$.
}
\end{ex}

\begin{ex}
Đoạn thẳng AB dài 32 cm. Đoạn CD có độ dài bằng độ dài AB giảm đi 4 lần. CD dài bao nhiêu xăng-ti-mét?
\shortans{8}
\loigiai{
$32 : 4 = 8$ cm.
}
\end{ex}

\newpage

\TaoPhieuBaiTap{Chủ đề 4 -- Phép nhân, phép chia trong phạm vi 100}{28}{Bài toán giải bằng hai bước tính}{https://loptoanthaylinh.vn/bai-tap/toan-3/bai-28}
\begin{ex}
Hàng trên có 8 bông hoa, hàng dưới có nhiều hơn hàng trên 4 bông hoa. Hỏi cả hai hàng có bao nhiêu bông hoa?
\shortans{20}
\loigiai{
Hàng dưới: $8 + 4 = 12$ bông. Cả hai hàng: $8 + 12 = 20$ bông.
}
\end{ex}

\begin{ex}
Ngăn trên có 15 quyển sách, ngăn dưới có nhiều hơn ngăn trên 7 quyển sách. Hỏi cả hai ngăn có bao nhiêu quyển sách?
\shortans{37}
\loigiai{
Ngăn dưới: $15 + 7 = 22$ quyển. Cả hai ngăn: $15 + 22 = 37$ quyển.
}
\end{ex}

\begin{ex}
Thùng thứ nhất có 30 lít dầu, thùng thứ hai có ít hơn thùng thứ nhất 8 lít dầu. Hỏi cả hai thùng có bao nhiêu lít dầu?
\shortans{52}
\loigiai{
Thùng thứ hai: $30 - 8 = 22$ lít. Cả hai thùng: $30 + 22 = 52$ lít.
}
\end{ex}

\begin{ex}
Bạn Mai gấp được 14 con hạc, bạn Lan gấp được ít hơn bạn Mai 5 con hạc. Hỏi cả hai bạn gấp được bao nhiêu con hạc?
\shortans{23}
\loigiai{
Lan gấp: $14 - 5 = 9$ con. Cả hai bạn: $14 + 9 = 23$ con.
}
\end{ex}

\begin{ex}
Bao gạo to nặng 45 kg, bao gạo bé nhẹ hơn bao gạo to 15 kg. Hỏi cả hai bao nặng bao nhiêu ki-lô-gam?
\shortans{75}
\loigiai{
Bao bé: $45 - 15 = 30$ kg. Cả hai bao: $45 + 30 = 75$ kg.
}
\end{ex}

\begin{ex}
Đội Một trồng được 35 cây, đội Hai trồng được gấp 2 lần đội Một. Hỏi cả hai đội trồng được bao nhiêu cây?
\shortans{105}
\loigiai{
Đội Hai: $35 \times 2 = 70$ cây. Cả hai đội: $35 + 70 = 105$ cây.
}
\end{ex}

\begin{ex}
Buổi sáng cửa hàng bán được 12 chiếc xe đạp, buổi chiều bán được số xe gấp 3 lần buổi sáng. Hỏi cả ngày cửa hàng bán được bao nhiêu chiếc xe đạp?
\shortans{48}
\loigiai{
Buổi chiều: $12 \times 3 = 36$ chiếc. Cả ngày: $12 + 36 = 48$ chiếc.
}
\end{ex}

\begin{ex}
Một con vịt nặng 2 kg, một con ngỗng nặng gấp 3 lần con vịt. Hỏi cả con vịt và con ngỗng nặng bao nhiêu ki-lô-gam?
\shortans{8}
\loigiai{
Con ngỗng: $2 \times 3 = 6$ kg. Cả hai con: $2 + 6 = 8$ kg.
}
\end{ex}

\begin{ex}
Lớp 3A có 18 học sinh nữ, số học sinh nam bằng số học sinh nữ giảm đi 2 lần. Hỏi lớp 3A có tất cả bao nhiêu học sinh?
\shortans{27}
\loigiai{
Học sinh nam: $18 : 2 = 9$ bạn. Tổng số học sinh: $18 + 9 = 27$ học sinh.
}
\end{ex}

\begin{ex}
Thùng to đựng 40 lít nước ngọt, thùng nhỏ đựng số lít nước ngọt bằng thùng to giảm đi 4 lần. Hỏi cả hai thùng đựng bao nhiêu lít nước ngọt?
\shortans{50}
\loigiai{
Thùng nhỏ: $40 : 4 = 10$ lít. Cả hai thùng: $40 + 10 = 50$ lít.
}
\end{ex}

\begin{ex}
Mẹ mua 5 hộp bánh, mỗi hộp có 4 chiếc bánh. Mẹ mua thêm 6 chiếc bánh rời. Hỏi mẹ đã mua tất cả bao nhiêu chiếc bánh?
\shortans{26}
\loigiai{
Số bánh trong hộp: $5 \times 4 = 20$ chiếc. Tổng số bánh: $20 + 6 = 26$ chiếc.
}
\end{ex}

\begin{ex}
Có 6 can dầu, mỗi can chứa 5 lít dầu. Người ta đã dùng hết 12 lít dầu. Hỏi còn lại bao nhiêu lít dầu?
\shortans{18}
\loigiai{
Tổng số dầu ban đầu: $6 \times 5 = 30$ lít. Còn lại: $30 - 12 = 18$ lít.
}
\end{ex}

\begin{ex}
Một cuộn vải dài 50 m. Người ta may 6 bộ quần áo, mỗi bộ hết 3 m vải. Hỏi cuộn vải còn lại bao nhiêu mét?
\shortans{32}
\loigiai{
Đã dùng: $6 \times 3 = 18$ m. Còn lại: $50 - 18 = 32$ m.
}
\end{ex}

\begin{ex}
Cô giáo có 35 quyển vở, cô thưởng cho 6 bạn, mỗi bạn 4 quyển. Hỏi cô giáo còn lại bao nhiêu quyển vở?
\shortans{11}
\loigiai{
Đã thưởng: $6 \times 4 = 24$ quyển. Còn lại: $35 - 24 = 11$ quyển.
}
\end{ex}

\begin{ex}
Bạn An có 4 phong kẹo, mỗi phong có 8 chiếc kẹo. An cho em 10 chiếc kẹo. Hỏi An còn lại bao nhiêu chiếc kẹo?
\shortans{22}
\loigiai{
Tổng số kẹo: $4 \times 8 = 32$ chiếc. Còn lại: $32 - 10 = 22$ chiếc.
}
\end{ex}

\begin{ex}
Có 3 tổ, mỗi tổ có 8 bạn học sinh. Sau đó có thêm 4 bạn chuyển vào. Hỏi có tất cả bao nhiêu bạn học sinh?
\shortans{28}
\loigiai{
Ban đầu: $3 \times 8 = 24$ bạn. Sau đó: $24 + 4 = 28$ bạn.
}
\end{ex}

\begin{ex}
Mẹ mua 2 túi táo, mỗi túi có 9 quả táo. Mẹ biếu bà 8 quả. Hỏi mẹ còn lại bao nhiêu quả táo?
\shortans{10}
\loigiai{
Tổng số táo: $2 \times 9 = 18$ quả. Còn lại: $18 - 8 = 10$ quả.
}
\end{ex}

\begin{ex}
Một cửa hàng nhập về 4 bao gạo, mỗi bao nặng 20 kg. Cửa hàng đã bán được 35 kg. Hỏi cửa hàng còn lại bao nhiêu ki-lô-gam gạo?
\shortans{45}
\loigiai{
Tổng số gạo: $4 \times 20 = 80$ kg. Còn lại: $80 - 35 = 45$ kg.
}
\end{ex}

\begin{ex}
Nam có 15 viên bi xanh và 25 viên bi đỏ. Nam chia đều tất cả số bi đó vào 5 túi. Hỏi mỗi túi có bao nhiêu viên bi?
\shortans{8}
\loigiai{
Tổng số bi: $15 + 25 = 40$ viên. Mỗi túi: $40 : 5 = 8$ viên bi.
}
\end{ex}

\begin{ex}
Có 18 bạn nam và 14 bạn nữ cùng xếp đều thành 4 hàng. Hỏi mỗi hàng có bao nhiêu bạn?
\shortans{8}
\loigiai{
Tổng số bạn: $18 + 14 = 32$ bạn. Mỗi hàng: $32 : 4 = 8$ bạn.
}
\end{ex}

\begin{ex}
Một người thu được 42 quả dừa ở cây thứ nhất và 30 quả ở cây thứ hai. Người đó đem tất cả xếp đều vào 8 sọt. Hỏi mỗi sọt có bao nhiêu quả dừa?
\shortans{9}
\loigiai{
Tổng số dừa: $42 + 30 = 72$ quả. Mỗi sọt: $72 : 8 = 9$ quả dừa.
}
\end{ex}

\begin{ex}
Mẹ có 45 kg thóc nếp và 27 kg thóc tẻ. Mẹ chia đều vào 8 bao. Hỏi mỗi bao có bao nhiêu ki-lô-gam thóc?
\shortans{9}
\loigiai{
Tổng số thóc: $45 + 27 = 72$ kg. Mỗi bao: $72 : 8 = 9$ kg.
}
\end{ex}

\begin{ex}
Khối 3 có 2 lớp, lớp 3A có 32 học sinh, lớp 3B có 28 học sinh. Tất cả các học sinh xếp đều thành 6 hàng. Hỏi mỗi hàng có bao nhiêu học sinh?
\shortans{10}
\loigiai{
Tổng số học sinh: $32 + 28 = 60$ học sinh. Mỗi hàng: $60 : 6 = 10$ bạn.
}
\end{ex}

\begin{ex}
Một mảnh vườn hình chữ nhật có chiều rộng 6 m, chiều dài gấp 3 lần chiều rộng. Chu vi mảnh vườn đó là bao nhiêu mét?
\shortans{48}
\loigiai{
Chiều dài: $6 \times 3 = 18$ m. Chu vi: $(18 + 6) \times 2 = 48$ m.
}
\end{ex}

\begin{ex}
Một hình vuông có cạnh dài 6 cm. Người ta tăng mỗi cạnh thêm 3 cm. Chu vi hình vuông mới là bao nhiêu xăng-ti-mét?
\shortans{36}
\loigiai{
Cạnh mới: $6 + 3 = 9$ cm. Chu vi mới: $9 \times 4 = 36$ cm.
}
\end{ex}

\newpage

\TaoPhieuBaiTap{Chủ đề 4 -- Phép nhân, phép chia trong phạm vi 100}{29}{Luyện tập chung phép tính}{https://loptoanthaylinh.vn/bai-tap/toan-3/bai-29}
\begin{ex}
Tính: $14 \times 4 = \dots$
\shortans{56}
\loigiai{
$14 \times 4 = 56$.
}
\end{ex}

\begin{ex}
Tính: $23 \times 4 = \dots$
\shortans{92}
\loigiai{
$23 \times 4 = 92$.
}
\end{ex}

\begin{ex}
Tính: $68 : 2 = \dots$
\shortans{34}
\loigiai{
$68 : 2 = 34$.
}
\end{ex}

\begin{ex}
Tính: $96 : 3 = \dots$
\shortans{32}
\loigiai{
$96 : 3 = 32$.
}
\end{ex}

\begin{ex}
Tính: $85 : 5 = \dots$
\shortans{17}
\loigiai{
$85 : 5 = 17$.
}
\end{ex}

\begin{ex}
Tính: $72 : 4 = \dots$
\shortans{18}
\loigiai{
$72 : 4 = 18$.
}
\end{ex}

\begin{ex}
Tìm số dư trong phép chia $47 : 6$.
\shortans{5}
\loigiai{
$47 = 6 \times 7 + 5$, dư 5.
}
\end{ex}

\begin{ex}
Tìm số dư trong phép chia $59 : 8$.
\shortans{3}
\loigiai{
$59 = 8 \times 7 + 3$, dư 3.
}
\end{ex}

\begin{ex}
Gấp 15 lên 3 lần được bao nhiêu?
\shortans{45}
\loigiai{
$15 \times 3 = 45$.
}
\end{ex}

\begin{ex}
Giảm 48 đi 4 lần được bao nhiêu?
\shortans{12}
\loigiai{
$48 : 4 = 12$.
}
\end{ex}

\begin{ex}
Tính giá trị biểu thức: $13 \times 5 - 25$.
\shortans{40}
\loigiai{
$65 - 25 = 40$.
}
\end{ex}

\begin{ex}
Tính giá trị biểu thức: $92 : 4 + 37$.
\shortans{60}
\loigiai{
$23 + 37 = 60$.
}
\end{ex}

\begin{ex}
Tính giá trị biểu thức: $(24 + 16) : 5$.
\shortans{8}
\loigiai{
$40 : 5 = 8$.
}
\end{ex}

\begin{ex}
Tính giá trị biểu thức: $60 - 45 : 3$.
\shortans{45}
\loigiai{
$60 - 15 = 45$.
}
\end{ex}

\begin{ex}
Một cuộn dây dài 64 m cắt thành 4 đoạn bằng nhau. Hỏi mỗi đoạn dài bao nhiêu mét?
\shortans{16}
\loigiai{
$64 : 4 = 16$ m.
}
\end{ex}

\begin{ex}
Mẹ mua 14 quả cam, mẹ mua số quả quýt gấp 3 lần số quả cam. Hỏi mẹ mua bao nhiêu quả quýt?
\shortans{42}
\loigiai{
$14 \times 3 = 42$ quả quýt.
}
\end{ex}

\begin{ex}
Một cửa hàng có 90 kg gạo, đã bán đi một phần ba số gạo đó. Hỏi cửa hàng đã bán bao nhiêu ki-lô-gam gạo?
\shortans{30}
\loigiai{
$90 : 3 = 30$ kg.
}
\end{ex}

\begin{ex}
Lớp 3C có 35 học sinh chia đều thành 5 nhóm. Hỏi mỗi nhóm có bao nhiêu học sinh?
\shortans{7}
\loigiai{
$35 : 5 = 7$ học sinh.
}
\end{ex}

\begin{ex}
Một thùng đựng 54 lít mật ong, người ta rót đều vào 6 can. Hỏi mỗi can có bao nhiêu lít mật ong?
\shortans{9}
\loigiai{
$54 : 6 = 9$ lít.
}
\end{ex}

\begin{ex}
Trong phép chia có số chia là 7, thương là 9 và số dư là 4. Số bị chia là bao nhiêu?
\shortans{67}
\loigiai{
$9 \times 7 + 4 = 67$.
}
\end{ex}

\begin{ex}
Có 28 học sinh nam và 20 học sinh nữ cùng xếp thành 6 hàng đều nhau. Mỗi hàng có bao nhiêu học sinh?
\shortans{8}
\loigiai{
$(28 + 20) : 6 = 48 : 6 = 8$ học sinh.
}
\end{ex}

\begin{ex}
Một hình chữ nhật có chiều rộng 7 cm, chiều dài gấp 2 lần chiều rộng. Chu vi hình chữ nhật đó là bao nhiêu xăng-ti-mét?
\shortans{42}
\loigiai{
Chiều dài = 14 cm. Chu vi = $(14 + 7) \times 2 = 42$ cm.
}
\end{ex}

\begin{ex}
Tìm $x$: $x \times 4 = 96$.
\shortans{24}
\loigiai{
$x = 96 : 4 = 24$.
}
\end{ex}

\begin{ex}
Tìm $x$: $x : 3 = 28$.
\shortans{84}
\loigiai{
$x = 28 \times 3 = 84$.
}
\end{ex}

\begin{ex}
Tính: $15 \times 4 + 15 = \dots$
\shortans{75}
\loigiai{
$60 + 15 = 75$.
}
\end{ex}

\newpage


% =========================================================================
% Một số đơn vị đo độ dài, khối lượng, dung tích, nhiệt độ
% =========================================================================
\section{Một số đơn vị đo độ dài, khối lượng, dung tích, nhiệt độ}

\TaoPhieuBaiTap{Chủ đề 5 -- Một số đơn vị đo độ dài, khối lượng, dung tích, nhiệt độ}{30}{Mi-li-mét}{https://loptoanthaylinh.vn/bai-tap/toan-3/bai-30}
\begin{ex}
Đổi: $1\text{ cm} = \dots\text{ mm}$. Điền số thích hợp.
\shortans{10}
\loigiai{
$1\text{ cm} = 10\text{ mm}$.
}
\end{ex}

\begin{ex}
Đổi: $1\text{ dm} = \dots\text{ mm}$. Điền số thích hợp.
\shortans{100}
\loigiai{
$1\text{ dm} = 100\text{ mm}$.
}
\end{ex}

\begin{ex}
Đổi: $1\text{ m} = \dots\text{ mm}$. Điền số thích hợp.
\shortans{1000}
\loigiai{
$1\text{ m} = 1000\text{ mm}$.
}
\end{ex}

\begin{ex}
Đổi: $5\text{ cm} = \dots\text{ mm}$. Điền số thích hợp.
\shortans{50}
\loigiai{
$5\text{ cm} = 50\text{ mm}$.
}
\end{ex}

\begin{ex}
Đổi: $8\text{ cm } 4\text{ mm} = \dots\text{ mm}$. Điền số thích hợp.
\shortans{84}
\loigiai{
$8\text{ cm } 4\text{ mm} = 80 + 4 = 84\text{ mm}$.
}
\end{ex}

\begin{ex}
Đổi: $12\text{ cm } 5\text{ mm} = \dots\text{ mm}$. Điền số thích hợp.
\shortans{125}
\loigiai{
$120 + 5 = 125\text{ mm}$.
}
\end{ex}

\begin{ex}
Đổi: $70\text{ mm} = \dots\text{ cm}$. Điền số thích hợp.
\shortans{7}
\loigiai{
$70\text{ mm} = 7\text{ cm}$.
}
\end{ex}

\begin{ex}
Tính: $25\text{ mm} + 45\text{ mm} = \dots\text{ mm}$. Điền số thích hợp.
\shortans{70}
\loigiai{
$25 + 45 = 70$.
}
\end{ex}

\begin{ex}
Tính: $90\text{ mm} - 38\text{ mm} = \dots\text{ mm}$. Điền số thích hợp.
\shortans{52}
\loigiai{
$90 - 38 = 52$.
}
\end{ex}

\begin{ex}
Tính: $15\text{ mm} \times 4 = \dots\text{ mm}$. Điền số thích hợp.
\shortans{60}
\loigiai{
$15 \times 4 = 60$.
}
\end{ex}

\begin{ex}
Tính: $72\text{ mm} : 8 = \dots\text{ mm}$. Điền số thích hợp.
\shortans{9}
\loigiai{
$72 : 8 = 9$.
}
\end{ex}

\begin{ex}
Một con cào cào dài 35 mm, một con châu chấu dài 42 mm. Con châu chấu dài hơn con cào cào bao nhiêu mi-li-mét?
\shortans{7}
\loigiai{
$42 - 35 = 7$ mm.
}
\end{ex}

\begin{ex}
Một chiếc đinh dài 4 cm. Chiếc đinh đó dài bao nhiêu mi-li-mét?
\shortans{40}
\loigiai{
$4\text{ cm} = 40\text{ mm}$.
}
\end{ex}

\begin{ex}
Bề dày của một cuốn sách toán là 8 mm. Bề dày của 5 cuốn sách xếp chồng lên nhau là bao nhiêu mi-li-mét?
\shortans{40}
\loigiai{
$8 \times 5 = 40$ mm.
}
\end{ex}

\begin{ex}
Một đoạn dây đồng dài 6 cm được cắt bớt 25 mm. Đoạn dây còn lại dài bao nhiêu mi-li-mét?
\shortans{35}
\loigiai{
$6\text{ cm} = 60\text{ mm}$; $60 - 25 = 35$ mm.
}
\end{ex}

\begin{ex}
Một cục tẩy dài 45 mm, gọt bớt 12 mm. Cục tẩy còn lại dài bao nhiêu mi-li-mét?
\shortans{33}
\loigiai{
$45 - 12 = 33$ mm.
}
\end{ex}

\begin{ex}
Bút chì của Nam dài 14 cm. Tính theo đơn vị mi-li-mét bút chì dài bao nhiêu?
\shortans{140}
\loigiai{
$14 \times 10 = 140$ mm.
}
\end{ex}

\begin{ex}
Con kiến bò được quãng đường dài 8 cm 5 mm. Hỏi con kiến đã bò được bao nhiêu mi-li-mét?
\shortans{85}
\loigiai{
$80 + 5 = 85$ mm.
}
\end{ex}

\begin{ex}
Một sợi chỉ dài 10 cm, bị đứt một đoạn dài 35 mm. Đoạn chỉ còn lại dài bao nhiêu mi-li-mét?
\shortans{65}
\loigiai{
$100 - 35 = 65$ mm.
}
\end{ex}

\begin{ex}
Một hình tam giác có ba cạnh đều dài 30 mm. Chu vi hình tam giác đó là bao nhiêu mi-li-mét?
\shortans{90}
\loigiai{
$30 \times 3 = 90$ mm.
}
\end{ex}

\begin{ex}
Một hình vuông có cạnh dài 25 mm. Chu vi hình vuông đó là bao nhiêu mi-li-mét?
\shortans{100}
\loigiai{
$25 \times 4 = 100$ mm.
}
\end{ex}

\begin{ex}
Một hình chữ nhật có chiều dài 40 mm, chiều rộng 25 mm. Chu vi hình chữ nhật đó là bao nhiêu mi-li-mét?
\shortans{130}
\loigiai{
$(40 + 25) \times 2 = 130$ mm.
}
\end{ex}

\begin{ex}
Đổi: $3\text{ dm} = \dots\text{ mm}$. Điền số thích hợp.
\shortans{300}
\loigiai{
$3\text{ dm} = 300\text{ mm}$.
}
\end{ex}

\begin{ex}
Tính: $24\text{ mm} \times 3 = \dots\text{ mm}$. Điền số thích hợp.
\shortans{72}
\loigiai{
$24 \times 3 = 72$.
}
\end{ex}

\begin{ex}
Tính: $80\text{ mm} : 4 = \dots\text{ mm}$. Điền số thích hợp.
\shortans{20}
\loigiai{
$80 : 4 = 20$.
}
\end{ex}

\newpage

\TaoPhieuBaiTap{Chủ đề 5 -- Một số đơn vị đo độ dài, khối lượng, dung tích, nhiệt độ}{31}{Gam}{https://loptoanthaylinh.vn/bai-tap/toan-3/bai-31}
\begin{ex}
Đổi: $1\text{ kg} = \dots\text{ g}$. Điền số thích hợp.
\shortans{1000}
\loigiai{
$1\text{ kg} = 1000\text{ g}$.
}
\end{ex}

\begin{ex}
Tính: $350\text{ g} + 250\text{ g} = \dots\text{ g}$. Điền số thích hợp.
\shortans{600}
\loigiai{
$350 + 250 = 600$.
}
\end{ex}

\begin{ex}
Tính: $800\text{ g} - 450\text{ g} = \dots\text{ g}$. Điền số thích hợp.
\shortans{350}
\loigiai{
$800 - 450 = 350$.
}
\end{ex}

\begin{ex}
Tính: $120\text{ g} \times 3 = \dots\text{ g}$. Điền số thích hợp.
\shortans{360}
\loigiai{
$120 \times 3 = 360$.
}
\end{ex}

\begin{ex}
Tính: $500\text{ g} : 5 = \dots\text{ g}$. Điền số thích hợp.
\shortans{100}
\loigiai{
$500 : 5 = 100$.
}
\end{ex}

\begin{ex}
Một quả táo cân nặng 180 g, một quả cam cân nặng 220 g. Cả quả táo và quả cam cân nặng bao nhiêu gam?
\shortans{400}
\loigiai{
$180 + 220 = 400$ g.
}
\end{ex}

\begin{ex}
Một hộp sữa bột cân nặng 900 g, trong đó vỏ hộp nặng 100 g. Hỏi sữa trong hộp nặng bao nhiêu gam?
\shortans{800}
\loigiai{
$900 - 100 = 800$ g.
}
\end{ex}

\begin{ex}
Mỗi gói bánh cân nặng 150 g. Hỏi 4 gói bánh như thế cân nặng bao nhiêu gam?
\shortans{600}
\loigiai{
$150 \times 4 = 600$ g.
}
\end{ex}

\begin{ex}
Một gói bột ngọt cân nặng 450 g. Người ta đã dùng 200 g. Hỏi gói bột ngọt còn lại bao nhiêu gam?
\shortans{250}
\loigiai{
$450 - 200 = 250$ g.
}
\end{ex}

\begin{ex}
Mẹ mua 1 kg đường và đã dùng 350 g để làm bánh. Hỏi mẹ còn lại bao nhiêu gam đường?
\shortans{650}
\loigiai{
$1000 - 350 = 650$ g.
}
\end{ex}

\begin{ex}
Có 5 gói trà nặng như nhau, tổng cân nặng là 500 g. Hỏi mỗi gói trà cân nặng bao nhiêu gam?
\shortans{100}
\loigiai{
$500 : 5 = 100$ g.
}
\end{ex}

\begin{ex}
Một gói mì tôm nặng 85 g. Hỏi 2 gói mì tôm như thế nặng bao nhiêu gam?
\shortans{170}
\loigiai{
$85 \times 2 = 170$ g.
}
\end{ex}

\begin{ex}
Tính: $250\text{ g} \times 2 + 100\text{ g} = \dots\text{ g}$.
\shortans{600}
\loigiai{
$500 + 100 = 600$.
}
\end{ex}

\begin{ex}
Tính: $1000\text{ g} - 400\text{ g} : 2 = \dots\text{ g}$.
\shortans{800}
\loigiai{
$1000 - 200 = 800$.
}
\end{ex}

\begin{ex}
Một con gà nặng 1800 g, một con vịt nặng 1500 g. Con gà nặng hơn con vịt bao nhiêu gam?
\shortans{300}
\loigiai{
$1800 - 1500 = 300$ g.
}
\end{ex}

\begin{ex}
Mẹ mua 3 gói kẹo, mỗi gói nặng 200 g và 1 gói bánh nặng 300 g. Mẹ mua tất cả bao nhiêu gam kẹo và bánh?
\shortans{900}
\loigiai{
$3 \times 200 + 300 = 900$ g.
}
\end{ex}

\begin{ex}
Điền số thích hợp: $450\text{ g} + \dots\text{ g} = 1000\text{ g}$.
\shortans{550}
\loigiai{
$1000 - 450 = 550$.
}
\end{ex}

\begin{ex}
Điền số thích hợp: $1\text{ kg} - \dots\text{ g} = 750\text{ g}$.
\shortans{250}
\loigiai{
$1000 - 750 = 250$.
}
\end{ex}

\begin{ex}
Một quả dưa hấu nặng 2 kg 500 g. Quả dưa hấu nặng bao nhiêu gam?
\shortans{2500}
\loigiai{
$2000 + 500 = 2500$ g.
}
\end{ex}

\begin{ex}
Có 800 g bột mì chia đều vào 4 túi. Mỗi túi đựng bao nhiêu gam bột mì?
\shortans{200}
\loigiai{
$800 : 4 = 200$ g.
}
\end{ex}

\begin{ex}
Một gói đường nặng 500 g. Hỏi 2 gói đường như thế nặng bao nhiêu ki-lô-gam?
\shortans{1}
\loigiai{
$500 \times 2 = 1000\text{ g} = 1\text{ kg}$.
}
\end{ex}

\begin{ex}
Một quả trứng gà nặng khoảng 50 g. Hỏi 6 quả trứng gà nặng khoảng bao nhiêu gam?
\shortans{300}
\loigiai{
$50 \times 6 = 300$ g.
}
\end{ex}

\begin{ex}
Bác An có 700 g hạt giống, đã gieo 420 g. Bác An còn lại bao nhiêu gam hạt giống?
\shortans{280}
\loigiai{
$700 - 420 = 280$ g.
}
\end{ex}

\begin{ex}
Tính: $60\text{ g} \times 5 = \dots\text{ g}$.
\shortans{300}
\loigiai{
$60 \times 5 = 300$.
}
\end{ex}

\begin{ex}
Tính: $900\text{ g} : 3 = \dots\text{ g}$.
\shortans{300}
\loigiai{
$900 : 3 = 300$.
}
\end{ex}

\newpage

\TaoPhieuBaiTap{Chủ đề 5 -- Một số đơn vị đo độ dài, khối lượng, dung tích, nhiệt độ}{32}{Mi-li-lít}{https://loptoanthaylinh.vn/bai-tap/toan-3/bai-32}
\begin{ex}
Đổi: $1\text{ l} = \dots\text{ ml}$. Điền số thích hợp.
\shortans{1000}
\loigiai{
$1\text{ l} = 1000\text{ ml}$.
}
\end{ex}

\begin{ex}
Tính: $250\text{ ml} + 150\text{ ml} = \dots\text{ ml}$.
\shortans{400}
\loigiai{
$250 + 150 = 400$.
}
\end{ex}

\begin{ex}
Tính: $750\text{ ml} - 300\text{ ml} = \dots\text{ ml}$.
\shortans{450}
\loigiai{
$750 - 300 = 450$.
}
\end{ex}

\begin{ex}
Tính: $120\text{ ml} \times 4 = \dots\text{ ml}$.
\shortans{480}
\loigiai{
$120 \times 4 = 480$.
}
\end{ex}

\begin{ex}
Tính: $600\text{ ml} : 3 = \dots\text{ ml}$.
\shortans{200}
\loigiai{
$600 : 3 = 200$.
}
\end{ex}

\begin{ex}
Một chai nước khoáng có dung tích 500 ml. Hỏi 2 chai nước như thế có tất cả bao nhiêu mi-li-lít nước?
\shortans{1000}
\loigiai{
$500 \times 2 = 1000$ ml.
}
\end{ex}

\begin{ex}
Một hộp sữa chua chứa 100 ml sữa. Hỏi 6 hộp sữa chua như thế chứa tất cả bao nhiêu mi-li-lít?
\shortans{600}
\loigiai{
$100 \times 6 = 600$ ml.
}
\end{ex}

\begin{ex}
Một bình chứa 1 lít nước, người ta đã rót ra 350 ml. Hỏi trong bình còn lại bao nhiêu mi-li-lít nước?
\shortans{650}
\loigiai{
$1000 - 350 = 650$ ml.
}
\end{ex}

\begin{ex}
Một thìa si-rô ho chứa 5 ml thuốc. Bác sĩ chỉ định uống mỗi ngày 3 thìa. Hỏi mỗi ngày bé uống bao nhiêu mi-li-lít si-rô?
\shortans{15}
\loigiai{
$5 \times 3 = 15$ ml.
}
\end{ex}

\begin{ex}
Một bình nước có 800 ml nước. Rót đều vào 4 cốc nước. Hỏi mỗi cốc có bao nhiêu mi-li-lít nước?
\shortans{200}
\loigiai{
$800 : 4 = 200$ ml.
}
\end{ex}

\begin{ex}
Tính: $350\text{ ml} + 650\text{ ml} = \dots\text{ ml}$.
\shortans{1000}
\loigiai{
$350 + 650 = 1000$.
}
\end{ex}

\begin{ex}
Tính: $1000\text{ ml} - 450\text{ ml} = \dots\text{ ml}$.
\shortans{550}
\loigiai{
$1000 - 450 = 550$.
}
\end{ex}

\begin{ex}
Điền số thích hợp: $400\text{ ml} + \dots\text{ ml} = 1000\text{ ml}$.
\shortans{600}
\loigiai{
$1000 - 400 = 600$.
}
\end{ex}

\begin{ex}
Một ca chứa 300 ml nước, thêm vào ca 250 ml nước nữa. Trong ca có bao nhiêu mi-li-lít nước?
\shortans{550}
\loigiai{
$300 + 250 = 550$ ml.
}
\end{ex}

\begin{ex}
Một chai mật ong có 750 ml. Mẹ dùng hết 250 ml để làm thuốc. Trong chai còn lại bao nhiêu mi-li-lít mật ong?
\shortans{500}
\loigiai{
$750 - 250 = 500$ ml.
}
\end{ex}

\begin{ex}
Mỗi ngày bạn Nam uống 4 cốc nước, mỗi cốc chứa 250 ml. Hỏi mỗi ngày Nam uống bao nhiêu mi-li-lít nước?
\shortans{1000}
\loigiai{
$4 \times 250 = 1000$ ml.
}
\end{ex}

\begin{ex}
Một lọ thuốc nhỏ mắt có dung tích 10 ml. Hỏi 5 lọ thuốc như thế có dung tích bao nhiêu mi-li-lít?
\shortans{50}
\loigiai{
$10 \times 5 = 50$ ml.
}
\end{ex}

\begin{ex}
Người ta rót 900 ml nước trái cây vào 9 chiếc cốc bằng nhau. Mỗi chiếc cốc chứa bao nhiêu mi-li-lít nước trái cây?
\shortans{100}
\loigiai{
$900 : 9 = 100$ ml.
}
\end{ex}

\begin{ex}
Trong bình có 500 ml nước, rót thêm vào 200 ml nước nữa rồi chia đều vào 7 cốc. Mỗi cốc có bao nhiêu mi-li-lít nước?
\shortans{100}
\loigiai{
Tổng lượng nước: $500 + 200 = 700$ ml. Mỗi cốc: $700 : 7 = 100$ ml.
}
\end{ex}

\begin{ex}
Tính: $150\text{ ml} \times 5 = \dots\text{ ml}$.
\shortans{750}
\loigiai{
$150 \times 5 = 750$.
}
\end{ex}

\begin{ex}
Tính: $450\text{ ml} : 9 = \dots\text{ ml}$.
\shortans{50}
\loigiai{
$450 : 9 = 50$.
}
\end{ex}

\begin{ex}
Một chai dầu ăn có dung tích 1 lít. Rót ra một can nhỏ 200 ml và một bát nhỏ 50 ml. Trong chai còn lại bao nhiêu mi-li-lít?
\shortans{750}
\loigiai{
$1000 - (200 + 50) = 750$ ml.
}
\end{ex}

\begin{ex}
2 chai nước khoáng mỗi chai 500 ml tương đương với bao nhiêu lít nước?
\shortans{1}
\loigiai{
$500 \times 2 = 1000\text{ ml} = 1\text{ l}$.
}
\end{ex}

\begin{ex}
Tính: $80\text{ ml} \times 6 = \dots\text{ ml}$.
\shortans{480}
\loigiai{
$80 \times 6 = 480$.
}
\end{ex}

\begin{ex}
Tính: $360\text{ ml} : 4 = \dots\text{ ml}$.
\shortans{90}
\loigiai{
$360 : 4 = 90$.
}
\end{ex}

\newpage

\TaoPhieuBaiTap{Chủ đề 5 -- Một số đơn vị đo độ dài, khối lượng, dung tích, nhiệt độ}{33}{Nhiệt độ. Đơn vị đo nhiệt độ}{https://loptoanthaylinh.vn/bai-tap/toan-3/bai-33}
\begin{ex}
Đơn vị đo nhiệt độ thường dùng ở Việt Nam là độ gì? (Độ C ghi 1, Độ F ghi 2).
\shortans{1}
\loigiai{
Đơn vị thường dùng là độ C ($^\circ$C).
}
\end{ex}

\begin{ex}
Nhiệt kế đang chỉ $30^\circ$C. Nếu nhiệt độ tăng thêm $5^\circ$C thì nhiệt kế chỉ bao nhiêu độ C?
\shortans{35}
\loigiai{
$30 + 5 = 35^\circ$C.
}
\end{ex}

\begin{ex}
Nhiệt kế đang chỉ $28^\circ$C. Vào ban đêm nhiệt độ giảm $4^\circ$C. Nhiệt độ ban đêm là bao nhiêu độ C?
\shortans{24}
\loigiai{
$28 - 4 = 24^\circ$C.
}
\end{ex}

\begin{ex}
Nhiệt độ cơ thể của một người bình thường khỏe mạnh khoảng bao nhiêu độ C?
\shortans{37}
\loigiai{
Nhiệt độ người khỏe mạnh khoảng $37^\circ$C.
}
\end{ex}

\begin{ex}
Nước đá đang tan có nhiệt độ là bao nhiêu độ C?
\shortans{0}
\loigiai{
Nước đá tan ở $0^\circ$C.
}
\end{ex}

\begin{ex}
Nước sôi ở điều kiện tiêu chuẩn có nhiệt độ là bao nhiêu độ C?
\shortans{100}
\loigiai{
Nước sôi ở $100^\circ$C.
}
\end{ex}

\begin{ex}
Nhiệt độ phòng lúc sáng sớm là $22^\circ$C, lúc giữa trưa là $32^\circ$C. Nhiệt độ giữa trưa cao hơn lúc sáng sớm bao nhiêu độ C?
\shortans{10}
\loigiai{
$32 - 22 = 10^\circ$C.
}
\end{ex}

\begin{ex}
Hôm qua nhiệt độ ở Sa Pa là $12^\circ$C, hôm nay giảm đi $3^\circ$C. Nhiệt độ ở Sa Pa hôm nay là bao nhiêu độ C?
\shortans{9}
\loigiai{
$12 - 3 = 9^\circ$C.
}
\end{ex}

\begin{ex}
Nhiệt độ tại Hà Nội là $28^\circ$C, tại TP. Hồ Chí Minh là $34^\circ$C. TP. Hồ Chí Minh nóng hơn Hà Nội bao nhiêu độ C?
\shortans{6}
\loigiai{
$34 - 28 = 6^\circ$C.
}
\end{ex}

\begin{ex}
Dụng cụ nào dùng để đo nhiệt độ? (Nhiệt kế ghi 1, thước kẻ ghi 2).
\shortans{1}
\loigiai{
Nhiệt kế dùng để đo nhiệt độ.
}
\end{ex}

\begin{ex}
Nhiệt độ lúc 8 giờ là $25^\circ$C. Đến 12 giờ nhiệt độ tăng thêm $6^\circ$C. Hỏi lúc 12 giờ nhiệt độ là bao nhiêu độ C?
\shortans{31}
\loigiai{
$25 + 6 = 31^\circ$C.
}
\end{ex}

\begin{ex}
Nhiệt độ ngăn mát tủ lạnh là $5^\circ$C. Ngăn đá có nhiệt độ thấp hơn ngăn mát $23^\circ$C. Hỏi nhiệt độ ngăn đá là âm bao nhiêu độ C? (Điền số tự nhiên thể hiện độ chênh lệch là 23).
\shortans{23}
\loigiai{
Chênh lệch nhiệt độ là $23^\circ$C.
}
\end{ex}

\begin{ex}
Một bạn bị sốt có thân nhiệt là $39^\circ$C. Nhiệt độ này cao hơn thân nhiệt bình thường ($37^\circ$C) bao nhiêu độ C?
\shortans{2}
\loigiai{
$39 - 37 = 2^\circ$C.
}
\end{ex}

\begin{ex}
Nhiệt độ hôm nay là $30^\circ$C, dự báo ngày mai nhiệt độ tăng $2^\circ$C. Nhiệt độ ngày mai dự kiến là bao nhiêu độ C?
\shortans{32}
\loigiai{
$30 + 2 = 32^\circ$C.
}
\end{ex}

\begin{ex}
Khi đo nhiệt độ nước tắm cho em bé, nhiệt kế chỉ $38^\circ$C. Sau một lúc nhiệt độ giảm $3^\circ$C. Hỏi nhiệt độ lúc sau là bao nhiêu độ C?
\shortans{35}
\loigiai{
$38 - 3 = 35^\circ$C.
}
\end{ex}

\begin{ex}
Nhiệt độ buổi sáng là $20^\circ$C, buổi trưa tăng thêm $8^\circ$C, buổi chiều giảm đi $4^\circ$C so với buổi trưa. Nhiệt độ buổi chiều là bao nhiêu độ C?
\shortans{24}
\loigiai{
$20 + 8 - 4 = 24^\circ$C.
}
\end{ex}

\begin{ex}
Tính: $25^\circ\text{C} + 12^\circ\text{C} = \dots^\circ\text{C}$.
\shortans{37}
\loigiai{
$25 + 12 = 37$.
}
\end{ex}

\begin{ex}
Tính: $36^\circ\text{C} - 15^\circ\text{C} = \dots^\circ\text{C}$.
\shortans{21}
\loigiai{
$36 - 15 = 21$.
}
\end{ex}

\begin{ex}
Nhiệt độ ngoài trời mùa đông ở Mẫu Sơn có thể xuống tới $2^\circ$C. Nếu tăng thêm $15^\circ$C thì được bao nhiêu độ C?
\shortans{17}
\loigiai{
$2 + 15 = 17^\circ$C.
}
\end{ex}

\begin{ex}
Bác sĩ đo nhiệt độ của bệnh nhân lúc 6 giờ là $38^\circ$C, sau khi uống thuốc lúc 8 giờ nhiệt độ giảm xuống còn $37^\circ$C. Nhiệt độ giảm bao nhiêu độ C?
\shortans{1}
\loigiai{
$38 - 37 = 1^\circ$C.
}
\end{ex}

\begin{ex}
Nhiệt độ ở Đà Lạt là $18^\circ$C, ở Nha Trang là $29^\circ$C. Nha Trang có nhiệt độ cao hơn Đà Lạt bao nhiêu độ C?
\shortans{11}
\loigiai{
$29 - 18 = 11^\circ$C.
}
\end{ex}

\begin{ex}
Tính: $16^\circ\text{C} + 18^\circ\text{C} = \dots^\circ\text{C}$.
\shortans{34}
\loigiai{
$16 + 18 = 34$.
}
\end{ex}

\begin{ex}
Tính: $40^\circ\text{C} - 16^\circ\text{C} = \dots^\circ\text{C}$.
\shortans{24}
\loigiai{
$40 - 16 = 24$.
}
\end{ex}

\begin{ex}
Nhiệt độ một ngày hè ở Hà Nội cao nhất là $38^\circ$C, thấp nhất là $28^\circ$C. Chênh lệch nhiệt độ trong ngày là bao nhiêu độ C?
\shortans{10}
\loigiai{
$38 - 28 = 10^\circ$C.
}
\end{ex}

\begin{ex}
Nhiệt độ nước trà nóng là $85^\circ$C. Để nguội nhiệt độ giảm đi $40^\circ$C. Nhiệt độ lúc sau là bao nhiêu độ C?
\shortans{45}
\loigiai{
$85 - 40 = 45^\circ$C.
}
\end{ex}

\newpage

\TaoPhieuBaiTap{Chủ đề 5 -- Một số đơn vị đo độ dài, khối lượng, dung tích, nhiệt độ}{34}{Thực hành trải nghiệm đơn vị đo}{https://loptoanthaylinh.vn/bai-tap/toan-3/bai-34}
\begin{ex}
Đoạn que dài 15 cm. Bạn Nam đo được bao nhiêu mi-li-mét?
\shortans{150}
\loigiai{
$15\text{ cm} = 150\text{ mm}$.
}
\end{ex}

\begin{ex}
Một quả chuối nặng khoảng 120 g. Hỏi 3 quả chuối như thế nặng khoảng bao nhiêu gam?
\shortans{360}
\loigiai{
$120 \times 3 = 360$ g.
}
\end{ex}

\begin{ex}
Rót hết chai nước 500 ml vào các cốc loại 100 ml. Rót được bao nhiêu cốc như thế?
\shortans{5}
\loigiai{
$500 : 100 = 5$ cốc.
}
\end{ex}

\begin{ex}
Nhiệt kế chỉ $26^\circ$C. Sau 1 giờ nhiệt độ tăng $3^\circ$C. Nhiệt kế chỉ bao nhiêu độ C?
\shortans{29}
\loigiai{
$26 + 3 = 29^\circ$C.
}
\end{ex}

\begin{ex}
Chiều dài một chiếc đinh ghim là 25 mm. Bốn chiếc đinh ghim xếp nối tiếp dài bao nhiêu mi-li-mét?
\shortans{100}
\loigiai{
$25 \times 4 = 100$ mm.
}
\end{ex}

\begin{ex}
Một gói muối nặng 500 g, một gói bột ngọt nặng 200 g. Cả hai gói nặng bao nhiêu gam?
\shortans{700}
\loigiai{
$500 + 200 = 700$ g.
}
\end{ex}

\begin{ex}
Một hộp sữa tươi có dung tích 180 ml. Bé uống 2 hộp thì đã uống bao nhiêu mi-li-lít sữa?
\shortans{360}
\loigiai{
$180 \times 2 = 360$ ml.
}
\end{ex}

\begin{ex}
Nhiệt độ nước trong cốc đá là $4^\circ$C. Rót thêm nước nóng làm nhiệt độ tăng lên $20^\circ$C. Nhiệt độ mới là bao nhiêu độ C?
\shortans{24}
\loigiai{
$4 + 20 = 24^\circ$C.
}
\end{ex}

\begin{ex}
Một cuốn vở ô li có bề dày 5 mm. Một chồng vở gồm 8 cuốn có bề dày bao nhiêu mi-li-mét?
\shortans{40}
\loigiai{
$5 \times 8 = 40$ mm.
}
\end{ex}

\begin{ex}
Bác thợ may dùng 40 cm vải để may túi. Chiều dài mảnh vải tính theo mi-li-mét là bao nhiêu?
\shortans{400}
\loigiai{
$40\text{ cm} = 400\text{ mm}$.
}
\end{ex}

\begin{ex}
Cân nặng của bạn An là 32 kg. Bạn Bình nhẹ hơn bạn An 3 kg. Bạn Bình nặng bao nhiêu ki-lô-gam?
\shortans{29}
\loigiai{
$32 - 3 = 29$ kg.
}
\end{ex}

\begin{ex}
Bình nước có dung tích 2 lít. Đổi sang mi-li-lít được bao nhiêu mi-li-lít?
\shortans{2000}
\loigiai{
$2\text{ l} = 2000\text{ ml}$.
}
\end{ex}

\begin{ex}
Người ta rót 400 ml nước vào bình đã có sẵn 350 ml nước. Lượng nước trong bình bây giờ là bao nhiêu mi-li-lít?
\shortans{750}
\loigiai{
$350 + 400 = 750$ ml.
}
\end{ex}

\begin{ex}
Một quả bóng bàn nặng khoảng 3 g. Hỏi 10 quả bóng bàn như thế nặng bao nhiêu gam?
\shortans{30}
\loigiai{
$3 \times 10 = 30$ g.
}
\end{ex}

\begin{ex}
Một thanh sô-cô-la nặng 100 g chia đều cho 5 bạn. Mỗi bạn được bao nhiêu gam sô-cô-la?
\shortans{20}
\loigiai{
$100 : 5 = 20$ g.
}
\end{ex}

\begin{ex}
Đổi: $4\text{ cm } 8\text{ mm} = \dots\text{ mm}$.
\shortans{48}
\loigiai{
$40 + 8 = 48$ mm.
}
\end{ex}

\begin{ex}
Đổi: $1\text{ l } - 300\text{ ml} = \dots\text{ ml}$.
\shortans{700}
\loigiai{
$1000 - 300 = 700$ ml.
}
\end{ex}

\begin{ex}
Đổi: $1\text{ kg } - 600\text{ g} = \dots\text{ g}$.
\shortans{400}
\loigiai{
$1000 - 600 = 400$ g.
}
\end{ex}

\begin{ex}
Nhiệt độ cao nhất trong ngày là $35^\circ$C, thấp nhất là $25^\circ$C. Chênh lệch là bao nhiêu độ C?
\shortans{10}
\loigiai{
$35 - 25 = 10^\circ$C.
}
\end{ex}

\begin{ex}
Một chiếc bút dạ dài 135 mm, chiếc bút bi ngắn hơn 15 mm. Chiếc bút bi dài bao nhiêu mi-li-mét?
\shortans{120}
\loigiai{
$135 - 15 = 120$ mm.
}
\end{ex}

\begin{ex}
Hộp bánh quy nặng 350 g, hộp kẹo nặng 250 g. Hộp bánh quy nặng hơn hộp kẹo bao nhiêu gam?
\shortans{100}
\loigiai{
$350 - 250 = 100$ g.
}
\end{ex}

\begin{ex}
Có 600 ml dầu ăn, dùng 150 ml để rán bánh chưng. Còn lại bao nhiêu mi-li-lít dầu ăn?
\shortans{450}
\loigiai{
$600 - 150 = 450$ ml.
}
\end{ex}

\begin{ex}
Nhiệt độ phòng tăng từ $21^\circ$C lên $27^\circ$C. Đã tăng bao nhiêu độ C?
\shortans{6}
\loigiai{
$27 - 21 = 6^\circ$C.
}
\end{ex}

\begin{ex}
Đoạn thẳng AB dài 50 mm, đoạn thẳng CD dài 7 cm. Đoạn CD dài hơn đoạn AB bao nhiêu mi-li-mét?
\shortans{20}
\loigiai{
$7\text{ cm} = 70\text{ mm}$; $70 - 50 = 20$ mm.
}
\end{ex}

\begin{ex}
Tính: $15\text{ g} \times 6 = \dots\text{ g}$.
\shortans{90}
\loigiai{
$15 \times 6 = 90$.
}
\end{ex}

\newpage

\TaoPhieuBaiTap{Chủ đề 5 -- Một số đơn vị đo độ dài, khối lượng, dung tích, nhiệt độ}{35}{Luyện tập chung đo lường}{https://loptoanthaylinh.vn/bai-tap/toan-3/bai-35}
\begin{ex}
Đổi: $6\text{ cm } 5\text{ mm} = \dots\text{ mm}$.
\shortans{65}
\loigiai{
$60 + 5 = 65$ mm.
}
\end{ex}

\begin{ex}
Đổi: $2\text{ dm } 4\text{ mm} = \dots\text{ mm}$.
\shortans{204}
\loigiai{
$200 + 4 = 204$ mm.
}
\end{ex}

\begin{ex}
Tính: $350\text{ g} + 450\text{ g} = \dots\text{ g}$.
\shortans{800}
\loigiai{
$350 + 450 = 800$.
}
\end{ex}

\begin{ex}
Tính: $1000\text{ g} - 320\text{ g} = \dots\text{ g}$.
\shortans{680}
\loigiai{
$1000 - 320 = 680$.
}
\end{ex}

\begin{ex}
Tính: $250\text{ ml} \times 3 = \dots\text{ ml}$.
\shortans{750}
\loigiai{
$250 \times 3 = 750$.
}
\end{ex}

\begin{ex}
Tính: $800\text{ ml} : 4 = \dots\text{ ml}$.
\shortans{200}
\loigiai{
$800 : 4 = 200$.
}
\end{ex}

\begin{ex}
Một con thỏ cân nặng 2500 g. Con thỏ cân nặng bao nhiêu ki-lô-gam và gam? (Chỉ ghi số gam lẻ ngoài 2 kg).
\shortans{500}
\loigiai{
$2500\text{ g} = 2\text{ kg } 500\text{ g}$.
}
\end{ex}

\begin{ex}
Tính: $14\text{ mm} \times 5 = \dots\text{ mm}$.
\shortans{70}
\loigiai{
$14 \times 5 = 70$ mm.
}
\end{ex}

\begin{ex}
Tính: $96\text{ mm} : 3 = \dots\text{ mm}$.
\shortans{32}
\loigiai{
$96 : 3 = 32$ mm.
}
\end{ex}

\begin{ex}
Một chai sữa có 450 ml, em uống 150 ml. Chai sữa còn lại bao nhiêu mi-li-lít?
\shortans{300}
\loigiai{
$450 - 150 = 300$ ml.
}
\end{ex}

\begin{ex}
Nhiệt độ sáng nay là $19^\circ$C, đến trưa tăng lên $27^\circ$C. Nhiệt độ đã tăng bao nhiêu độ C?
\shortans{8}
\loigiai{
$27 - 19 = 8^\circ$C.
}
\end{ex}

\begin{ex}
Mẹ mua 4 gói mì chính, mỗi gói nặng 100 g. Mẹ đã mua bao nhiêu gam mì chính?
\shortans{400}
\loigiai{
$100 \times 4 = 400$ g.
}
\end{ex}

\begin{ex}
Một chai nước mắm có 750 ml, một chai giấm có 250 ml. Cả hai chai có bao nhiêu mi-li-lít?
\shortans{1000}
\loigiai{
$750 + 250 = 1000$ ml.
}
\end{ex}

\begin{ex}
Một hình chữ nhật có chiều dài 50 mm, chiều rộng 30 mm. Chu vi hình chữ nhật đó là bao nhiêu mi-li-mét?
\shortans{160}
\loigiai{
$(50 + 30) \times 2 = 160$ mm.
}
\end{ex}

\begin{ex}
Một hình vuông có cạnh dài 20 mm. Chu vi hình vuông đó là bao nhiêu mi-li-mét?
\shortans{80}
\loigiai{
$20 \times 4 = 80$ mm.
}
\end{ex}

\begin{ex}
Nhiệt độ nước nóng là $70^\circ$C, sau 30 phút giảm đi $25^\circ$C. Nhiệt độ lúc sau là bao nhiêu độ C?
\shortans{45}
\loigiai{
$70 - 25 = 45^\circ$C.
}
\end{ex}

\begin{ex}
Có 1 kg bột làm bánh, mẹ dùng hết 650 g. Mẹ còn lại bao nhiêu gam bột?
\shortans{350}
\loigiai{
$1000 - 650 = 350$ g.
}
\end{ex}

\begin{ex}
Một gói kẹo nặng 125 g. Hỏi 4 gói kẹo như vậy nặng bao nhiêu gam?
\shortans{500}
\loigiai{
$125 \times 4 = 500$ g.
}
\end{ex}

\begin{ex}
Rót 900 ml nước cam vào các ly, mỗi ly 150 ml. Rót được bao nhiêu ly?
\shortans{6}
\loigiai{
$900 : 150 = 6$ ly.
}
\end{ex}

\begin{ex}
Đoạn thẳng AB dài 12 cm, đoạn thẳng CD dài 75 mm. Đoạn AB dài hơn đoạn CD bao nhiêu mi-li-mét?
\shortans{45}
\loigiai{
$120 - 75 = 45$ mm.
}
\end{ex}

\begin{ex}
Tính: $28^\circ\text{C} + 9^\circ\text{C} = \dots^\circ\text{C}$.
\shortans{37}
\loigiai{
$28 + 9 = 37$.
}
\end{ex}

\begin{ex}
Tính: $35^\circ\text{C} - 17^\circ\text{C} = \dots^\circ\text{C}$.
\shortans{18}
\loigiai{
$35 - 17 = 18$.
}
\end{ex}

\begin{ex}
Có 3 bình mật ong, mỗi bình chứa 200 ml. Tổng số mật ong là bao nhiêu mi-li-lít?
\shortans{600}
\loigiai{
$200 \times 3 = 600$ ml.
}
\end{ex}

\begin{ex}
Một chiếc bút chì dài 18 cm, một chiếc kẹp giấy dài 30 mm. Chiếc bút chì dài gấp mấy lần chiếc kẹp giấy?
\shortans{6}
\loigiai{
$18\text{ cm} = 180\text{ mm}$; $180 : 30 = 6$ lần.
}
\end{ex}

\begin{ex}
Tính: $1000\text{ mm} - 650\text{ mm} = \dots\text{ mm}$.
\shortans{350}
\loigiai{
$1000 - 650 = 350$.
}
\end{ex}

\newpage


% =========================================================================
% Phép nhân, phép chia trong phạm vi 1 000
% =========================================================================
\section{Phép nhân, phép chia trong phạm vi 1 000}

\TaoPhieuBaiTap{Chủ đề 6 -- Phép nhân, phép chia trong phạm vi 1 000}{36}{Nhân số có ba chữ số với số có một chữ số}{https://loptoanthaylinh.vn/bai-tap/toan-3/bai-36}
\begin{ex}
Tính: $123 \times 2 = \dots$
\shortans{246}
\loigiai{
$123 \times 2 = 246$.
}
\end{ex}

\begin{ex}
Tính: $214 \times 2 = \dots$
\shortans{428}
\loigiai{
$214 \times 2 = 428$.
}
\end{ex}

\begin{ex}
Tính: $321 \times 3 = \dots$
\shortans{963}
\loigiai{
$321 \times 3 = 963$.
}
\end{ex}

\begin{ex}
Tính: $102 \times 4 = \dots$
\shortans{408}
\loigiai{
$102 \times 4 = 408$.
}
\end{ex}

\begin{ex}
Tính: $115 \times 3 = \dots$
\shortans{345}
\loigiai{
$115 \times 3 = 345$.
}
\end{ex}

\begin{ex}
Tính: $126 \times 3 = \dots$
\shortans{378}
\loigiai{
$126 \times 3 = 378$.
}
\end{ex}

\begin{ex}
Tính: $142 \times 4 = \dots$
\shortans{568}
\loigiai{
$142 \times 4 = 568$.
}
\end{ex}

\begin{ex}
Tính: $165 \times 2 = \dots$
\shortans{330}
\loigiai{
$165 \times 2 = 330$.
}
\end{ex}

\begin{ex}
Tính: $175 \times 3 = \dots$
\shortans{525}
\loigiai{
$175 \times 3 = 525$.
}
\end{ex}

\begin{ex}
Tính: $184 \times 2 = \dots$
\shortans{368}
\loigiai{
$184 \times 2 = 368$.
}
\end{ex}

\begin{ex}
Tính: $218 \times 4 = \dots$
\shortans{872}
\loigiai{
$218 \times 4 = 872$.
}
\end{ex}

\begin{ex}
Tính: $245 \times 3 = \dots$
\shortans{735}
\loigiai{
$245 \times 3 = 735$.
}
\end{ex}

\begin{ex}
Tính: $152 \times 5 = \dots$
\shortans{760}
\loigiai{
$152 \times 5 = 760$.
}
\end{ex}

\begin{ex}
Tính: $160 \times 6 = \dots$
\shortans{960}
\loigiai{
$160 \times 6 = 960$.
}
\end{ex}

\begin{ex}
Mỗi chuyến xe buýt chở được 125 hành khách. Hỏi 3 chuyến xe như thế chở được bao nhiêu hành khách?
\shortans{375}
\loigiai{
$125 \times 3 = 375$ hành khách.
}
\end{ex}

\begin{ex}
Một cửa hàng có 4 thùng dầu, mỗi thùng chứa 115 lít dầu. Hỏi cửa hàng có tất cả bao nhiêu lít dầu?
\shortans{460}
\loigiai{
$115 \times 4 = 460$ lít.
}
\end{ex}

\begin{ex}
Mỗi cuộn vải dài 140 m. Hỏi 5 cuộn vải như thế dài bao nhiêu mét?
\shortans{700}
\loigiai{
$140 \times 5 = 700$ m.
}
\end{ex}

\begin{ex}
Một trường tiểu học mua 3 bộ bàn ghế, mỗi bộ giá 250 nghìn đồng. Hỏi phải trả bao nhiêu nghìn đồng?
\shortans{750}
\loigiai{
$250 \times 3 = 750$ nghìn đồng.
}
\end{ex}

\begin{ex}
Tính giá trị biểu thức: $132 \times 3 + 120$.
\shortans{516}
\loigiai{
$396 + 120 = 516$.
}
\end{ex}

\begin{ex}
Tính giá trị biểu thức: $215 \times 4 - 360$.
\shortans{500}
\loigiai{
$860 - 360 = 500$.
}
\end{ex}

\begin{ex}
Tìm $x$, biết: $x : 3 = 214$.
\shortans{642}
\loigiai{
$x = 214 \times 3 = 642$.
}
\end{ex}

\begin{ex}
Tìm $x$, biết: $x : 4 = 152$.
\shortans{608}
\loigiai{
$x = 152 \times 4 = 608$.
}
\end{ex}

\begin{ex}
Mỗi bao xi măng nặng 50 kg. Hỏi 18 bao xi măng như thế nặng bao nhiêu ki-lô-gam?
\shortans{900}
\loigiai{
$50 \times 18 = 900$ kg.
}
\end{ex}

\begin{ex}
Một trang trại nuôi vịt có 2 chuồng, mỗi chuồng có 325 con vịt. Trang trại có bao nhiêu con vịt?
\shortans{650}
\loigiai{
$325 \times 2 = 650$ con vịt.
}
\end{ex}

\begin{ex}
Tính: $111 \times 9 = \dots$
\shortans{999}
\loigiai{
$111 \times 9 = 999$.
}
\end{ex}

\newpage

\TaoPhieuBaiTap{Chủ đề 6 -- Phép nhân, phép chia trong phạm vi 1 000}{37}{Chia số có ba chữ số cho số có một chữ số}{https://loptoanthaylinh.vn/bai-tap/toan-3/bai-37}
\begin{ex}
Tính: $248 : 2 = \dots$
\shortans{124}
\loigiai{
$248 : 2 = 124$.
}
\end{ex}

\begin{ex}
Tính: $369 : 3 = \dots$
\shortans{123}
\loigiai{
$369 : 3 = 123$.
}
\end{ex}

\begin{ex}
Tính: $488 : 4 = \dots$
\shortans{122}
\loigiai{
$488 : 4 = 122$.
}
\end{ex}

\begin{ex}
Tính: $555 : 5 = \dots$
\shortans{111}
\loigiai{
$555 : 5 = 111$.
}
\end{ex}

\begin{ex}
Tính: $426 : 3 = \dots$
\shortans{142}
\loigiai{
$426 : 3 = 142$.
}
\end{ex}

\begin{ex}
Tính: $515 : 5 = \dots$
\shortans{103}
\loigiai{
$515 : 5 = 103$.
}
\end{ex}

\begin{ex}
Tính: $624 : 6 = \dots$
\shortans{104}
\loigiai{
$624 : 6 = 104$.
}
\end{ex}

\begin{ex}
Tính: $735 : 7 = \dots$
\shortans{105}
\loigiai{
$735 : 7 = 105$.
}
\end{ex}

\begin{ex}
Tính: $848 : 8 = \dots$
\shortans{106}
\loigiai{
$848 : 8 = 106$.
}
\end{ex}

\begin{ex}
Tính: $927 : 9 = \dots$
\shortans{103}
\loigiai{
$927 : 9 = 103$.
}
\end{ex}

\begin{ex}
Tính: $456 : 4 = \dots$
\shortans{114}
\loigiai{
$456 : 4 = 114$.
}
\end{ex}

\begin{ex}
Tính: $565 : 5 = \dots$
\shortans{113}
\loigiai{
$565 : 5 = 113$.
}
\end{ex}

\begin{ex}
Tính: $672 : 6 = \dots$
\shortans{112}
\loigiai{
$672 : 6 = 112$.
}
\end{ex}

\begin{ex}
Tính: $784 : 7 = \dots$
\shortans{112}
\loigiai{
$784 : 7 = 112$.
}
\end{ex}

\begin{ex}
Tính: $856 : 8 = \dots$
\shortans{107}
\loigiai{
$856 : 8 = 107$.
}
\end{ex}

\begin{ex}
Có 345 kg đường chia đều vào 3 bao. Hỏi mỗi bao có bao nhiêu ki-lô-gam đường?
\shortans{115}
\loigiai{
$345 : 3 = 115$ kg.
}
\end{ex}

\begin{ex}
Một cuộn dây dài 484 m cắt thành 4 đoạn bằng nhau. Hỏi mỗi đoạn dài bao nhiêu mét?
\shortans{121}
\loigiai{
$484 : 4 = 121$ m.
}
\end{ex}

\begin{ex}
Một trường có 636 học sinh chia đều vào 6 khối. Hỏi mỗi khối có bao nhiêu học sinh?
\shortans{106}
\loigiai{
$636 : 6 = 106$ học sinh.
}
\end{ex}

\begin{ex}
Tìm số dư trong phép chia $547 : 5$.
\shortans{2}
\loigiai{
$547 = 5 \times 109 + 2$, dư 2.
}
\end{ex}

\begin{ex}
Tìm thương trong phép chia $648 : 6$.
\shortans{108}
\loigiai{
$648 : 6 = 108$.
}
\end{ex}

\begin{ex}
Tìm $x$, biết: $x \times 4 = 528$.
\shortans{132}
\loigiai{
$x = 528 : 4 = 132$.
}
\end{ex}

\begin{ex}
Tìm $x$, biết: $3 \times x = 726$.
\shortans{242}
\loigiai{
$x = 726 : 3 = 242$.
}
\end{ex}

\begin{ex}
Tính giá trị biểu thức: $840 : 4 + 40$.
\shortans{250}
\loigiai{
$210 + 40 = 250$.
}
\end{ex}

\begin{ex}
Tính giá trị biểu thức: $963 : 3 - 121$.
\shortans{200}
\loigiai{
$321 - 121 = 200$.
}
\end{ex}

\begin{ex}
Có 750 lít dầu chia đều vào 5 thùng. Mỗi thùng chứa bao nhiêu lít dầu?
\shortans{150}
\loigiai{
$750 : 5 = 150$ lít.
}
\end{ex}

\newpage

\TaoPhieuBaiTap{Chủ đề 6 -- Phép nhân, phép chia trong phạm vi 1 000}{38}{Biểu thức số. Tính giá trị của biểu thức số}{https://loptoanthaylinh.vn/bai-tap/toan-3/bai-38}
\begin{ex}
Tính giá trị của biểu thức: $125 + 45 - 70$.
\shortans{100}
\loigiai{
$170 - 70 = 100$.
}
\end{ex}

\begin{ex}
Tính giá trị của biểu thức: $80 - 20 + 35$.
\shortans{95}
\loigiai{
$60 + 35 = 95$.
}
\end{ex}

\begin{ex}
Tính giá trị của biểu thức: $24 \times 2 : 4$.
\shortans{12}
\loigiai{
$48 : 4 = 12$.
}
\end{ex}

\begin{ex}
Tính giá trị của biểu thức: $96 : 3 \times 2$.
\shortans{64}
\loigiai{
$32 \times 2 = 64$.
}
\end{ex}

\begin{ex}
Tính giá trị của biểu thức: $50 + 20 \times 3$.
\shortans{110}
\loigiai{
Nhân trước cộng sau: $50 + 60 = 110$.
}
\end{ex}

\begin{ex}
Tính giá trị của biểu thức: $90 - 40 : 2$.
\shortans{70}
\loigiai{
Chia trước trừ sau: $90 - 20 = 70$.
}
\end{ex}

\begin{ex}
Tính giá trị của biểu thức: $150 - 30 \times 4$.
\shortans{30}
\loigiai{
$150 - 120 = 30$.
}
\end{ex}

\begin{ex}
Tính giá trị của biểu thức: $200 + 400 : 5$.
\shortans{280}
\loigiai{
$200 + 80 = 280$.
}
\end{ex}

\begin{ex}
Tính giá trị của biểu thức: $(45 + 15) \times 3$.
\shortans{180}
\loigiai{
Trong ngoặc trước: $60 \times 3 = 180$.
}
\end{ex}

\begin{ex}
Tính giá trị của biểu thức: $(84 - 24) : 6$.
\shortans{10}
\loigiai{
Trong ngoặc trước: $60 : 6 = 10$.
}
\end{ex}

\begin{ex}
Tính giá trị của biểu thức: $120 : (2 + 4)$.
\shortans{20}
\loigiai{
$120 : 6 = 20$.
}
\end{ex}

\begin{ex}
Tính giá trị của biểu thức: $40 \times (10 - 6)$.
\shortans{160}
\loigiai{
$40 \times 4 = 160$.
}
\end{ex}

\begin{ex}
Tính giá trị của biểu thức: $35 + 25 \times 2 - 10$.
\shortans{75}
\loigiai{
$35 + 50 - 10 = 75$.
}
\end{ex}

\begin{ex}
Tính giá trị của biểu thức: $100 - 60 : 2 + 15$.
\shortans{85}
\loigiai{
$100 - 30 + 15 = 85$.
}
\end{ex}

\begin{ex}
Tính giá trị của biểu thức: $(30 + 50) : 4 + 12$.
\shortans{32}
\loigiai{
$80 : 4 + 12 = 20 + 12 = 32$.
}
\end{ex}

\begin{ex}
Tính giá trị của biểu thức: $250 - (60 + 40) \times 2$.
\shortans{50}
\loigiai{
$250 - 100 \times 2 = 250 - 200 = 50$.
}
\end{ex}

\begin{ex}
Tính giá trị của biểu thức: $48 : 4 + 36 : 3$.
\shortans{24}
\loigiai{
$12 + 12 = 24$.
}
\end{ex}

\begin{ex}
Tính giá trị của biểu thức: $15 \times 4 - 40 : 2$.
\shortans{40}
\loigiai{
$60 - 20 = 40$.
}
\end{ex}

\begin{ex}
Tính giá trị của biểu thức: $70 + 30 : (5 + 5)$.
\shortans{73}
\loigiai{
$70 + 30 : 10 = 70 + 3 = 73$.
}
\end{ex}

\begin{ex}
Tính giá trị của biểu thức: $(100 - 20) : (4 + 4)$.
\shortans{10}
\loigiai{
$80 : 8 = 10$.
}
\end{ex}

\begin{ex}
Tìm giá trị của $A$ biết $A = 12 \times 5 + 40$.
\shortans{100}
\loigiai{
$A = 60 + 40 = 100$.
}
\end{ex}

\begin{ex}
Tìm giá trị của $B$ biết $B = 180 : 3 - 25$.
\shortans{35}
\loigiai{
$B = 60 - 25 = 35$.
}
\end{ex}

\begin{ex}
Điền số thích hợp: $(35 + \dots) \times 2 = 100$.
\shortans{15}
\loigiai{
$35 + x = 50 \Rightarrow x = 15$.
}
\end{ex}

\begin{ex}
Điền số thích hợp: $80 - \dots : 3 = 60$.
\shortans{60}
\loigiai{
$x : 3 = 20 \Rightarrow x = 60$.
}
\end{ex}

\begin{ex}
Tính giá trị của biểu thức: $240 : 6 \times 5$.
\shortans{200}
\loigiai{
$40 \times 5 = 200$.
}
\end{ex}

\newpage

\TaoPhieuBaiTap{Chủ đề 6 -- Phép nhân, phép chia trong phạm vi 1 000}{39}{So sánh số lớn gấp mấy lần số bé}{https://loptoanthaylinh.vn/bai-tap/toan-3/bai-39}
\begin{ex}
Số 24 gấp số 6 mấy lần?
\shortans{4}
\loigiai{
$24 : 6 = 4$ lần.
}
\end{ex}

\begin{ex}
Số 35 gấp số 7 mấy lần?
\shortans{5}
\loigiai{
$35 : 7 = 5$ lần.
}
\end{ex}

\begin{ex}
Số 48 gấp số 8 mấy lần?
\shortans{6}
\loigiai{
$48 : 8 = 6$ lần.
}
\end{ex}

\begin{ex}
Số 54 gấp số 9 mấy lần?
\shortans{6}
\loigiai{
$54 : 9 = 6$ lần.
}
\end{ex}

\begin{ex}
Số 72 gấp số 8 mấy lần?
\shortans{9}
\loigiai{
$72 : 8 = 9$ lần.
}
\end{ex}

\begin{ex}
Số 63 gấp số 7 mấy lần?
\shortans{9}
\loigiai{
$63 : 7 = 9$ lần.
}
\end{ex}

\begin{ex}
Số 80 gấp số 10 mấy lần?
\shortans{8}
\loigiai{
$80 : 10 = 8$ lần.
}
\end{ex}

\begin{ex}
Số 90 gấp số 30 mấy lần?
\shortans{3}
\loigiai{
$90 : 30 = 3$ lần.
}
\end{ex}

\begin{ex}
Số 100 gấp số 20 mấy lần?
\shortans{5}
\loigiai{
$100 : 20 = 5$ lần.
}
\end{ex}

\begin{ex}
Đoạn thẳng AB dài 18 cm, đoạn CD dài 6 cm. Đoạn AB dài gấp mấy lần đoạn CD?
\shortans{3}
\loigiai{
$18 : 6 = 3$ lần.
}
\end{ex}

\begin{ex}
Mẹ 36 tuổi, con 6 tuổi. Tuổi mẹ gấp mấy lần tuổi con?
\shortans{6}
\loigiai{
$36 : 6 = 6$ lần.
}
\end{ex}

\begin{ex}
Bố 40 tuổi, con 8 tuổi. Tuổi bố gấp mấy lần tuổi con?
\shortans{5}
\loigiai{
$40 : 8 = 5$ lần.
}
\end{ex}

\begin{ex}
Bà 70 tuổi, cháu 7 tuổi. Tuổi bà gấp mấy lần tuổi cháu?
\shortans{10}
\loigiai{
$70 : 7 = 10$ lần.
}
\end{ex}

\begin{ex}
Bao gạo to nặng 50 kg, bao gạo nhỏ nặng 10 kg. Bao to nặng gấp mấy lần bao nhỏ?
\shortans{5}
\loigiai{
$50 : 10 = 5$ lần.
}
\end{ex}

\begin{ex}
Bể thứ nhất chứa 64 lít nước, bể thứ hai chứa 8 lít nước. Bể thứ nhất gấp mấy lần bể thứ hai?
\shortans{8}
\loigiai{
$64 : 8 = 8$ lần.
}
\end{ex}

\begin{ex}
Quyển sách dày 96 trang, quyển truyện tranh dày 24 trang. Quyển sách dày gấp mấy lần quyển truyện?
\shortans{4}
\loigiai{
$96 : 24 = 4$ lần.
}
\end{ex}

\begin{ex}
Lớp 3A có 32 học sinh, trong đó có 8 bạn học sinh xuất sắc. Số học sinh cả lớp gấp mấy lần số học sinh xuất sắc?
\shortans{4}
\loigiai{
$32 : 8 = 4$ lần.
}
\end{ex}

\begin{ex}
Một sợi dây dài 45 m, cắt lấy một đoạn dài 5 m. Sợi dây ban đầu dài gấp mấy lần đoạn đã cắt?
\shortans{9}
\loigiai{
$45 : 5 = 9$ lần.
}
\end{ex}

\begin{ex}
Hùng có 42 viên bi, Dũng có 7 viên bi. Số bi của Hùng gấp mấy lần số bi của Dũng?
\shortans{6}
\loigiai{
$42 : 7 = 6$ lần.
}
\end{ex}

\begin{ex}
Vườn nhà Mai có 48 cây bưởi và 6 cây cam. Số cây bưởi gấp mấy lần số cây cam?
\shortans{8}
\loigiai{
$48 : 6 = 8$ lần.
}
\end{ex}

\begin{ex}
Một mảnh vườn hình chữ nhật có chiều dài 24 m, chiều rộng 6 m. Chiều dài gấp mấy lần chiều rộng?
\shortans{4}
\loigiai{
$24 : 6 = 4$ lần.
}
\end{ex}

\begin{ex}
Số lớn là 70, số bé là 14. Số lớn gấp mấy lần số bé?
\shortans{5}
\loigiai{
$70 : 14 = 5$ lần.
}
\end{ex}

\begin{ex}
Số lớn là 84, số bé là 12. Số lớn gấp mấy lần số bé?
\shortans{7}
\loigiai{
$84 : 12 = 7$ lần.
}
\end{ex}

\begin{ex}
Một hình vuông có chu vi 36 cm, cạnh 9 cm. Chu vi gấp mấy lần cạnh?
\shortans{4}
\loigiai{
$36 : 9 = 4$ lần.
}
\end{ex}

\begin{ex}
Một ngày có 24 giờ. 24 giờ gấp mấy lần khoảng thời gian 4 giờ?
\shortans{6}
\loigiai{
$24 : 4 = 6$ lần.
}
\end{ex}

\newpage

\TaoPhieuBaiTap{Chủ đề 6 -- Phép nhân, phép chia trong phạm vi 1 000}{40}{Luyện tập chung biểu thức}{https://loptoanthaylinh.vn/bai-tap/toan-3/bai-40}
\begin{ex}
Tính: $142 \times 3 = \dots$
\shortans{426}
\loigiai{
$142 \times 3 = 426$.
}
\end{ex}

\begin{ex}
Tính: $215 \times 4 = \dots$
\shortans{860}
\loigiai{
$215 \times 4 = 860$.
}
\end{ex}

\begin{ex}
Tính: $486 : 2 = \dots$
\shortans{243}
\loigiai{
$486 : 2 = 243$.
}
\end{ex}

\begin{ex}
Tính: $729 : 9 = \dots$
\shortans{81}
\loigiai{
$729 : 9 = 81$.
}
\end{ex}

\begin{ex}
Tính giá trị biểu thức: $120 + 30 \times 4$.
\shortans{240}
\loigiai{
$120 + 120 = 240$.
}
\end{ex}

\begin{ex}
Tính giá trị biểu thức: $250 - 150 : 5$.
\shortans{220}
\loigiai{
$250 - 30 = 220$.
}
\end{ex}

\begin{ex}
Tính giá trị biểu thức: $(48 + 32) : 8$.
\shortans{10}
\loigiai{
$80 : 8 = 10$.
}
\end{ex}

\begin{ex}
Tính giá trị biểu thức: $16 \times (15 - 10)$.
\shortans{80}
\loigiai{
$16 \times 5 = 80$.
}
\end{ex}

\begin{ex}
Số 64 gấp số 8 mấy lần?
\shortans{8}
\loigiai{
$64 : 8 = 8$ lần.
}
\end{ex}

\begin{ex}
Số 81 gấp số 9 mấy lần?
\shortans{9}
\loigiai{
$81 : 9 = 9$ lần.
}
\end{ex}

\begin{ex}
Mẹ 45 tuổi, con 9 tuổi. Tuổi mẹ gấp mấy lần tuổi con?
\shortans{5}
\loigiai{
$45 : 9 = 5$ lần.
}
\end{ex}

\begin{ex}
Có 3 thùng dầu, mỗi thùng chứa 120 lít dầu. Hỏi cả 3 thùng có tất cả bao nhiêu lít dầu?
\shortans{360}
\loigiai{
$120 \times 3 = 360$ lít.
}
\end{ex}

\begin{ex}
Một cuộn dây dài 636 m được chia thành 6 đoạn bằng nhau. Hỏi mỗi đoạn dài bao nhiêu mét?
\shortans{106}
\loigiai{
$636 : 6 = 106$ m.
}
\end{ex}

\begin{ex}
Tính số dư của phép chia: $435 : 4$.
\shortans{3}
\loigiai{
$435 = 4 \times 108 + 3$, dư 3.
}
\end{ex}

\begin{ex}
Tìm $x$: $x : 3 = 145$.
\shortans{435}
\loigiai{
$x = 145 \times 3 = 435$.
}
\end{ex}

\begin{ex}
Tìm $x$: $4 \times x = 484$.
\shortans{121}
\loigiai{
$x = 484 : 4 = 121$.
}
\end{ex}

\begin{ex}
Một bao gạo nặng 60 kg, bao đậu xanh nặng 15 kg. Bao gạo nặng gấp mấy lần bao đậu xanh?
\shortans{4}
\loigiai{
$60 : 15 = 4$ lần.
}
\end{ex}

\begin{ex}
Tính: $118 \times 5 = \dots$
\shortans{590}
\loigiai{
$118 \times 5 = 590$.
}
\end{ex}

\begin{ex}
Tính: $848 : 4 = \dots$
\shortans{212}
\loigiai{
$848 : 4 = 212$.
}
\end{ex}

\begin{ex}
Tính giá trị biểu thức: $300 - 80 \times 3$.
\shortans{60}
\loigiai{
$300 - 240 = 60$.
}
\end{ex}

\begin{ex}
Tính giá trị biểu thức: $100 + 400 : 4$.
\shortans{200}
\loigiai{
$100 + 100 = 200$.
}
\end{ex}

\begin{ex}
Một rổ có 36 quả xoài và 6 quả ổi. Số xoài gấp mấy lần số ổi?
\shortans{6}
\loigiai{
$36 : 6 = 6$ lần.
}
\end{ex}

\begin{ex}
Một cửa hàng bán được 124 kg gạo trong ngày thứ nhất, ngày thứ hai bán được gấp đôi ngày thứ nhất. Ngày thứ hai bán được bao nhiêu ki-lô-gam?
\shortans{248}
\loigiai{
$124 \times 2 = 248$ kg.
}
\end{ex}

\begin{ex}
Một tấm vải dài 250 m, người ta bán đi 100 m rồi chia đều phần còn lại thành 3 phần. Mỗi phần dài bao nhiêu mét?
\shortans{50}
\loigiai{
Còn lại: $250 - 100 = 150$ m. Mỗi phần: $150 : 3 = 50$ m.
}
\end{ex}

\begin{ex}
Tính: $125 \times 4 = \dots$
\shortans{500}
\loigiai{
$125 \times 4 = 500$.
}
\end{ex}

\newpage


% =========================================================================
% Ôn tập học kì 1
% =========================================================================
\section{Ôn tập học kì 1}

\TaoPhieuBaiTap{Chủ đề 7 -- Ôn tập học kì 1}{41}{Ôn tập phép nhân, phép chia trong phạm vi 1000}{https://loptoanthaylinh.vn/bai-tap/toan-3/bai-41}
\begin{ex}
Tính: $16 \times 4 = \dots$
\shortans{64}
\loigiai{
$16 \times 4 = 64$.
}
\end{ex}

\begin{ex}
Tính: $25 \times 3 = \dots$
\shortans{75}
\loigiai{
$25 \times 3 = 75$.
}
\end{ex}

\begin{ex}
Tính: $132 \times 3 = \dots$
\shortans{396}
\loigiai{
$132 \times 3 = 396$.
}
\end{ex}

\begin{ex}
Tính: $214 \times 4 = \dots$
\shortans{856}
\loigiai{
$214 \times 4 = 856$.
}
\end{ex}

\begin{ex}
Tính: $84 : 4 = \dots$
\shortans{21}
\loigiai{
$84 : 4 = 21$.
}
\end{ex}

\begin{ex}
Tính: $96 : 3 = \dots$
\shortans{32}
\loigiai{
$96 : 3 = 32$.
}
\end{ex}

\begin{ex}
Tính: $484 : 4 = \dots$
\shortans{121}
\loigiai{
$484 : 4 = 121$.
}
\end{ex}

\begin{ex}
Tính: $639 : 3 = \dots$
\shortans{213}
\loigiai{
$639 : 3 = 213$.
}
\end{ex}

\begin{ex}
Tính: $550 : 5 = \dots$
\shortans{110}
\loigiai{
$550 : 5 = 110$.
}
\end{ex}

\begin{ex}
Tìm thương trong phép chia: $742 : 7$.
\shortans{106}
\loigiai{
$742 : 7 = 106$.
}
\end{ex}

\begin{ex}
Tìm số dư trong phép chia: $65 : 7$.
\shortans{2}
\loigiai{
$65 = 7 \times 9 + 2$, dư 2.
}
\end{ex}

\begin{ex}
Tìm số dư trong phép chia: $487 : 4$.
\shortans{3}
\loigiai{
$487 = 4 \times 121 + 3$, dư 3.
}
\end{ex}

\begin{ex}
Tìm $x$, biết: $x \times 5 = 125$.
\shortans{25}
\loigiai{
$x = 125 : 5 = 25$.
}
\end{ex}

\begin{ex}
Tìm $x$, biết: $x : 3 = 112$.
\shortans{336}
\loigiai{
$x = 112 \times 3 = 336$.
}
\end{ex}

\begin{ex}
Một người nuôi 4 lồng gà, mỗi lồng có 12 con gà. Hỏi có tất cả bao nhiêu con gà?
\shortans{48}
\loigiai{
$12 \times 4 = 48$ con gà.
}
\end{ex}

\begin{ex}
Có 72 lít mật ong rót đều vào 8 can. Hỏi mỗi can có bao nhiêu lít mật ong?
\shortans{9}
\loigiai{
$72 : 8 = 9$ lít.
}
\end{ex}

\begin{ex}
Một cuộn vải dài 180 m, đã may hết 60 m. Số vải còn lại chia đều cho 4 người. Mỗi người được bao nhiêu mét vải?
\shortans{30}
\loigiai{
Còn lại: $180 - 60 = 120$ m. Mỗi người: $120 : 4 = 30$ m.
}
\end{ex}

\begin{ex}
Tính: $14 \times 5 + 30 = \dots$
\shortans{100}
\loigiai{
$70 + 30 = 100$.
}
\end{ex}

\begin{ex}
Tính: $88 : 4 - 12 = \dots$
\shortans{10}
\loigiai{
$22 - 12 = 10$.
}
\end{ex}

\begin{ex}
Tính: $105 \times 4 = \dots$
\shortans{420}
\loigiai{
$105 \times 4 = 420$.
}
\end{ex}

\begin{ex}
Một trang trại có 3 chuồng lợn, mỗi chuồng có 45 con lợn. Trang trại có tất cả bao nhiêu con lợn?
\shortans{135}
\loigiai{
$45 \times 3 = 135$ con lợn.
}
\end{ex}

\begin{ex}
Có 96 quyển vở chia đều cho 8 bạn học sinh xuất sắc. Mỗi bạn nhận được bao nhiêu quyển vở?
\shortans{12}
\loigiai{
$96 : 8 = 12$ quyển vở.
}
\end{ex}

\begin{ex}
Một số chia cho 6 được thương là 14 và dư 5. Tìm số đó.
\shortans{89}
\loigiai{
$14 \times 6 + 5 = 84 + 5 = 89$.
}
\end{ex}

\begin{ex}
Một số chia cho 4 được thương là 125 và không có dư. Tìm số đó.
\shortans{500}
\loigiai{
$125 \times 4 = 500$.
}
\end{ex}

\begin{ex}
Tính: $600 : 3 : 2 = \dots$
\shortans{100}
\loigiai{
$200 : 2 = 100$.
}
\end{ex}

\newpage

\TaoPhieuBaiTap{Chủ đề 7 -- Ôn tập học kì 1}{42}{Ôn tập biểu thức số}{https://loptoanthaylinh.vn/bai-tap/toan-3/bai-42}
\begin{ex}
Tính giá trị của biểu thức: $45 + 55 - 30$.
\shortans{70}
\loigiai{
$100 - 30 = 70$.
}
\end{ex}

\begin{ex}
Tính giá trị của biểu thức: $120 - 40 + 60$.
\shortans{140}
\loigiai{
$80 + 60 = 140$.
}
\end{ex}

\begin{ex}
Tính giá trị của biểu thức: $15 \times 4 : 2$.
\shortans{30}
\loigiai{
$60 : 2 = 30$.
}
\end{ex}

\begin{ex}
Tính giá trị của biểu thức: $90 : 3 \times 5$.
\shortans{150}
\loigiai{
$30 \times 5 = 150$.
}
\end{ex}

\begin{ex}
Tính giá trị của biểu thức: $40 + 20 \times 4$.
\shortans{120}
\loigiai{
$40 + 80 = 120$.
}
\end{ex}

\begin{ex}
Tính giá trị của biểu thức: $100 - 60 : 3$.
\shortans{80}
\loigiai{
$100 - 20 = 80$.
}
\end{ex}

\begin{ex}
Tính giá trị của biểu thức: $150 + 250 : 5$.
\shortans{200}
\loigiai{
$150 + 50 = 200$.
}
\end{ex}

\begin{ex}
Tính giá trị của biểu thức: $200 - 45 \times 2$.
\shortans{110}
\loigiai{
$200 - 90 = 110$.
}
\end{ex}

\begin{ex}
Tính giá trị của biểu thức: $(60 + 40) \times 4$.
\shortans{400}
\loigiai{
$100 \times 4 = 400$.
}
\end{ex}

\begin{ex}
Tính giá trị của biểu thức: $(150 - 30) : 4$.
\shortans{30}
\loigiai{
$120 : 4 = 30$.
}
\end{ex}

\begin{ex}
Tính giá trị của biểu thức: $180 : (3 + 6)$.
\shortans{20}
\loigiai{
$180 : 9 = 20$.
}
\end{ex}

\begin{ex}
Tính giá trị của biểu thức: $25 \times (12 - 8)$.
\shortans{100}
\loigiai{
$25 \times 4 = 100$.
}
\end{ex}

\begin{ex}
Tính giá trị của biểu thức: $60 + 40 : 2 + 10$.
\shortans{90}
\loigiai{
$60 + 20 + 10 = 90$.
}
\end{ex}

\begin{ex}
Tính giá trị của biểu thức: $200 - (50 + 30) : 2$.
\shortans{160}
\loigiai{
$200 - 80 : 2 = 200 - 40 = 160$.
}
\end{ex}

\begin{ex}
Tính giá trị của biểu thức: $(20 + 30) \times (10 - 6)$.
\shortans{200}
\loigiai{
$50 \times 4 = 200$.
}
\end{ex}

\begin{ex}
Tính giá trị của biểu thức: $140 : 2 - 35$.
\shortans{35}
\loigiai{
$70 - 35 = 35$.
}
\end{ex}

\begin{ex}
Tính giá trị của biểu thức: $36 : 4 + 45 : 5$.
\shortans{18}
\loigiai{
$9 + 9 = 18$.
}
\end{ex}

\begin{ex}
Tính giá trị của biểu thức: $18 \times 3 - 24 : 4$.
\shortans{48}
\loigiai{
$54 - 6 = 48$.
}
\end{ex}

\begin{ex}
Tìm $x$, biết: $x + 20 \times 3 = 100$.
\shortans{40}
\loigiai{
$x + 60 = 100 \Rightarrow x = 40$.
}
\end{ex}

\begin{ex}
Tìm $x$, biết: $x - 60 : 2 = 50$.
\shortans{80}
\loigiai{
$x - 30 = 50 \Rightarrow x = 80$.
}
\end{ex}

\begin{ex}
Tính giá trị của biểu thức: $(24 + 16) \times 5$.
\shortans{200}
\loigiai{
$40 \times 5 = 200$.
}
\end{ex}

\begin{ex}
Tính giá trị của biểu thức: $300 - 150 : 3$.
\shortans{250}
\loigiai{
$300 - 50 = 250$.
}
\end{ex}

\begin{ex}
Tính giá trị của biểu thức: $80 : (2 \times 4)$.
\shortans{10}
\loigiai{
$80 : 8 = 10$.
}
\end{ex}

\begin{ex}
Tính giá trị của biểu thức: $(50 + 70) : 6 + 15$.
\shortans{35}
\loigiai{
$120 : 6 + 15 = 20 + 15 = 35$.
}
\end{ex}

\begin{ex}
Tính: $500 - 100 \times 3 = \dots$
\shortans{200}
\loigiai{
$500 - 300 = 200$.
}
\end{ex}

\newpage

\TaoPhieuBaiTap{Chủ đề 7 -- Ôn tập học kì 1}{43}{Ôn tập hình học và đo lường}{https://loptoanthaylinh.vn/bai-tap/toan-3/bai-43}
\begin{ex}
Đoạn thẳng AB dài 16 cm, M là trung điểm của AB. Độ dài đoạn AM là bao nhiêu xăng-ti-mét?
\shortans{8}
\loigiai{
$16 : 2 = 8$ cm.
}
\end{ex}

\begin{ex}
Một hình tròn có bán kính 9 cm. Đường kính hình tròn đó là bao nhiêu xăng-ti-mét?
\shortans{18}
\loigiai{
$9 \times 2 = 18$ cm.
}
\end{ex}

\begin{ex}
Một hình tròn có đường kính 30 cm. Bán kính hình tròn đó là bao nhiêu xăng-ti-mét?
\shortans{15}
\loigiai{
$30 : 2 = 15$ cm.
}
\end{ex}

\begin{ex}
Hình vuông có cạnh dài 6 cm. Chu vi của hình vuông đó là bao nhiêu xăng-ti-mét?
\shortans{24}
\loigiai{
$6 \times 4 = 24$ cm.
}
\end{ex}

\begin{ex}
Hình chữ nhật có chiều dài 12 cm, chiều rộng 8 cm. Chu vi của hình chữ nhật đó là bao nhiêu xăng-ti-mét?
\shortans{40}
\loigiai{
$(12 + 8) \times 2 = 40$ cm.
}
\end{ex}

\begin{ex}
Một hình tam giác có độ dài 3 cạnh là 9 cm, 12 cm, 15 cm. Chu vi tam giác đó là bao nhiêu xăng-ti-mét?
\shortans{36}
\loigiai{
$9 + 12 + 15 = 36$ cm.
}
\end{ex}

\begin{ex}
Một khối lập phương có bao nhiêu mặt?
\shortans{6}
\loigiai{
Có 6 mặt.
}
\end{ex}

\begin{ex}
Một khối lập phương có bao nhiêu cạnh?
\shortans{12}
\loigiai{
Có 12 cạnh.
}
\end{ex}

\begin{ex}
Một khối hộp chữ nhật có bao nhiêu đỉnh?
\shortans{8}
\loigiai{
Có 8 đỉnh.
}
\end{ex}

\begin{ex}
Đổi: $7\text{ cm } 4\text{ mm} = \dots\text{ mm}$.
\shortans{74}
\loigiai{
$70 + 4 = 74$ mm.
}
\end{ex}

\begin{ex}
Đổi: $1\text{ kg } = \dots\text{ g}$.
\shortans{1000}
\loigiai{
$1\text{ kg} = 1000\text{ g}$.
}
\end{ex}

\begin{ex}
Đổi: $1\text{ l } = \dots\text{ ml}$.
\shortans{1000}
\loigiai{
$1\text{ l} = 1000\text{ ml}$.
}
\end{ex}

\begin{ex}
Một chai sữa chứa 250 ml, 4 chai như thế chứa bao nhiêu mi-li-lít sữa?
\shortans{1000}
\loigiai{
$250 \times 4 = 1000$ ml.
}
\end{ex}

\begin{ex}
Một quả dưa nặng 1 kg 400 g. Quả dưa đó nặng bao nhiêu gam?
\shortans{1400}
\loigiai{
$1000 + 400 = 1400$ g.
}
\end{ex}

\begin{ex}
Nhiệt độ phòng là $24^\circ$C, sau đó tăng thêm $5^\circ$C. Nhiệt độ lúc sau là bao nhiêu độ C?
\shortans{29}
\loigiai{
$24 + 5 = 29^\circ$C.
}
\end{ex}

\begin{ex}
Một hình vuông có chu vi 32 cm. Độ dài mỗi cạnh của hình vuông là bao nhiêu xăng-ti-mét?
\shortans{8}
\loigiai{
$32 : 4 = 8$ cm.
}
\end{ex}

\begin{ex}
Một hình chữ nhật có chu vi 50 cm, chiều dài 15 cm. Chiều rộng hình chữ nhật đó là bao nhiêu xăng-ti-mét?
\shortans{10}
\loigiai{
Nửa chu vi = 25 cm; chiều rộng = $25 - 15 = 10$ cm.
}
\end{ex}

\begin{ex}
Dùng ê-ke để kiểm tra, hình chữ nhật có bao nhiêu góc vuông?
\shortans{4}
\loigiai{
Có 4 góc vuông.
}
\end{ex}

\begin{ex}
Một khối lập phương có cạnh dài 4 cm. Tổng độ dài tất cả 12 cạnh là bao nhiêu xăng-ti-mét?
\shortans{48}
\loigiai{
$4 \times 12 = 48$ cm.
}
\end{ex}

\begin{ex}
Đổi: $3\text{ dm } 5\text{ mm} = \dots\text{ mm}$.
\shortans{305}
\loigiai{
$300 + 5 = 305$ mm.
}
\end{ex}

\begin{ex}
Bác thợ rèn thanh sắt dài 1 m thành một khung hình vuông. Mỗi cạnh của khung hình vuông dài bao nhiêu xăng-ti-mét?
\shortans{25}
\loigiai{
$1\text{ m} = 100\text{ cm}$; $100 : 4 = 25$ cm.
}
\end{ex}

\begin{ex}
Một gói mì chính nặng 450 g, mẹ đã dùng 150 g. Còn lại bao nhiêu gam?
\shortans{300}
\loigiai{
$450 - 150 = 300$ g.
}
\end{ex}

\begin{ex}
Nhiệt kế chỉ $36^\circ$C. Vào ban đêm nhiệt độ giảm $7^\circ$C. Ban đêm nhiệt độ là bao nhiêu độ C?
\shortans{29}
\loigiai{
$36 - 7 = 29^\circ$C.
}
\end{ex}

\begin{ex}
Tính chu vi hình tam giác đều có cạnh 11 cm theo xăng-ti-mét.
\shortans{33}
\loigiai{
$11 \times 3 = 33$ cm.
}
\end{ex}

\begin{ex}
Một hình tứ giác có 4 cạnh đều bằng 7 cm. Chu vi của tứ giác đó là bao nhiêu xăng-ti-mét?
\shortans{28}
\loigiai{
$7 \times 4 = 28$ cm.
}
\end{ex}

\newpage

\TaoPhieuBaiTap{Chủ đề 7 -- Ôn tập học kì 1}{44}{Ôn tập chung học kì 1}{https://loptoanthaylinh.vn/bai-tap/toan-3/bai-44}
\begin{ex}
Số liền trước của 800 là số nào?
\shortans{799}
\loigiai{
$800 - 1 = 799$.
}
\end{ex}

\begin{ex}
Số liền sau của 699 là số nào?
\shortans{700}
\loigiai{
$699 + 1 = 700$.
}
\end{ex}

\begin{ex}
Tính: $356 + 243 = \dots$
\shortans{599}
\loigiai{
$356 + 243 = 599$.
}
\end{ex}

\begin{ex}
Tính: $785 - 342 = \dots$
\shortans{443}
\loigiai{
$785 - 342 = 443$.
}
\end{ex}

\begin{ex}
Tính: $124 \times 3 = \dots$
\shortans{372}
\loigiai{
$124 \times 3 = 372$.
}
\end{ex}

\begin{ex}
Tính: $648 : 2 = \dots$
\shortans{324}
\loigiai{
$648 : 2 = 324$.
}
\end{ex}

\begin{ex}
Tính giá trị biểu thức: $40 + 60 : 2$.
\shortans{70}
\loigiai{
$40 + 30 = 70$.
}
\end{ex}

\begin{ex}
Tính giá trị biểu thức: $(50 + 30) \times 3$.
\shortans{240}
\loigiai{
$80 \times 3 = 240$.
}
\end{ex}

\begin{ex}
Số 42 gấp số 7 mấy lần?
\shortans{6}
\loigiai{
$42 : 7 = 6$ lần.
}
\end{ex}

\begin{ex}
Số 72 gấp số 9 mấy lần?
\shortans{8}
\loigiai{
$72 : 9 = 8$ lần.
}
\end{ex}

\begin{ex}
Tìm $\frac{1}{6}$ của 54 quả cam.
\shortans{9}
\loigiai{
$54 : 6 = 9$ quả cam.
}
\end{ex}

\begin{ex}
Tìm $\frac{1}{8}$ của 48 kg gạo.
\shortans{6}
\loigiai{
$48 : 8 = 6$ kg gạo.
}
\end{ex}

\begin{ex}
Một hình chữ nhật có chiều dài 15 cm, chiều rộng 5 cm. Chu vi hình chữ nhật đó là bao nhiêu xăng-ti-mét?
\shortans{40}
\loigiai{
$(15 + 5) \times 2 = 40$ cm.
}
\end{ex}

\begin{ex}
Một hình vuông có cạnh 9 cm. Chu vi hình vuông đó là bao nhiêu xăng-ti-mét?
\shortans{36}
\loigiai{
$9 \times 4 = 36$ cm.
}
\end{ex}

\begin{ex}
Hình tròn có bán kính 7 cm thì đường kính là bao nhiêu xăng-ti-mét?
\shortans{14}
\loigiai{
$7 \times 2 = 14$ cm.
}
\end{ex}

\begin{ex}
Một thùng chứa 120 lít dầu, người ta bán đi $\frac{1}{3}$ số dầu đó. Hỏi đã bán bao nhiêu lít dầu?
\shortans{40}
\loigiai{
$120 : 3 = 40$ lít.
}
\end{ex}

\begin{ex}
Một cửa hàng ngày thứ nhất bán được 135 kg gạo, ngày thứ hai bán gấp đôi ngày thứ nhất. Cả hai ngày bán được bao nhiêu ki-lô-gam gạo?
\shortans{405}
\loigiai{
Ngày 2: $135 \times 2 = 270$ kg. Cả hai ngày: $135 + 270 = 405$ kg.
}
\end{ex}

\begin{ex}
Khối lập phương có bao nhiêu đỉnh?
\shortans{8}
\loigiai{
Có 8 đỉnh.
}
\end{ex}

\begin{ex}
Khối hộp chữ nhật có bao nhiêu mặt?
\shortans{6}
\loigiai{
Có 6 mặt.
}
\end{ex}

\begin{ex}
Đổi: $5\text{ cm } 6\text{ mm} = \dots\text{ mm}$.
\shortans{56}
\loigiai{
$50 + 6 = 56$ mm.
}
\end{ex}

\begin{ex}
Đổi: $1\text{ kg} - 250\text{ g} = \dots\text{ g}$.
\shortans{750}
\loigiai{
$1000 - 250 = 750$ g.
}
\end{ex}

\begin{ex}
Đổi: $1\text{ l} - 450\text{ ml} = \dots\text{ ml}$.
\shortans{550}
\loigiai{
$1000 - 450 = 550$ ml.
}
\end{ex}

\begin{ex}
Nhiệt độ trong phòng học là $25^\circ$C, ngoài trời là $33^\circ$C. Ngoài trời nóng hơn trong phòng bao nhiêu độ C?
\shortans{8}
\loigiai{
$33 - 25 = 8^\circ$C.
}
\end{ex}

\begin{ex}
Tìm $x$: $x \times 4 = 360$.
\shortans{90}
\loigiai{
$x = 360 : 4 = 90$.
}
\end{ex}

\begin{ex}
Tìm $x$: $x : 3 = 120 + 10$.
\shortans{390}
\loigiai{
$x : 3 = 130 \Rightarrow x = 130 \times 3 = 390$.
}
\end{ex}

\newpage


% =========================================================================
% Các số đến 10 000
% =========================================================================
\section{Các số đến 10 000}

\TaoPhieuBaiTap{Chủ đề 8 -- Các số đến 10 000}{45}{Các số có bốn chữ số. Số 10 000}{https://loptoanthaylinh.vn/bai-tap/toan-3/bai-45}
\begin{ex}
Số gồm 3 nghìn, 5 trăm, 2 chục và 7 đơn vị viết là gì?
\shortans{3527}
\loigiai{
Số đó là 3527.
}
\end{ex}

\begin{ex}
Số 6841 có chữ số hàng nghìn là bao nhiêu?
\shortans{6}
\loigiai{
Chữ số hàng nghìn là 6.
}
\end{ex}

\begin{ex}
Số 7054 có chữ số hàng trăm là bao nhiêu?
\shortans{0}
\loigiai{
Chữ số hàng trăm là 0.
}
\end{ex}

\begin{ex}
Số liền sau của số 9999 là số nào?
\shortans{10000}
\loigiai{
Số liền sau của 9999 là $9999 + 1 = 10000$.
}
\end{ex}

\begin{ex}
Số liền trước của số 5000 là số nào?
\shortans{4999}
\loigiai{
Số liền trước của 5000 là $5000 - 1 = 4999$.
}
\end{ex}

\begin{ex}
Số lớn nhất có bốn chữ số là số nào?
\shortans{9999}
\loigiai{
Số lớn nhất có bốn chữ số là 9999.
}
\end{ex}

\begin{ex}
Số bé nhất có bốn chữ số là số nào?
\shortans{1000}
\loigiai{
Số bé nhất có bốn chữ số là 1000.
}
\end{ex}

\begin{ex}
Số tròn nghìn liền sau của 4000 là số nào?
\shortans{5000}
\loigiai{
Số tròn nghìn liền sau của 4000 là 5000.
}
\end{ex}

\begin{ex}
Số tròn trăm liền trước của 6300 là số nào?
\shortans{6200}
\loigiai{
Số tròn trăm liền trước của 6300 là 6200.
}
\end{ex}

\begin{ex}
Số tròn chục liền sau của 7820 là số nào?
\shortans{7830}
\loigiai{
Số tròn chục liền sau là 7830.
}
\end{ex}

\begin{ex}
Viết số thích hợp: $4000 + 500 + 60 + 3 = \dots$
\shortans{4563}
\loigiai{
$4000 + 500 + 60 + 3 = 4563$.
}
\end{ex}

\begin{ex}
Viết số thích hợp: $8000 + 20 + 9 = \dots$
\shortans{8029}
\loigiai{
$8000 + 20 + 9 = 8029$.
}
\end{ex}

\begin{ex}
Viết số thích hợp: $9000 + 700 = \dots$
\shortans{9700}
\loigiai{
$9000 + 700 = 9700$.
}
\end{ex}

\begin{ex}
Số 8425 có giá trị của chữ số 4 là bao nhiêu?
\shortans{400}
\loigiai{
Chữ số 4 ở hàng trăm có giá trị là 400.
}
\end{ex}

\begin{ex}
Số 5913 có giá trị của chữ số 5 là bao nhiêu?
\shortans{5000}
\loigiai{
Chữ số 5 ở hàng nghìn có giá trị là 5000.
}
\end{ex}

\begin{ex}
Cho các chữ số 3, 0, 8, 5. Số lớn nhất có bốn chữ số khác nhau lập từ các chữ số trên là số nào?
\shortans{8530}
\loigiai{
Sắp xếp giảm dần: 8530.
}
\end{ex}

\begin{ex}
Cho các chữ số 1, 0, 4, 7. Số bé nhất có bốn chữ số khác nhau lập từ các chữ số trên là số nào?
\shortans{1047}
\loigiai{
Chữ số hàng nghìn khác 0 nhỏ nhất là 1, tiếp theo là 0, 4, 7: 1047.
}
\end{ex}

\begin{ex}
Có bao nhiêu số tròn nghìn từ 1000 đến 10 000?
\shortans{10}
\loigiai{
Các số: 1000, 2000, ..., 10000 (10 số).
}
\end{ex}

\begin{ex}
Điền số thích hợp: $8999 < \dots < 9001$.
\shortans{9000}
\loigiai{
Số ở giữa là 9000.
}
\end{ex}

\begin{ex}
Số liền sau của số 4329 là số nào?
\shortans{4330}
\loigiai{
$4329 + 1 = 4330$.
}
\end{ex}

\begin{ex}
Số liền trước của số 7100 là số nào?
\shortans{7099}
\loigiai{
$7100 - 1 = 7099$.
}
\end{ex}

\begin{ex}
Tổng của số lớn nhất có bốn chữ số và số bé nhất có một chữ số là bao nhiêu?
\shortans{9999}
\loigiai{
$9999 + 0 = 9999$.
}
\end{ex}

\begin{ex}
Hiệu của số bé nhất có bốn chữ số và số lớn nhất có ba chữ số là bao nhiêu?
\shortans{1}
\loigiai{
$1000 - 999 = 1$.
}
\end{ex}

\begin{ex}
Số 6666 có bao nhiêu chữ số 6?
\shortans{4}
\loigiai{
Có 4 chữ số 6.
}
\end{ex}

\begin{ex}
Giá trị của chữ số 9 trong số 9234 gấp mấy lần giá trị của chữ số 9 trong số 2394?
\shortans{100}
\loigiai{
9000 gấp 90 là $9000 : 90 = 100$ lần.
}
\end{ex}

\newpage

\TaoPhieuBaiTap{Chủ đề 8 -- Các số đến 10 000}{46}{So sánh các số trong phạm vi 10 000}{https://loptoanthaylinh.vn/bai-tap/toan-3/bai-46}
\begin{ex}
Trong các số $4520, 4250, 5420, 5240$, số nào lớn nhất?
\shortans{5420}
\loigiai{
5420 là số lớn nhất.
}
\end{ex}

\begin{ex}
Trong các số $7890, 7980, 7098, 7908$, số nào bé nhất?
\shortans{7098}
\loigiai{
7098 là số bé nhất.
}
\end{ex}

\begin{ex}
Trong các số $9102, 9012, 9201, 9120$, số lớn thứ nhì là số nào?
\shortans{9120}
\loigiai{
Thứ tự giảm dần: 9201, 9120, 9102, 9012. Số lớn thứ nhì là 9120.
}
\end{ex}

\begin{ex}
Tìm $x$ biết $x$ là số tròn nghìn thỏa mãn: $6000 < x < 8000$.
\shortans{7000}
\loigiai{
Số tròn nghìn duy nhất là 7000.
}
\end{ex}

\begin{ex}
Tìm số tự nhiên $y$ biết: $4567 < y < 4569$.
\shortans{4568}
\loigiai{
Số ở giữa là 4568.
}
\end{ex}

\begin{ex}
Cho 4 số: 3400, 3040, 3004, 3404. Hiệu của số lớn nhất và số bé nhất là bao nhiêu?
\shortans{400}
\loigiai{
Lớn nhất: 3404, bé nhất: 3004. Hiệu = $3404 - 3004 = 400$.
}
\end{ex}

\begin{ex}
Có bao nhiêu số tự nhiên có bốn chữ số lớn hơn 9995?
\shortans{4}
\loigiai{
Các số: 9996, 9997, 9998, 9999 (4 số).
}
\end{ex}

\begin{ex}
Có bao nhiêu số tự nhiên có bốn chữ số nhỏ hơn 1005?
\shortans{5}
\loigiai{
Các số: 1000, 1001, 1002, 1003, 1004 (5 số).
}
\end{ex}

\begin{ex}
Tìm số lớn nhất có 4 chữ số khác nhau lập từ các chữ số 9, 8, 7, 6.
\shortans{9876}
\loigiai{
Số đó là 9876.
}
\end{ex}

\begin{ex}
Tìm số bé nhất có 4 chữ số khác nhau lập từ các chữ số 1, 2, 3, 4.
\shortans{1234}
\loigiai{
Số đó là 1234.
}
\end{ex}

\begin{ex}
Điền chữ số thích hợp vào dấu * để: $56*8 > 5688$. Điền chữ số duy nhất thỏa mãn.
\shortans{9}
\loigiai{
Vì $5698 > 5688$, nên chữ số đó là 9.
}
\end{ex}

\begin{ex}
Điền chữ số thích hợp vào dấu * để: $7*35 < 7135$.
\shortans{0}
\loigiai{
Hàng trăm phải nhỏ hơn 1 nên $* = 0$.
}
\end{ex}

\begin{ex}
Khoảng cách giữa hai số 8500 và 8000 là bao nhiêu?
\shortans{500}
\loigiai{
$8500 - 8000 = 500$.
}
\end{ex}

\begin{ex}
Số 3500 kém số 5000 bao nhiêu đơn vị?
\shortans{1500}
\loigiai{
$5000 - 3500 = 1500$.
}
\end{ex}

\begin{ex}
Số 7200 hơn số 6800 bao nhiêu đơn vị?
\shortans{400}
\loigiai{
$7200 - 6800 = 400$.
}
\end{ex}

\begin{ex}
Đoàn Một chở 2500 kg hàng, Đoàn Hai chở 3200 kg hàng. Đoàn Hai chở nhiều hơn Đoàn Một bao nhiêu ki-lô-gam?
\shortans{700}
\loigiai{
$3200 - 2500 = 700$ kg.
}
\end{ex}

\begin{ex}
Thôn Đoài có 4520 người, thôn Đông có 4510 người. Thôn Đoài đông hơn thôn Đông bao nhiêu người?
\shortans{10}
\loigiai{
$4520 - 4510 = 10$ người.
}
\end{ex}

\begin{ex}
Bác An thu hoạch 2400 bắp ngô, bác Ba thu hoạch 2100 bắp ngô. Bác An nhiều hơn bác Ba bao nhiêu bắp ngô?
\shortans{300}
\loigiai{
$2400 - 2100 = 300$ bắp.
}
\end{ex}

\begin{ex}
Tìm $x$ là số tròn trăm thỏa mãn $4800 < x < 5000$.
\shortans{4900}
\loigiai{
Số đó là 4900.
}
\end{ex}

\begin{ex}
Tìm $x$ là số tròn chục thỏa mãn $6710 < x < 6730$.
\shortans{6720}
\loigiai{
Số đó là 6720.
}
\end{ex}

\begin{ex}
Trong dãy số tự nhiên từ 1000 đến 1010 có bao nhiêu số?
\shortans{11}
\loigiai{
$1010 - 1000 + 1 = 11$ số.
}
\end{ex}

\begin{ex}
Có bao nhiêu số có 4 chữ số tận cùng bằng 0 từ 1000 đến 1050?
\shortans{6}
\loigiai{
1000, 1010, 1020, 1030, 1040, 1050 (6 số).
}
\end{ex}

\begin{ex}
Số bé nhất có bốn chữ số giống nhau là số nào?
\shortans{1111}
\loigiai{
Số đó là 1111.
}
\end{ex}

\begin{ex}
Số lớn nhất có bốn chữ số giống nhau là số nào?
\shortans{9999}
\loigiai{
Số đó là 9999.
}
\end{ex}

\begin{ex}
Hiệu của số lớn nhất có 4 chữ số giống nhau và số bé nhất có 4 chữ số giống nhau là bao nhiêu?
\shortans{8888}
\loigiai{
$9999 - 1111 = 8888$.
}
\end{ex}

\newpage

\TaoPhieuBaiTap{Chủ đề 8 -- Các số đến 10 000}{47}{Làm quen với chữ số La Mã}{https://loptoanthaylinh.vn/bai-tap/toan-3/bai-47}
\begin{ex}
Số La Mã V biểu thị số tự nhiên nào?
\shortans{5}
\loigiai{
V là 5.
}
\end{ex}

\begin{ex}
Số La Mã X biểu thị số tự nhiên nào?
\shortans{10}
\loigiai{
X là 10.
}
\end{ex}

\begin{ex}
Số La Mã I biểu thị số tự nhiên nào?
\shortans{1}
\loigiai{
I là 1.
}
\end{ex}

\begin{ex}
Số La Mã IV biểu thị số tự nhiên nào?
\shortans{4}
\loigiai{
IV là 4.
}
\end{ex}

\begin{ex}
Số La Mã VI biểu thị số tự nhiên nào?
\shortans{6}
\loigiai{
VI là 6.
}
\end{ex}

\begin{ex}
Số La Mã IX biểu thị số tự nhiên nào?
\shortans{9}
\loigiai{
IX là 9.
}
\end{ex}

\begin{ex}
Số La Mã XI biểu thị số tự nhiên nào?
\shortans{11}
\loigiai{
XI là 11.
}
\end{ex}

\begin{ex}
Số La Mã XII biểu thị số tự nhiên nào?
\shortans{12}
\loigiai{
XII là 12.
}
\end{ex}

\begin{ex}
Số La Mã XIV biểu thị số tự nhiên nào?
\shortans{14}
\loigiai{
XIV là 14.
}
\end{ex}

\begin{ex}
Số La Mã XV biểu thị số tự nhiên nào?
\shortans{15}
\loigiai{
XV là 15.
}
\end{ex}

\begin{ex}
Số La Mã XIX biểu thị số tự nhiên nào?
\shortans{19}
\loigiai{
XIX là 19.
}
\end{ex}

\begin{ex}
Số La Mã XX biểu thị số tự nhiên nào?
\shortans{20}
\loigiai{
XX là 20.
}
\end{ex}

\begin{ex}
Số La Mã XXI biểu thị số tự nhiên nào?
\shortans{21}
\loigiai{
XXI là 21.
}
\end{ex}

\begin{ex}
Số La Mã VII biểu thị số tự nhiên nào?
\shortans{7}
\loigiai{
VII là 7.
}
\end{ex}

\begin{ex}
Số La Mã VIII biểu thị số tự nhiên nào?
\shortans{8}
\loigiai{
VIII là 8.
}
\end{ex}

\begin{ex}
Số La Mã XIII biểu thị số tự nhiên nào?
\shortans{13}
\loigiai{
XIII là 13.
}
\end{ex}

\begin{ex}
Số La Mã XVI biểu thị số tự nhiên nào?
\shortans{16}
\loigiai{
XVI là 16.
}
\end{ex}

\begin{ex}
Số La Mã XVII biểu thị số tự nhiên nào?
\shortans{17}
\loigiai{
XVII là 17.
}
\end{ex}

\begin{ex}
Số La Mã XVIII biểu thị số tự nhiên nào?
\shortans{18}
\loigiai{
XVIII là 18.
}
\end{ex}

\begin{ex}
Đồng hồ La Mã có kim ngắn chỉ vào số IX, kim dài chỉ vào số XII. Hỏi đồng hồ chỉ mấy giờ đúng?
\shortans{9}
\loigiai{
IX là 9, đồng hồ chỉ 9 giờ đúng.
}
\end{ex}

\begin{ex}
Đồng hồ La Mã có kim ngắn chỉ vào số IV, kim dài chỉ vào số XII. Đồng hồ chỉ mấy giờ đúng?
\shortans{4}
\loigiai{
IV là 4, đồng hồ chỉ 4 giờ đúng.
}
\end{ex}

\begin{ex}
Tính giá trị bằng số tự nhiên: V + VII = \dots
\shortans{12}
\loigiai{
$5 + 7 = 12$.
}
\end{ex}

\begin{ex}
Tính giá trị bằng số tự nhiên: X + IX = \dots
\shortans{19}
\loigiai{
$10 + 9 = 19$.
}
\end{ex}

\begin{ex}
Tính giá trị bằng số tự nhiên: XX - VIII = \dots
\shortans{12}
\loigiai{
$20 - 8 = 12$.
}
\end{ex}

\begin{ex}
Dùng que tính để xếp số La Mã VII cần bao nhiêu que tính?
\shortans{4}
\loigiai{
Chữ V cần 2 que, mỗi chữ I cần 1 que: $2 + 1 + 1 = 4$ que tính.
}
\end{ex}

\newpage

\TaoPhieuBaiTap{Chủ đề 8 -- Các số đến 10 000}{48}{Làm tròn số đến hàng chục, hàng trăm}{https://loptoanthaylinh.vn/bai-tap/toan-3/bai-48}
\begin{ex}
Làm tròn số 43 đến hàng chục ta được số nào?
\shortans{40}
\loigiai{
Chữ số hàng đơn vị là 3 < 5 nên làm tròn thành 40.
}
\end{ex}

\begin{ex}
Làm tròn số 78 đến hàng chục ta được số nào?
\shortans{80}
\loigiai{
Chữ số hàng đơn vị là 8 $\ge$ 5 nên làm tròn thành 80.
}
\end{ex}

\begin{ex}
Làm tròn số 65 đến hàng chục ta được số nào?
\shortans{70}
\loigiai{
Chữ số hàng đơn vị là 5 nên làm tròn lên 70.
}
\end{ex}

\begin{ex}
Làm tròn số 234 đến hàng chục ta được số nào?
\shortans{230}
\loigiai{
Chữ số hàng đơn vị là 4 nên làm tròn thành 230.
}
\end{ex}

\begin{ex}
Làm tròn số 586 đến hàng chục ta được số nào?
\shortans{590}
\loigiai{
Chữ số hàng đơn vị là 6 nên làm tròn lên 590.
}
\end{ex}

\begin{ex}
Làm tròn số 315 đến hàng chục ta được số nào?
\shortans{320}
\loigiai{
Chữ số hàng đơn vị là 5 nên làm tròn lên 320.
}
\end{ex}

\begin{ex}
Làm tròn số 320 đến hàng trăm ta được số nào?
\shortans{300}
\loigiai{
Chữ số hàng chục là 2 < 5 nên làm tròn thành 300.
}
\end{ex}

\begin{ex}
Làm tròn số 670 đến hàng trăm ta được số nào?
\shortans{700}
\loigiai{
Chữ số hàng chục là 7 $\ge$ 5 nên làm tròn lên 700.
}
\end{ex}

\begin{ex}
Làm tròn số 450 đến hàng trăm ta được số nào?
\shortans{500}
\loigiai{
Chữ số hàng chục là 5 nên làm tròn lên 500.
}
\end{ex}

\begin{ex}
Làm tròn số 1248 đến hàng trăm ta được số nào?
\shortans{1200}
\loigiai{
Chữ số hàng chục là 4 < 5 nên làm tròn thành 1200.
}
\end{ex}

\begin{ex}
Làm tròn số 3682 đến hàng trăm ta được số nào?
\shortans{3700}
\loigiai{
Chữ số hàng chục là 8 $\ge$ 5 nên làm tròn lên 3700.
}
\end{ex}

\begin{ex}
Làm tròn số 5450 đến hàng trăm ta được số nào?
\shortans{5500}
\loigiai{
Chữ số hàng chục là 5 nên làm tròn lên 5500.
}
\end{ex}

\begin{ex}
Làm tròn số 7819 đến hàng chục ta được số nào?
\shortans{7820}
\loigiai{
Hàng đơn vị là 9 nên làm tròn lên 7820.
}
\end{ex}

\begin{ex}
Làm tròn số 9004 đến hàng chục ta được số nào?
\shortans{9000}
\loigiai{
Hàng đơn vị là 4 nên làm tròn xuống 9000.
}
\end{ex}

\begin{ex}
Làm tròn số 8975 đến hàng trăm ta được số nào?
\shortans{9000}
\loigiai{
Hàng chục là 7 nên làm tròn lên: $8900 + 100 = 9000$.
}
\end{ex}

\begin{ex}
Có bao nhiêu số tự nhiên có hai chữ số khi làm tròn đến hàng chục thì được số 50?
\shortans{10}
\loigiai{
Các số từ 45 đến 54 (gồm 10 số).
}
\end{ex}

\begin{ex}
Số nhỏ nhất khi làm tròn đến hàng chục được số 80 là số nào?
\shortans{75}
\loigiai{
Số nhỏ nhất là 75.
}
\end{ex}

\begin{ex}
Số lớn nhất khi làm tròn đến hàng chục được số 80 là số nào?
\shortans{84}
\loigiai{
Số lớn nhất là 84.
}
\end{ex}

\begin{ex}
Số nhỏ nhất khi làm tròn đến hàng trăm được số 600 là số nào?
\shortans{550}
\loigiai{
Số nhỏ nhất là 550.
}
\end{ex}

\begin{ex}
Số lớn nhất khi làm tròn đến hàng trăm được số 600 là số nào?
\shortans{649}
\loigiai{
Số lớn nhất là 649.
}
\end{ex}

\begin{ex}
Quãng đường từ nhà An đến trường dài 782 m. Làm tròn đến hàng trăm thì khoảng bao nhiêu mét?
\shortans{800}
\loigiai{
Làm tròn 782 đến hàng trăm được 800 m.
}
\end{ex}

\begin{ex}
Một cửa hàng bán được 148 quyển vở. Làm tròn đến hàng chục là bao nhiêu quyển vở?
\shortans{150}
\loigiai{
Làm tròn 148 đến hàng chục được 150.
}
\end{ex}

\begin{ex}
Cân nặng một con bê là 85 kg. Làm tròn đến hàng chục là bao nhiêu ki-lô-gam?
\shortans{90}
\loigiai{
Làm tròn lên thành 90 kg.
}
\end{ex}

\begin{ex}
Làm tròn số 349 đến hàng trăm ta được số nào?
\shortans{300}
\loigiai{
Hàng chục là 4 < 5 nên làm tròn thành 300.
}
\end{ex}

\begin{ex}
Làm tròn số 351 đến hàng trăm ta được số nào?
\shortans{400}
\loigiai{
Hàng chục là 5 nên làm tròn lên 400.
}
\end{ex}

\newpage

\TaoPhieuBaiTap{Chủ đề 8 -- Các số đến 10 000}{49}{Luyện tập chung số có bốn chữ số}{https://loptoanthaylinh.vn/bai-tap/toan-3/bai-49}
\begin{ex}
Số 4725 gồm mấy nghìn?
\shortans{4}
\loigiai{
Chữ số hàng nghìn là 4.
}
\end{ex}

\begin{ex}
Số liền sau của 6999 là số nào?
\shortans{7000}
\loigiai{
$6999 + 1 = 7000$.
}
\end{ex}

\begin{ex}
Số liền trước của 8000 là số nào?
\shortans{7999}
\loigiai{
$8000 - 1 = 7999$.
}
\end{ex}

\begin{ex}
Trong các số $6450, 6540, 6405, 6504$, số lớn nhất là số nào?
\shortans{6540}
\loigiai{
Số lớn nhất là 6540.
}
\end{ex}

\begin{ex}
Số La Mã XVIII tương ứng với số tự nhiên nào?
\shortans{18}
\loigiai{
XVIII là 18.
}
\end{ex}

\begin{ex}
Số La Mã XIV tương ứng với số tự nhiên nào?
\shortans{14}
\loigiai{
XIV là 14.
}
\end{ex}

\begin{ex}
Làm tròn số 564 đến hàng chục được số nào?
\shortans{560}
\loigiai{
Làm tròn xuống 560.
}
\end{ex}

\begin{ex}
Làm tròn số 750 đến hàng trăm được số nào?
\shortans{800}
\loigiai{
Làm tròn lên 800.
}
\end{ex}

\begin{ex}
Viết số thích hợp: $7000 + 300 + 4 = \dots$
\shortans{7304}
\loigiai{
$7000 + 300 + 4 = 7304$.
}
\end{ex}

\begin{ex}
Tìm $x$ là số tròn nghìn biết $4000 < x < 6000$.
\shortans{5000}
\loigiai{
Số đó là 5000.
}
\end{ex}

\begin{ex}
Cho 4 chữ số 2, 0, 8, 9. Số lớn nhất có 4 chữ số ghép từ cả 4 chữ số này là số nào?
\shortans{9820}
\loigiai{
Số đó là 9820.
}
\end{ex}

\begin{ex}
Cho 4 chữ số 1, 0, 5, 8. Số bé nhất có 4 chữ số khác nhau ghép từ 4 chữ số này là số nào?
\shortans{1058}
\loigiai{
Số đó là 1058.
}
\end{ex}

\begin{ex}
Khoảng cách giữa hai số 9000 và 8700 là bao nhiêu?
\shortans{300}
\loigiai{
$9000 - 8700 = 300$.
}
\end{ex}

\begin{ex}
Đồng hồ La Mã chỉ vào số VIII nghĩa là mấy giờ?
\shortans{8}
\loigiai{
VIII là 8 giờ.
}
\end{ex}

\begin{ex}
Số 5555 có bao nhiêu chữ số 5?
\shortans{4}
\loigiai{
Có 4 chữ số 5.
}
\end{ex}

\begin{ex}
Làm tròn số 8263 đến hàng trăm ta được số nào?
\shortans{8300}
\loigiai{
Hàng chục là 6 nên làm tròn lên 8300.
}
\end{ex}

\begin{ex}
Làm tròn số 4096 đến hàng chục ta được số nào?
\shortans{4100}
\loigiai{
Hàng đơn vị là 6 làm tròn lên: $4090 + 10 = 4100$.
}
\end{ex}

\begin{ex}
Số tự nhiên nhỏ nhất có 4 chữ số khác nhau là số nào?
\shortans{1023}
\loigiai{
Số đó là 1023.
}
\end{ex}

\begin{ex}
Số tự nhiên lớn nhất có 4 chữ số khác nhau là số nào?
\shortans{9876}
\loigiai{
Số đó là 9876.
}
\end{ex}

\begin{ex}
Tính: XXI - IX = \dots (kết quả bằng số tự nhiên).
\shortans{12}
\loigiai{
$21 - 9 = 12$.
}
\end{ex}

\begin{ex}
Tính: VI $\times$ III = \dots (kết quả bằng số tự nhiên).
\shortans{18}
\loigiai{
$6 \times 3 = 18$.
}
\end{ex}

\begin{ex}
Số lớn nhất có 4 chữ số có tích các chữ số bằng 1 là số nào?
\shortans{1111}
\loigiai{
Tất cả các chữ số phải là 1: 1111.
}
\end{ex}

\begin{ex}
Có bao nhiêu số tròn trăm từ 2000 đến 3000?
\shortans{11}
\loigiai{
2000, 2100, ..., 3000 (11 số).
}
\end{ex}

\begin{ex}
Làm tròn số 649 đến hàng trăm được số nào?
\shortans{600}
\loigiai{
Hàng chục là 4 nên làm tròn thành 600.
}
\end{ex}

\begin{ex}
Làm tròn số 650 đến hàng trăm được số nào?
\shortans{700}
\loigiai{
Hàng chục là 5 nên làm tròn lên 700.
}
\end{ex}

\newpage


% =========================================================================
% Chu vi, diện tích một số hình phẳng
% =========================================================================
\section{Chu vi, diện tích một số hình phẳng}

\TaoPhieuBaiTap{Chủ đề 9 -- Chu vi, diện tích một số hình phẳng}{50}{Chu vi hình tam giác, tứ giác, hình chữ nhật, hình vuông}{https://loptoanthaylinh.vn/bai-tap/toan-3/bai-50}
\begin{ex}
Một hình tam giác có độ dài các cạnh là 15 cm, 20 cm, 25 cm. Chu vi hình tam giác đó là bao nhiêu xăng-ti-mét?
\shortans{60}
\loigiai{
$15 + 20 + 25 = 60$ cm.
}
\end{ex}

\begin{ex}
Một hình tam giác đều có cạnh 14 cm. Chu vi của nó là bao nhiêu xăng-ti-mét?
\shortans{42}
\loigiai{
$14 \times 3 = 42$ cm.
}
\end{ex}

\begin{ex}
Một hình tứ giác có độ dài 4 cạnh là 10 cm, 12 cm, 15 cm, 18 cm. Chu vi hình tứ giác đó là bao nhiêu xăng-ti-mét?
\shortans{55}
\loigiai{
$10 + 12 + 15 + 18 = 55$ cm.
}
\end{ex}

\begin{ex}
Một hình chữ nhật có chiều dài 25 cm, chiều rộng 15 cm. Chu vi của hình chữ nhật đó là bao nhiêu xăng-ti-mét?
\shortans{80}
\loigiai{
$(25 + 15) \times 2 = 80$ cm.
}
\end{ex}

\begin{ex}
Một hình chữ nhật có chiều dài 30 cm, chiều rộng 20 cm. Nửa chu vi hình chữ nhật đó là bao nhiêu xăng-ti-mét?
\shortans{50}
\loigiai{
$30 + 20 = 50$ cm.
}
\end{ex}

\begin{ex}
Một hình vuông có cạnh dài 15 cm. Chu vi hình vuông đó là bao nhiêu xăng-ti-mét?
\shortans{60}
\loigiai{
$15 \times 4 = 60$ cm.
}
\end{ex}

\begin{ex}
Một hình vuông có cạnh dài 25 cm. Chu vi hình vuông đó là bao nhiêu xăng-ti-mét?
\shortans{100}
\loigiai{
$25 \times 4 = 100$ cm.
}
\end{ex}

\begin{ex}
Một hình vuông có chu vi 48 cm. Cạnh của hình vuông đó dài bao nhiêu xăng-ti-mét?
\shortans{12}
\loigiai{
$48 : 4 = 12$ cm.
}
\end{ex}

\begin{ex}
Một hình vuông có chu vi 80 cm. Cạnh của hình vuông đó dài bao nhiêu xăng-ti-mét?
\shortans{20}
\loigiai{
$80 : 4 = 20$ cm.
}
\end{ex}

\begin{ex}
Một hình chữ nhật có chu vi 64 cm, chiều rộng là 12 cm. Chiều dài hình chữ nhật đó là bao nhiêu xăng-ti-mét?
\shortans{20}
\loigiai{
Nửa chu vi = $64 : 2 = 32$ cm. Chiều dài = $32 - 12 = 20$ cm.
}
\end{ex}

\begin{ex}
Một hình chữ nhật có nửa chu vi là 40 cm, chiều dài là 25 cm. Chiều rộng hình chữ nhật đó là bao nhiêu xăng-ti-mét?
\shortans{15}
\loigiai{
$40 - 25 = 15$ cm.
}
\end{ex}

\begin{ex}
Một khu đất hình vuông có cạnh dài 50 m. Chu vi khu đất đó là bao nhiêu mét?
\shortans{200}
\loigiai{
$50 \times 4 = 200$ m.
}
\end{ex}

\begin{ex}
Một sân trường hình chữ nhật có chiều dài 60 m, chiều rộng 40 m. Chu vi sân trường đó là bao nhiêu mét?
\shortans{200}
\loigiai{
$(60 + 40) \times 2 = 200$ m.
}
\end{ex}

\begin{ex}
Uốn một đoạn dây thép thành một hình tam giác có 3 cạnh bằng nhau, mỗi cạnh dài 8 cm. Đoạn dây thép dài bao nhiêu xăng-ti-mét?
\shortans{24}
\loigiai{
$8 \times 3 = 24$ cm.
}
\end{ex}

\begin{ex}
Uốn một đoạn dây thép dài 36 cm thành một hình vuông. Cạnh của hình vuông đó dài bao nhiêu xăng-ti-mét?
\shortans{9}
\loigiai{
$36 : 4 = 9$ cm.
}
\end{ex}

\begin{ex}
Một hình chữ nhật có chiều rộng 8 cm, chiều dài gấp 3 lần chiều rộng. Chu vi hình chữ nhật đó là bao nhiêu xăng-ti-mét?
\shortans{64}
\loigiai{
Chiều dài = 24 cm. Chu vi = $(24 + 8) \times 2 = 64$ cm.
}
\end{ex}

\begin{ex}
Một hình chữ nhật có chiều dài 20 cm, chiều rộng bằng nửa chiều dài. Chu vi hình chữ nhật đó là bao nhiêu xăng-ti-mét?
\shortans{60}
\loigiai{
Chiều rộng = 10 cm. Chu vi = $(20 + 10) \times 2 = 60$ cm.
}
\end{ex}

\begin{ex}
Ghép hai hình vuông cạnh 5 cm lại thành một hình chữ nhật. Chu vi hình chữ nhật tạo thành là bao nhiêu xăng-ti-mét?
\shortans{30}
\loigiai{
Chiều dài = 10 cm, rộng = 5 cm. Chu vi = $(10 + 5) \times 2 = 30$ cm.
}
\end{ex}

\begin{ex}
Ghép ba hình vuông cạnh 4 cm liền nhau thành một hình chữ nhật. Chu vi hình chữ nhật đó là bao nhiêu xăng-ti-mét?
\shortans{32}
\loigiai{
Dài = 12 cm, rộng = 4 cm. Chu vi = $(12 + 4) \times 2 = 32$ cm.
}
\end{ex}

\begin{ex}
Một mảnh vườn hình chữ nhật có chu vi 100 m. Nửa chu vi mảnh vườn đó là bao nhiêu mét?
\shortans{50}
\loigiai{
$100 : 2 = 50$ m.
}
\end{ex}

\begin{ex}
Một hình tam giác có chu vi 45 cm, biết hai cạnh dài 12 cm và 18 cm. Cạnh thứ ba dài bao nhiêu xăng-ti-mét?
\shortans{15}
\loigiai{
$45 - (12 + 18) = 15$ cm.
}
\end{ex}

\begin{ex}
Một tứ giác có 4 cạnh bằng nhau và chu vi bằng 44 cm. Độ dài mỗi cạnh là bao nhiêu xăng-ti-mét?
\shortans{11}
\loigiai{
$44 : 4 = 11$ cm.
}
\end{ex}

\begin{ex}
Một khung ảnh hình vuông cạnh 18 cm. Bạn Mai dán nẹp xung quanh khung ảnh hết bao nhiêu xăng-ti-mét nẹp?
\shortans{72}
\loigiai{
$18 \times 4 = 72$ cm.
}
\end{ex}

\begin{ex}
Nếu cạnh của một hình vuông tăng lên 2 cm thì chu vi hình vuông đó tăng thêm bao nhiêu xăng-ti-mét?
\shortans{8}
\loigiai{
$2 \times 4 = 8$ cm.
}
\end{ex}

\begin{ex}
Nếu cạnh hình vuông gấp lên 3 lần thì chu vi của nó gấp lên mấy lần?
\shortans{3}
\loigiai{
Chu vi gấp lên 3 lần.
}
\end{ex}

\newpage

\TaoPhieuBaiTap{Chủ đề 9 -- Chu vi, diện tích một số hình phẳng}{51}{Diện tích của một hình. Xăng-ti-mét vuông}{https://loptoanthaylinh.vn/bai-tap/toan-3/bai-51}
\begin{ex}
Đơn vị đo diện tích xăng-ti-mét vuông viết tắt là gì? (Viết số: $\text{cm}^2$ có số mũ là bao nhiêu?).
\shortans{2}
\loigiai{
Số mũ là 2.
}
\end{ex}

\begin{ex}
Một hình vuông có cạnh dài 1 cm. Diện tích của hình vuông đó là bao nhiêu xăng-ti-mét vuông?
\shortans{1}
\loigiai{
$1 \times 1 = 1\text{ cm}^2$.
}
\end{ex}

\begin{ex}
Một hình chữ nhật được ghép bởi 6 ô vuông nhỏ, mỗi ô vuông có diện tích $1\text{ cm}^2$. Diện tích hình chữ nhật đó là bao nhiêu xăng-ti-mét vuông?
\shortans{6}
\loigiai{
$6 \times 1 = 6\text{ cm}^2$.
}
\end{ex}

\begin{ex}
Một tấm bìa được tạo bởi 15 ô vuông nhỏ diện tích $1\text{ cm}^2$. Diện tích tấm bìa là bao nhiêu xăng-ti-mét vuông?
\shortans{15}
\loigiai{
$15\text{ cm}^2$.
}
\end{ex}

\begin{ex}
Tính: $25\text{ cm}^2 + 35\text{ cm}^2 = \dots\text{ cm}^2$.
\shortans{60}
\loigiai{
$25 + 35 = 60$.
}
\end{ex}

\begin{ex}
Tính: $90\text{ cm}^2 - 42\text{ cm}^2 = \dots\text{ cm}^2$.
\shortans{48}
\loigiai{
$90 - 42 = 48$.
}
\end{ex}

\begin{ex}
Tính: $15\text{ cm}^2 \times 4 = \dots\text{ cm}^2$.
\shortans{60}
\loigiai{
$15 \times 4 = 60$.
}
\end{ex}

\begin{ex}
Tính: $84\text{ cm}^2 : 4 = \dots\text{ cm}^2$.
\shortans{21}
\loigiai{
$84 : 4 = 21$.
}
\end{ex}

\begin{ex}
Một miếng bìa diện tích $45\text{ cm}^2$, người ta cắt đi $18\text{ cm}^2$. Miếng bìa còn lại bao nhiêu xăng-ti-mét vuông?
\shortans{27}
\loigiai{
$45 - 18 = 27\text{ cm}^2$.
}
\end{ex}

\begin{ex}
Một hình gồm 8 ô vuông, mỗi ô vuông có diện tích $2\text{ cm}^2$. Diện tích cả hình là bao nhiêu xăng-ti-mét vuông?
\shortans{16}
\loigiai{
$8 \times 2 = 16\text{ cm}^2$.
}
\end{ex}

\begin{ex}
Tính: $30\text{ cm}^2 + 20\text{ cm}^2 \times 2$.
\shortans{70}
\loigiai{
$30 + 40 = 70$.
}
\end{ex}

\begin{ex}
Tính: $(50\text{ cm}^2 - 20\text{ cm}^2) : 3$.
\shortans{10}
\loigiai{
$30 : 3 = 10$.
}
\end{ex}

\begin{ex}
Một hình chữ nhật có diện tích $24\text{ cm}^2$ được chia thành 4 phần bằng nhau. Mỗi phần có diện tích bao nhiêu xăng-ti-mét vuông?
\shortans{6}
\loigiai{
$24 : 4 = 6\text{ cm}^2$.
}
\end{ex}

\begin{ex}
Một hình tam giác có diện tích $18\text{ cm}^2$, một hình tứ giác có diện tích $25\text{ cm}^2$. Tổng diện tích hai hình là bao nhiêu xăng-ti-mét vuông?
\shortans{43}
\loigiai{
$18 + 25 = 43\text{ cm}^2$.
}
\end{ex}

\begin{ex}
Hình A có diện tích $36\text{ cm}^2$, hình B có diện tích $12\text{ cm}^2$. Diện tích hình A gấp mấy lần diện tích hình B?
\shortans{3}
\loigiai{
$36 : 12 = 3$ lần.
}
\end{ex}

\begin{ex}
Hình M có diện tích $40\text{ cm}^2$, hình N có diện tích kém hình M là $15\text{ cm}^2$. Diện tích hình N là bao nhiêu xăng-ti-mét vuông?
\shortans{25}
\loigiai{
$40 - 15 = 25\text{ cm}^2$.
}
\end{ex}

\begin{ex}
Tính: $12\text{ cm}^2 \times 5 = \dots\text{ cm}^2$.
\shortans{60}
\loigiai{
$12 \times 5 = 60$.
}
\end{ex}

\begin{ex}
Tính: $100\text{ cm}^2 : 5 = \dots\text{ cm}^2$.
\shortans{20}
\loigiai{
$100 : 5 = 20$.
}
\end{ex}

\begin{ex}
Một tờ giấy màu đỏ diện tích $50\text{ cm}^2$, tờ giấy màu xanh diện tích $70\text{ cm}^2$. Cả hai tờ giấy có diện tích bao nhiêu xăng-ti-mét vuông?
\shortans{120}
\loigiai{
$50 + 70 = 120\text{ cm}^2$.
}
\end{ex}

\begin{ex}
Một con tem hình chữ nhật có diện tích $6\text{ cm}^2$. Hỏi 5 con tem như thế có tổng diện tích bao nhiêu xăng-ti-mét vuông?
\shortans{30}
\loigiai{
$6 \times 5 = 30\text{ cm}^2$.
}
\end{ex}

\begin{ex}
Tấm bìa hình vuông gồm 16 ô vuông nhỏ mỗi ô diện tích $1\text{ cm}^2$. Hỏi diện tích tấm bìa là bao nhiêu xăng-ti-mét vuông?
\shortans{16}
\loigiai{
$16 \times 1 = 16\text{ cm}^2$.
}
\end{ex}

\begin{ex}
Tính: $14\text{ cm}^2 + 36\text{ cm}^2 : 4$.
\shortans{23}
\loigiai{
$14 + 9 = 23$.
}
\end{ex}

\begin{ex}
Tính: $60\text{ cm}^2 - 15\text{ cm}^2 \times 2$.
\shortans{30}
\loigiai{
$60 - 30 = 30$.
}
\end{ex}

\begin{ex}
Một miếng dán sticker có diện tích $8\text{ cm}^2$. Một trang vở dán được 9 miếng như vậy. Tổng diện tích các miếng sticker là bao nhiêu xăng-ti-mét vuông?
\shortans{72}
\loigiai{
$8 \times 9 = 72\text{ cm}^2$.
}
\end{ex}

\begin{ex}
Hình vuông cạnh 1 cm thì chu vi là 4 cm, còn diện tích là bao nhiêu xăng-ti-mét vuông?
\shortans{1}
\loigiai{
$1 \times 1 = 1\text{ cm}^2$.
}
\end{ex}

\newpage

\TaoPhieuBaiTap{Chủ đề 9 -- Chu vi, diện tích một số hình phẳng}{52}{Diện tích hình chữ nhật, diện tích hình vuông}{https://loptoanthaylinh.vn/bai-tap/toan-3/bai-52}
\begin{ex}
Một hình chữ nhật có chiều dài 8 cm, chiều rộng 5 cm. Tính diện tích hình chữ nhật đó theo xăng-ti-mét vuông.
\shortans{40}
\loigiai{
$8 \times 5 = 40\text{ cm}^2$.
}
\end{ex}

\begin{ex}
Một hình chữ nhật có chiều dài 9 cm, chiều rộng 6 cm. Tính diện tích hình chữ nhật đó theo xăng-ti-mét vuông.
\shortans{54}
\loigiai{
$9 \times 6 = 54\text{ cm}^2$.
}
\end{ex}

\begin{ex}
Một hình chữ nhật có chiều dài 12 cm, chiều rộng 4 cm. Diện tích của hình chữ nhật đó là bao nhiêu xăng-ti-mét vuông?
\shortans{48}
\loigiai{
$12 \times 4 = 48\text{ cm}^2$.
}
\end{ex}

\begin{ex}
Một hình vuông có cạnh dài 6 cm. Tính diện tích hình vuông đó theo xăng-ti-mét vuông.
\shortans{36}
\loigiai{
$6 \times 6 = 36\text{ cm}^2$.
}
\end{ex}

\begin{ex}
Một hình vuông có cạnh dài 8 cm. Tính diện tích hình vuông đó theo xăng-ti-mét vuông.
\shortans{64}
\loigiai{
$8 \times 8 = 64\text{ cm}^2$.
}
\end{ex}

\begin{ex}
Một hình vuông có cạnh dài 9 cm. Tính diện tích hình vuông đó theo xăng-ti-mét vuông.
\shortans{81}
\loigiai{
$9 \times 9 = 81\text{ cm}^2$.
}
\end{ex}

\begin{ex}
Một hình chữ nhật có diện tích $35\text{ cm}^2$, chiều dài là 7 cm. Chiều rộng của hình chữ nhật đó là bao nhiêu xăng-ti-mét?
\shortans{5}
\loigiai{
$35 : 7 = 5$ cm.
}
\end{ex}

\begin{ex}
Một hình chữ nhật có diện tích $48\text{ cm}^2$, chiều rộng là 6 cm. Chiều dài của hình chữ nhật đó là bao nhiêu xăng-ti-mét?
\shortans{8}
\loigiai{
$48 : 6 = 8$ cm.
}
\end{ex}

\begin{ex}
Một hình vuông có chu vi 24 cm. Cạnh của hình vuông đó dài bao nhiêu xăng-ti-mét?
\shortans{6}
\loigiai{
$24 : 4 = 6$ cm.
}
\end{ex}

\begin{ex}
Một hình vuông có chu vi 28 cm. Tính diện tích hình vuông đó theo xăng-ti-mét vuông.
\shortans{49}
\loigiai{
Cạnh = $28 : 4 = 7$ cm. Diện tích = $7 \times 7 = 49\text{ cm}^2$.
}
\end{ex}

\begin{ex}
Một hình vuông có chu vi 36 cm. Diện tích hình vuông đó là bao nhiêu xăng-ti-mét vuông?
\shortans{81}
\loigiai{
Cạnh = $36 : 4 = 9$ cm. Diện tích = $9 \times 9 = 81\text{ cm}^2$.
}
\end{ex}

\begin{ex}
Một hình vuông có chu vi 20 cm. Diện tích hình vuông đó là bao nhiêu xăng-ti-mét vuông?
\shortans{25}
\loigiai{
Cạnh = $20 : 4 = 5$ cm. Diện tích = $5 \times 5 = 25\text{ cm}^2$.
}
\end{ex}

\begin{ex}
Một hình chữ nhật có chiều rộng 4 cm, chiều dài gấp 3 lần chiều rộng. Diện tích hình chữ nhật đó là bao nhiêu xăng-ti-mét vuông?
\shortans{48}
\loigiai{
Chiều dài = $4 \times 3 = 12$ cm. Diện tích = $12 \times 4 = 48\text{ cm}^2$.
}
\end{ex}

\begin{ex}
Một hình chữ nhật có nửa chu vi là 15 cm, chiều dài là 9 cm. Diện tích hình chữ nhật đó là bao nhiêu xăng-ti-mét vuông?
\shortans{54}
\loigiai{
Chiều rộng = $15 - 9 = 6$ cm. Diện tích = $9 \times 6 = 54\text{ cm}^2$.
}
\end{ex}

\begin{ex}
Một tấm thiệp hình chữ nhật có chiều dài 15 cm, chiều rộng 10 cm. Diện tích tấm thiệp là bao nhiêu xăng-ti-mét vuông?
\shortans{150}
\loigiai{
$15 \times 10 = 150\text{ cm}^2$.
}
\end{ex}

\begin{ex}
Một viên gạch men hình vuông có cạnh dài 20 cm. Diện tích viên gạch men là bao nhiêu xăng-ti-mét vuông?
\shortans{400}
\loigiai{
$20 \times 20 = 400\text{ cm}^2$.
}
\end{ex}

\begin{ex}
Một hình vuông có diện tích $16\text{ cm}^2$. Độ dài cạnh của hình vuông đó là bao nhiêu xăng-ti-mét?
\shortans{4}
\loigiai{
Vì $4 \times 4 = 16$ nên cạnh bằng 4 cm.
}
\end{ex}

\begin{ex}
Một hình vuông có diện tích $100\text{ cm}^2$. Cạnh hình vuông dài bao nhiêu xăng-ti-mét?
\shortans{10}
\loigiai{
Vì $10 \times 10 = 100$ nên cạnh bằng 10 cm.
}
\end{ex}

\begin{ex}
Một hình chữ nhật có chu vi 30 cm, chiều dài 10 cm. Diện tích hình chữ nhật đó là bao nhiêu xăng-ti-mét vuông?
\shortans{50}
\loigiai{
Nửa chu vi = 15 cm; rộng = $15 - 10 = 5$ cm. Diện tích = $10 \times 5 = 50\text{ cm}^2$.
}
\end{ex}

\begin{ex}
Một hình chữ nhật có chiều dài 14 cm, chiều rộng bằng nửa chiều dài. Diện tích hình chữ nhật đó là bao nhiêu xăng-ti-mét vuông?
\shortans{98}
\loigiai{
Chiều rộng = 7 cm. Diện tích = $14 \times 7 = 98\text{ cm}^2$.
}
\end{ex}

\begin{ex}
Ghép hai hình vuông cạnh 4 cm thành một hình chữ nhật. Diện tích hình chữ nhật tạo thành là bao nhiêu xăng-ti-mét vuông?
\shortans{32}
\loigiai{
Diện tích 1 hình vuông = $4 \times 4 = 16\text{ cm}^2$. Hai hình = $16 \times 2 = 32\text{ cm}^2$.
}
\end{ex}

\begin{ex}
Một miếng bìa hình chữ nhật dài 18 cm, rộng 5 cm. Người ta cắt ra một hình vuông cạnh 5 cm. Diện tích phần bìa còn lại là bao nhiêu xăng-ti-mét vuông?
\shortans{65}
\loigiai{
Diện tích ban đầu = $18 \times 5 = 90\text{ cm}^2$. Cắt đi = $5 \times 5 = 25\text{ cm}^2$. Còn lại: $90 - 25 = 65\text{ cm}^2$.
}
\end{ex}

\begin{ex}
Nếu gấp cạnh của hình vuông lên 2 lần thì diện tích hình vuông gấp lên mấy lần?
\shortans{4}
\loigiai{
Diện tích tăng lên $2 \times 2 = 4$ lần.
}
\end{ex}

\begin{ex}
Một khung ảnh hình vuông cạnh 12 cm. Diện tích mặt kính của khung ảnh là bao nhiêu xăng-ti-mét vuông?
\shortans{144}
\loigiai{
$12 \times 12 = 144\text{ cm}^2$.
}
\end{ex}

\begin{ex}
Tính diện tích hình chữ nhật có chiều dài 20 cm và chiều rộng 7 cm theo xăng-ti-mét vuông.
\shortans{140}
\loigiai{
$20 \times 7 = 140\text{ cm}^2$.
}
\end{ex}

\newpage

\TaoPhieuBaiTap{Chủ đề 9 -- Chu vi, diện tích một số hình phẳng}{53}{Luyện tập chung chu vi và diện tích}{https://loptoanthaylinh.vn/bai-tap/toan-3/bai-53}
\begin{ex}
Một hình vuông có cạnh 7 cm. Chu vi hình vuông đó là bao nhiêu xăng-ti-mét?
\shortans{28}
\loigiai{
$7 \times 4 = 28$ cm.
}
\end{ex}

\begin{ex}
Một hình vuông có cạnh 7 cm. Diện tích hình vuông đó là bao nhiêu xăng-ti-mét vuông?
\shortans{49}
\loigiai{
$7 \times 7 = 49\text{ cm}^2$.
}
\end{ex}

\begin{ex}
Một hình chữ nhật có chiều dài 11 cm, chiều rộng 6 cm. Chu vi của nó là bao nhiêu xăng-ti-mét?
\shortans{34}
\loigiai{
$(11 + 6) \times 2 = 34$ cm.
}
\end{ex}

\begin{ex}
Một hình chữ nhật có chiều dài 11 cm, chiều rộng 6 cm. Diện tích của nó là bao nhiêu xăng-ti-mét vuông?
\shortans{66}
\loigiai{
$11 \times 6 = 66\text{ cm}^2$.
}
\end{ex}

\begin{ex}
Một hình tam giác đều có chu vi 36 cm. Cạnh của hình tam giác đó dài bao nhiêu xăng-ti-mét?
\shortans{12}
\loigiai{
$36 : 3 = 12$ cm.
}
\end{ex}

\begin{ex}
Một hình tứ giác có 4 cạnh đều bằng 9 cm. Chu vi tứ giác đó là bao nhiêu xăng-ti-mét?
\shortans{36}
\loigiai{
$9 \times 4 = 36$ cm.
}
\end{ex}

\begin{ex}
Một hình vuông có diện tích bằng diện tích hình chữ nhật dài 8 cm, rộng 2 cm. Cạnh của hình vuông đó là bao nhiêu xăng-ti-mét?
\shortans{4}
\loigiai{
Diện tích chữ nhật = $8 \times 2 = 16\text{ cm}^2$. Cạnh hình vuông = 4 cm vì $4 \times 4 = 16$.
}
\end{ex}

\begin{ex}
Một hình vuông có chu vi bằng chu vi hình chữ nhật dài 10 cm, rộng 6 cm. Cạnh hình vuông đó là bao nhiêu xăng-ti-mét?
\shortans{8}
\loigiai{
Chu vi chữ nhật = $(10 + 6) \times 2 = 32$ cm. Cạnh hình vuông = $32 : 4 = 8$ cm.
}
\end{ex}

\begin{ex}
Tính: $45\text{ cm}^2 + 55\text{ cm}^2 = \dots\text{ cm}^2$.
\shortans{100}
\loigiai{
$45 + 55 = 100$.
}
\end{ex}

\begin{ex}
Tính: $120\text{ cm}^2 - 45\text{ cm}^2 = \dots\text{ cm}^2$.
\shortans{75}
\loigiai{
$120 - 45 = 75$.
}
\end{ex}

\begin{ex}
Tính: $25\text{ cm}^2 \times 3 = \dots\text{ cm}^2$.
\shortans{75}
\loigiai{
$25 \times 3 = 75$.
}
\end{ex}

\begin{ex}
Tính: $96\text{ cm}^2 : 4 = \dots\text{ cm}^2$.
\shortans{24}
\loigiai{
$96 : 4 = 24$.
}
\end{ex}

\begin{ex}
Một tấm thảm hình vuông có cạnh 3 m. Chu vi tấm thảm đó là bao nhiêu mét?
\shortans{12}
\loigiai{
$3 \times 4 = 12$ m.
}
\end{ex}

\begin{ex}
Một sân bóng chuyền hình chữ nhật dài 18 m, rộng 9 m. Chu vi sân bóng chuyền là bao nhiêu mét?
\shortans{54}
\loigiai{
$(18 + 9) \times 2 = 54$ m.
}
\end{ex}

\begin{ex}
Một hình chữ nhật có chu vi 40 cm, chiều rộng là 8 cm. Diện tích hình chữ nhật đó là bao nhiêu xăng-ti-mét vuông?
\shortans{96}
\loigiai{
Nửa chu vi = 20 cm; dài = $20 - 8 = 12$ cm. Diện tích = $12 \times 8 = 96\text{ cm}^2$.
}
\end{ex}

\begin{ex}
Một hình vuông có chu vi 44 cm. Cạnh của nó dài bao nhiêu xăng-ti-mét?
\shortans{11}
\loigiai{
$44 : 4 = 11$ cm.
}
\end{ex}

\begin{ex}
Một hình vuông có cạnh 11 cm. Diện tích của nó là bao nhiêu xăng-ti-mét vuông?
\shortans{121}
\loigiai{
$11 \times 11 = 121\text{ cm}^2$.
}
\end{ex}

\begin{ex}
Nửa chu vi hình chữ nhật là 24 cm, chiều dài gấp 3 lần chiều rộng. Chiều rộng hình chữ nhật đó là bao nhiêu xăng-ti-mét?
\shortans{6}
\loigiai{
Tổng số phần: $1 + 3 = 4$ phần. Chiều rộng = $24 : 4 = 6$ cm.
}
\end{ex}

\begin{ex}
Một hình chữ nhật có diện tích $72\text{ cm}^2$, chiều dài 9 cm. Chiều rộng hình chữ nhật đó là bao nhiêu xăng-ti-mét?
\shortans{8}
\loigiai{
$72 : 9 = 8$ cm.
}
\end{ex}

\begin{ex}
Cạnh của hình vuông là 5 cm. Nếu tăng cạnh thêm 3 cm thì diện tích tăng thêm bao nhiêu xăng-ti-mét vuông?
\shortans{39}
\loigiai{
Cạnh mới = 8 cm. Diện tích mới = $8 \times 8 = 64\text{ cm}^2$. Tăng thêm: $64 - 25 = 39\text{ cm}^2$.
}
\end{ex}

\begin{ex}
Một bức tranh hình vuông có diện tích $81\text{ cm}^2$. Chu vi bức tranh đó là bao nhiêu xăng-ti-mét?
\shortans{36}
\loigiai{
Cạnh = 9 cm. Chu vi = $9 \times 4 = 36$ cm.
}
\end{ex}

\begin{ex}
Một khu đất hình chữ nhật có chiều dài 25 m, chiều rộng 12 m. Chu vi khu đất đó là bao nhiêu mét?
\shortans{74}
\loigiai{
$(25 + 12) \times 2 = 74$ m.
}
\end{ex}

\begin{ex}
Một hình chữ nhật có nửa chu vi là 18 cm, chiều dài hơn chiều rộng 4 cm. Chiều dài là bao nhiêu xăng-ti-mét?
\shortans{11}
\loigiai{
Chiều dài = $(18 + 4) : 2 = 11$ cm.
}
\end{ex}

\begin{ex}
Một hình chữ nhật có nửa chu vi là 18 cm, chiều dài hơn chiều rộng 4 cm. Diện tích của nó là bao nhiêu xăng-ti-mét vuông?
\shortans{77}
\loigiai{
Dài = 11 cm, rộng = 7 cm. Diện tích = $11 \times 7 = 77\text{ cm}^2$.
}
\end{ex}

\begin{ex}
Tính: $15\text{ cm}^2 \times 6 = \dots\text{ cm}^2$.
\shortans{90}
\loigiai{
$15 \times 6 = 90$.
}
\end{ex}

\newpage


% =========================================================================
% Cộng, trừ, nhân, chia trong phạm vi 10 000
% =========================================================================
\section{Cộng, trừ, nhân, chia trong phạm vi 10 000}

\TaoPhieuBaiTap{Chủ đề 10 -- Cộng, trừ, nhân, chia trong phạm vi 10 000}{54}{Phép cộng trong phạm vi 10 000}{https://loptoanthaylinh.vn/bai-tap/toan-3/bai-54}
\begin{ex}
Tính: $3524 + 2143 = \dots$
\shortans{5667}
\loigiai{
$3524 + 2143 = 5667$.
}
\end{ex}

\begin{ex}
Tính: $4816 + 2153 = \dots$
\shortans{6969}
\loigiai{
$4816 + 2153 = 6969$.
}
\end{ex}

\begin{ex}
Tính: $5248 + 3619 = \dots$
\shortans{8867}
\loigiai{
$5248 + 3619 = 8867$.
}
\end{ex}

\begin{ex}
Tính: $2745 + 4538 = \dots$
\shortans{7283}
\loigiai{
$2745 + 4538 = 7283$.
}
\end{ex}

\begin{ex}
Tính: $6382 + 1849 = \dots$
\shortans{8231}
\loigiai{
$6382 + 1849 = 8231$.
}
\end{ex}

\begin{ex}
Tính: $4500 + 3500 = \dots$
\shortans{8000}
\loigiai{
$4500 + 3500 = 8000$.
}
\end{ex}

\begin{ex}
Tính: $7200 + 1800 = \dots$
\shortans{9000}
\loigiai{
$7200 + 1800 = 9000$.
}
\end{ex}

\begin{ex}
Tính: $3860 + 2140 = \dots$
\shortans{6000}
\loigiai{
$3860 + 2140 = 6000$.
}
\end{ex}

\begin{ex}
Tính: $5412 + 735 = \dots$
\shortans{6147}
\loigiai{
$5412 + 735 = 6147$.
}
\end{ex}

\begin{ex}
Tính: $6084 + 879 = \dots$
\shortans{6963}
\loigiai{
$6084 + 879 = 6963$.
}
\end{ex}

\begin{ex}
Tính giá trị biểu thức: $3400 + 1600 + 2500$.
\shortans{7500}
\loigiai{
$5000 + 2500 = 7500$.
}
\end{ex}

\begin{ex}
Tính giá trị biểu thức: $4215 + 1380 + 405$.
\shortans{6000}
\loigiai{
$4215 + 1380 = 5595$; $5595 + 405 = 6000$.
}
\end{ex}

\begin{ex}
Một đội công nhân tuần thứ nhất trải được 2150 m đường, tuần thứ hai trải được 2450 m đường. Cả hai tuần đội đó trải được bao nhiêu mét đường?
\shortans{4600}
\loigiai{
$2150 + 2450 = 4600$ m.
}
\end{ex}

\begin{ex}
Kho thứ nhất có 3650 kg thóc, kho thứ hai có 4280 kg thóc. Cả hai kho có tất cả bao nhiêu ki-lô-gam thóc?
\shortans{7930}
\loigiai{
$3650 + 4280 = 7930$ kg.
}
\end{ex}

\begin{ex}
Trường Tiểu học Kim Đồng có 1240 học sinh nam và 1180 học sinh nữ. Trường có tất cả bao nhiêu học sinh?
\shortans{2420}
\loigiai{
$1240 + 1180 = 2420$ học sinh.
}
\end{ex}

\begin{ex}
Điền số thích hợp: $x - 3400 = 2600$.
\shortans{6000}
\loigiai{
$x = 2600 + 3400 = 6000$.
}
\end{ex}

\begin{ex}
Điền số thích hợp: $x - 1850 = 3150$.
\shortans{5000}
\loigiai{
$x = 3150 + 1850 = 5000$.
}
\end{ex}

\begin{ex}
Tìm tổng của số lớn nhất có bốn chữ số và số bé nhất có bốn chữ số.
\shortans{10999}
\loigiai{
Số lớn nhất có 4 chữ số: 9999; số bé nhất có 4 chữ số: 1000. $9999 + 1000 = 10999$.
}
\end{ex}

\begin{ex}
Một xe buýt chở 38 người, một tàu hỏa chở 1250 người. Cả hai phương tiện chở bao nhiêu người?
\shortans{1288}
\loigiai{
$1250 + 38 = 1288$ người.
}
\end{ex}

\begin{ex}
Tính nhanh: $2500 + 1700 + 3500$.
\shortans{7700}
\loigiai{
$2500 + 3500 + 1700 = 6000 + 1700 = 7700$.
}
\end{ex}

\begin{ex}
Quãng đường từ điểm A đến điểm B dài 3420 m, từ B đến C dài 2580 m. Quãng đường từ A đến C qua B dài bao nhiêu mét?
\shortans{6000}
\loigiai{
$3420 + 2580 = 6000$ m.
}
\end{ex}

\begin{ex}
Một vườn ươm có 3800 cây keo và 4500 cây bạch đàn. Vườn có tất cả bao nhiêu cây?
\shortans{8300}
\loigiai{
$3800 + 4500 = 8300$ cây.
}
\end{ex}

\begin{ex}
Tính: $5000 + 4999 = \dots$
\shortans{9999}
\loigiai{
$5000 + 4999 = 9999$.
}
\end{ex}

\begin{ex}
Tính: $7895 + 105 = \dots$
\shortans{8000}
\loigiai{
$7895 + 105 = 8000$.
}
\end{ex}

\begin{ex}
Tính: $6240 + 3760 = \dots$
\shortans{10000}
\loigiai{
$6240 + 3760 = 10000$.
}
\end{ex}

\newpage

\TaoPhieuBaiTap{Chủ đề 10 -- Cộng, trừ, nhân, chia trong phạm vi 10 000}{55}{Phép trừ trong phạm vi 10 000}{https://loptoanthaylinh.vn/bai-tap/toan-3/bai-55}
\begin{ex}
Tính: $7685 - 3421 = \dots$
\shortans{4264}
\loigiai{
$7685 - 3421 = 4264$.
}
\end{ex}

\begin{ex}
Tính: $8956 - 4528 = \dots$
\shortans{4428}
\loigiai{
$8956 - 4528 = 4428$.
}
\end{ex}

\begin{ex}
Tính: $6352 - 2718 = \dots$
\shortans{3634}
\loigiai{
$6352 - 2718 = 3634$.
}
\end{ex}

\begin{ex}
Tính: $5240 - 1835 = \dots$
\shortans{3405}
\loigiai{
$5240 - 1835 = 3405$.
}
\end{ex}

\begin{ex}
Tính: $8000 - 3500 = \dots$
\shortans{4500}
\loigiai{
$8000 - 3500 = 4500$.
}
\end{ex}

\begin{ex}
Tính: $10000 - 4000 = \dots$
\shortans{6000}
\loigiai{
$10000 - 4000 = 6000$.
}
\end{ex}

\begin{ex}
Tính: $7000 - 2450 = \dots$
\shortans{4550}
\loigiai{
$7000 - 2450 = 4550$.
}
\end{ex}

\begin{ex}
Tính: $9000 - 4380 = \dots$
\shortans{4620}
\loigiai{
$9000 - 4380 = 4620$.
}
\end{ex}

\begin{ex}
Tính: $6540 - 890 = \dots$
\shortans{5650}
\loigiai{
$6540 - 890 = 5650$.
}
\end{ex}

\begin{ex}
Tính: $4000 - 125 = \dots$
\shortans{3875}
\loigiai{
$4000 - 125 = 3875$.
}
\end{ex}

\begin{ex}
Một cửa hàng có 5000 kg xi măng, đã bán được 2450 kg. Cửa hàng còn lại bao nhiêu ki-lô-gam xi măng?
\shortans{2550}
\loigiai{
$5000 - 2450 = 2550$ kg.
}
\end{ex}

\begin{ex}
Một đội trồng rừng dự định trồng 8500 cây thông, đã trồng được 5600 cây. Hỏi đội còn phải trồng bao nhiêu cây nữa?
\shortans{2900}
\loigiai{
$8500 - 5600 = 2900$ cây.
}
\end{ex}

\begin{ex}
Một bể chứa 6000 lít nước, người ta đã bơm ra tưới rau 2350 lít nước. Trong bể còn lại bao nhiêu lít nước?
\shortans{3650}
\loigiai{
$6000 - 2350 = 3650$ lít.
}
\end{ex}

\begin{ex}
Tìm $x$, biết: $x + 3500 = 8000$.
\shortans{4500}
\loigiai{
$x = 8000 - 3500 = 4500$.
}
\end{ex}

\begin{ex}
Tìm $x$, biết: $2450 + x = 7000$.
\shortans{4550}
\loigiai{
$x = 7000 - 2450 = 4550$.
}
\end{ex}

\begin{ex}
Tìm $x$, biết: $8500 - x = 3200$.
\shortans{5300}
\loigiai{
$x = 8500 - 3200 = 5300$.
}
\end{ex}

\begin{ex}
Tìm $x$, biết: $10000 - x = 6500$.
\shortans{3500}
\loigiai{
$x = 10000 - 6500 = 3500$.
}
\end{ex}

\begin{ex}
Hiệu của số lớn nhất có bốn chữ số và số tròn nghìn bé nhất có bốn chữ số là bao nhiêu?
\shortans{8999}
\loigiai{
$9999 - 1000 = 8999$.
}
\end{ex}

\begin{ex}
Tính giá trị biểu thức: $8500 - 2500 - 1500$.
\shortans{4500}
\loigiai{
$6000 - 1500 = 4500$.
}
\end{ex}

\begin{ex}
Tính giá trị biểu thức: $10000 - (3500 + 4500)$.
\shortans{2000}
\loigiai{
$10000 - 8000 = 2000$.
}
\end{ex}

\begin{ex}
Một thư viện có 9840 cuốn sách, đã cho học sinh mượn 3420 cuốn. Thư viện còn lại bao nhiêu cuốn sách?
\shortans{6420}
\loigiai{
$9840 - 3420 = 6420$ cuốn.
}
\end{ex}

\begin{ex}
Năm nay một trang trại thu được 7500 quả trứng vịt, nhiều hơn năm ngoái 1800 quả. Năm ngoái trang trại thu được bao nhiêu quả trứng vịt?
\shortans{5700}
\loigiai{
$7500 - 1800 = 5700$ quả.
}
\end{ex}

\begin{ex}
Quãng đường từ A đến B dài 4500 m. Một người đã đi được 2800 m. Người đó còn phải đi bao nhiêu mét nữa để đến B?
\shortans{1700}
\loigiai{
$4500 - 2800 = 1700$ m.
}
\end{ex}

\begin{ex}
Tính: $9500 - 5500 = \dots$
\shortans{4000}
\loigiai{
$9500 - 5500 = 4000$.
}
\end{ex}

\begin{ex}
Tính: $7340 - 340 = \dots$
\shortans{7000}
\loigiai{
$7340 - 340 = 7000$.
}
\end{ex}

\newpage

\TaoPhieuBaiTap{Chủ đề 10 -- Cộng, trừ, nhân, chia trong phạm vi 10 000}{56}{Nhân số có bốn chữ số với số có một chữ số}{https://loptoanthaylinh.vn/bai-tap/toan-3/bai-56}
\begin{ex}
Tính: $1234 \times 2 = \dots$
\shortans{2468}
\loigiai{
$1234 \times 2 = 2468$.
}
\end{ex}

\begin{ex}
Tính: $2131 \times 3 = \dots$
\shortans{6393}
\loigiai{
$2131 \times 3 = 6393$.
}
\end{ex}

\begin{ex}
Tính: $1023 \times 3 = \dots$
\shortans{3069}
\loigiai{
$1023 \times 3 = 3069$.
}
\end{ex}

\begin{ex}
Tính: $1215 \times 4 = \dots$
\shortans{4860}
\loigiai{
$1215 \times 4 = 4860$.
}
\end{ex}

\begin{ex}
Tính: $1420 \times 3 = \dots$
\shortans{4260}
\loigiai{
$1420 \times 3 = 4260$.
}
\end{ex}

\begin{ex}
Tính: $1512 \times 5 = \dots$
\shortans{7560}
\loigiai{
$1512 \times 5 = 7560$.
}
\end{ex}

\begin{ex}
Tính: $1624 \times 4 = \dots$
\shortans{6496}
\loigiai{
$1624 \times 4 = 6496$.
}
\end{ex}

\begin{ex}
Tính: $2118 \times 3 = \dots$
\shortans{6354}
\loigiai{
$2118 \times 3 = 6354$.
}
\end{ex}

\begin{ex}
Tính: $2415 \times 4 = \dots$
\shortans{9660}
\loigiai{
$2415 \times 4 = 9660$.
}
\end{ex}

\begin{ex}
Tính: $1800 \times 5 = \dots$
\shortans{9000}
\loigiai{
$1800 \times 5 = 9000$.
}
\end{ex}

\begin{ex}
Tính: $1250 \times 6 = \dots$
\shortans{7500}
\loigiai{
$1250 \times 6 = 7500$.
}
\end{ex}

\begin{ex}
Tính: $1111 \times 8 = \dots$
\shortans{8888}
\loigiai{
$1111 \times 8 = 8888$.
}
\end{ex}

\begin{ex}
Tính: $1005 \times 7 = \dots$
\shortans{7035}
\loigiai{
$1005 \times 7 = 7035$.
}
\end{ex}

\begin{ex}
Tính: $1408 \times 6 = \dots$
\shortans{8448}
\loigiai{
$1408 \times 6 = 8448$.
}
\end{ex}

\begin{ex}
Một xe tải chở 1250 viên gạch. Hỏi 4 xe tải như thế chở được bao nhiêu viên gạch?
\shortans{5000}
\loigiai{
$1250 \times 4 = 5000$ viên gạch.
}
\end{ex}

\begin{ex}
Mỗi cuộn dây thép nặng 1500 g. Hỏi 6 cuộn dây thép như thế nặng bao nhiêu gam?
\shortans{9000}
\loigiai{
$1500 \times 6 = 9000$ g.
}
\end{ex}

\begin{ex}
Một trang trại nuôi bò có 4 khu, mỗi khu có 1150 con bò. Trang trại có tất cả bao nhiêu con bò?
\shortans{4600}
\loigiai{
$1150 \times 4 = 4600$ con bò.
}
\end{ex}

\begin{ex}
Một công ty mua 3 chiếc máy in, mỗi chiếc giá 2300 nghìn đồng. Công ty phải trả bao nhiêu nghìn đồng?
\shortans{6900}
\loigiai{
$2300 \times 3 = 6900$ nghìn đồng.
}
\end{ex}

\begin{ex}
Tìm $x$, biết: $x : 4 = 1520$.
\shortans{6080}
\loigiai{
$x = 1520 \times 4 = 6080$.
}
\end{ex}

\begin{ex}
Tìm $x$, biết: $x : 3 = 2410$.
\shortans{7230}
\loigiai{
$x = 2410 \times 3 = 7230$.
}
\end{ex}

\begin{ex}
Tính giá trị biểu thức: $1200 \times 4 + 1500$.
\shortans{6300}
\loigiai{
$4800 + 1500 = 6300$.
}
\end{ex}

\begin{ex}
Tính giá trị biểu thức: $2100 \times 3 - 3500$.
\shortans{2800}
\loigiai{
$6300 - 3500 = 2800$.
}
\end{ex}

\begin{ex}
Một hình vuông có cạnh dài 1250 m. Chu vi hình vuông đó là bao nhiêu mét?
\shortans{5000}
\loigiai{
$1250 \times 4 = 5000$ m.
}
\end{ex}

\begin{ex}
Mỗi bao gạo nặng 50 kg. Một kho có 140 bao gạo như vậy, tổng số ki-lô-gam gạo là bao nhiêu?
\shortans{7000}
\loigiai{
$50 \times 140 = 7000$ kg.
}
\end{ex}

\begin{ex}
Tính: $2500 \times 4 = \dots$
\shortans{10000}
\loigiai{
$2500 \times 4 = 10000$.
}
\end{ex}

\newpage

\TaoPhieuBaiTap{Chủ đề 10 -- Cộng, trừ, nhân, chia trong phạm vi 10 000}{57}{Chia số có bốn chữ số cho số có một chữ số}{https://loptoanthaylinh.vn/bai-tap/toan-3/bai-57}
\begin{ex}
Tính: $4826 : 2 = \dots$
\shortans{2413}
\loigiai{
$4826 : 2 = 2413$.
}
\end{ex}

\begin{ex}
Tính: $6936 : 3 = \dots$
\shortans{2312}
\loigiai{
$6936 : 3 = 2312$.
}
\end{ex}

\begin{ex}
Tính: $8448 : 4 = \dots$
\shortans{2112}
\loigiai{
$8448 : 4 = 2112$.
}
\end{ex}

\begin{ex}
Tính: $5550 : 5 = \dots$
\shortans{1110}
\loigiai{
$5550 : 5 = 1110$.
}
\end{ex}

\begin{ex}
Tính: $6248 : 2 = \dots$
\shortans{3124}
\loigiai{
$6248 : 2 = 3124$.
}
\end{ex}

\begin{ex}
Tính: $7535 : 5 = \dots$
\shortans{1507}
\loigiai{
$7535 : 5 = 1507$.
}
\end{ex}

\begin{ex}
Tính: $6480 : 6 = \dots$
\shortans{1080}
\loigiai{
$6480 : 6 = 1080$.
}
\end{ex}

\begin{ex}
Tính: $8470 : 7 = \dots$
\shortans{1210}
\loigiai{
$8470 : 7 = 1210$.
}
\end{ex}

\begin{ex}
Tính: $9632 : 8 = \dots$
\shortans{1204}
\loigiai{
$9632 : 8 = 1204$.
}
\end{ex}

\begin{ex}
Tính: $9018 : 9 = \dots$
\shortans{1002}
\loigiai{
$9018 : 9 = 1002$.
}
\end{ex}

\begin{ex}
Tính: $7254 : 6 = \dots$
\shortans{1209}
\loigiai{
$7254 : 6 = 1209$.
}
\end{ex}

\begin{ex}
Tính: $8540 : 5 = \dots$
\shortans{1708}
\loigiai{
$8540 : 5 = 1708$.
}
\end{ex}

\begin{ex}
Tính: $9849 : 7 = \dots$
\shortans{1407}
\loigiai{
$9849 : 7 = 1407$.
}
\end{ex}

\begin{ex}
Có 4860 kg thóc chia đều vào 4 kho. Hỏi mỗi kho nhận được bao nhiêu ki-lô-gam thóc?
\shortans{1215}
\loigiai{
$4860 : 4 = 1215$ kg.
}
\end{ex}

\begin{ex}
Một cuộn dây thép dài 6390 m cắt thành 3 đoạn bằng nhau. Mỗi đoạn dài bao nhiêu mét?
\shortans{2130}
\loigiai{
$6390 : 3 = 2130$ m.
}
\end{ex}

\begin{ex}
Một trường học có 1248 học sinh chia đều vào 4 khối lớp. Mỗi khối có bao nhiêu học sinh?
\shortans{312}
\loigiai{
$1248 : 4 = 312$ học sinh.
}
\end{ex}

\begin{ex}
Tìm số dư trong phép chia: $5627 : 5$.
\shortans{2}
\loigiai{
$5627 = 5 \times 1125 + 2$, dư 2.
}
\end{ex}

\begin{ex}
Tìm thương trong phép chia: $6845 : 4$.
\shortans{1711}
\loigiai{
$6845 = 4 \times 1711 + 1$, thương là 1711.
}
\end{ex}

\begin{ex}
Tìm $x$, biết: $x \times 3 = 4560$.
\shortans{1520}
\loigiai{
$x = 4560 : 3 = 1520$.
}
\end{ex}

\begin{ex}
Tìm $x$, biết: $5 \times x = 7650$.
\shortans{1530}
\loigiai{
$x = 7650 : 5 = 1530$.
}
\end{ex}

\begin{ex}
Tính giá trị biểu thức: $8400 : 4 + 1200$.
\shortans{3300}
\loigiai{
$2100 + 1200 = 3300$.
}
\end{ex}

\begin{ex}
Tính giá trị biểu thức: $9600 : 3 - 1800$.
\shortans{1400}
\loigiai{
$3200 - 1800 = 1400$.
}
\end{ex}

\begin{ex}
Người ta đóng 7200 hộp sữa vào 8 thùng lớn bằng nhau. Mỗi thùng chứa bao nhiêu hộp sữa?
\shortans{900}
\loigiai{
$7200 : 8 = 900$ hộp sữa.
}
\end{ex}

\begin{ex}
Một mảnh đất hình vuông có chu vi 4840 m. Độ dài mỗi cạnh mảnh đất là bao nhiêu mét?
\shortans{1210}
\loigiai{
$4840 : 4 = 1210$ m.
}
\end{ex}

\begin{ex}
Tính: $10000 : 2 = \dots$
\shortans{5000}
\loigiai{
$10000 : 2 = 5000$.
}
\end{ex}

\newpage

\TaoPhieuBaiTap{Chủ đề 10 -- Cộng, trừ, nhân, chia trong phạm vi 10 000}{58}{Luyện tập chung phép tính 10 000}{https://loptoanthaylinh.vn/bai-tap/toan-3/bai-58}
\begin{ex}
Tính: $4520 + 3280 = \dots$
\shortans{7800}
\loigiai{
$4520 + 3280 = 7800$.
}
\end{ex}

\begin{ex}
Tính: $8950 - 4520 = \dots$
\shortans{4430}
\loigiai{
$8950 - 4520 = 4430$.
}
\end{ex}

\begin{ex}
Tính: $1420 \times 4 = \dots$
\shortans{5680}
\loigiai{
$1420 \times 4 = 5680$.
}
\end{ex}

\begin{ex}
Tính: $7535 : 5 = \dots$
\shortans{1507}
\loigiai{
$7535 : 5 = 1507$.
}
\end{ex}

\begin{ex}
Tính giá trị biểu thức: $2400 + 1200 \times 3$.
\shortans{6000}
\loigiai{
$2400 + 3600 = 6000$.
}
\end{ex}

\begin{ex}
Tính giá trị biểu thức: $8000 - 4800 : 4$.
\shortans{6800}
\loigiai{
$8000 - 1200 = 6800$.
}
\end{ex}

\begin{ex}
Tính giá trị biểu thức: $(3500 + 1500) : 5$.
\shortans{1000}
\loigiai{
$5000 : 5 = 1000$.
}
\end{ex}

\begin{ex}
Tìm $x$: $x + 2450 = 6000$.
\shortans{3550}
\loigiai{
$x = 6000 - 2450 = 3550$.
}
\end{ex}

\begin{ex}
Tìm $x$: $x - 3800 = 4200$.
\shortans{8000}
\loigiai{
$x = 4200 + 3800 = 8000$.
}
\end{ex}

\begin{ex}
Tìm $x$: $x \times 4 = 8480$.
\shortans{2120}
\loigiai{
$x = 8480 : 4 = 2120$.
}
\end{ex}

\begin{ex}
Tìm $x$: $x : 5 = 1420$.
\shortans{7100}
\loigiai{
$x = 1420 \times 5 = 7100$.
}
\end{ex}

\begin{ex}
Một cửa hàng ngày đầu bán được 1250 kg gạo, ngày sau bán được gấp đôi ngày đầu. Cả hai ngày bán được bao nhiêu ki-lô-gam gạo?
\shortans{3750}
\loigiai{
Ngày sau: $1250 \times 2 = 2500$ kg. Cả hai ngày: $1250 + 2500 = 3750$ kg.
}
\end{ex}

\begin{ex}
Có 4 thùng nước, mỗi thùng chứa 1500 lít. Đã dùng 3500 lít. Còn lại bao nhiêu lít nước?
\shortans{2500}
\loigiai{
Tổng: $1500 \times 4 = 6000$ lít. Còn lại: $6000 - 3500 = 2500$ lít.
}
\end{ex}

\begin{ex}
Chu vi hình vuông có cạnh 1400 cm là bao nhiêu mét?
\shortans{56}
\loigiai{
$1400 \times 4 = 5600\text{ cm} = 56\text{ m}$.
}
\end{ex}

\begin{ex}
Một hình chữ nhật có chu vi 8000 m, chiều dài là 2500 m. Chiều rộng hình chữ nhật đó là bao nhiêu mét?
\shortans{1500}
\loigiai{
Nửa chu vi = 4000 m; rộng = $4000 - 2500 = 1500$ m.
}
\end{ex}

\begin{ex}
Một số chia cho 6 được thương là 1250 và dư 3. Tìm số đó.
\shortans{7503}
\loigiai{
$1250 \times 6 + 3 = 7500 + 3 = 7503$.
}
\end{ex}

\begin{ex}
Tính số dư của phép chia: $6843 : 5$.
\shortans{3}
\loigiai{
$6843 = 5 \times 1368 + 3$, dư 3.
}
\end{ex}

\begin{ex}
Tìm tổng của số lớn nhất có 4 chữ số khác nhau và số bé nhất có 4 chữ số khác nhau.
\shortans{10899}
\loigiai{
$9876 + 1023 = 10899$.
}
\end{ex}

\begin{ex}
Một phân xưởng may ngày thứ nhất may được 1120 bộ quần áo, ngày thứ hai may được 1380 bộ. Hỏi cả hai ngày may được bao nhiêu bộ?
\shortans{2500}
\loigiai{
$1120 + 1380 = 2500$ bộ.
}
\end{ex}

\begin{ex}
Tính nhanh: $3600 + 2400 + 4000$.
\shortans{10000}
\loigiai{
$6000 + 4000 = 10000$.
}
\end{ex}

\begin{ex}
Tính: $9999 - 5555 = \dots$
\shortans{4444}
\loigiai{
$9999 - 5555 = 4444$.
}
\end{ex}

\begin{ex}
Tính: $1500 \times 6 = \dots$
\shortans{9000}
\loigiai{
$1500 \times 6 = 9000$.
}
\end{ex}

\begin{ex}
Tính: $9600 : 8 = \dots$
\shortans{1200}
\loigiai{
$9600 : 8 = 1200$.
}
\end{ex}

\begin{ex}
Bác An nuôi 1200 con gà, bác Bình nuôi gấp 3 lần bác An. Cả hai bác nuôi bao nhiêu con gà?
\shortans{4800}
\loigiai{
Bình: $1200 \times 3 = 3600$. Cả hai: $1200 + 3600 = 4800$ con.
}
\end{ex}

\begin{ex}
Tính: $10000 - 3500 \times 2 = \dots$
\shortans{3000}
\loigiai{
$10000 - 7000 = 3000$.
}
\end{ex}

\newpage


% =========================================================================
% Các số đến 100 000
% =========================================================================
\section{Các số đến 100 000}

\TaoPhieuBaiTap{Chủ đề 11 -- Các số đến 100 000}{59}{Các số có năm chữ số. Số 100 000}{https://loptoanthaylinh.vn/bai-tap/toan-3/bai-59}
\begin{ex}
Số gồm 5 chục nghìn, 3 nghìn, 6 trăm, 2 chục và 8 đơn vị viết là gì?
\shortans{53628}
\loigiai{
Số đó là 53628.
}
\end{ex}

\begin{ex}
Số 74152 có chữ số hàng chục nghìn là bao nhiêu?
\shortans{7}
\loigiai{
Chữ số hàng chục nghìn là 7.
}
\end{ex}

\begin{ex}
Số 80425 có chữ số hàng nghìn là bao nhiêu?
\shortans{0}
\loigiai{
Chữ số hàng nghìn là 0.
}
\end{ex}

\begin{ex}
Số liền sau của số 99 999 là số nào?
\shortans{100000}
\loigiai{
Số liền sau của 99 999 là 100 000.
}
\end{ex}

\begin{ex}
Số liền trước của số 50 000 là số nào?
\shortans{49999}
\loigiai{
$50000 - 1 = 49999$.
}
\end{ex}

\begin{ex}
Số lớn nhất có năm chữ số là số nào?
\shortans{99999}
\loigiai{
Số lớn nhất có 5 chữ số là 99999.
}
\end{ex}

\begin{ex}
Số bé nhất có năm chữ số là số nào?
\shortans{10000}
\loigiai{
Số bé nhất có 5 chữ số là 10000.
}
\end{ex}

\begin{ex}
Số tròn chục nghìn liền sau của 60 000 là số nào?
\shortans{70000}
\loigiai{
Số tròn chục nghìn liền sau là 70 000.
}
\end{ex}

\begin{ex}
Số tròn nghìn liền trước của 45 000 là số nào?
\shortans{44000}
\loigiai{
Số tròn nghìn liền trước là 44 000.
}
\end{ex}

\begin{ex}
Số tròn chục liền sau của 81 240 là số nào?
\shortans{81250}
\loigiai{
Số đó là 81250.
}
\end{ex}

\begin{ex}
Viết số thích hợp: $60000 + 4000 + 300 + 20 + 5 = \dots$
\shortans{64325}
\loigiai{
Số đó là 64325.
}
\end{ex}

\begin{ex}
Viết số thích hợp: $80000 + 500 + 40 + 7 = \dots$
\shortans{80547}
\loigiai{
Số đó là 80547.
}
\end{ex}

\begin{ex}
Viết số thích hợp: $90000 + 8000 + 6 = \dots$
\shortans{98006}
\loigiai{
Số đó là 98006.
}
\end{ex}

\begin{ex}
Số 75 412 có giá trị của chữ số 7 là bao nhiêu?
\shortans{70000}
\loigiai{
Chữ số 7 ở hàng chục nghìn có giá trị 70 000.
}
\end{ex}

\begin{ex}
Số 38 650 có giá trị của chữ số 8 là bao nhiêu?
\shortans{8000}
\loigiai{
Chữ số 8 ở hàng nghìn có giá trị 8000.
}
\end{ex}

\begin{ex}
Cho các chữ số 4, 1, 9, 0, 7. Số lớn nhất có năm chữ số khác nhau lập từ các chữ số trên là số nào?
\shortans{97410}
\loigiai{
Sắp xếp giảm dần: 97410.
}
\end{ex}

\begin{ex}
Cho các chữ số 5, 2, 0, 8, 3. Số bé nhất có năm chữ số khác nhau lập từ các chữ số trên là số nào?
\shortans{20358}
\loigiai{
Chữ số đầu khác 0 nhỏ nhất là 2, tiếp theo là 0, 3, 5, 8: 20358.
}
\end{ex}

\begin{ex}
Có bao nhiêu số tròn chục nghìn từ 10 000 đến 100 000?
\shortans{10}
\loigiai{
10000, 20000, ..., 100000 (10 số).
}
\end{ex}

\begin{ex}
Điền số thích hợp: $69999 < \dots < 70001$.
\shortans{70000}
\loigiai{
Số ở giữa là 70000.
}
\end{ex}

\begin{ex}
Số liền sau của số 45 899 là số nào?
\shortans{45900}
\loigiai{
$45899 + 1 = 45900$.
}
\end{ex}

\begin{ex}
Số liền trước của số 32 000 là số nào?
\shortans{31999}
\loigiai{
$32000 - 1 = 31999$.
}
\end{ex}

\begin{ex}
Hiệu của số bé nhất có năm chữ số và số lớn nhất có bốn chữ số là bao nhiêu?
\shortans{1}
\loigiai{
$10000 - 9999 = 1$.
}
\end{ex}

\begin{ex}
Số 88 888 có bao nhiêu chữ số 8?
\shortans{5}
\loigiai{
Có 5 chữ số 8.
}
\end{ex}

\begin{ex}
Trong số 65 321, chữ số hàng trăm là chữ số nào?
\shortans{3}
\loigiai{
Chữ số hàng trăm là 3.
}
\end{ex}

\begin{ex}
Giá trị của chữ số 5 ở hàng chục nghìn gấp mấy lần giá trị của chữ số 5 ở hàng trăm?
\shortans{100}
\loigiai{
$50000 : 500 = 100$ lần.
}
\end{ex}

\newpage

\TaoPhieuBaiTap{Chủ đề 11 -- Các số đến 100 000}{60}{So sánh các số trong phạm vi 100 000}{https://loptoanthaylinh.vn/bai-tap/toan-3/bai-60}
\begin{ex}
Trong các số $45600, 46500, 54600, 56400$, số nào lớn nhất?
\shortans{56400}
\loigiai{
56400 là số lớn nhất.
}
\end{ex}

\begin{ex}
Trong các số $87900, 89700, 87090, 89070$, số nào bé nhất?
\shortans{87090}
\loigiai{
87090 là số bé nhất.
}
\end{ex}

\begin{ex}
Tìm số lớn thứ nhì trong các số: $71200, 72100, 70999, 72050$.
\shortans{72050}
\loigiai{
Giảm dần: 72100, 72050, 71200, 70999. Lớn thứ nhì là 72050.
}
\end{ex}

\begin{ex}
Tìm $x$ là số tròn chục nghìn thỏa mãn $40000 < x < 60000$.
\shortans{50000}
\loigiai{
Số đó là 50000.
}
\end{ex}

\begin{ex}
Tìm $y$ là số tự nhiên biết: $78954 < y < 78956$.
\shortans{78955}
\loigiai{
Số đó là 78955.
}
\end{ex}

\begin{ex}
Cho 4 số: 50000, 45000, 48000, 52000. Hiệu giữa số lớn nhất và số bé nhất là bao nhiêu?
\shortans{7000}
\loigiai{
Lớn nhất: 52000, bé nhất: 45000. Hiệu = $52000 - 45000 = 7000$.
}
\end{ex}

\begin{ex}
Có bao nhiêu số tự nhiên có năm chữ số lớn hơn 99 996?
\shortans{3}
\loigiai{
99997, 99998, 99999 (3 số).
}
\end{ex}

\begin{ex}
Có bao nhiêu số tự nhiên có năm chữ số nhỏ hơn 10 004?
\shortans{4}
\loigiai{
10000, 10001, 10002, 10003 (4 số).
}
\end{ex}

\begin{ex}
Điền chữ số thích hợp vào dấu * để: $68*50 > 68850$.
\shortans{9}
\loigiai{
Hàng trăm phải lớn hơn 8 nên $* = 9$.
}
\end{ex}

\begin{ex}
Điền chữ số thích hợp vào dấu * để: $4*500 < 41500$.
\shortans{0}
\loigiai{
Hàng nghìn phải nhỏ hơn 1 nên $* = 0$.
}
\end{ex}

\begin{ex}
Khoảng cách giữa hai số 60 000 và 55 000 là bao nhiêu?
\shortans{5000}
\loigiai{
$60000 - 55000 = 5000$.
}
\end{ex}

\begin{ex}
Số 75 000 hơn số 72 000 bao nhiêu đơn vị?
\shortans{3000}
\loigiai{
$75000 - 72000 = 3000$.
}
\end{ex}

\begin{ex}
Số 84 000 kém số 90 000 bao nhiêu đơn vị?
\shortans{6000}
\loigiai{
$90000 - 84000 = 6000$.
}
\end{ex}

\begin{ex}
Tỉnh A có 85 400 người, tỉnh B có 92 100 người. Tỉnh B nhiều hơn tỉnh A bao nhiêu người?
\shortans{6700}
\loigiai{
$92100 - 85400 = 6700$ người.
}
\end{ex}

\begin{ex}
Một nhà máy năm nay sản xuất 48 500 sản phẩm, nhiều hơn năm ngoái 3500 sản phẩm. Năm ngoái nhà máy sản xuất bao nhiêu sản phẩm?
\shortans{45000}
\loigiai{
$48500 - 3500 = 45000$ sản phẩm.
}
\end{ex}

\begin{ex}
Tìm số tự nhiên lớn nhất có 5 chữ số khác nhau.
\shortans{98765}
\loigiai{
Số đó là 98765.
}
\end{ex}

\begin{ex}
Tìm số tự nhiên bé nhất có 5 chữ số khác nhau.
\shortans{10234}
\loigiai{
Số đó là 10234.
}
\end{ex}

\begin{ex}
Tìm $x$ là số tròn nghìn thỏa mãn $38000 < x < 40000$.
\shortans{39000}
\loigiai{
Số đó là 39000.
}
\end{ex}

\begin{ex}
Tìm $x$ là số tròn trăm thỏa mãn $54800 < x < 55000$.
\shortans{54900}
\loigiai{
Số đó là 54900.
}
\end{ex}

\begin{ex}
Trong dãy số tự nhiên từ 20 000 đến 20 010 có bao nhiêu số?
\shortans{11}
\loigiai{
$20010 - 20000 + 1 = 11$ số.
}
\end{ex}

\begin{ex}
Số bé nhất có năm chữ số giống nhau là số nào?
\shortans{11111}
\loigiai{
Số đó là 11111.
}
\end{ex}

\begin{ex}
Số lớn nhất có năm chữ số giống nhau là số nào?
\shortans{99999}
\loigiai{
Số đó là 99999.
}
\end{ex}

\begin{ex}
Hiệu của số lớn nhất có 5 chữ số giống nhau và số bé nhất có 5 chữ số giống nhau là bao nhiêu?
\shortans{88888}
\loigiai{
$99999 - 11111 = 88888$.
}
\end{ex}

\begin{ex}
Số nào lớn hơn: 68 700 hay 68 070? (Hiệu giữa số lớn và số bé là bao nhiêu?).
\shortans{630}
\loigiai{
$68700 - 68070 = 630$.
}
\end{ex}

\begin{ex}
Số 100 000 có bao nhiêu chữ số 0?
\shortans{5}
\loigiai{
Có 5 chữ số 0.
}
\end{ex}

\newpage

\TaoPhieuBaiTap{Chủ đề 11 -- Các số đến 100 000}{61}{Làm tròn số đến hàng nghìn, hàng chục nghìn}{https://loptoanthaylinh.vn/bai-tap/toan-3/bai-61}
\begin{ex}
Làm tròn số 4200 đến hàng nghìn ta được số nào?
\shortans{4000}
\loigiai{
Hàng trăm là 2 < 5 nên làm tròn thành 4000.
}
\end{ex}

\begin{ex}
Làm tròn số 7800 đến hàng nghìn ta được số nào?
\shortans{8000}
\loigiai{
Hàng trăm là 8 $\ge$ 5 nên làm tròn lên 8000.
}
\end{ex}

\begin{ex}
Làm tròn số 5500 đến hàng nghìn ta được số nào?
\shortans{6000}
\loigiai{
Hàng trăm là 5 nên làm tròn lên 6000.
}
\end{ex}

\begin{ex}
Làm tròn số 14 300 đến hàng nghìn ta được số nào?
\shortans{14000}
\loigiai{
Hàng trăm là 3 nên làm tròn thành 14 000.
}
\end{ex}

\begin{ex}
Làm tròn số 26 750 đến hàng nghìn ta được số nào?
\shortans{27000}
\loigiai{
Hàng trăm là 7 nên làm tròn lên 27 000.
}
\end{ex}

\begin{ex}
Làm tròn số 38 500 đến hàng nghìn ta được số nào?
\shortans{39000}
\loigiai{
Hàng trăm là 5 nên làm tròn lên 39 000.
}
\end{ex}

\begin{ex}
Làm tròn số 34 000 đến hàng chục nghìn ta được số nào?
\shortans{30000}
\loigiai{
Hàng nghìn là 4 < 5 nên làm tròn thành 30 000.
}
\end{ex}

\begin{ex}
Làm tròn số 67 000 đến hàng chục nghìn ta được số nào?
\shortans{70000}
\loigiai{
Hàng nghìn là 7 $\ge$ 5 nên làm tròn lên 70 000.
}
\end{ex}

\begin{ex}
Làm tròn số 45 000 đến hàng chục nghìn ta được số nào?
\shortans{50000}
\loigiai{
Hàng nghìn là 5 nên làm tròn lên 50 000.
}
\end{ex}

\begin{ex}
Làm tròn số 82 450 đến hàng chục nghìn ta được số nào?
\shortans{80000}
\loigiai{
Hàng nghìn là 2 nên làm tròn thành 80 000.
}
\end{ex}

\begin{ex}
Làm tròn số 95 800 đến hàng chục nghìn ta được số nào?
\shortans{100000}
\loigiai{
Hàng nghìn là 5 nên làm tròn lên 100 000.
}
\end{ex}

\begin{ex}
Làm tròn số 56 120 đến hàng nghìn ta được số nào?
\shortans{56000}
\loigiai{
Hàng trăm là 1 nên làm tròn thành 56 000.
}
\end{ex}

\begin{ex}
Làm tròn số 73 890 đến hàng nghìn ta được số nào?
\shortans{74000}
\loigiai{
Hàng trăm là 8 nên làm tròn lên 74 000.
}
\end{ex}

\begin{ex}
Số nhỏ nhất khi làm tròn đến hàng nghìn được 40 000 là số nào?
\shortans{39500}
\loigiai{
Số nhỏ nhất là 39 500.
}
\end{ex}

\begin{ex}
Số lớn nhất khi làm tròn đến hàng nghìn được 40 000 là số nào?
\shortans{40499}
\loigiai{
Số lớn nhất là 40 499.
}
\end{ex}

\begin{ex}
Số nhỏ nhất khi làm tròn đến hàng chục nghìn được 60 000 là số nào?
\shortans{55000}
\loigiai{
Số nhỏ nhất là 55 000.
}
\end{ex}

\begin{ex}
Số lớn nhất khi làm tròn đến hàng chục nghìn được 60 000 là số nào?
\shortans{64999}
\loigiai{
Số lớn nhất là 64 999.
}
\end{ex}

\begin{ex}
Một sân vận động có 42 150 khán giả. Làm tròn đến hàng nghìn là khoảng bao nhiêu khán giả?
\shortans{42000}
\loigiai{
Làm tròn thành 42 000.
}
\end{ex}

\begin{ex}
Một công ty xuất khẩu 78 600 tấn gạo. Làm tròn đến hàng chục nghìn là khoảng bao nhiêu tấn gạo?
\shortans{80000}
\loigiai{
Làm tròn lên 80 000 tấn.
}
\end{ex}

\begin{ex}
Dân số một xã là 14 520 người. Làm tròn đến hàng nghìn là bao nhiêu người?
\shortans{15000}
\loigiai{
Làm tròn lên 15 000.
}
\end{ex}

\begin{ex}
Quãng đường xe chạy được 54 800 km. Làm tròn đến hàng chục nghìn là bao nhiêu ki-lô-mét?
\shortans{50000}
\loigiai{
Làm tròn thành 50 000 km.
}
\end{ex}

\begin{ex}
Làm tròn số 24 999 đến hàng chục nghìn ta được số nào?
\shortans{20000}
\loigiai{
Hàng nghìn là 4 nên làm tròn thành 20 000.
}
\end{ex}

\begin{ex}
Làm tròn số 25 000 đến hàng chục nghìn ta được số nào?
\shortans{30000}
\loigiai{
Hàng nghìn là 5 nên làm tròn lên 30 000.
}
\end{ex}

\begin{ex}
Làm tròn số 99 600 đến hàng nghìn ta được số nào?
\shortans{100000}
\loigiai{
Làm tròn lên 100 000.
}
\end{ex}

\begin{ex}
Làm tròn số 89 400 đến hàng nghìn ta được số nào?
\shortans{89000}
\loigiai{
Hàng trăm là 4 nên làm tròn thành 89 000.
}
\end{ex}

\newpage

\TaoPhieuBaiTap{Chủ đề 11 -- Các số đến 100 000}{62}{Luyện tập chung số có năm chữ số}{https://loptoanthaylinh.vn/bai-tap/toan-3/bai-62}
\begin{ex}
Số 84 520 gồm mấy chục nghìn?
\shortans{8}
\loigiai{
Chữ số hàng chục nghìn là 8.
}
\end{ex}

\begin{ex}
Số liền sau của 89 999 là số nào?
\shortans{90000}
\loigiai{
$89999 + 1 = 90000$.
}
\end{ex}

\begin{ex}
Số liền trước của 100 000 là số nào?
\shortans{99999}
\loigiai{
$100000 - 1 = 99999$.
}
\end{ex}

\begin{ex}
Trong các số $75400, 74500, 75040, 74050$, số lớn nhất là số nào?
\shortans{75400}
\loigiai{
75400 là số lớn nhất.
}
\end{ex}

\begin{ex}
Làm tròn số 58 400 đến hàng nghìn được số nào?
\shortans{58000}
\loigiai{
Hàng trăm là 4 nên làm tròn thành 58 000.
}
\end{ex}

\begin{ex}
Làm tròn số 65 000 đến hàng chục nghìn được số nào?
\shortans{70000}
\loigiai{
Hàng nghìn là 5 nên làm tròn lên 70 000.
}
\end{ex}

\begin{ex}
Viết số thích hợp: $70000 + 5000 + 400 + 3 = \dots$
\shortans{75403}
\loigiai{
Số đó là 75403.
}
\end{ex}

\begin{ex}
Tìm $x$ là số tròn chục nghìn thỏa mãn $70000 < x < 90000$.
\shortans{80000}
\loigiai{
Số đó là 80000.
}
\end{ex}

\begin{ex}
Cho 5 chữ số 1, 0, 4, 6, 8. Số lớn nhất có 5 chữ số khác nhau ghép từ các chữ số trên là số nào?
\shortans{86410}
\loigiai{
Sắp xếp giảm dần: 86410.
}
\end{ex}

\begin{ex}
Cho 5 chữ số 2, 0, 3, 5, 7. Số bé nhất có 5 chữ số khác nhau ghép từ các chữ số trên là số nào?
\shortans{20357}
\loigiai{
Số đó là 20357.
}
\end{ex}

\begin{ex}
Khoảng cách giữa hai số 80 000 và 76 000 là bao nhiêu?
\shortans{4000}
\loigiai{
$80000 - 76000 = 4000$.
}
\end{ex}

\begin{ex}
Số 65 000 hơn số 61 500 bao nhiêu đơn vị?
\shortans{3500}
\loigiai{
$65000 - 61500 = 3500$.
}
\end{ex}

\begin{ex}
Số tự nhiên nhỏ nhất có 5 chữ số khác nhau là số nào?
\shortans{10234}
\loigiai{
Số đó là 10234.
}
\end{ex}

\begin{ex}
Số tự nhiên lớn nhất có 5 chữ số khác nhau là số nào?
\shortans{98765}
\loigiai{
Số đó là 98765.
}
\end{ex}

\begin{ex}
Làm tròn số 47 620 đến hàng nghìn ta được số nào?
\shortans{48000}
\loigiai{
Hàng trăm là 6 nên làm tròn lên 48 000.
}
\end{ex}

\begin{ex}
Làm tròn số 83 200 đến hàng chục nghìn ta được số nào?
\shortans{80000}
\loigiai{
Hàng nghìn là 3 nên làm tròn thành 80 000.
}
\end{ex}

\begin{ex}
Có bao nhiêu số tròn chục nghìn có 5 chữ số?
\shortans{9}
\loigiai{
10000, 20000, ..., 90000 (9 số).
}
\end{ex}

\begin{ex}
Điền số thích hợp: $49999 < \dots < 50001$.
\shortans{50000}
\loigiai{
Số ở giữa là 50000.
}
\end{ex}

\begin{ex}
Số 77 777 có bao nhiêu chữ số 7?
\shortans{5}
\loigiai{
Có 5 chữ số 7.
}
\end{ex}

\begin{ex}
Tính: $50000 + 40000 = \dots$
\shortans{90000}
\loigiai{
$50000 + 40000 = 90000$.
}
\end{ex}

\begin{ex}
Tính: $100000 - 30000 = \dots$
\shortans{70000}
\loigiai{
$100000 - 30000 = 70000$.
}
\end{ex}

\begin{ex}
Tính: $20000 \times 4 = \dots$
\shortans{80000}
\loigiai{
$20000 \times 4 = 80000$.
}
\end{ex}

\begin{ex}
Tính: $90000 : 3 = \dots$
\shortans{30000}
\loigiai{
$90000 : 3 = 30000$.
}
\end{ex}

\begin{ex}
Làm tròn số 94 800 đến hàng chục nghìn được số nào?
\shortans{90000}
\loigiai{
Hàng nghìn là 4 nên làm tròn thành 90 000.
}
\end{ex}

\begin{ex}
Làm tròn số 95 000 đến hàng chục nghìn được số nào?
\shortans{100000}
\loigiai{
Hàng nghìn là 5 nên làm tròn lên 100 000.
}
\end{ex}

\newpage


% =========================================================================
% Cộng, trừ trong phạm vi 100 000
% =========================================================================
\section{Cộng, trừ trong phạm vi 100 000}

\TaoPhieuBaiTap{Chủ đề 12 -- Cộng, trừ trong phạm vi 100 000}{63}{Phép cộng trong phạm vi 100 000}{https://loptoanthaylinh.vn/bai-tap/toan-3/bai-63}
\begin{ex}
Tính: $35240 + 21430 = \dots$
\shortans{56670}
\loigiai{
$35240 + 21430 = 56670$.
}
\end{ex}

\begin{ex}
Tính: $48160 + 21530 = \dots$
\shortans{69690}
\loigiai{
$48160 + 21530 = 69690$.
}
\end{ex}

\begin{ex}
Tính: $52480 + 36190 = \dots$
\shortans{88670}
\loigiai{
$52480 + 36190 = 88670$.
}
\end{ex}

\begin{ex}
Tính: $27450 + 45380 = \dots$
\shortans{72830}
\loigiai{
$27450 + 45380 = 72830$.
}
\end{ex}

\begin{ex}
Tính: $63820 + 18490 = \dots$
\shortans{82310}
\loigiai{
$63820 + 18490 = 82310$.
}
\end{ex}

\begin{ex}
Tính: $40000 + 30000 = \dots$
\shortans{70000}
\loigiai{
$40000 + 30000 = 70000$.
}
\end{ex}

\begin{ex}
Tính: $75000 + 15000 = \dots$
\shortans{90000}
\loigiai{
$75000 + 15000 = 90000$.
}
\end{ex}

\begin{ex}
Tính: $64000 + 36000 = \dots$
\shortans{100000}
\loigiai{
$64000 + 36000 = 100000$.
}
\end{ex}

\begin{ex}
Tính: $54120 + 7350 = \dots$
\shortans{61470}
\loigiai{
$54120 + 7350 = 61470$.
}
\end{ex}

\begin{ex}
Tính: $82400 + 17600 = \dots$
\shortans{100000}
\loigiai{
$82400 + 17600 = 100000$.
}
\end{ex}

\begin{ex}
Tính giá trị biểu thức: $35000 + 25000 + 18000$.
\shortans{78000}
\loigiai{
$60000 + 18000 = 78000$.
}
\end{ex}

\begin{ex}
Tính giá trị biểu thức: $42000 + 18000 + 40000$.
\shortans{100000}
\loigiai{
$60000 + 40000 = 100000$.
}
\end{ex}

\begin{ex}
Một công ty tháng Một thu được 32 000 USD, tháng Hai thu được 45 000 USD. Cả hai tháng thu được bao nhiêu USD?
\shortans{77000}
\loigiai{
$32000 + 45000 = 77000$ USD.
}
\end{ex}

\begin{ex}
Một kho chứa 48 500 kg thóc, nhập thêm 21 500 kg thóc nữa. Trong kho có tất cả bao nhiêu ki-lô-gam thóc?
\shortans{70000}
\loigiai{
$48500 + 21500 = 70000$ kg.
}
\end{ex}

\begin{ex}
Huyện A có 45 600 người, huyện B có 52 300 người. Cả hai huyện có bao nhiêu người?
\shortans{97900}
\loigiai{
$45600 + 52300 = 97900$ người.
}
\end{ex}

\begin{ex}
Tìm $x$, biết: $x - 35000 = 45000$.
\shortans{80000}
\loigiai{
$x = 45000 + 35000 = 80000$.
}
\end{ex}

\begin{ex}
Tìm $x$, biết: $x - 28000 = 62000$.
\shortans{90000}
\loigiai{
$x = 62000 + 28000 = 90000$.
}
\end{ex}

\begin{ex}
Tính nhanh: $25000 + 34000 + 25000$.
\shortans{84000}
\loigiai{
$50000 + 34000 = 84000$.
}
\end{ex}

\begin{ex}
Một vườn ươm có 28 000 cây giống, đợt hai ươm thêm 35 000 cây giống. Cả hai đợt ươm được bao nhiêu cây?
\shortans{63000}
\loigiai{
$28000 + 35000 = 63000$ cây.
}
\end{ex}

\begin{ex}
Một bể nước có sẵn 25 000 lít nước, bơm thêm vào 45 000 lít nước. Bể có bao nhiêu lít nước?
\shortans{70000}
\loigiai{
$25000 + 45000 = 70000$ lít.
}
\end{ex}

\begin{ex}
Tính: $99000 + 1000 = \dots$
\shortans{100000}
\loigiai{
$99000 + 1000 = 100000$.
}
\end{ex}

\begin{ex}
Tính: $55555 + 44444 = \dots$
\shortans{99999}
\loigiai{
$55555 + 44444 = 99999$.
}
\end{ex}

\begin{ex}
Tổng của số lớn nhất có năm chữ số và số bé nhất có một chữ số là bao nhiêu?
\shortans{99999}
\loigiai{
$99999 + 0 = 99999$.
}
\end{ex}

\begin{ex}
Tính: $37200 + 42800 = \dots$
\shortans{80000}
\loigiai{
$37200 + 42800 = 80000$.
}
\end{ex}

\begin{ex}
Tính: $61500 + 28500 = \dots$
\shortans{90000}
\loigiai{
$61500 + 28500 = 90000$.
}
\end{ex}

\newpage

\TaoPhieuBaiTap{Chủ đề 12 -- Cộng, trừ trong phạm vi 100 000}{64}{Phép trừ trong phạm vi 100 000}{https://loptoanthaylinh.vn/bai-tap/toan-3/bai-64}
\begin{ex}
Tính: $76850 - 34210 = \dots$
\shortans{42640}
\loigiai{
$76850 - 34210 = 42640$.
}
\end{ex}

\begin{ex}
Tính: $89560 - 45280 = \dots$
\shortans{44280}
\loigiai{
$89560 - 45280 = 44280$.
}
\end{ex}

\begin{ex}
Tính: $63520 - 27180 = \dots$
\shortans{36340}
\loigiai{
$63520 - 27180 = 36340$.
}
\end{ex}

\begin{ex}
Tính: $100000 - 40000 = \dots$
\shortans{60000}
\loigiai{
$100000 - 40000 = 60000$.
}
\end{ex}

\begin{ex}
Tính: $100000 - 65000 = \dots$
\shortans{35000}
\loigiai{
$100000 - 65000 = 35000$.
}
\end{ex}

\begin{ex}
Tính: $80000 - 35000 = \dots$
\shortans{45000}
\loigiai{
$80000 - 35000 = 45000$.
}
\end{ex}

\begin{ex}
Tính: $90000 - 42500 = \dots$
\shortans{47500}
\loigiai{
$90000 - 42500 = 47500$.
}
\end{ex}

\begin{ex}
Tính: $75000 - 28000 = \dots$
\shortans{47000}
\loigiai{
$75000 - 28000 = 47000$.
}
\end{ex}

\begin{ex}
Tính: $60000 - 15800 = \dots$
\shortans{44200}
\loigiai{
$60000 - 15800 = 44200$.
}
\end{ex}

\begin{ex}
Tính: $50000 - 4999 = \dots$
\shortans{45001}
\loigiai{
$50000 - 4999 = 45001$.
}
\end{ex}

\begin{ex}
Một nông trường có 85 000 con gà, đã xuất chuồng 38 000 con. Nông trường còn lại bao nhiêu con gà?
\shortans{47000}
\loigiai{
$85000 - 38000 = 47000$ con.
}
\end{ex}

\begin{ex}
Một hồ nước có 90 000 lít nước, tháo bớt 34 500 lít để vệ sinh. Trong hồ còn lại bao nhiêu lít nước?
\shortans{55500}
\loigiai{
$90000 - 34500 = 55500$ lít.
}
\end{ex}

\begin{ex}
Một công ty đặt mục tiêu sản xuất 100 000 sản phẩm, đã sản xuất được 76 000 sản phẩm. Công ty còn phải sản xuất bao nhiêu sản phẩm nữa?
\shortans{24000}
\loigiai{
$100000 - 76000 = 24000$ sản phẩm.
}
\end{ex}

\begin{ex}
Tìm $x$, biết: $x + 45000 = 90000$.
\shortans{45000}
\loigiai{
$x = 90000 - 45000 = 45000$.
}
\end{ex}

\begin{ex}
Tìm $x$, biết: $32000 + x = 80000$.
\shortans{48000}
\loigiai{
$x = 80000 - 32000 = 48000$.
}
\end{ex}

\begin{ex}
Tìm $x$, biết: $100000 - x = 55000$.
\shortans{45000}
\loigiai{
$x = 100000 - 55000 = 45000$.
}
\end{ex}

\begin{ex}
Tính giá trị biểu thức: $95000 - 35000 - 20000$.
\shortans{40000}
\loigiai{
$60000 - 20000 = 40000$.
}
\end{ex}

\begin{ex}
Tính giá trị biểu thức: $100000 - (45000 + 35000)$.
\shortans{20000}
\loigiai{
$100000 - 80000 = 20000$.
}
\end{ex}

\begin{ex}
Một khu rừng có 72 000 cây xanh, đợt bão làm đổ 8500 cây. Khu rừng còn lại bao nhiêu cây xanh?
\shortans{63500}
\loigiai{
$72000 - 8500 = 63500$ cây.
}
\end{ex}

\begin{ex}
Quãng đường từ tỉnh A đến tỉnh B dài 65 000 m. Một đoàn tàu đã đi được 42 000 m. Đoàn tàu còn cách B bao nhiêu mét?
\shortans{23000}
\loigiai{
$65000 - 42000 = 23000$ m.
}
\end{ex}

\begin{ex}
Hiệu của 100 000 và 1 là bao nhiêu?
\shortans{99999}
\loigiai{
$100000 - 1 = 99999$.
}
\end{ex}

\begin{ex}
Tính: $88888 - 33333 = \dots$
\shortans{55555}
\loigiai{
$88888 - 33333 = 55555$.
}
\end{ex}

\begin{ex}
Tính: $70000 - 25000 = \dots$
\shortans{45000}
\loigiai{
$70000 - 25000 = 45000$.
}
\end{ex}

\begin{ex}
Tính: $92000 - 32000 = \dots$
\shortans{60000}
\loigiai{
$92000 - 32000 = 60000$.
}
\end{ex}

\begin{ex}
Tính: $84600 - 4600 = \dots$
\shortans{80000}
\loigiai{
$84600 - 4600 = 80000$.
}
\end{ex}

\newpage

\TaoPhieuBaiTap{Chủ đề 12 -- Cộng, trừ trong phạm vi 100 000}{65}{Luyện tập chung cộng trừ 100 000}{https://loptoanthaylinh.vn/bai-tap/toan-3/bai-65}
\begin{ex}
Tính: $45000 + 35000 = \dots$
\shortans{80000}
\loigiai{
$45000 + 35000 = 80000$.
}
\end{ex}

\begin{ex}
Tính: $90000 - 45000 = \dots$
\shortans{45000}
\loigiai{
$90000 - 45000 = 45000$.
}
\end{ex}

\begin{ex}
Tính: $62000 + 38000 = \dots$
\shortans{100000}
\loigiai{
$62000 + 38000 = 100000$.
}
\end{ex}

\begin{ex}
Tính: $100000 - 72000 = \dots$
\shortans{28000}
\loigiai{
$100000 - 72000 = 28000$.
}
\end{ex}

\begin{ex}
Tính giá trị biểu thức: $40000 + 30000 - 25000$.
\shortans{45000}
\loigiai{
$70000 - 25000 = 45000$.
}
\end{ex}

\begin{ex}
Tính giá trị biểu thức: $100000 - (30000 + 40000)$.
\shortans{30000}
\loigiai{
$100000 - 70000 = 30000$.
}
\end{ex}

\begin{ex}
Tìm $x$: $x + 28000 = 75000$.
\shortans{47000}
\loigiai{
$x = 75000 - 28000 = 47000$.
}
\end{ex}

\begin{ex}
Tìm $x$: $x - 34000 = 56000$.
\shortans{90000}
\loigiai{
$x = 56000 + 34000 = 90000$.
}
\end{ex}

\begin{ex}
Tìm $x$: $80000 - x = 32000$.
\shortans{48000}
\loigiai{
$x = 80000 - 32000 = 48000$.
}
\end{ex}

\begin{ex}
Một cửa hàng ngày thứ nhất thu được 25 000 nghìn đồng, ngày thứ hai thu được 35 000 nghìn đồng. Cả hai ngày thu được bao nhiêu nghìn đồng?
\shortans{60000}
\loigiai{
$25000 + 35000 = 60000$ nghìn đồng.
}
\end{ex}

\begin{ex}
Một xe tải chở 45 000 kg xi măng, đã giao cho công trường 18 000 kg. Trên xe còn lại bao nhiêu ki-lô-gam?
\shortans{27000}
\loigiai{
$45000 - 18000 = 27000$ kg.
}
\end{ex}

\begin{ex}
Quãng đường xe khách đi trong 2 ngày: ngày một đi 32 000 m, ngày hai đi 48 000 m. Cả hai ngày xe đi được bao nhiêu mét?
\shortans{80000}
\loigiai{
$32000 + 48000 = 80000$ m.
}
\end{ex}

\begin{ex}
Tổng của số tròn chục nghìn lớn nhất có 5 chữ số và số tròn chục nghìn bé nhất có 5 chữ số là bao nhiêu?
\shortans{100000}
\loigiai{
$90000 + 10000 = 100000$.
}
\end{ex}

\begin{ex}
Hiệu của 100 000 và số tròn chục nghìn lớn nhất có 5 chữ số là bao nhiêu?
\shortans{10000}
\loigiai{
$100000 - 90000 = 10000$.
}
\end{ex}

\begin{ex}
Tính: $35400 + 24600 = \dots$
\shortans{60000}
\loigiai{
$35400 + 24600 = 60000$.
}
\end{ex}

\begin{ex}
Tính: $78500 - 28500 = \dots$
\shortans{50000}
\loigiai{
$78500 - 28500 = 50000$.
}
\end{ex}

\begin{ex}
Kho thóc có 68 000 kg, đã xuất đi 25 000 kg rồi lại nhập thêm 17 000 kg. Kho thóc hiện có bao nhiêu ki-lô-gam?
\shortans{60000}
\loigiai{
$68000 - 25000 + 17000 = 60000$ kg.
}
\end{ex}

\begin{ex}
Tính nhanh: $45000 + 27000 + 15000$.
\shortans{87000}
\loigiai{
$45000 + 15000 + 27000 = 60000 + 27000 = 87000$.
}
\end{ex}

\begin{ex}
Tính: $99999 - 88888 = \dots$
\shortans{11111}
\loigiai{
$99999 - 88888 = 11111$.
}
\end{ex}

\begin{ex}
Tính: $50000 + 50000 = \dots$
\shortans{100000}
\loigiai{
$50000 + 50000 = 100000$.
}
\end{ex}

\begin{ex}
Tính: $100000 - 10000 = \dots$
\shortans{90000}
\loigiai{
$100000 - 10000 = 90000$.
}
\end{ex}

\begin{ex}
Một thư viện có 54 000 cuốn sách, trường quyên góp thêm 16 000 cuốn sách. Hỏi thư viện có tất cả bao nhiêu cuốn sách?
\shortans{70000}
\loigiai{
$54000 + 16000 = 70000$ cuốn.
}
\end{ex}

\begin{ex}
Một trại nuôi cá thu hoạch 42 000 kg cá tra và 28 000 kg cá basa. Tổng sản lượng cá là bao nhiêu ki-lô-gam?
\shortans{70000}
\loigiai{
$42000 + 28000 = 70000$ kg.
}
\end{ex}

\begin{ex}
Tính: $82000 - 41000 = \dots$
\shortans{41000}
\loigiai{
$82000 - 41000 = 41000$.
}
\end{ex}

\begin{ex}
Tính: $33333 + 66666 = \dots$
\shortans{99999}
\loigiai{
$33333 + 66666 = 99999$.
}
\end{ex}

\newpage


% =========================================================================
% Xem đồng hồ. tháng -- năm. tiền việt nam
% =========================================================================
\section{Xem đồng hồ. tháng -- năm. tiền việt nam}

\TaoPhieuBaiTap{Chủ đề 13 -- Xem đồng hồ. tháng -- năm. tiền việt nam}{66}{Xem đồng hồ}{https://loptoanthaylinh.vn/bai-tap/toan-3/bai-66}
\begin{ex}
Một năm thông thường có bao nhiêu tháng?
\shortans{12}
\loigiai{
Một năm có 12 tháng.
}
\end{ex}

\begin{ex}
Một năm thông thường không nhuận có bao nhiêu ngày?
\shortans{365}
\loigiai{
Năm thường có 365 ngày.
}
\end{ex}

\begin{ex}
Một năm nhuận có bao nhiêu ngày?
\shortans{366}
\loigiai{
Năm nhuận có 366 ngày.
}
\end{ex}

\begin{ex}
Tháng 1 có bao nhiêu ngày?
\shortans{31}
\loigiai{
Tháng 1 có 31 ngày.
}
\end{ex}

\begin{ex}
Tháng 2 của năm không nhuận có bao nhiêu ngày?
\shortans{28}
\loigiai{
Tháng 2 năm thường có 28 ngày.
}
\end{ex}

\begin{ex}
Tháng 2 của năm nhuận có bao nhiêu ngày?
\shortans{29}
\loigiai{
Tháng 2 năm nhuận có 29 ngày.
}
\end{ex}

\begin{ex}
Tháng 4 có bao nhiêu ngày?
\shortans{30}
\loigiai{
Tháng 4 có 30 ngày.
}
\end{ex}

\begin{ex}
Tháng 7 có bao nhiêu ngày?
\shortans{31}
\loigiai{
Tháng 7 có 31 ngày.
}
\end{ex}

\begin{ex}
Tháng 8 có bao nhiêu ngày?
\shortans{31}
\loigiai{
Tháng 8 có 31 ngày.
}
\end{ex}

\begin{ex}
Tháng 9 có bao nhiêu ngày?
\shortans{30}
\loigiai{
Tháng 9 có 30 ngày.
}
\end{ex}

\begin{ex}
Tháng 11 có bao nhiêu ngày?
\shortans{30}
\loigiai{
Tháng 11 có 30 ngày.
}
\end{ex}

\begin{ex}
Tháng 12 có bao nhiêu ngày?
\shortans{31}
\loigiai{
Tháng 12 có 31 ngày.
}
\end{ex}

\begin{ex}
Trong một năm có bao nhiêu tháng có 31 ngày?
\shortans{7}
\loigiai{
Các tháng: 1, 3, 5, 7, 8, 10, 12 (7 tháng).
}
\end{ex}

\begin{ex}
Trong một năm có bao nhiêu tháng có 30 ngày?
\shortans{4}
\loigiai{
Các tháng: 4, 6, 9, 11 (4 tháng).
}
\end{ex}

\begin{ex}
Một ngày có bao nhiêu giờ?
\shortans{24}
\loigiai{
Một ngày có 24 giờ.
}
\end{ex}

\begin{ex}
Một giờ có bao nhiêu phút?
\shortans{60}
\loigiai{
Một giờ có 60 phút.
}
\end{ex}

\begin{ex}
Một phút có bao nhiêu giây?
\shortans{60}
\loigiai{
Một phút có 60 giây.
}
\end{ex}

\begin{ex}
Đồng hồ chỉ 8 giờ kém 15 phút, tức là 7 giờ bao nhiêu phút?
\shortans{45}
\loigiai{
8 giờ kém 15 phút là 7 giờ 45 phút.
}
\end{ex}

\begin{ex}
Đồng hồ chỉ 3 giờ kém 10 phút, tức là 2 giờ bao nhiêu phút?
\shortans{50}
\loigiai{
2 giờ 50 phút.
}
\end{ex}

\begin{ex}
Đồng hồ chỉ 9 giờ rưỡi, tức là 9 giờ bao nhiêu phút?
\shortans{30}
\loigiai{
9 giờ rưỡi là 9 giờ 30 phút.
}
\end{ex}

\begin{ex}
Một tuần lễ có bao nhiêu ngày?
\shortans{7}
\loigiai{
Một tuần lễ có 7 ngày.
}
\end{ex}

\begin{ex}
Nếu ngày 1 tháng 5 là thứ Năm thì ngày 8 tháng 5 là thứ mấy? (Thứ Hai ghi 2, Ba ghi 3, Tư ghi 4, Năm ghi 5,...).
\shortans{5}
\loigiai{
Sau đúng 7 ngày vẫn là thứ Năm (5).
}
\end{ex}

\begin{ex}
Nếu ngày 2 tháng 9 là thứ Ba thì ngày 9 tháng 9 là thứ mấy? (Điền số tương ứng).
\shortans{3}
\loigiai{
Sau 7 ngày là thứ Ba (3).
}
\end{ex}

\begin{ex}
Từ 8 giờ sáng đến 11 giờ trưa cùng ngày có bao nhiêu giờ?
\shortans{3}
\loigiai{
$11 - 8 = 3$ giờ.
}
\end{ex}

\begin{ex}
Từ 14 giờ đến 18 giờ có bao nhiêu giờ?
\shortans{4}
\loigiai{
$18 - 14 = 4$ giờ.
}
\end{ex}

\newpage

\TaoPhieuBaiTap{Chủ đề 13 -- Xem đồng hồ. tháng -- năm. tiền việt nam}{67}{Tháng -- Năm}{https://loptoanthaylinh.vn/bai-tap/toan-3/bai-67}
\begin{ex}
Buổi học bắt đầu lúc 7 giờ 30 phút và kết thúc lúc 11 giờ cùng ngày. Buổi học kéo dài bao nhiêu phút?
\shortans{210}
\loigiai{
Từ 7h30 đến 11h là 3 giờ 30 phút = $3 \times 60 + 30 = 210$ phút.
}
\end{ex}

\begin{ex}
Hùng đi từ nhà lúc 6 giờ 45 phút và đến trường lúc 7 giờ. Hùng đi hết bao nhiêu phút?
\shortans{15}
\loigiai{
$7\text{h} - 6\text{h}45 = 15$ phút.
}
\end{ex}

\begin{ex}
Trận bóng đá bắt đầu lúc 15 giờ và kéo dài 90 phút. Trận đấu kết thúc lúc mấy giờ chiều và bao nhiêu phút? (Chỉ ghi số phút lẻ sau 16 giờ).
\shortans{30}
\loigiai{
15 giờ + 90 phút = 16 giờ 30 phút.
}
\end{ex}

\begin{ex}
Một bộ phim kéo dài từ 19 giờ 15 phút đến 21 giờ 15 phút. Bộ phim dài bao nhiêu giờ?
\shortans{2}
\loigiai{
$21 - 19 = 2$ giờ.
}
\end{ex}

\begin{ex}
Nếu ngày 20 tháng 11 là thứ Tư thì ngày 27 tháng 11 là thứ mấy? (Điền số thứ).
\shortans{4}
\loigiai{
Thứ Tư (4).
}
\end{ex}

\begin{ex}
Nếu ngày 8 tháng 3 là thứ Bảy thì ngày 1 tháng 3 là thứ mấy? (Điền số thứ).
\shortans{7}
\loigiai{
Lùi 7 ngày vẫn là thứ Bảy (7).
}
\end{ex}

\begin{ex}
Tháng 10 có bao nhiêu ngày?
\shortans{31}
\loigiai{
Tháng 10 có 31 ngày.
}
\end{ex}

\begin{ex}
Tháng 6 có bao nhiêu ngày?
\shortans{30}
\loigiai{
Tháng 6 có 30 ngày.
}
\end{ex}

\begin{ex}
Một tuần có 7 ngày. Hỏi 4 tuần có bao nhiêu ngày?
\shortans{28}
\loigiai{
$7 \times 4 = 28$ ngày.
}
\end{ex}

\begin{ex}
Bạn Lan ngủ từ 21 giờ đêm hôm trước đến 6 giờ sáng hôm sau. Bạn Lan đã ngủ bao nhiêu giờ?
\shortans{9}
\loigiai{
Từ 21h đến 24h là 3 giờ; $3 + 6 = 9$ giờ.
}
\end{ex}

\begin{ex}
Tàu hỏa xuất phát lúc 6 giờ sáng và chạy hết 8 giờ. Tàu đến ga đích lúc mấy giờ?
\shortans{14}
\loigiai{
$6 + 8 = 14$ giờ.
}
\end{ex}

\begin{ex}
Mẹ đi chợ từ 7 giờ 15 phút đến 8 giờ. Mẹ đi chợ hết bao nhiêu phút?
\shortans{45}
\loigiai{
$60 - 15 = 45$ phút.
}
\end{ex}

\begin{ex}
Một tiết học kéo dài 40 phút, bắt đầu lúc 8 giờ. Tiết học kết thúc lúc 8 giờ bao nhiêu phút?
\shortans{40}
\loigiai{
8 giờ 40 phút.
}
\end{ex}

\begin{ex}
Một ngày bác Nam làm việc 8 giờ. Một tuần bác làm việc 5 ngày. Hỏi trong một tuần bác Nam làm việc bao nhiêu giờ?
\shortans{40}
\loigiai{
$8 \times 5 = 40$ giờ.
}
\end{ex}

\begin{ex}
Nếu hôm nay là ngày 28 tháng 2 của một năm không nhuận thì ngày mai là ngày 1 tháng mấy?
\shortans{3}
\loigiai{
Tháng 3.
}
\end{ex}

\begin{ex}
Nếu hôm nay là ngày 31 tháng 12 thì ngày mai là ngày 1 tháng mấy?
\shortans{1}
\loigiai{
Ngày 1 tháng 1 năm sau.
}
\end{ex}

\begin{ex}
Tháng 3 có 31 ngày. Hỏi tháng 3 nhiều hơn tháng 4 bao nhiêu ngày?
\shortans{1}
\loigiai{
$31 - 30 = 1$ ngày.
}
\end{ex}

\begin{ex}
2 giờ bằng bao nhiêu phút?
\shortans{120}
\loigiai{
$2 \times 60 = 120$ phút.
}
\end{ex}

\begin{ex}
3 giờ bằng bao nhiêu phút?
\shortans{180}
\loigiai{
$3 \times 60 = 180$ phút.
}
\end{ex}

\begin{ex}
180 phút bằng bao nhiêu giờ?
\shortans{3}
\loigiai{
$180 : 60 = 3$ giờ.
}
\end{ex}

\begin{ex}
2 ngày có bao nhiêu giờ?
\shortans{48}
\loigiai{
$24 \times 2 = 48$ giờ.
}
\end{ex}

\begin{ex}
3 tuần lễ có bao nhiêu ngày?
\shortans{21}
\loigiai{
$7 \times 3 = 21$ ngày.
}
\end{ex}

\begin{ex}
Đồng hồ chỉ 10 giờ kém 20 phút. Hỏi còn bao nhiêu phút nữa thì đến 10 giờ đúng?
\shortans{20}
\loigiai{
Còn 20 phút nữa.
}
\end{ex}

\begin{ex}
Kim phút quay được 1 vòng quanh mặt đồng hồ thì hết bao nhiêu phút?
\shortans{60}
\loigiai{
Hết 60 phút (1 giờ).
}
\end{ex}

\begin{ex}
Kim giờ quay được 1 vòng quanh mặt đồng hồ thì hết bao nhiêu giờ?
\shortans{12}
\loigiai{
Hết 12 giờ.
}
\end{ex}

\newpage

\TaoPhieuBaiTap{Chủ đề 13 -- Xem đồng hồ. tháng -- năm. tiền việt nam}{68}{Tiền Việt Nam}{https://loptoanthaylinh.vn/bai-tap/toan-3/bai-68}
\begin{ex}
Đổi: 1 tờ 10 000 đồng đổi được bao nhiêu tờ 5 000 đồng?
\shortans{2}
\loigiai{
$10000 : 5000 = 2$ tờ.
}
\end{ex}

\begin{ex}
Đổi: 1 tờ 20 000 đồng đổi được bao nhiêu tờ 10 000 đồng?
\shortans{2}
\loigiai{
$20000 : 10000 = 2$ tờ.
}
\end{ex}

\begin{ex}
Đổi: 1 tờ 50 000 đồng đổi được bao nhiêu tờ 10 000 đồng?
\shortans{5}
\loigiai{
$50000 : 10000 = 5$ tờ.
}
\end{ex}

\begin{ex}
Đổi: 1 tờ 100 000 đồng đổi được bao nhiêu tờ 50 000 đồng?
\shortans{2}
\loigiai{
$100000 : 50000 = 2$ tờ.
}
\end{ex}

\begin{ex}
Đổi: 1 tờ 100 000 đồng đổi được bao nhiêu tờ 20 000 đồng?
\shortans{5}
\loigiai{
$100000 : 20000 = 5$ tờ.
}
\end{ex}

\begin{ex}
Đổi: 1 tờ 100 000 đồng đổi được bao nhiêu tờ 10 000 đồng?
\shortans{10}
\loigiai{
$100000 : 10000 = 10$ tờ.
}
\end{ex}

\begin{ex}
Bạn Mai có 3 tờ 10 000 đồng và 1 tờ 5 000 đồng. Mai có tất cả bao nhiêu đồng?
\shortans{35000}
\loigiai{
$3 \times 10000 + 5000 = 35000$ đồng.
}
\end{ex}

\begin{ex}
Bạn Nam mua một chiếc bút chì giá 5 000 đồng và một quyển vở giá 8 000 đồng. Nam phải trả bao nhiêu đồng?
\shortans{13000}
\loigiai{
$5000 + 8000 = 13000$ đồng.
}
\end{ex}

\begin{ex}
Một cục tẩy giá 4 000 đồng, một chiếc thước kẻ giá 6 000 đồng. Cả hai món đồ giá bao nhiêu đồng?
\shortans{10000}
\loigiai{
$4000 + 6000 = 10000$ đồng.
}
\end{ex}

\begin{ex}
Bạn An đưa cô bán hàng tờ 50 000 đồng để mua một hộp màu giá 35 000 đồng. Cô bán hàng phải trả lại An bao nhiêu đồng?
\shortans{15000}
\loigiai{
$50000 - 35000 = 15000$ đồng.
}
\end{ex}

\begin{ex}
Mẹ mua một chiếc nón lá giá 40 000 đồng và đưa người bán tờ 100 000 đồng. Người bán phải trả lại mẹ bao nhiêu đồng?
\shortans{60000}
\loigiai{
$100000 - 40000 = 60000$ đồng.
}
\end{ex}

\begin{ex}
Có 4 tờ 20 000 đồng. Tổng số tiền là bao nhiêu đồng?
\shortans{80000}
\loigiai{
$20000 \times 4 = 80000$ đồng.
}
\end{ex}

\begin{ex}
Có 2 tờ 50 000 đồng. Tổng số tiền là bao nhiêu đồng?
\shortans{100000}
\loigiai{
$50000 \times 2 = 100000$ đồng.
}
\end{ex}

\begin{ex}
Mua 5 quyển vở, mỗi quyển giá 6 000 đồng. Tổng số tiền mua vở là bao nhiêu đồng?
\shortans{30000}
\loigiai{
$6000 \times 5 = 30000$ đồng.
}
\end{ex}

\begin{ex}
Mua 3 chiếc bút bi, mỗi chiếc giá 4 000 đồng. Tổng số tiền mua bút là bao nhiêu đồng?
\shortans{12000}
\loigiai{
$4000 \times 3 = 12000$ đồng.
}
\end{ex}

\begin{ex}
Bạn Bình có 50 000 đồng, mua một cái kéo giá 12 000 đồng và hồ dán giá 8 000 đồng. Bình còn lại bao nhiêu đồng?
\shortans{30000}
\loigiai{
$50000 - (12000 + 8000) = 30000$ đồng.
}
\end{ex}

\begin{ex}
Một chiếc ba lô giá 85 000 đồng. Chị có 100 000 đồng thì sau khi mua ba lô chị còn lại bao nhiêu đồng?
\shortans{15000}
\loigiai{
$100000 - 85000 = 15000$ đồng.
}
\end{ex}

\begin{ex}
Giá tiền của 1 chiếc kem là 10 000 đồng. Mua 4 chiếc kem hết bao nhiêu đồng?
\shortans{40000}
\loigiai{
$10000 \times 4 = 40000$ đồng.
}
\end{ex}

\begin{ex}
Cô giáo mua 6 hộp phấn, mỗi hộp giá 5 000 đồng. Cô giáo đưa tờ 50 000 đồng thì được trả lại bao nhiêu đồng?
\shortans{20000}
\loigiai{
Tiền mua: $6 \times 5000 = 30000$ đồng. Trả lại: $50000 - 30000 = 20000$ đồng.
}
\end{ex}

\begin{ex}
Một gói kẹo giá 15 000 đồng, một gói bánh đắt hơn gói kẹo 10 000 đồng. Gói bánh giá bao nhiêu đồng?
\shortans{25000}
\loigiai{
$15000 + 10000 = 25000$ đồng.
}
\end{ex}

\begin{ex}
Trong các tờ tiền: 10 000 đồng, 20 000 đồng, 50 000 đồng, 100 000 đồng. Tờ có mệnh giá lớn nhất có mệnh giá là bao nhiêu nghìn đồng?
\shortans{100}
\loigiai{
100 000 đồng = 100 nghìn đồng.
}
\end{ex}

\begin{ex}
Bạn Hải nuôi lợn đất tiết kiệm được 2 tờ 20 000 đồng và 3 tờ 10 000 đồng. Hải tiết kiệm được bao nhiêu đồng?
\shortans{70000}
\loigiai{
$2 \times 20000 + 3 \times 10000 = 70000$ đồng.
}
\end{ex}

\begin{ex}
Hòa có 45 000 đồng. Hòa muốn mua một quyển truyện giá 60 000 đồng. Hòa còn thiếu bao nhiêu đồng?
\shortans{15000}
\loigiai{
$60000 - 45000 = 15000$ đồng.
}
\end{ex}

\begin{ex}
Tính tổng giá trị của: 1 tờ 50 000 đồng, 2 tờ 20 000 đồng và 1 tờ 10 000 đồng theo đơn vị đồng.
\shortans{100000}
\loigiai{
$50000 + 40000 + 10000 = 100000$ đồng.
}
\end{ex}

\begin{ex}
5 chiếc kẹp tóc giá 25 000 đồng. Một chiếc kẹp tóc giá bao nhiêu đồng?
\shortans{5000}
\loigiai{
$25000 : 5 = 5000$ đồng.
}
\end{ex}

\newpage

\TaoPhieuBaiTap{Chủ đề 13 -- Xem đồng hồ. tháng -- năm. tiền việt nam}{69}{Luyện tập chung thời gian và tiền tệ}{https://loptoanthaylinh.vn/bai-tap/toan-3/bai-69}
\begin{ex}
Tháng 5 có bao nhiêu ngày?
\shortans{31}
\loigiai{
Tháng 5 có 31 ngày.
}
\end{ex}

\begin{ex}
Tháng 6 có bao nhiêu ngày?
\shortans{30}
\loigiai{
Tháng 6 có 30 ngày.
}
\end{ex}

\begin{ex}
Một năm thường có bao nhiêu ngày?
\shortans{365}
\loigiai{
Có 365 ngày.
}
\end{ex}

\begin{ex}
Một ngày có bao nhiêu giờ?
\shortans{24}
\loigiai{
Có 24 giờ.
}
\end{ex}

\begin{ex}
Đồng hồ chỉ 14 giờ tức là mấy giờ chiều?
\shortans{2}
\loigiai{
$14 - 12 = 2$ giờ chiều.
}
\end{ex}

\begin{ex}
Đồng hồ chỉ 20 giờ tức là mấy giờ tối?
\shortans{8}
\loigiai{
$20 - 12 = 8$ giờ tối.
}
\end{ex}

\begin{ex}
Một giờ có bao nhiêu phút?
\shortans{60}
\loigiai{
Có 60 phút.
}
\end{ex}

\begin{ex}
Một tờ 50 000 đồng đổi được bao nhiêu tờ 10 000 đồng?
\shortans{5}
\loigiai{
$50000 : 10000 = 5$ tờ.
}
\end{ex}

\begin{ex}
Có 5 tờ 20 000 đồng. Tổng số tiền là bao nhiêu đồng?
\shortans{100000}
\loigiai{
$20000 \times 5 = 100000$ đồng.
}
\end{ex}

\begin{ex}
Mẹ mua một bó hoa giá 35 000 đồng và một thiệp chúc mừng giá 15 000 đồng. Mẹ phải trả bao nhiêu đồng?
\shortans{50000}
\loigiai{
$35000 + 15000 = 50000$ đồng.
}
\end{ex}

\begin{ex}
An có 100 000 đồng, mua một quả bóng đá giá 75 000 đồng. An còn lại bao nhiêu đồng?
\shortans{25000}
\loigiai{
$100000 - 75000 = 25000$ đồng.
}
\end{ex}

\begin{ex}
Một buổi xem xiếc bắt đầu lúc 19 giờ 30 phút và kết thúc lúc 21 giờ cùng ngày. Buổi biểu diễn kéo dài bao nhiêu phút?
\shortans{90}
\loigiai{
Kéo dài 1 giờ 30 phút = 90 phút.
}
\end{ex}

\begin{ex}
Nếu ngày 10 tháng 10 là thứ Sáu thì ngày 17 tháng 10 là thứ mấy? (Điền số thứ).
\shortans{6}
\loigiai{
Sau 7 ngày vẫn là thứ Sáu (6).
}
\end{ex}

\begin{ex}
Nếu ngày 1 tháng 6 là Chủ nhật thì ngày 8 tháng 6 là thứ mấy? (Chủ nhật quy ước là 8).
\shortans{8}
\loigiai{
Sau 7 ngày là Chủ nhật (8).
}
\end{ex}

\begin{ex}
Trong một năm có bao nhiêu tháng có 30 ngày?
\shortans{4}
\loigiai{
Có 4 tháng (4, 6, 9, 11).
}
\end{ex}

\begin{ex}
Trong một năm có bao nhiêu tháng có 31 ngày?
\shortans{7}
\loigiai{
Có 7 tháng.
}
\end{ex}

\begin{ex}
Một bộ đồ chơi xếp hình giá 65 000 đồng. Bạn Nam có 50 000 đồng. Nam còn thiếu bao nhiêu đồng?
\shortans{15000}
\loigiai{
$65000 - 50000 = 15000$ đồng.
}
\end{ex}

\begin{ex}
Mua 4 gói bánh, mỗi gói giá 12 000 đồng. Phải trả bao nhiêu đồng?
\shortans{48000}
\loigiai{
$12000 \times 4 = 48000$ đồng.
}
\end{ex}

\begin{ex}
Có 60 000 đồng chia đều cho 3 anh em. Mỗi người được bao nhiêu đồng?
\shortans{20000}
\loigiai{
$60000 : 3 = 20000$ đồng.
}
\end{ex}

\begin{ex}
Từ 9 giờ sáng đến 12 giờ trưa có bao nhiêu giờ?
\shortans{3}
\loigiai{
$12 - 9 = 3$ giờ.
}
\end{ex}

\begin{ex}
Nửa giờ bằng bao nhiêu phút?
\shortans{30}
\loigiai{
$60 : 2 = 30$ phút.
}
\end{ex}

\begin{ex}
Một phần tư giờ bằng bao nhiêu phút?
\shortans{15}
\loigiai{
$60 : 4 = 15$ phút.
}
\end{ex}

\begin{ex}
Đổi 2 tờ 50 000 đồng lấy các tờ 20 000 đồng thì được bao nhiêu tờ?
\shortans{5}
\loigiai{
$100000 : 20000 = 5$ tờ.
}
\end{ex}

\begin{ex}
Mẹ mua cá hết 45 000 đồng, mua rau hết 15 000 đồng. Mẹ đưa cô bán hàng 100 000 đồng thì được trả lại bao nhiêu đồng?
\shortans{40000}
\loigiai{
Tổng tiền mua = 60 000 đồng. Trả lại = $100000 - 60000 = 40000$ đồng.
}
\end{ex}

\begin{ex}
Tháng 2 của năm nhuận có bao nhiêu ngày?
\shortans{29}
\loigiai{
Có 29 ngày.
}
\end{ex}

\newpage


% =========================================================================
% Nhân, chia trong phạm vi 100 000
% =========================================================================
\section{Nhân, chia trong phạm vi 100 000}

\TaoPhieuBaiTap{Chủ đề 14 -- Nhân, chia trong phạm vi 100 000}{70}{Nhân số có năm chữ số với số có một chữ số}{https://loptoanthaylinh.vn/bai-tap/toan-3/bai-70}
\begin{ex}
Tính: $12341 \times 2 = \dots$
\shortans{24682}
\loigiai{
$12341 \times 2 = 24682$.
}
\end{ex}

\begin{ex}
Tính: $21321 \times 3 = \dots$
\shortans{63963}
\loigiai{
$21321 \times 3 = 63963$.
}
\end{ex}

\begin{ex}
Tính: $10214 \times 4 = \dots$
\shortans{40856}
\loigiai{
$10214 \times 4 = 40856$.
}
\end{ex}

\begin{ex}
Tính: $11215 \times 3 = \dots$
\shortans{33645}
\loigiai{
$11215 \times 3 = 33645$.
}
\end{ex}

\begin{ex}
Tính: $14205 \times 2 = \dots$
\shortans{28410}
\loigiai{
$14205 \times 2 = 28410$.
}
\end{ex}

\begin{ex}
Tính: $12150 \times 4 = \dots$
\shortans{48600}
\loigiai{
$12150 \times 4 = 48600$.
}
\end{ex}

\begin{ex}
Tính: $15000 \times 5 = \dots$
\shortans{75000}
\loigiai{
$15000 \times 5 = 75000$.
}
\end{ex}

\begin{ex}
Tính: $16000 \times 4 = \dots$
\shortans{64000}
\loigiai{
$16000 \times 4 = 64000$.
}
\end{ex}

\begin{ex}
Tính: $18000 \times 5 = \dots$
\shortans{90000}
\loigiai{
$18000 \times 5 = 90000$.
}
\end{ex}

\begin{ex}
Tính: $25000 \times 3 = \dots$
\shortans{75000}
\loigiai{
$25000 \times 3 = 75000$.
}
\end{ex}

\begin{ex}
Tính: $20000 \times 4 = \dots$
\shortans{80000}
\loigiai{
$20000 \times 4 = 80000$.
}
\end{ex}

\begin{ex}
Tính: $25000 \times 4 = \dots$
\shortans{100000}
\loigiai{
$25000 \times 4 = 100000$.
}
\end{ex}

\begin{ex}
Tính: $11111 \times 9 = \dots$
\shortans{99999}
\loigiai{
$11111 \times 9 = 99999$.
}
\end{ex}

\begin{ex}
Tính: $13250 \times 3 = \dots$
\shortans{39750}
\loigiai{
$13250 \times 3 = 39750$.
}
\end{ex}

\begin{ex}
Mỗi chuyến xe chở được 12 500 kg hàng. Hỏi 4 chuyến xe như thế chở được bao nhiêu ki-lô-gam hàng?
\shortans{50000}
\loigiai{
$12500 \times 4 = 50000$ kg.
}
\end{ex}

\begin{ex}
Một công ty mỗi tháng tiết kiệm được 15 000 nghìn đồng. Sau 6 tháng công ty tiết kiệm được bao nhiêu nghìn đồng?
\shortans{90000}
\loigiai{
$15000 \times 6 = 90000$ nghìn đồng.
}
\end{ex}

\begin{ex}
Một hình vuông có cạnh dài 15 000 m. Chu vi hình vuông đó là bao nhiêu mét?
\shortans{60000}
\loigiai{
$15000 \times 4 = 60000$ m.
}
\end{ex}

\begin{ex}
Tính giá trị biểu thức: $12000 \times 3 + 24000$.
\shortans{60000}
\loigiai{
$36000 + 24000 = 60000$.
}
\end{ex}

\begin{ex}
Tính giá trị biểu thức: $18000 \times 4 - 32000$.
\shortans{40000}
\loigiai{
$72000 - 32000 = 40000$.
}
\end{ex}

\begin{ex}
Tìm $x$, biết: $x : 3 = 14200$.
\shortans{42600}
\loigiai{
$x = 14200 \times 3 = 42600$.
}
\end{ex}

\begin{ex}
Tìm $x$, biết: $x : 4 = 21000$.
\shortans{84000}
\loigiai{
$x = 21000 \times 4 = 84000$.
}
\end{ex}

\begin{ex}
Mỗi bao xi măng nặng 50 kg. Một đoàn tàu chở 1200 bao xi măng. Tổng khối lượng xi măng là bao nhiêu ki-lô-gam?
\shortans{60000}
\loigiai{
$50 \times 1200 = 60000$ kg.
}
\end{ex}

\begin{ex}
Một huyện có 3 xã vùng cao, mỗi xã được cấp 18 000 cuốn vở. Cả 3 xã được cấp bao nhiêu cuốn vở?
\shortans{54000}
\loigiai{
$18000 \times 3 = 54000$ cuốn.
}
\end{ex}

\begin{ex}
Tính: $10500 \times 6 = \dots$
\shortans{63000}
\loigiai{
$10500 \times 6 = 63000$.
}
\end{ex}

\begin{ex}
Tính: $16200 \times 5 = \dots$
\shortans{81000}
\loigiai{
$16200 \times 5 = 81000$.
}
\end{ex}

\newpage

\TaoPhieuBaiTap{Chủ đề 14 -- Nhân, chia trong phạm vi 100 000}{71}{Chia số có năm chữ số cho số có một chữ số}{https://loptoanthaylinh.vn/bai-tap/toan-3/bai-71}
\begin{ex}
Tính: $48264 : 2 = \dots$
\shortans{24132}
\loigiai{
$48264 : 2 = 24132$.
}
\end{ex}

\begin{ex}
Tính: $69369 : 3 = \dots$
\shortans{23123}
\loigiai{
$69369 : 3 = 23123$.
}
\end{ex}

\begin{ex}
Tính: $84488 : 4 = \dots$
\shortans{21122}
\loigiai{
$84488 : 4 = 21122$.
}
\end{ex}

\begin{ex}
Tính: $55550 : 5 = \dots$
\shortans{11110}
\loigiai{
$55550 : 5 = 11110$.
}
\end{ex}

\begin{ex}
Tính: $64820 : 2 = \dots$
\shortans{32410}
\loigiai{
$64820 : 2 = 32410$.
}
\end{ex}

\begin{ex}
Tính: $75350 : 5 = \dots$
\shortans{15070}
\loigiai{
$75350 : 5 = 15070$.
}
\end{ex}

\begin{ex}
Tính: $60000 : 3 = \dots$
\shortans{20000}
\loigiai{
$60000 : 3 = 20000$.
}
\end{ex}

\begin{ex}
Tính: $80000 : 4 = \dots$
\shortans{20000}
\loigiai{
$80000 : 4 = 20000$.
}
\end{ex}

\begin{ex}
Tính: $100000 : 2 = \dots$
\shortans{50000}
\loigiai{
$100000 : 2 = 50000$.
}
\end{ex}

\begin{ex}
Tính: $100000 : 5 = \dots$
\shortans{20000}
\loigiai{
$100000 : 5 = 20000$.
}
\end{ex}

\begin{ex}
Tính: $96000 : 3 = \dots$
\shortans{32000}
\loigiai{
$96000 : 3 = 32000$.
}
\end{ex}

\begin{ex}
Tính: $84000 : 6 = \dots$
\shortans{14000}
\loigiai{
$84000 : 6 = 14000$.
}
\end{ex}

\begin{ex}
Tính: $91000 : 7 = \dots$
\shortans{13000}
\loigiai{
$91000 : 7 = 13000$.
}
\end{ex}

\begin{ex}
Tính: $72000 : 8 = \dots$
\shortans{9000}
\loigiai{
$72000 : 8 = 9000$.
}
\end{ex}

\begin{ex}
Tính: $81000 : 9 = \dots$
\shortans{9000}
\loigiai{
$81000 : 9 = 9000$.
}
\end{ex}

\begin{ex}
Có 48 000 lít dầu chia đều vào 4 bể. Mỗi bể chứa bao nhiêu lít dầu?
\shortans{12000}
\loigiai{
$48000 : 4 = 12000$ lít.
}
\end{ex}

\begin{ex}
Một cuộn dây cáp dài 65 000 m cắt đều thành 5 đoạn. Mỗi đoạn dài bao nhiêu mét?
\shortans{13000}
\loigiai{
$65000 : 5 = 13000$ m.
}
\end{ex}

\begin{ex}
Một nông trường thu hoạch được 72 000 kg cam và chia đều vào 6 xe tải. Mỗi xe chở bao nhiêu ki-lô-gam cam?
\shortans{12000}
\loigiai{
$72000 : 6 = 12000$ kg.
}
\end{ex}

\begin{ex}
Tìm $x$, biết: $x \times 4 = 48480$.
\shortans{12120}
\loigiai{
$x = 48480 : 4 = 12120$.
}
\end{ex}

\begin{ex}
Tìm $x$, biết: $3 \times x = 96360$.
\shortans{32120}
\loigiai{
$x = 96360 : 3 = 32120$.
}
\end{ex}

\begin{ex}
Tính giá trị biểu thức: $64000 : 4 + 14000$.
\shortans{30000}
\loigiai{
$16000 + 14000 = 30000$.
}
\end{ex}

\begin{ex}
Tính giá trị biểu thức: $85000 : 5 - 7000$.
\shortans{10000}
\loigiai{
$17000 - 7000 = 10000$.
}
\end{ex}

\begin{ex}
Một mảnh đất hình vuông có chu vi 44 800 m. Mỗi cạnh mảnh đất dài bao nhiêu mét?
\shortans{11200}
\loigiai{
$44800 : 4 = 11200$ m.
}
\end{ex}

\begin{ex}
Tìm số dư trong phép chia: $50003 : 5$.
\shortans{3}
\loigiai{
$50003 = 5 \times 10000 + 3$, dư 3.
}
\end{ex}

\begin{ex}
Tính: $99999 : 9 = \dots$
\shortans{11111}
\loigiai{
$99999 : 9 = 11111$.
}
\end{ex}

\newpage

\TaoPhieuBaiTap{Chủ đề 14 -- Nhân, chia trong phạm vi 100 000}{72}{Luyện tập chung nhân chia 100 000}{https://loptoanthaylinh.vn/bai-tap/toan-3/bai-72}
\begin{ex}
Tính: $12000 \times 4 = \dots$
\shortans{48000}
\loigiai{
$12000 \times 4 = 48000$.
}
\end{ex}

\begin{ex}
Tính: $15000 \times 3 = \dots$
\shortans{45000}
\loigiai{
$15000 \times 3 = 45000$.
}
\end{ex}

\begin{ex}
Tính: $24000 \times 2 = \dots$
\shortans{48000}
\loigiai{
$24000 \times 2 = 48000$.
}
\end{ex}

\begin{ex}
Tính: $68000 : 2 = \dots$
\shortans{34000}
\loigiai{
$68000 : 2 = 34000$.
}
\end{ex}

\begin{ex}
Tính: $96000 : 3 = \dots$
\shortans{32000}
\loigiai{
$96000 : 3 = 32000$.
}
\end{ex}

\begin{ex}
Tính: $85000 : 5 = \dots$
\shortans{17000}
\loigiai{
$85000 : 5 = 17000$.
}
\end{ex}

\begin{ex}
Tính giá trị biểu thức: $14000 \times 3 + 18000$.
\shortans{60000}
\loigiai{
$42000 + 18000 = 60000$.
}
\end{ex}

\begin{ex}
Tính giá trị biểu thức: $75000 : 3 - 15000$.
\shortans{10000}
\loigiai{
$25000 - 15000 = 10000$.
}
\end{ex}

\begin{ex}
Tính giá trị biểu thức: $(20000 + 30000) \times 2$.
\shortans{100000}
\loigiai{
$50000 \times 2 = 100000$.
}
\end{ex}

\begin{ex}
Tính giá trị biểu thức: $(90000 - 30000) : 3$.
\shortans{20000}
\loigiai{
$60000 : 3 = 20000$.
}
\end{ex}

\begin{ex}
Tìm $x$: $x \times 5 = 75000$.
\shortans{15000}
\loigiai{
$x = 75000 : 5 = 15000$.
}
\end{ex}

\begin{ex}
Tìm $x$: $x : 4 = 16000$.
\shortans{64000}
\loigiai{
$x = 16000 \times 4 = 64000$.
}
\end{ex}

\begin{ex}
Một cửa hàng nhập về 3 xe hàng, mỗi xe chở 14 000 kg xi măng. Cửa hàng nhập tất cả bao nhiêu ki-lô-gam xi măng?
\shortans{42000}
\loigiai{
$14000 \times 3 = 42000$ kg.
}
\end{ex}

\begin{ex}
Có 60 000 gói mì tôm chia đều vào 5 thùng lớn. Mỗi thùng có bao nhiêu gói mì?
\shortans{12000}
\loigiai{
$60000 : 5 = 12000$ gói.
}
\end{ex}

\begin{ex}
Một hình vuông có cạnh 18 000 m. Chu vi hình vuông đó là bao nhiêu mét?
\shortans{72000}
\loigiai{
$18000 \times 4 = 72000$ m.
}
\end{ex}

\begin{ex}
Một hình vuông có chu vi 88 000 m. Độ dài mỗi cạnh của hình vuông là bao nhiêu mét?
\shortans{22000}
\loigiai{
$88000 : 4 = 22000$ m.
}
\end{ex}

\begin{ex}
Tính số dư trong phép chia: $65004 : 5$.
\shortans{4}
\loigiai{
$65004 = 5 \times 13000 + 4$, dư 4.
}
\end{ex}

\begin{ex}
Tìm thương trong phép chia: $84000 : 4$.
\shortans{21000}
\loigiai{
$84000 : 4 = 21000$.
}
\end{ex}

\begin{ex}
Tính: $25000 \times 2 + 50000 = \dots$
\shortans{100000}
\loigiai{
$50000 + 50000 = 100000$.
}
\end{ex}

\begin{ex}
Tính: $100000 : 4 = \dots$
\shortans{25000}
\loigiai{
$100000 : 4 = 25000$.
}
\end{ex}

\begin{ex}
Một người nuôi 4 chuồng vịt, mỗi chuồng có 11 500 con. Người đó có tất cả bao nhiêu con vịt?
\shortans{46000}
\loigiai{
$11500 \times 4 = 46000$ con.
}
\end{ex}

\begin{ex}
Có 72 000 viên gạch lát nền chia đều vào 8 khu nhà. Mỗi khu nhà nhận được bao nhiêu viên gạch?
\shortans{9000}
\loigiai{
$72000 : 8 = 9000$ viên.
}
\end{ex}

\begin{ex}
Tính: $15000 \times 4 - 20000 = \dots$
\shortans{40000}
\loigiai{
$60000 - 20000 = 40000$.
}
\end{ex}

\begin{ex}
Tính: $80000 : 8 \times 5 = \dots$
\shortans{50000}
\loigiai{
$10000 \times 5 = 50000$.
}
\end{ex}

\begin{ex}
Tính: $11111 \times 5 = \dots$
\shortans{55555}
\loigiai{
$11111 \times 5 = 55555$.
}
\end{ex}

\newpage


% =========================================================================
% Làm quen với yếu tố thống kê, xác suất
% =========================================================================
\section{Làm quen với yếu tố thống kê, xác suất}

\TaoPhieuBaiTap{Chủ đề 15 -- Làm quen với yếu tố thống kê, xác suất}{73}{Thu thập, phân loại, ghi chép số liệu}{https://loptoanthaylinh.vn/bai-tap/toan-3/bai-73}
\begin{ex}
Lớp 3A có 15 bạn thích màu đỏ, 12 bạn thích màu xanh, 8 bạn thích màu vàng. Hỏi có bao nhiêu bạn thích màu đỏ?
\shortans{15}
\loigiai{
Có 15 bạn thích màu đỏ.
}
\end{ex}

\begin{ex}
Theo số liệu trên, tổng số học sinh thích cả 3 màu là bao nhiêu?
\shortans{35}
\loigiai{
$15 + 12 + 8 = 35$ bạn.
}
\end{ex}

\begin{ex}
Số bạn thích màu đỏ nhiều hơn số bạn thích màu vàng bao nhiêu bạn?
\shortans{7}
\loigiai{
$15 - 8 = 7$ bạn.
}
\end{ex}

\begin{ex}
Trong một cuộc khảo sát về thú cưng: 18 bạn thích chó, 14 bạn thích mèo, 6 bạn thích thỏ. Thú cưng nào được nhiều bạn thích nhất? (Chó ghi 1, mèo ghi 2, thỏ ghi 3).
\shortans{1}
\loigiai{
Chó có 18 bạn thích (nhiều nhất).
}
\end{ex}

\begin{ex}
Tổng số bạn tham gia khảo sát thú cưng trên là bao nhiêu?
\shortans{38}
\loigiai{
$18 + 14 + 6 = 38$ bạn.
}
\end{ex}

\begin{ex}
Một tuần bạn Lan ghi chép số trang sách đọc được mỗi ngày: Thứ Hai: 10, Thứ Ba: 12, Thứ Tư: 15, Thứ Năm: 8, Thứ Sáu: 15. Ngày nào Lan đọc được ít trang sách nhất? (Ghi số thứ, ví dụ Thứ Năm ghi 5).
\shortans{5}
\loigiai{
Thứ Năm (8 trang, ít nhất).
}
\end{ex}

\begin{ex}
Tổng số trang sách Lan đọc được từ Thứ Hai đến Thứ Sáu là bao nhiêu trang?
\shortans{60}
\loigiai{
$10 + 12 + 15 + 8 + 15 = 60$ trang.
}
\end{ex}

\begin{ex}
Số trang sách Lan đọc nhiều nhất trong một ngày là bao nhiêu?
\shortans{15}
\loigiai{
Nhiều nhất là 15 trang.
}
\end{ex}

\begin{ex}
Bảng theo dõi số cây trồng của 4 tổ: Tổ 1: 25 cây, Tổ 2: 30 cây, Tổ 3: 28 cây, Tổ 4: 32 cây. Tổ trồng nhiều cây nhất là tổ mấy?
\shortans{4}
\loigiai{
Tổ 4 trồng được 32 cây.
}
\end{ex}

\begin{ex}
Tổng số cây cả 4 tổ trồng được là bao nhiêu cây?
\shortans{115}
\loigiai{
$25 + 30 + 28 + 32 = 115$ cây.
}
\end{ex}

\begin{ex}
Tổ 4 trồng nhiều hơn Tổ 1 bao nhiêu cây?
\shortans{7}
\loigiai{
$32 - 25 = 7$ cây.
}
\end{ex}

\begin{ex}
Bảng ghi số học sinh giỏi của 3 lớp: Lớp 3A: 12 bạn, Lớp 3B: 15 bạn, Lớp 3C: 10 bạn. Cả 3 lớp có bao nhiêu học sinh giỏi?
\shortans{37}
\loigiai{
$12 + 15 + 10 = 37$ bạn.
}
\end{ex}

\begin{ex}
Lớp có số học sinh giỏi ít nhất có bao nhiêu bạn?
\shortans{10}
\loigiai{
Lớp 3C có 10 bạn.
}
\end{ex}

\begin{ex}
Số học sinh giỏi lớp 3B nhiều hơn lớp 3C bao nhiêu bạn?
\shortans{5}
\loigiai{
$15 - 10 = 5$ bạn.
}
\end{ex}

\begin{ex}
Số lượng bóng các màu trong rổ: Đỏ: 8 quả, Xanh: 12 quả, Vàng: 10 quả. Tổng số bóng trong rổ là bao nhiêu quả?
\shortans{30}
\loigiai{
$8 + 12 + 10 = 30$ quả.
}
\end{ex}

\begin{ex}
Số bóng xanh gấp số bóng đỏ bao nhiêu lần? (Viết dưới dạng phân số tối giản hoặc tỉ số: lấy số bóng xanh chia cho 4 được mấy?).
\shortans{3}
\loigiai{
$12 : 4 = 3$.
}
\end{ex}

\begin{ex}
Hiệu giữa số bóng xanh và số bóng đỏ là bao nhiêu quả?
\shortans{4}
\loigiai{
$12 - 8 = 4$ quả.
}
\end{ex}

\begin{ex}
Số bông hoa 4 bạn hái được: Mai: 9 bông, Lan: 12 bông, Cúc: 8 bông, Đào: 11 bông. Ai hái được ít hoa nhất? (Mai ghi 1, Lan ghi 2, Cúc ghi 3, Đào ghi 4).
\shortans{3}
\loigiai{
Cúc hái được 8 bông (ít nhất).
}
\end{ex}

\begin{ex}
Tổng số hoa 4 bạn hái được là bao nhiêu bông?
\shortans{40}
\loigiai{
$9 + 12 + 8 + 11 = 40$ bông.
}
\end{ex}

\begin{ex}
Số điểm 10 của bạn Minh trong tuần: Thứ Hai: 2, Thứ Ba: 3, Thứ Tư: 1, Thứ Năm: 2, Thứ Sáu: 4. Tổng số điểm 10 Minh đạt được là bao nhiêu?
\shortans{12}
\loigiai{
$2 + 3 + 1 + 2 + 4 = 12$ điểm 10.
}
\end{ex}

\begin{ex}
Ngày Minh đạt nhiều điểm 10 nhất có bao nhiêu điểm 10?
\shortans{4}
\loigiai{
Có 4 điểm 10 vào thứ Sáu.
}
\end{ex}

\begin{ex}
Theo bảng số liệu: Khối 1 có 120 học sinh, Khối 2 có 130 học sinh, Khối 3 có 140 học sinh. Cả 3 khối có bao nhiêu học sinh?
\shortans{390}
\loigiai{
$120 + 130 + 140 = 390$ học sinh.
}
\end{ex}

\begin{ex}
Khối 3 nhiều hơn Khối 1 bao nhiêu học sinh?
\shortans{20}
\loigiai{
$140 - 120 = 20$ học sinh.
}
\end{ex}

\begin{ex}
Bảng số liệu bán hàng ngày Chủ nhật: Bán được 45 cái bút, 60 quyển vở, 20 cục tẩy. Tổng số món đồ đã bán là bao nhiêu?
\shortans{125}
\loigiai{
$45 + 60 + 20 = 125$ món.
}
\end{ex}

\begin{ex}
Số quyển vở bán được nhiều hơn số cái bút bao nhiêu cái?
\shortans{15}
\loigiai{
$60 - 45 = 15$.
}
\end{ex}

\newpage

\TaoPhieuBaiTap{Chủ đề 15 -- Làm quen với yếu tố thống kê, xác suất}{74}{Khả năng xảy ra của một sự kiện}{https://loptoanthaylinh.vn/bai-tap/toan-3/bai-74}
\begin{ex}
Gieo một con xúc xắc 6 mặt (từ 1 đến 6 chấm). Có thể xuất hiện mặt 7 chấm không? (Có ghi 1, Không ghi 0).
\shortans{0}
\loigiai{
Xúc xắc chỉ có từ 1 đến 6 chấm nên không thể có 7 chấm.
}
\end{ex}

\begin{ex}
Gieo một đồng xu có hai mặt Sấp (S) và Ngửa (N). Có bao nhiêu khả năng có thể xảy ra về mặt xuất hiện?
\shortans{2}
\loigiai{
Có 2 khả năng: mặt Sấp hoặc mặt Ngửa.
}
\end{ex}

\begin{ex}
Trong hộp chỉ có 5 viên bi đỏ. Lấy ngẫu nhiên 1 viên bi thì chắc chắn lấy được bi màu gì? (Bi đỏ ghi 1, Bi xanh ghi 2).
\shortans{1}
\loigiai{
Chắc chắn lấy được bi đỏ.
}
\end{ex}

\begin{ex}
Trong hộp có 3 viên bi đỏ và 3 viên bi xanh. Lấy ngẫu nhiên 1 viên bi, có thể lấy được viên bi vàng không? (Có ghi 1, Không thể ghi 0).
\shortans{0}
\loigiai{
Không thể lấy được bi vàng.
}
\end{ex}

\begin{ex}
Trong hộp có 4 viên bi xanh và 2 viên bi đỏ. Sự kiện 'Lấy được viên bi xanh' là: (Chắc chắn ghi 1, Có thể ghi 2, Không thể ghi 3).
\shortans{2}
\loigiai{
Có thể lấy được bi xanh (hoặc bi đỏ).
}
\end{ex}

\begin{ex}
Mặt trời mọc ở hướng Đông vào mỗi sáng là sự kiện gì? (Chắc chắn ghi 1, Có thể ghi 2, Không thể ghi 3).
\shortans{1}
\loigiai{
Chắc chắn xảy ra.
}
\end{ex}

\begin{ex}
Một con cá biết bay trên trời như chim là sự kiện gì? (Chắc chắn ghi 1, Có thể ghi 2, Không thể ghi 3).
\shortans{3}
\loigiai{
Không thể xảy ra.
}
\end{ex}

\begin{ex}
Ngày mai trời có mưa là sự kiện gì? (Chắc chắn ghi 1, Có thể ghi 2, Không thể ghi 3).
\shortans{2}
\loigiai{
Có thể xảy ra.
}
\end{ex}

\begin{ex}
Một hộp có 10 chiếc kẹo, trong đó có 10 chiếc kẹo cam. Lấy ngẫu nhiên 1 chiếc kẹo thì chắc chắn lấy được kẹo gì? (Kẹo cam ghi 1, kẹo dâu ghi 2).
\shortans{1}
\loigiai{
Chắc chắn là kẹo cam.
}
\end{ex}

\begin{ex}
Tung một con xúc xắc, sự kiện 'Số chấm xuất hiện nhỏ hơn 7' là sự kiện gì? (Chắc chắn ghi 1, Có thể ghi 2, Không thể ghi 3).
\shortans{1}
\loigiai{
Chắc chắn, vì xúc xắc có từ 1 đến 6 chấm (đều < 7).
}
\end{ex}

\begin{ex}
Tung một con xúc xắc, sự kiện 'Số chấm xuất hiện bằng 0' là sự kiện gì? (Chắc chắn ghi 1, Có thể ghi 2, Không thể ghi 3).
\shortans{3}
\loigiai{
Không thể, xúc xắc không có mặt 0 chấm.
}
\end{ex}

\begin{ex}
Tung một con xúc xắc, sự kiện 'Số chấm xuất hiện là số chẵn' là sự kiện gì? (Chắc chắn ghi 1, Có thể ghi 2, Không thể ghi 3).
\shortans{2}
\loigiai{
Có thể xảy ra (mặt 2, 4, hoặc 6).
}
\end{ex}

\begin{ex}
Có bao nhiêu mặt là số chẵn trên một con xúc xắc 6 mặt?
\shortans{3}
\loigiai{
Có 3 mặt chẵn: 2, 4, 6.
}
\end{ex}

\begin{ex}
Có bao nhiêu mặt là số lẻ trên một con xúc xắc 6 mặt?
\shortans{3}
\loigiai{
Có 3 mặt lẻ: 1, 3, 5.
}
\end{ex}

\begin{ex}
Trong túi có 5 quả bóng xanh và 5 quả bóng đỏ. Không nhìn vào túi, phải lấy ra ít nhất bao nhiêu quả bóng để chắc chắn có ít nhất 1 quả bóng đỏ?
\shortans{6}
\loigiai{
Trường hợp xấu nhất lấy phải 5 quả bóng xanh, cần lấy quả thứ 6 để chắc chắn có bóng đỏ.
}
\end{ex}

\begin{ex}
Trong hộp có 4 viên bi vàng và 6 viên bi đỏ. Cần lấy ra ít nhất bao nhiêu viên bi để chắc chắn có ít nhất 1 viên bi vàng?
\shortans{7}
\loigiai{
Trường hợp xấu nhất lấy hết 6 viên bi đỏ, lấy thêm 1 viên nữa: $6 + 1 = 7$ viên.
}
\end{ex}

\begin{ex}
Trong hộp có 10 viên bi xanh. Lấy ra 2 viên bi thì có thể có viên bi đỏ nào không? (Có ghi 1, Không ghi 0).
\shortans{0}
\loigiai{
Không thể có bi đỏ.
}
\end{ex}

\begin{ex}
Quay một vòng quay chia thành 4 phần bằng nhau có ghi số 1, 2, 3, 4. Có bao nhiêu khả năng mũi tên dừng lại ở một ô số?
\shortans{4}
\loigiai{
Có 4 khả năng.
}
\end{ex}

\begin{ex}
Vòng quay có 4 số: 1, 2, 3, 4. Mũi tên có thể chỉ vào số 5 không? (Có ghi 1, Không ghi 0).
\shortans{0}
\loigiai{
Không thể chỉ vào số 5.
}
\end{ex}

\begin{ex}
Rút ngẫu nhiên một tấm thẻ từ tập thẻ đánh số từ 1 đến 10. Có thể rút được thẻ ghi số 8 không? (Có ghi 1, Không ghi 0).
\shortans{1}
\loigiai{
Có thể rút được thẻ số 8.
}
\end{ex}

\begin{ex}
Rút ngẫu nhiên một thẻ từ tập thẻ từ 1 đến 10. Sự kiện rút được thẻ số 12 là sự kiện gì? (Chắc chắn ghi 1, Có thể ghi 2, Không thể ghi 3).
\shortans{3}
\loigiai{
Không thể xảy ra.
}
\end{ex}

\begin{ex}
Rút ngẫu nhiên một thẻ từ 1 đến 10. Có bao nhiêu thẻ mang số chẵn?
\shortans{5}
\loigiai{
Các số 2, 4, 6, 8, 10 (5 thẻ).
}
\end{ex}

\begin{ex}
Rút ngẫu nhiên một thẻ từ 1 đến 10. Có bao nhiêu thẻ mang số lẻ?
\shortans{5}
\loigiai{
Các số 1, 3, 5, 7, 9 (5 thẻ).
}
\end{ex}

\begin{ex}
Bạn Nam tung đồng xu 10 lần thì số lần mặt Ngửa xuất hiện có thể lớn hơn 10 không? (Có ghi 1, Không ghi 0).
\shortans{0}
\loigiai{
Tung 10 lần thì tối đa chỉ có 10 lần ngửa.
}
\end{ex}

\begin{ex}
Hai bạn chơi oẳn tù tì. Trong một ván chơi có mấy kết quả có thể xảy ra đối với bạn thứ nhất? (Thắng, thua, hòa).
\shortans{3}
\loigiai{
Có 3 kết quả: Thắng, Thua hoặc Hòa.
}
\end{ex}

\newpage

\TaoPhieuBaiTap{Chủ đề 15 -- Làm quen với yếu tố thống kê, xác suất}{75}{Thực hành trải nghiệm thống kê}{https://loptoanthaylinh.vn/bai-tap/toan-3/bai-75}
\begin{ex}
Kiểm đếm số xe đi qua cổng trường trong 10 phút: 12 xe đạp, 25 xe máy, 8 xe ô tô. Có bao nhiêu xe máy đi qua?
\shortans{25}
\loigiai{
Có 25 xe máy.
}
\end{ex}

\begin{ex}
Tổng số xe các loại đi qua cổng trường trong 10 phút trên là bao nhiêu xe?
\shortans{45}
\loigiai{
$12 + 25 + 8 = 45$ xe.
}
\end{ex}

\begin{ex}
Loại phương tiện nào đi qua ít nhất? (Xe đạp ghi 1, xe máy ghi 2, ô tô ghi 3).
\shortans{3}
\loigiai{
Ô tô có 8 xe (ít nhất).
}
\end{ex}

\begin{ex}
Xe máy nhiều hơn xe đạp bao nhiêu chiếc?
\shortans{13}
\loigiai{
$25 - 12 = 13$ xe.
}
\end{ex}

\begin{ex}
Một nhóm học sinh đo chiều cao của 4 bạn: An: 132 cm, Bình: 135 cm, Cường: 130 cm, Dũng: 138 cm. Bạn cao nhất cao bao nhiêu xăng-ti-mét?
\shortans{138}
\loigiai{
Dũng cao nhất: 138 cm.
}
\end{ex}

\begin{ex}
Bạn thấp nhất cao bao nhiêu xăng-ti-mét?
\shortans{130}
\loigiai{
Cường thấp nhất: 130 cm.
}
\end{ex}

\begin{ex}
Bạn cao nhất cao hơn bạn thấp nhất bao nhiêu xăng-ti-mét?
\shortans{8}
\loigiai{
$138 - 130 = 8$ cm.
}
\end{ex}

\begin{ex}
Trong 4 bạn trên, có bao nhiêu bạn có chiều cao từ 132 cm trở lên?
\shortans{3}
\loigiai{
An (132), Bình (135), Dũng (138): có 3 bạn.
}
\end{ex}

\begin{ex}
Kiểm đếm số lượng hoa quả mẹ mua: 10 quả cam, 15 quả quýt, 5 quả xoài. Quýt nhiều hơn xoài bao nhiêu quả?
\shortans{10}
\loigiai{
$15 - 5 = 10$ quả.
}
\end{ex}

\begin{ex}
Tổng số quả mẹ mua là bao nhiêu quả?
\shortans{30}
\loigiai{
$10 + 15 + 5 = 30$ quả.
}
\end{ex}

\begin{ex}
Tung đồng xu 20 lần, thấy có 12 lần xuất hiện mặt Ngửa. Hỏi có bao nhiêu lần xuất hiện mặt Sấp?
\shortans{8}
\loigiai{
$20 - 12 = 8$ lần.
}
\end{ex}

\begin{ex}
Gieo xúc xắc 30 lần, thấy có 6 lần xuất hiện mặt 6 chấm. Số lần xuất hiện các mặt còn lại là bao nhiêu lần?
\shortans{24}
\loigiai{
$30 - 6 = 24$ lần.
}
\end{ex}

\begin{ex}
Bảng điểm kiểm tra môn Toán của tổ 1: có 3 điểm 10, 4 điểm 9, 2 điểm 8, 1 điểm 7. Tổ 1 có bao nhiêu học sinh?
\shortans{10}
\loigiai{
$3 + 4 + 2 + 1 = 10$ học sinh.
}
\end{ex}

\begin{ex}
Có bao nhiêu bạn đạt điểm 9 trở lên trong tổ 1?
\shortans{7}
\loigiai{
$3 + 4 = 7$ bạn.
}
\end{ex}

\begin{ex}
Điểm số có nhiều bạn đạt được nhất là điểm mấy?
\shortans{9}
\loigiai{
Điểm 9 có 4 bạn đạt (nhiều nhất).
}
\end{ex}

\begin{ex}
Điểm số có ít bạn đạt được nhất là điểm mấy?
\shortans{7}
\loigiai{
Điểm 7 chỉ có 1 bạn đạt.
}
\end{ex}

\begin{ex}
Theo dõi số ngày mưa trong tháng 7 ở một địa phương: có 12 ngày mưa. Hỏi có bao nhiêu ngày không mưa trong tháng 7?
\shortans{19}
\loigiai{
Tháng 7 có 31 ngày: $31 - 12 = 19$ ngày không mưa.
}
\end{ex}

\begin{ex}
Số lượng sách các môn trong tủ sách lớp: Toán: 45 cuốn, Tiếng Việt: 50 cuốn, Tự nhiên và Xã hội: 30 cuốn. Tổng số sách là bao nhiêu cuốn?
\shortans{125}
\loigiai{
$45 + 50 + 30 = 125$ cuốn.
}
\end{ex}

\begin{ex}
Sách Tiếng Việt nhiều hơn sách Tự nhiên và Xã hội bao nhiêu cuốn?
\shortans{20}
\loigiai{
$50 - 30 = 20$ cuốn.
}
\end{ex}

\begin{ex}
Trong hộp có 8 viên kẹo đỏ và 12 viên kẹo xanh. Lấy ra ngẫu nhiên 1 viên kẹo, có thể lấy được kẹo màu vàng không? (Có ghi 1, Không ghi 0).
\shortans{0}
\loigiai{
Không thể có kẹo vàng.
}
\end{ex}

\begin{ex}
Tung con xúc xắc 1 lần, khả năng xuất hiện mặt có số chấm lớn hơn 6 là mấy? (Ghi 0 nếu không thể, ghi 1 nếu chắc chắn).
\shortans{0}
\loigiai{
Không thể có mặt lớn hơn 6.
}
\end{ex}

\begin{ex}
Một bình chứa 5 quả cầu đỏ, 5 quả cầu xanh. Cần lấy ít nhất mấy quả cầu để chắc chắn có quả cầu xanh?
\shortans{6}
\loigiai{
$5 + 1 = 6$ quả.
}
\end{ex}

\begin{ex}
Đo nhiệt độ lúc 12 giờ trưa trong 3 ngày liên tiếp ghi được: $32^\circ$C, $35^\circ$C, $33^\circ$C. Nhiệt độ cao nhất trong 3 ngày là bao nhiêu độ C?
\shortans{35}
\loigiai{
$35^\circ$C là cao nhất.
}
\end{ex}

\begin{ex}
Chênh lệch giữa ngày nóng nhất và ngày mát nhất trong 3 ngày đó là bao nhiêu độ C?
\shortans{3}
\loigiai{
$35 - 32 = 3^\circ$C.
}
\end{ex}

\begin{ex}
Theo bảng số liệu rác thải nhựa thu gom được: Khối 3: 40 kg, Khối 4: 55 kg, Khối 5: 65 kg. Khối 5 thu gom nhiều hơn Khối 3 bao nhiêu ki-lô-gam?
\shortans{25}
\loigiai{
$65 - 40 = 25$ kg.
}
\end{ex}

\newpage


% =========================================================================
% Ôn tập cuối năm
% =========================================================================
\section{Ôn tập cuối năm}

\TaoPhieuBaiTap{Chủ đề 16 -- Ôn tập cuối năm}{76}{Ôn tập các số trong phạm vi 100 000}{https://loptoanthaylinh.vn/bai-tap/toan-3/bai-76}
\begin{ex}
Số gồm 8 chục nghìn, 5 nghìn, 3 trăm, 4 chục và 2 đơn vị viết là gì?
\shortans{85342}
\loigiai{
Số đó là 85342.
}
\end{ex}

\begin{ex}
Số 74 520 có chữ số hàng chục nghìn là bao nhiêu?
\shortans{7}
\loigiai{
Chữ số hàng chục nghìn là 7.
}
\end{ex}

\begin{ex}
Số liền sau của số 99 999 là số nào?
\shortans{100000}
\loigiai{
$99999 + 1 = 100000$.
}
\end{ex}

\begin{ex}
Số liền trước của số 70 000 là số nào?
\shortans{69999}
\loigiai{
$70000 - 1 = 69999$.
}
\end{ex}

\begin{ex}
Làm tròn số 48 650 đến hàng nghìn ta được số nào?
\shortans{49000}
\loigiai{
Làm tròn lên 49 000.
}
\end{ex}

\begin{ex}
Làm tròn số 73 200 đến hàng chục nghìn ta được số nào?
\shortans{70000}
\loigiai{
Làm tròn thành 70 000.
}
\end{ex}

\begin{ex}
Số La Mã XIX biểu thị số tự nhiên nào?
\shortans{19}
\loigiai{
XIX là 19.
}
\end{ex}

\begin{ex}
Số La Mã XIV biểu thị số tự nhiên nào?
\shortans{14}
\loigiai{
XIV là 14.
}
\end{ex}

\begin{ex}
Trong các số $45600, 46500, 54600, 56400$, số nào bé nhất?
\shortans{45600}
\loigiai{
45600 là số bé nhất.
}
\end{ex}

\begin{ex}
Trong các số $82450, 84250, 82540, 84520$, số nào lớn nhất?
\shortans{84520}
\loigiai{
84520 là số lớn nhất.
}
\end{ex}

\begin{ex}
Số lớn nhất có 5 chữ số khác nhau là số nào?
\shortans{98765}
\loigiai{
Số đó là 98765.
}
\end{ex}

\begin{ex}
Số bé nhất có 5 chữ số khác nhau là số nào?
\shortans{10234}
\loigiai{
Số đó là 10234.
}
\end{ex}

\begin{ex}
Số 55 555 có bao nhiêu chữ số 5?
\shortans{5}
\loigiai{
Có 5 chữ số 5.
}
\end{ex}

\begin{ex}
Viết số thích hợp: $50000 + 7000 + 200 + 4 = \dots$
\shortans{57204}
\loigiai{
Số đó là 57204.
}
\end{ex}

\begin{ex}
Khoảng cách giữa hai số 60 000 và 58 500 là bao nhiêu?
\shortans{1500}
\loigiai{
$60000 - 58500 = 1500$.
}
\end{ex}

\begin{ex}
Có bao nhiêu số tròn chục nghìn từ 10 000 đến 100 000?
\shortans{10}
\loigiai{
10 số.
}
\end{ex}

\begin{ex}
Số La Mã XX tương ứng số tự nhiên nào?
\shortans{20}
\loigiai{
XX là 20.
}
\end{ex}

\begin{ex}
Đồng hồ La Mã chỉ vào số XII nghĩa là mấy giờ?
\shortans{12}
\loigiai{
XII là 12 giờ.
}
\end{ex}

\begin{ex}
Làm tròn số 850 đến hàng trăm được số nào?
\shortans{900}
\loigiai{
Làm tròn lên 900.
}
\end{ex}

\begin{ex}
Làm tròn số 342 đến hàng chục được số nào?
\shortans{340}
\loigiai{
Làm tròn thành 340.
}
\end{ex}

\begin{ex}
Điền số thích hợp: $79999 < \dots < 80001$.
\shortans{80000}
\loigiai{
Số ở giữa là 80000.
}
\end{ex}

\begin{ex}
Giá trị của chữ số 6 trong số 65 432 là bao nhiêu?
\shortans{60000}
\loigiai{
Chữ số 6 ở hàng chục nghìn có giá trị 60 000.
}
\end{ex}

\begin{ex}
Hiệu của số lớn nhất có 5 chữ số và số bé nhất có 5 chữ số là bao nhiêu?
\shortans{89999}
\loigiai{
$99999 - 10000 = 89999$.
}
\end{ex}

\begin{ex}
Số 100 000 đọc là một trăm nghìn, có bao nhiêu chữ số 0?
\shortans{5}
\loigiai{
Có 5 chữ số 0.
}
\end{ex}

\begin{ex}
Tìm $x$ là số tròn nghìn biết $68000 < x < 70000$.
\shortans{69000}
\loigiai{
Số đó là 69000.
}
\end{ex}

\newpage

\TaoPhieuBaiTap{Chủ đề 16 -- Ôn tập cuối năm}{77}{Ôn tập phép cộng, phép trừ trong phạm vi 100 000}{https://loptoanthaylinh.vn/bai-tap/toan-3/bai-77}
\begin{ex}
Tính: $45200 + 34600 = \dots$
\shortans{79800}
\loigiai{
$45200 + 34600 = 79800$.
}
\end{ex}

\begin{ex}
Tính: $56340 + 28190 = \dots$
\shortans{84530}
\loigiai{
$56340 + 28190 = 84530$.
}
\end{ex}

\begin{ex}
Tính: $78650 - 34210 = \dots$
\shortans{44440}
\loigiai{
$78650 - 34210 = 44440$.
}
\end{ex}

\begin{ex}
Tính: $82500 - 36800 = \dots$
\shortans{45700}
\loigiai{
$82500 - 36800 = 45700$.
}
\end{ex}

\begin{ex}
Tính: $100000 - 45000 = \dots$
\shortans{55000}
\loigiai{
$100000 - 45000 = 55000$.
}
\end{ex}

\begin{ex}
Tính: $60000 + 40000 = \dots$
\shortans{100000}
\loigiai{
$60000 + 40000 = 100000$.
}
\end{ex}

\begin{ex}
Tính giá trị biểu thức: $35000 + 25000 - 40000$.
\shortans{20000}
\loigiai{
$60000 - 40000 = 20000$.
}
\end{ex}

\begin{ex}
Tính giá trị biểu thức: $100000 - (35000 + 25000)$.
\shortans{40000}
\loigiai{
$100000 - 60000 = 40000$.
}
\end{ex}

\begin{ex}
Tìm $x$: $x + 32000 = 80000$.
\shortans{48000}
\loigiai{
$x = 80000 - 32000 = 48000$.
}
\end{ex}

\begin{ex}
Tìm $x$: $x - 45000 = 35000$.
\shortans{80000}
\loigiai{
$x = 35000 + 45000 = 80000$.
}
\end{ex}

\begin{ex}
Tìm $x$: $90000 - x = 25000$.
\shortans{65000}
\loigiai{
$x = 90000 - 25000 = 65000$.
}
\end{ex}

\begin{ex}
Một huyện năm ngoái thu hoạch 42 500 tấn thóc, năm nay thu hoạch được 48 300 tấn. Cả hai năm huyện thu hoạch bao nhiêu tấn thóc?
\shortans{90800}
\loigiai{
$42500 + 48300 = 90800$ tấn.
}
\end{ex}

\begin{ex}
Một kho có 75 000 kg đường, đã xuất bán 28 500 kg. Kho còn lại bao nhiêu ki-lô-gam đường?
\shortans{46500}
\loigiai{
$75000 - 28500 = 46500$ kg.
}
\end{ex}

\begin{ex}
Một công ty dự định sản xuất 80 000 sản phẩm, hiện tại đã làm được 52 000 sản phẩm. Cần làm thêm bao nhiêu sản phẩm nữa?
\shortans{28000}
\loigiai{
$80000 - 52000 = 28000$ sản phẩm.
}
\end{ex}

\begin{ex}
Tính nhanh: $27000 + 45000 + 23000$.
\shortans{95000}
\loigiai{
$27000 + 23000 + 45000 = 50000 + 45000 = 95000$.
}
\end{ex}

\begin{ex}
Tính: $99999 - 44444 = \dots$
\shortans{55555}
\loigiai{
$99999 - 44444 = 55555$.
}
\end{ex}

\begin{ex}
Tổng của số tròn chục nghìn bé nhất và lớn nhất có năm chữ số là bao nhiêu?
\shortans{100000}
\loigiai{
$10000 + 90000 = 100000$.
}
\end{ex}

\begin{ex}
Hiệu của 100 000 và 50 000 là bao nhiêu?
\shortans{50000}
\loigiai{
$100000 - 50000 = 50000$.
}
\end{ex}

\begin{ex}
Một sân bay ngày thứ nhất đón 35 400 lượt khách, ngày thứ hai đón 42 600 lượt khách. Cả hai ngày đón bao nhiêu lượt khách?
\shortans{78000}
\loigiai{
$35400 + 42600 = 78000$ lượt khách.
}
\end{ex}

\begin{ex}
Bác Nam có 65 000 nghìn đồng tiết kiệm, bác rút ra 20 000 nghìn đồng để sửa nhà. Bác còn lại bao nhiêu nghìn đồng?
\shortans{45000}
\loigiai{
$65000 - 20000 = 45000$ nghìn đồng.
}
\end{ex}

\begin{ex}
Tính: $38400 + 21600 = \dots$
\shortans{60000}
\loigiai{
$38400 + 21600 = 60000$.
}
\end{ex}

\begin{ex}
Tính: $72500 - 22500 = \dots$
\shortans{50000}
\loigiai{
$72500 - 22500 = 50000$.
}
\end{ex}

\begin{ex}
Tìm $x$: $x + 15500 = 60000$.
\shortans{44500}
\loigiai{
$x = 60000 - 15500 = 44500$.
}
\end{ex}

\begin{ex}
Tìm $x$: $x - 18000 = 52000$.
\shortans{70000}
\loigiai{
$x = 52000 + 18000 = 70000$.
}
\end{ex}

\begin{ex}
Tính: $100000 - 38000 = \dots$
\shortans{62000}
\loigiai{
$100000 - 38000 = 62000$.
}
\end{ex}

\newpage

\TaoPhieuBaiTap{Chủ đề 16 -- Ôn tập cuối năm}{78}{Ôn tập phép nhân, phép chia trong phạm vi 100 000}{https://loptoanthaylinh.vn/bai-tap/toan-3/bai-78}
\begin{ex}
Tính: $14200 \times 2 = \dots$
\shortans{28400}
\loigiai{
$14200 \times 2 = 28400$.
}
\end{ex}

\begin{ex}
Tính: $21300 \times 3 = \dots$
\shortans{63900}
\loigiai{
$21300 \times 3 = 63900$.
}
\end{ex}

\begin{ex}
Tính: $12500 \times 4 = \dots$
\shortans{50000}
\loigiai{
$12500 \times 4 = 50000$.
}
\end{ex}

\begin{ex}
Tính: $15000 \times 5 = \dots$
\shortans{75000}
\loigiai{
$15000 \times 5 = 75000$.
}
\end{ex}

\begin{ex}
Tính: $20000 \times 5 = \dots$
\shortans{100000}
\loigiai{
$20000 \times 5 = 100000$.
}
\end{ex}

\begin{ex}
Tính: $48600 : 2 = \dots$
\shortans{24300}
\loigiai{
$48600 : 2 = 24300$.
}
\end{ex}

\begin{ex}
Tính: $69300 : 3 = \dots$
\shortans{23100}
\loigiai{
$69300 : 3 = 23100$.
}
\end{ex}

\begin{ex}
Tính: $84800 : 4 = \dots$
\shortans{21200}
\loigiai{
$84800 : 4 = 21200$.
}
\end{ex}

\begin{ex}
Tính: $65000 : 5 = \dots$
\shortans{13000}
\loigiai{
$65000 : 5 = 13000$.
}
\end{ex}

\begin{ex}
Tính: $100000 : 2 = \dots$
\shortans{50000}
\loigiai{
$100000 : 2 = 50000$.
}
\end{ex}

\begin{ex}
Tính giá trị biểu thức: $12000 \times 4 + 15000$.
\shortans{63000}
\loigiai{
$48000 + 15000 = 63000$.
}
\end{ex}

\begin{ex}
Tính giá trị biểu thức: $80000 : 4 - 8000$.
\shortans{12000}
\loigiai{
$20000 - 8000 = 12000$.
}
\end{ex}

\begin{ex}
Tìm $x$: $x \times 3 = 45000$.
\shortans{15000}
\loigiai{
$x = 45000 : 3 = 15000$.
}
\end{ex}

\begin{ex}
Tìm $x$: $x : 4 = 18000$.
\shortans{72000}
\loigiai{
$x = 18000 \times 4 = 72000$.
}
\end{ex}

\begin{ex}
Mỗi xe tải chở được 15 000 viên gạch. Hỏi 4 xe tải chở được bao nhiêu viên gạch?
\shortans{60000}
\loigiai{
$15000 \times 4 = 60000$ viên.
}
\end{ex}

\begin{ex}
Có 75 000 lít nước chia đều vào 3 bể. Mỗi bể chứa bao nhiêu lít nước?
\shortans{25000}
\loigiai{
$75000 : 3 = 25000$ lít.
}
\end{ex}

\begin{ex}
Một hình vuông có cạnh dài 12 000 m. Chu vi hình vuông đó là bao nhiêu mét?
\shortans{48000}
\loigiai{
$12000 \times 4 = 48000$ m.
}
\end{ex}

\begin{ex}
Tìm số dư trong phép chia: $45003 : 5$.
\shortans{3}
\loigiai{
$45003 = 5 \times 9000 + 3$, dư 3.
}
\end{ex}

\begin{ex}
Tìm thương trong phép chia: $96000 : 3$.
\shortans{32000}
\loigiai{
$96000 : 3 = 32000$.
}
\end{ex}

\begin{ex}
Tính: $25000 \times 3 + 25000 = \dots$
\shortans{100000}
\loigiai{
$75000 + 25000 = 100000$.
}
\end{ex}

\begin{ex}
Tính: $100000 : 5 \times 3 = \dots$
\shortans{60000}
\loigiai{
$20000 \times 3 = 60000$.
}
\end{ex}

\begin{ex}
Có 5 thùng dầu, mỗi thùng chứa 16 000 lít dầu. Tổng lượng dầu là bao nhiêu lít?
\shortans{80000}
\loigiai{
$16000 \times 5 = 80000$ lít.
}
\end{ex}

\begin{ex}
Một kho có 84 000 bao xi măng chia đều cho 6 công trường. Mỗi công trường nhận bao nhiêu bao?
\shortans{14000}
\loigiai{
$84000 : 6 = 14000$ bao.
}
\end{ex}

\begin{ex}
Tính: $11111 \times 7 = \dots$
\shortans{77777}
\loigiai{
$11111 \times 7 = 77777$.
}
\end{ex}

\begin{ex}
Tính: $99999 : 9 = \dots$
\shortans{11111}
\loigiai{
$99999 : 9 = 11111$.
}
\end{ex}

\newpage

\TaoPhieuBaiTap{Chủ đề 16 -- Ôn tập cuối năm}{79}{Ôn tập hình học và đo lường}{https://loptoanthaylinh.vn/bai-tap/toan-3/bai-79}
\begin{ex}
Một hình vuông có cạnh dài 8 cm. Chu vi của hình vuông đó là bao nhiêu xăng-ti-mét?
\shortans{32}
\loigiai{
$8 \times 4 = 32$ cm.
}
\end{ex}

\begin{ex}
Một hình vuông có cạnh dài 8 cm. Diện tích của hình vuông đó là bao nhiêu xăng-ti-mét vuông?
\shortans{64}
\loigiai{
$8 \times 8 = 64\text{ cm}^2$.
}
\end{ex}

\begin{ex}
Một hình chữ nhật có chiều dài 12 cm, chiều rộng 7 cm. Chu vi hình chữ nhật đó là bao nhiêu xăng-ti-mét?
\shortans{38}
\loigiai{
$(12 + 7) \times 2 = 38$ cm.
}
\end{ex}

\begin{ex}
Một hình chữ nhật có chiều dài 12 cm, chiều rộng 7 cm. Diện tích hình chữ nhật đó là bao nhiêu xăng-ti-mét vuông?
\shortans{84}
\loigiai{
$12 \times 7 = 84\text{ cm}^2$.
}
\end{ex}

\begin{ex}
Một hình tròn có bán kính 8 cm. Đường kính của hình tròn đó là bao nhiêu xăng-ti-mét?
\shortans{16}
\loigiai{
$8 \times 2 = 16$ cm.
}
\end{ex}

\begin{ex}
Một hình tròn có đường kính 20 cm. Bán kính của hình tròn đó là bao nhiêu xăng-ti-mét?
\shortans{10}
\loigiai{
$20 : 2 = 10$ cm.
}
\end{ex}

\begin{ex}
Khối lập phương có bao nhiêu mặt?
\shortans{6}
\loigiai{
Có 6 mặt.
}
\end{ex}

\begin{ex}
Khối lập phương có bao nhiêu đỉnh?
\shortans{8}
\loigiai{
Có 8 đỉnh.
}
\end{ex}

\begin{ex}
Khối hộp chữ nhật có bao nhiêu cạnh?
\shortans{12}
\loigiai{
Có 12 cạnh.
}
\end{ex}

\begin{ex}
Đổi: $5\text{ m } 8\text{ cm} = \dots\text{ cm}$.
\shortans{508}
\loigiai{
$500 + 8 = 508$ cm.
}
\end{ex}

\begin{ex}
Đổi: $3\text{ km} = \dots\text{ m}$.
\shortans{3000}
\loigiai{
$3\text{ km} = 3000\text{ m}$.
}
\end{ex}

\begin{ex}
Đổi: $1\text{ kg} = \dots\text{ g}$.
\shortans{1000}
\loigiai{
$1\text{ kg} = 1000\text{ g}$.
}
\end{ex}

\begin{ex}
Đổi: $1\text{ l} = \dots\text{ ml}$.
\shortans{1000}
\loigiai{
$1\text{ l} = 1000\text{ ml}$.
}
\end{ex}

\begin{ex}
Một hình vuông có chu vi 40 cm. Cạnh của hình vuông đó dài bao nhiêu xăng-ti-mét?
\shortans{10}
\loigiai{
$40 : 4 = 10$ cm.
}
\end{ex}

\begin{ex}
Một hình vuông có diện tích $49\text{ cm}^2$. Cạnh của hình vuông đó dài bao nhiêu xăng-ti-mét?
\shortans{7}
\loigiai{
Vì $7 \times 7 = 49$ nên cạnh là 7 cm.
}
\end{ex}

\begin{ex}
Một hình chữ nhật có diện tích $60\text{ cm}^2$, chiều rộng 5 cm. Chiều dài hình chữ nhật đó là bao nhiêu xăng-ti-mét?
\shortans{12}
\loigiai{
$60 : 5 = 12$ cm.
}
\end{ex}

\begin{ex}
Nhiệt độ sôi của nước ở điều kiện chuẩn là bao nhiêu độ C?
\shortans{100}
\loigiai{
$100^\circ$C.
}
\end{ex}

\begin{ex}
Nhiệt độ nước đá đang tan là bao nhiêu độ C?
\shortans{0}
\loigiai{
$0^\circ$C.
}
\end{ex}

\begin{ex}
Một gói bánh nặng 250 g, 4 gói bánh như thế nặng bao nhiêu gam?
\shortans{1000}
\loigiai{
$250 \times 4 = 1000$ g.
}
\end{ex}

\begin{ex}
Chai dầu có 750 ml, dùng hết 250 ml. Chai dầu còn lại bao nhiêu mi-li-lít?
\shortans{500}
\loigiai{
$750 - 250 = 500$ ml.
}
\end{ex}

\begin{ex}
Đoạn thẳng AB dài 18 cm, M là trung điểm của AB. Độ dài đoạn AM là bao nhiêu xăng-ti-mét?
\shortans{9}
\loigiai{
$18 : 2 = 9$ cm.
}
\end{ex}

\begin{ex}
Một góc vuông có số đo là bao nhiêu độ?
\shortans{90}
\loigiai{
Góc vuông bằng $90^\circ$.
}
\end{ex}

\begin{ex}
Chu vi hình tam giác đều cạnh 15 cm là bao nhiêu xăng-ti-mét?
\shortans{45}
\loigiai{
$15 \times 3 = 45$ cm.
}
\end{ex}

\begin{ex}
Chu vi hình tứ giác có 4 cạnh đều bằng 12 cm là bao nhiêu xăng-ti-mét?
\shortans{48}
\loigiai{
$12 \times 4 = 48$ cm.
}
\end{ex}

\begin{ex}
Một mảnh vườn hình chữ nhật có chiều dài 20 m, rộng 10 m. Diện tích mảnh vườn là bao nhiêu mét vuông?
\shortans{200}
\loigiai{
$20 \times 10 = 200\text{ m}^2$.
}
\end{ex}

\newpage

\TaoPhieuBaiTap{Chủ đề 16 -- Ôn tập cuối năm}{80}{Ôn tập bảng số liệu và khả năng xảy ra}{https://loptoanthaylinh.vn/bai-tap/toan-3/bai-80}
\begin{ex}
Số bạn yêu thích môn thể thao của lớp 3A: Bóng đá: 18 bạn, Cầu lông: 12 bạn, Bơi lội: 10 bạn. Môn nào được nhiều bạn thích nhất? (Bóng đá ghi 1, Cầu lông ghi 2, Bơi lội ghi 3).
\shortans{1}
\loigiai{
Bóng đá có 18 bạn.
}
\end{ex}

\begin{ex}
Tổng số học sinh tham gia bình chọn môn thể thao trên là bao nhiêu bạn?
\shortans{40}
\loigiai{
$18 + 12 + 10 = 40$ bạn.
}
\end{ex}

\begin{ex}
Số bạn thích bóng đá nhiều hơn số bạn thích bơi lội bao nhiêu bạn?
\shortans{8}
\loigiai{
$18 - 10 = 8$ bạn.
}
\end{ex}

\begin{ex}
Bảng số cây trồng của 3 lớp: 3A: 45 cây, 3B: 50 cây, 3C: 40 cây. Cả 3 lớp trồng được bao nhiêu cây?
\shortans{135}
\loigiai{
$45 + 50 + 40 = 135$ cây.
}
\end{ex}

\begin{ex}
Lớp trồng được nhiều cây nhất trồng bao nhiêu cây?
\shortans{50}
\loigiai{
Lớp 3B trồng 50 cây.
}
\end{ex}

\begin{ex}
Gieo một con xúc xắc 6 mặt 1 lần. Khả năng xuất hiện mặt 8 chấm có thể xảy ra không? (Có ghi 1, Không ghi 0).
\shortans{0}
\loigiai{
Không thể có mặt 8 chấm.
}
\end{ex}

\begin{ex}
Gieo một đồng xu 1 lần. Có mấy kết quả có thể xảy ra? (Sấp, Ngửa).
\shortans{2}
\loigiai{
Có 2 khả năng.
}
\end{ex}

\begin{ex}
Trong hộp có 6 bi xanh, không có bi nào khác. Lấy ngẫu nhiên 1 viên bi thì chắc chắn lấy được bi màu gì? (Bi xanh ghi 1, Bi đỏ ghi 2).
\shortans{1}
\loigiai{
Chắc chắn là bi xanh.
}
\end{ex}

\begin{ex}
Trong hộp có 5 viên bi đỏ và 5 viên bi vàng. Lấy ngẫu nhiên 1 viên bi, có thể lấy được bi xanh không? (Có ghi 1, Không thể ghi 0).
\shortans{0}
\loigiai{
Không thể lấy được bi xanh.
}
\end{ex}

\begin{ex}
Tung xúc xắc 1 lần, sự kiện xuất hiện mặt có số chấm bé hơn 7 là sự kiện gì? (Chắc chắn ghi 1, Có thể ghi 2, Không thể ghi 3).
\shortans{1}
\loigiai{
Chắc chắn (tất cả các mặt từ 1 đến 6 đều < 7).
}
\end{ex}

\begin{ex}
Tung xúc xắc 1 lần, sự kiện xuất hiện mặt 3 chấm là sự kiện gì? (Chắc chắn ghi 1, Có thể ghi 2, Không thể ghi 3).
\shortans{2}
\loigiai{
Có thể xảy ra.
}
\end{ex}

\begin{ex}
Tung xúc xắc 1 lần, sự kiện xuất hiện mặt 0 chấm là sự kiện gì? (Chắc chắn ghi 1, Có thể ghi 2, Không thể ghi 3).
\shortans{3}
\loigiai{
Không thể xảy ra.
}
\end{ex}

\begin{ex}
Trong hộp có 7 viên bi đỏ và 3 viên bi xanh. Muốn chắc chắn lấy được ít nhất 1 viên bi xanh thì cần lấy ra ít nhất bao nhiêu viên bi?
\shortans{8}
\loigiai{
Trường hợp xấu nhất lấy hết 7 bi đỏ, lấy thêm 1 viên: $7 + 1 = 8$ viên.
}
\end{ex}

\begin{ex}
Trong túi có 6 quả cầu trắng và 4 quả cầu đen. Cần lấy ít nhất mấy quả cầu để chắc chắn có quả cầu trắng?
\shortans{5}
\loigiai{
$4 + 1 = 5$ quả.
}
\end{ex}

\begin{ex}
Số sách đọc trong tháng của 4 bạn: An: 4 cuốn, Bình: 6 cuốn, Chi: 5 cuốn, Dung: 7 cuốn. Ai đọc được nhiều sách nhất? (An: 1, Bình: 2, Chi: 3, Dung: 4).
\shortans{4}
\loigiai{
Dung đọc 7 cuốn (nhiều nhất).
}
\end{ex}

\begin{ex}
Tổng số sách cả 4 bạn đọc được là bao nhiêu cuốn?
\shortans{22}
\loigiai{
$4 + 6 + 5 + 7 = 22$ cuốn.
}
\end{ex}

\begin{ex}
Bình đọc nhiều hơn An bao nhiêu cuốn sách?
\shortans{2}
\loigiai{
$6 - 4 = 2$ cuốn.
}
\end{ex}

\begin{ex}
Tung đồng xu 30 lần, có 16 lần xuất hiện mặt Sấp. Có bao nhiêu lần xuất hiện mặt Ngửa?
\shortans{14}
\loigiai{
$30 - 16 = 14$ lần.
}
\end{ex}

\begin{ex}
Quay vòng quay có 6 ô số từ 1 đến 6. Có bao nhiêu khả năng mũi tên dừng lại ở ô số chẵn?
\shortans{3}
\loigiai{
Có 3 ô số chẵn: 2, 4, 6.
}
\end{ex}

\begin{ex}
Quay vòng quay có 6 ô số từ 1 đến 6. Có bao nhiêu khả năng dừng lại ở ô số lẻ?
\shortans{3}
\loigiai{
Có 3 ô số lẻ: 1, 3, 5.
}
\end{ex}

\begin{ex}
Hai bạn cùng chơi tung xúc xắc. Tổng số chấm lớn nhất có thể nhận được khi tung 2 con xúc xắc là bao nhiêu?
\shortans{12}
\loigiai{
$6 + 6 = 12$ chấm.
}
\end{ex}

\begin{ex}
Tổng số chấm nhỏ nhất có thể nhận được khi tung 2 con xúc xắc là bao nhiêu?
\shortans{2}
\loigiai{
$1 + 1 = 2$ chấm.
}
\end{ex}

\begin{ex}
Kiểm đếm số học sinh đạt điểm 10 môn Tiếng Việt của khối 3: Lớp 3A: 8 bạn, Lớp 3B: 12 bạn, Lớp 3C: 10 bạn. Cả khối có bao nhiêu bạn đạt điểm 10?
\shortans{30}
\loigiai{
$8 + 12 + 10 = 30$ bạn.
}
\end{ex}

\begin{ex}
Lớp 3B nhiều hơn lớp 3A bao nhiêu bạn đạt điểm 10?
\shortans{4}
\loigiai{
$12 - 8 = 4$ bạn.
}
\end{ex}

\begin{ex}
Trong một tuần lễ có bao nhiêu ngày?
\shortans{7}
\loigiai{
Có 7 ngày.
}
\end{ex}

\newpage

\TaoPhieuBaiTap{Chủ đề 16 -- Ôn tập cuối năm}{81}{Ôn tập chung cuối năm}{https://loptoanthaylinh.vn/bai-tap/toan-3/bai-81}
\begin{ex}
Số gồm 6 chục nghìn, 4 nghìn, 5 trăm và 2 đơn vị viết là gì?
\shortans{64502}
\loigiai{
Số đó là 64502.
}
\end{ex}

\begin{ex}
Số liền sau của số 89 999 là số nào?
\shortans{90000}
\loigiai{
$89999 + 1 = 90000$.
}
\end{ex}

\begin{ex}
Số liền trước của số 100 000 là số nào?
\shortans{99999}
\loigiai{
$100000 - 1 = 99999$.
}
\end{ex}

\begin{ex}
Tính: $35000 + 45000 = \dots$
\shortans{80000}
\loigiai{
$35000 + 45000 = 80000$.
}
\end{ex}

\begin{ex}
Tính: $90000 - 35000 = \dots$
\shortans{55000}
\loigiai{
$90000 - 35000 = 55000$.
}
\end{ex}

\begin{ex}
Tính: $16000 \times 4 = \dots$
\shortans{64000}
\loigiai{
$16000 \times 4 = 64000$.
}
\end{ex}

\begin{ex}
Tính: $84000 : 4 = \dots$
\shortans{21000}
\loigiai{
$84000 : 4 = 21000$.
}
\end{ex}

\begin{ex}
Tính giá trị biểu thức: $15000 \times 3 + 25000$.
\shortans{70000}
\loigiai{
$45000 + 25000 = 70000$.
}
\end{ex}

\begin{ex}
Tính giá trị biểu thức: $100000 - 60000 : 2$.
\shortans{70000}
\loigiai{
$100000 - 30000 = 70000$.
}
\end{ex}

\begin{ex}
Tìm $x$: $x \times 5 = 60000$.
\shortans{12000}
\loigiai{
$x = 60000 : 5 = 12000$.
}
\end{ex}

\begin{ex}
Tìm $x$: $x : 3 = 24000$.
\shortans{72000}
\loigiai{
$x = 24000 \times 3 = 72000$.
}
\end{ex}

\begin{ex}
Một hình chữ nhật có chiều dài 15 cm, chiều rộng 8 cm. Chu vi của nó là bao nhiêu xăng-ti-mét?
\shortans{46}
\loigiai{
$(15 + 8) \times 2 = 46$ cm.
}
\end{ex}

\begin{ex}
Một hình chữ nhật có chiều dài 15 cm, chiều rộng 8 cm. Diện tích của nó là bao nhiêu xăng-ti-mét vuông?
\shortans{120}
\loigiai{
$15 \times 8 = 120\text{ cm}^2$.
}
\end{ex}

\begin{ex}
Một hình vuông có chu vi 36 cm. Cạnh của hình vuông đó dài bao nhiêu xăng-ti-mét?
\shortans{9}
\loigiai{
$36 : 4 = 9$ cm.
}
\end{ex}

\begin{ex}
Một hình vuông có cạnh dài 9 cm. Diện tích của nó là bao nhiêu xăng-ti-mét vuông?
\shortans{81}
\loigiai{
$9 \times 9 = 81\text{ cm}^2$.
}
\end{ex}

\begin{ex}
Hình tròn có bán kính 6 cm thì đường kính là bao nhiêu xăng-ti-mét?
\shortans{12}
\loigiai{
$6 \times 2 = 12$ cm.
}
\end{ex}

\begin{ex}
Khối lập phương có bao nhiêu đỉnh?
\shortans{8}
\loigiai{
Có 8 đỉnh.
}
\end{ex}

\begin{ex}
Khối hộp chữ nhật có bao nhiêu mặt?
\shortans{6}
\loigiai{
Có 6 mặt.
}
\end{ex}

\begin{ex}
Đổi 1 tờ 100 000 đồng lấy các tờ 20 000 đồng được bao nhiêu tờ?
\shortans{5}
\loigiai{
$100000 : 20000 = 5$ tờ.
}
\end{ex}

\begin{ex}
Mua 4 chiếc bút bi giá 5 000 đồng mỗi chiếc. Đưa tờ 50 000 đồng thì được trả lại bao nhiêu đồng?
\shortans{30000}
\loigiai{
$50000 - 4 \times 5000 = 30000$ đồng.
}
\end{ex}

\begin{ex}
Tháng 8 có bao nhiêu ngày?
\shortans{31}
\loigiai{
Tháng 8 có 31 ngày.
}
\end{ex}

\begin{ex}
Tháng 9 có bao nhiêu ngày?
\shortans{30}
\loigiai{
Tháng 9 có 30 ngày.
}
\end{ex}

\begin{ex}
Một ngày có bao nhiêu giờ?
\shortans{24}
\loigiai{
Có 24 giờ.
}
\end{ex}

\begin{ex}
Làm tròn số 64 800 đến hàng nghìn ta được số nào?
\shortans{65000}
\loigiai{
Làm tròn lên 65 000.
}
\end{ex}

\begin{ex}
Số La Mã XV tương ứng với số tự nhiên nào?
\shortans{15}
\loigiai{
XV là 15.
}
\end{ex}

\newpage

\end{document}
